%%%%%%%%%%%%%%%%%%%%%%%%%%%%%%%%%%%%%%%%%%%%%%%%%%%%%%%%%%%%%%%%%
% Contents: Customising LaTeX output
%%%%%%%%%%%%%%%%%%%%%%%%%%%%%%%%%%%%%%%%%%%%%%%%%%%%%%%%%%%%%%%%%
\chapter{Personalización de \LaTeX}

\begin{intro}
Los documentos producidos mediante las órdenes que ha aprendido hasta este punto parecerán aceptables a una amplia audiencia.  Aunque no tienen un aspecto extraordinario, obedecen todas las reglas establecidas de composición correcta, lo que los hará fáciles de leer y plácidos a la vista.

Sin embargo, hay situaciones donde \LaTeX{} no proporciona una orden o entorno que cubra sus necesidades, o la salida producida por algunas órdenes existentes puede no satisfacer sus expectativas.

En este capítulo, se darán algunas pistas para enseñar a \LaTeX{} nuevos trucos y hacerle producir salidas con diferente aspecto del producido por omisión.
\end{intro}


\section{Nuevas órdenes, entornos y paquetes}

Puede haber notado que todas las órdenes que presento en este libro se componen en una caja, y que se muestran en el índice al final del libro.  En lugar de usar directamente las órdenes \LaTeX{} necesarias para conseguirlo, he creado un \wi{paquete} en que defino nuevas órdenes y entornos con este propósito.  Ahora puedo escribir simplemente:

\begin{example}
\begin{lscommand}
\ci{dum}
\end{lscommand}
\end{example}

En este ejemplo, estoy usando tanto un nuevo entorno llamado\\ \ei{lscommand}, que es responsable de dibujar la caja alrededor de la orden, y una nueva orden llamada \ci{ci}, que compone el nombre de la orden y hace la correspondiente entrada en el índice.  Puede comprobarlo buscando la orden \ci{dum} en el índice al final del libro, donde pude encontrar una entrada para \ci{dum}, apuntando a cada página donde he mencionado la orden \ci{dum}.

Si alguna vez decido que no me gusta que las órdenes se compongan en una caja, puedo simplemente cambiar la definición del entorno \texttt{lscommand} para crear un nuevo aspecto.  Esto es mucho más fácil que ir por todo el documento localizando todos los lugares en que he usado comandos \LaTeX{} genéricos para dibujar una caja alrededor de una palabra.


\subsection{Órdenes nuevas}

Para añadir sus órdenes nuevas, use la orden
\begin{lscommand}
\ci{newcommand}\verb|{|%
       \emph{nombre}\verb|}[|\emph{núm}\verb|]{|\emph{definición}\verb|}|
\end{lscommand}
%command. 
Básicamente, lo orden requiere dos argumentos: el \emph{nombre} de la orden que quiere crear, y la \emph{definición} de la orden.  El argumento \emph{núm} entre corchetes es opcional e indica el número de argumentos que toma la nueva orden (hasta 9 son posibles).  Si no se indica el valor es 0, es decir, no se permiten argumentos.

Los siguientes dos ejemplos deberían ayudarle a entender la idea. El primer ejemplo define una nueva orden llamada \ci{intc}.  Es la abreviatura de ``La introducción no-tan-corta a \LaTeXe''.  Tal orden podría ser útil si tuviera que escribir el título del libro una y otra vez.

\begin{example}
\newcommand{\intc}{La 
    introducción no-tan-corta a
    \LaTeXe}
Esto es ``\intc'' \ldots{} 
``\intc''
\end{example}

El siguiente ejemplo ilustra cómo definir una orden nueva que toma un argumento. Los caracteres \verb|#1| se sustituyen por el argumento indicado.  Si quisiera usar un segundo argumento, use \verb|#2| y así sucesivamente.

\begin{example}
\newcommand{\txsit}[1]
 {Esta es la Introducción
   \emph{#1}-corta a \LaTeXe}
% en el cuerpo del documento: 
\begin{itemize}
    \item \txsit{no-tan}
    \item \txsit{súper}
\end{itemize}
\end{example}

\LaTeX{} no le permitirá crear una nueva orden sobre una ya existente.  Pero hay una orden especial en el caso de que explícitamente quisiera reemplazarla: \ci{renewcommand}. Usa la misma sintaxis que la orden \verb|\newcommand|.

En ciertos casos puede querer usar la orden \ci{providecommand}.  Funciona como \ci{newcommand} y hace que la orden sea definida si aún no existe, pero no hace nada si ya estaba definida.

Hay algunos puntos que comentar sobre los espacios que siguen a las órdenes de \LaTeX{}.  Vea la página  \pageref{whitespace} para más información.

\subsection{Nuevos entornos}

Similar a la orden \verb|\newcommand|, hay una orden para crear sus propios entornos.  La orden \ci{newenvironment} usa la siguiente sintaxis:

\begin{lscommand}
\ci{newenvironment}\verb|{|%
       \emph{nombre}\verb|}[|\emph{núm}\verb|]{|%
       \emph{antes}\verb|}{|\emph{después}\verb|}|
\end{lscommand}

También \ci{newenvironment} puede tener un argumento opcional.  El material indicado en el argumento \emph{antes} se procesa antes de que se procese el texto del entorno.  El material en el argumento \emph{después} se procesa cuando se encuentra la orden \verb|\end{|\emph{nombre}\verb|}|.

El ejemplo siguiente ilustra el uso de la orden \ci{newenvironment}.
\begin{example}
\newenvironment{king}
 {\rule{1ex}{1ex}%
      \hspace{\stretch{1}}}
 {\hspace{\stretch{1}}%
      \rule{1ex}{1ex}}

\begin{king} 
Mis humildes ideas\ldots
\end{king}
\end{example}

El argumento \emph{núm} se usa igual que con la orden \verb|\newcommand|. \LaTeX{} se asegura de que usted no defina un entorno que ya existe; pero si quiere alguna vez cambiar un entorno existente, puede usar la orden \ci{renewenvironment}.  Usa la misma sintaxis que la orden \ci{newenvironment}.

La orden usada en este ejemplo se explicará más tarde.  Para la orden \ci{rule} véase la página \pageref{sec:rule}, para \ci{stretch} vaya a la página \pageref{cmd:stretch}, y puede hallar más información sobre \ci{hspace} en la página \pageref{sec:hspace}.

\subsection{Espacio extra}

Al crear un entorno nuevo puede hallar dificultades en el manejo del espacio adicional, que puede llegar a tener efectos fatales.  Por ejemplo, cuando quiera crear un entorno para títulos que suprima su propia sangría así como la del siguiente párrafo.  La orden \ci{ignorespaces} en el bloque de comienzo del entorno hará que éste prescinda de cualquier espacio tras ejecutar el bloque de comienzo. El bloque final requiere un poco más de cuidado porque tiene lugar un proceso especial al final del entorno.  La orden \ci{ignorespacesafterend} hará que \LaTeX{} ejecute \ci{ignorespaces} después de que el proceso especial tenga lugar.

\begin{example}
\newenvironment{simple}%
 {\noindent}%
 {\par\noindent}

\begin{simple}
Mire el espacio\\a la izquierda.
\end{simple}
También\\aquí.
\end{example}

\begin{example}
\newenvironment{correct}%
 {\noindent\ignorespaces}%
 {\par\noindent%
   \ignorespacesafterend}

\begin{correct}
Sin espacio\\a la izquierda.
\end{correct}
También\\aquí.
\end{example}

\subsection{Línea de órdenes \LaTeX}

Si trabaja en un sistema operativo estilo \textsc{posix} (GNU o \textsc{unix}), quizás use \ci{Makefile} para compilar sus documentos de \LaTeX{}.  Entonces podría ser interesante producir diferentes versiones del mismo documento llamando a \LaTeX{} con diversos parámetros en la línea de órdenes.  Si añade la siguiente estructura a su documento:

\begin{verbatim}
\usepackage{ifthen}
\ifthenelse{\equal{\blancoynegro}{verdadero}}{
  % modo "blanco y negro"; hacer algo..
}{
  % modo "color"; hacer algo diferente..
}
\end{verbatim}

Ahora puede llamar a \LaTeX{} así:
\begin{verbatim}
latex '\newcommand{\blancoynegro}{verdadero}

%English version
\newcommand{\TFRGB}[2]{#2} % #1 en francais #2 in english

\newcommand{\dft}{By default}
\newcommand{\SSCT}[2]{ \section{#2}}
\newcommand{\SbSSCT}[2]{ \subsection{#2}}
\newcommand{\Par}[2]{ \paragraph{#2}}
\newcommand{\SbSbSSCT}[2]{ \subsubsection{#2}}
\newcommand{\SSCTTC}[4]{ \section[#2]{#2}}
\newcommand{\SbSSCTTC}[4]{ \subsection[#2]{#2}}
\newcommand{\SbSbSSCTTC}[4]{ \subsubsection[#3]{#4}}

\newcommand{\maboite}[1]{\begin{center}  \tikz \draw node[draw,fill=yellow!20,inner sep=0.2cm,text centered,text width=.75\linewidth] {Load package : #1} ; \end{center} }

%\newcommand{\DW}[1]{\begin{center}  \tikz \draw node[draw,fill=red!50,inner sep=0.1cm,text centered] {\tiny don't work !} ; \end{center} }

\newcommand{\DW}[1]{\tikz[baseline=-1mm]  \draw node[draw,fill=red!50] {{\tiny don't work !}}; \index{\textbf{6 list of don't work }}} % for english version
%% Version française

\newcommand{\TFRGB}[2]{#1} % #1 en francais #2 in english
\newcommand{\dft}{Par défaut : }
\newcommand{\SSCT}[2]{\section{#1}}
\newcommand{\SbSSCT}[2]{\subsection{#1}}
\newcommand{\Par}[2]{ \paragraph{#1}}
\newcommand{\SbSbSSCT}[2]{\subsubsection{#1}}
\newcommand{\SSCTTC}[4]{\section[#1]{#2}}
\newcommand{\SbSSCTTC}[4]{\subsection[#1]{#2}}
\newcommand{\SbSbSSCTTC}[4]{\subsubsection[#1]{#2}}



\newcommand{\maboite}[1]{\begin{center}  \tikz \draw node[draw,fill=yellow!20,inner sep=0.2cm,text centered,text width=.75\linewidth] {Charger l'extension: #1} ; \end{center} }


\newcommand{\DW}[1]{\tikz[baseline=-1mm]  \draw node[draw,fill=red!50] {{\tiny ne fonctionne pas !}}; \index{\textbf{6 liste des non-fonctionnels }}} % pour la version française

 \documentclass[a4paper,10pt]{article}

 \usepackage{fontspec}
\usepackage[french,english]{babel}

%\TFRGB{\selectlanguage{french}}{\selectlanguage{english}}
\usepackage{tikzpeople}

\usepackage{amsmath,amsfonts,amssymb}

\usepackage{pdfpages}  


%\usepackage{pst-all}

\usepackage{graphicx} 
\usepackage{hyperref}

\usepackage{animate}
\usepackage{makeidx}
%\usepackage{wrapfig}
%\usepackage{tikz-dependency}
\usepackage{pgfplots} %<<<<<<<<<<<<<<<<<<<<<<<<<<<<< 
\usepackage{tikz}
\usepackage{tkz-tab}

 
%\usepgflibrary{shapes.callouts}
\usepackage{tikz-qtree}
\usepackage{tkz-tab}
\usepackage{csquotes}

  
\usetikzlibrary{angles}
\usetikzlibrary{arrows}

\usetikzlibrary{shadings}
\usetikzlibrary{calc}
\usetikzlibrary{backgrounds}
\usetikzlibrary{decorations.pathmorphing}

\usetikzlibrary{decorations.markings}
\usetikzlibrary{decorations.footprints}
\usetikzlibrary{decorations.shapes}
\usetikzlibrary{decorations.text}
\usetikzlibrary{decorations.fractals}
\usepgflibrary{shapes.geometric}
\usetikzlibrary{intersections}
\usetikzlibrary{scopes}
\usetikzlibrary{shapes.symbols}
\usetikzlibrary{shapes.arrows}
\usetikzlibrary{shapes.callouts}
\usetikzlibrary{shapes.misc}
\usepgflibrary{shapes.multipart}
\usetikzlibrary{plotmarks}
\usetikzlibrary{trees}
\usetikzlibrary{fadings}
\usetikzlibrary{arrows.meta}
\usetikzlibrary{bending}
\usetikzlibrary{fit}
%\usetikzlibrary{circuits}
\usetikzlibrary{circuits.ee.IEC}
\usetikzlibrary{circuits.logic.IEC}
\usetikzlibrary[circuits.logic.US]
\usetikzlibrary{circuits.logic.CDH}
%\usetikzlibrary{decorations}
\usetikzlibrary{shapes.gates.logic.IEC}
\usetikzlibrary{matrix}
\usetikzlibrary{chains}
%\usetikzlibrary{circuit.plc.sfc}
\usepackage{tikzsymbols}
\usetikzlibrary{datavisualization}
\usetikzlibrary{datavisualization.formats.functions}
%
\usepackage{tikzducks}

\usepackage{tikzrput}
\usepackage{pgfornament}

%\usetikzlibrary{babel}
\usetikzlibrary{math}
\usetikzlibrary{optics}
\usetikzlibrary{through}
\usetikzlibrary{turtle}
\usetikzlibrary{quotes}


\pgfplotsset{compat=1.8}
\usetikzlibrary{positioning}

\usepackage{geometry}
\geometry{a4paper,top={3cm}}

\usepackage{ifpdf}
\usepackage{ifluatex}
\usetikzlibrary{spy}




%====================================================================

\makeindex

\newcommand{\AC}[1]{\{#1\}}

\newcommand{\BS}[1]{$\backslash$#1}

\newcommand{\BSB}[1]{\textbf{\color{blue} {$\backslash$#1}}}


\newcommand{\BSR}[1]{\textbf{\color{red}  $\backslash$#1}}

%\newcommand{\RDDX}[2]{{\color{red}#1} \index{\textbf{3 Paramètres et options}!#2=#1}}


\newcommand{\RRR}[1]{\tikz[baseline=-1mm,inner sep=2pt]  \draw node[draw,fill=red!20] {{\footnotesize  PGFmanual section :  #1}} ; }

\newcommand{\RRP}[1]{\tikz[baseline=-1mm]  \draw node[draw,fill=red!20] {{\footnotesize  pgfplots section :  #1}} ; }

%\newcommand{\RRR}[1]{\tikz[baseline=-1mm]  \draw node[draw,fill=red!20] {{\footnotesize  PGFmanual section :  #1}} ;\index{\textbf{5 PGFmanual }!#1} }

\newcommand{\DFR}{ \tikzpicture[scale=.25]
\draw [fill=blue](0,0) rectangle (3,1.5);
\draw [fill=white](1,0) rectangle (2,1.5);
\draw [fill=red](2,0) rectangle (3,1.5);\endtikzpicture }

\newcommand{\DGE}{ \tikzpicture[scale=.25]
\draw [fill=yellow](0,0) rectangle (3,.5);
\draw [fill=red]((0,.5) rectangle (3,1);
\draw [fill=black](0,1) rectangle (3,1.5); \endtikzpicture }

\newcommand{\DGB}{ \tikzpicture[scale=.25]
\draw [fill=blue](0,0) rectangle (3,1.5);
\draw [white,line width=.1cm](0,0) -- (3,1.5);
\draw [white,line width=.1cm](0,1.5) -- (3,0);
\draw [white,line width=.1cm](1.5,0) -- (1.5,1.5);
\draw [white,line width=.1cm](0,0.75) -- (3,0.75);
\draw [red,line width=.05cm](0,0) -- (3,1.5);
\draw [red,line width=.05cm](0,1.5) -- (3,0);
\draw [red,line width=.05cm](1.5,0) -- (1.5,1.5);
\draw [red,line width=.05cm](0,0.75) -- (3,0.75);
\endtikzpicture }



\TFRGB{
\newcommand{\ESS}[1]{\textbf{\textbackslash begin\AC{#1}}\index{\textbf{1 Environnements}!#1}}

\newcommand{\BSS}[1]{\textbf{\textbackslash{#1}}\index{\textbf{2 Commandes}!#1 @\textbackslash{}#1}}


\newcommand{\DDD}[1]{{\color{red}  #1}\index{\textbf{3 Paramètres et options}!#1}}
\newcommand{\RDD}[1]{{\color{red}  #1}\index{\textbf{3 Paramètres et options}!#1}}
%4
\newcommand{\BDD}[1]{{\color{blue}  #1}\index{\textbf{4 Valeurs Tikz}!#1}}

\newcommand{\RDDX}[2]{{\color{red}#1} \index{\textbf{4 Valeurs Tikz}!#1 (#2)}}
%5
\newcommand{\FDD}[1]{{\color{red}  #1}\index{\textbf{5 Extrémit\'es}!#1}}
}
{
\newcommand{\ESS}[1]{\textbf{\textbackslash{#1}}\index{\textbf{1 Environments}!#1 @\textbackslash{}#1}}
%2
\newcommand{\BSS}[1]{\textbf{\textbackslash{#1}}\index{\textbf{2 Commands}!#1 @\textbackslash{}#1}}
%3
\newcommand{\RDD}[1]{{\color{red}  #1}\index{\textbf{3 Parameters and options}!#1}}

\newcommand{\DDD}[1]{{\color{red}  #1}\index{\textbf{3 Parameters and options}!#1}}

\newcommand{\RDDX}[2]{{\color{red}#1} \index{\textbf{4 Values Tikz}!#1 (#2)}}
%
\newcommand{\BDD}[1]{{\color{blue}  #1}\index{\textbf{4 Values Tikz}!#1}}

\newcommand{\FDD}[1]{{\color{red}  #1}\index{\textbf{5 Extremities}!#1}}
}

\newcommand{\rouge}[1] {{\color{red}  #1}}
\newcommand{\blll}[1] {{\color{blue}  #1}}



 \begin{document}


%\selectlanguage{english}
\selectlanguage{french}



%  
\author{{\Huge Jean Pierre Casteleyn } \\ {\Huge IUT Génie Thermique et \'Energie } \\ {\Huge Dunkerque, France }}

\DeclareFixedFont{\RM}{T1}{ptm}{b}{n}{2cm}

\DeclareFixedFont{\RMM}{T1}{ptm}{b}{n}{1cm}

\title{ {\RM Visual TikZ} \\ \vspace{1cm} {\RMM Version 0.66} }



\date{
\begin{center}
\begin{animateinline}[loop,autoplay]{12}%
 \multiframe{24}{iAngle=0+15,icol=0+5}{\begin{tikzpicture}[rotate=90]
    \draw  (0,0) node[fill=white,circle] {\includegraphics[width=4cm]{LogoIUT}}  (0,0) circle (1);
  \end{tikzpicture}} 
\end{animateinline}% 
\end{center}
{\LARGE \TFRGB{mis à jour le \today}{Updated on \today} 
}
}


\maketitle



 \begin{animateinline}[autoplay,loop]{12}%
 \multiframe{24}{iAngle=0+15,icol=0+5}{\begin{tikzpicture}
 [scale=1.8] %
   \draw[line width=0pt] (-2,-2) rectangle(6,2); %
   \draw  (0,0) node[fill=white,circle,rotate=\iAngle] {\includegraphics[width=2cm]{LogoIUT}}  (0,0) circle (1);
    \draw (0,0) circle (1);
    \coordinate (abc) at (${sqrt(9-sin(\iAngle)*sin(\iAngle))+cos(\iAngle)}*(1,0)$) ;
    \coordinate (xyz) at (\iAngle:1);
    \draw[ultra thick] (0,0) --(xyz); 
    \draw[ultra thick] (xyz) -- (abc) ;
    \fill[color=blue!\icol] (abc)++(0.5,-1) rectangle (5,1) ;
    \draw[ultra thick] (abc) ++(0,-1) rectangle ++(.5,2) ;
    \draw[ultra thick]  (1.5,1) -- (5,1) -- (5,-1) -- (1.5,-1);
    \fill[red] (xyz) circle (4pt);
    \fill[red] (abc) circle (4pt); 
  \end{tikzpicture}}
 \end{animateinline} 


 
\newpage
 
 
\TFRGB{
\textbf{Objectifs }: 

\begin{itemize}
\item Avoir une image par  commande ou par paramètre.
\item Avoir un texte réduit au strict minimum.
\item Etre le plus complet possible au fil de mises à jour régulières.
\item Garder la même structure que visuel pstricks
\end{itemize} 
}
{\textbf{Objectives }: 

\begin{itemize}
\item One image per command or parameter.
\item the minimum amount of text possible.
\item the most complete possible update after update.
\item keep the same structure as VisualPSTricks
\end{itemize}}


\vspace{1cm}

\TFRGB{
\textbf{Remarques }: Le code donné est minimal et ne sert qu'à montrer les commandes concernées. Les effets sont parfois exagérés pour bien les mettre en évidence. Pour en savoir plus, vous pouvez voir la documentation. Pour se faire j'ai indiqué le numéro de \tikz[baseline=-1mm]  \draw node[draw,fill=red!20] {Section de pgfmanual} ;
}
{\textbf{Remarks }:
Minimal code is given to show the effect of a command or a parameter. The effects are sometime exaggerated for clarity   .To consult the documentation, I have given the number of the  \tikz[baseline=-1mm]  \draw node[draw,fill=red!20] {Section in pgfmanual} ;
}
\vspace{1cm}


\TFRGB{
\textbf{Vous pouvez me contacter à}
 \href{mailto:jpcdk@yahoo.fr}{mon e-mail personnel} pour

\begin{itemize}
\item me signaler les erreurs que vous avez constatés (merci d'indiquer la page où vous l'avez constaté)
\item me faire part de vos commentaires, suggestions \dots
\end{itemize}}
{
\textbf{You can contact me at }
 \href{mailto:jpcdk@yahoo.fr}{my personal email} to

\begin{itemize}
\item let me know the mistakes found (please indicate the page)
\item give me your commentaries, your suggestions \dots
\end{itemize}}

\vspace{1cm}
\TFRGB{
\textbf{Quoi de neuf ! } :

\begin{itemize}
\item Ajout de la library  chains \pageref{lib-chains}
\item Ajout de la library  through \pageref{lib-through}
\item Ajout de la library  turtle \pageref{lib-turtle}
\item Ajout de la library positioning \pageref{lib-pos}
\item Ajout du module tikzsymbols \pageref{symbol}
\item mise à jour du module tikzducks \pageref{ducks}
\item mise à jour des modules shape \pageref{formes}
\end{itemize}

}
{
\textbf{What's new } :
\begin{itemize}
\item chains library added \pageref{lib-chains}
\item through library added \pageref{lib-through}
\item turtle library added \pageref{lib-turtle}
\item positioning library added \pageref{lib-pos}
\item Tikzsymbols package added \pageref{symbol}
\item Tikzducks package updated \pageref{ducks}
\item shapes packages updated \pageref{formes}
\end{itemize}
}



\vspace{1cm}
\textbf{Licence } :


This work may be distributed and/or modified under the conditions of the LaTeX Project Public License, either version 1.3 of this license or (at your option) any later version.

 The latest version of this license is in  http://www.latex-project.org/lppl.txt and version 1.3 or later is part of all distributions of LaTeX
version 2005/12/01 or later.

This work has the LPPL maintenance status `maintained'.

The Current Maintainer of this work is M. Jean Pierre Casteleyn.

\vspace{2cm}
\textbf{\TFRGB{Merci à }{Thanks to}}:

Till Tantau  ,
Alain Matthes ,
Jim Diamond ,
Falk Rühl ,
Axel Kielhorn ,
Nils Fleischhacker ,
Michel Fruchart ,
Ben Vitecek
 
\newpage


% \tableofcontents
%
%
%\newpage




\RRR{75-2 = Concept: Data Points and Data Formats}

\begin{tikzpicture}
\datavisualization [school book axes, visualize as smooth line]
data {
x, y
-1.5, 2.25
-1, 1
-.5, .25
0, 0
.5, .25
1, 1
1.5, 2.25
};
\end{tikzpicture}


\begin{tikzpicture}
\datavisualization [school book axes, visualize as smooth line]
data [format=function] {
var x : interval [-1.5:1.5] samples 7;
func y = \value x*\value x;
};
\end{tikzpicture}


\begin{tikzpicture}
\datavisualization [school book axes, visualize as smooth line]
data [format=function] {
var x : interval [-1.5:1.5] samples 3;
func y = \value x*\value x;
};
\end{tikzpicture}

Section 76 gives an in-depth coverage of the available data formats and explains how new data formats
can be defined.


\RRR{75-3 = Concept: Axes, Ticks, and Grids}


\begin{tikzpicture}
\datavisualization [
scientific axes,
x axis={length=3cm, ticks=few},
visualize as smooth line
]
data [format=function] {
var x : interval [-1.5:1.5] samples 20;
func y = \value x*\value x;
};
\end{tikzpicture}

\begin{tikzpicture}
\datavisualization [
scientific axes=clean,
x axis={length=3cm, ticks=few},
all axes={grid},
visualize as smooth line
]
data [format=function] {
var x : interval [-1.5:1.5] samples 7;
func y = \value x*\value x;
};
\end{tikzpicture}

Section 77 explains in more detail how axes, ticks, and grid lines can be chosen and configured.


\RRR {75-4 = Concept: Visualizers}

\begin{tikzpicture}
\datavisualization [
scientific axes=clean,
x axis={length=3cm, ticks=few},
visualize as smooth line
]
data [format=function] {
var x : interval [-1.5:1.5] samples 7;
func y = \value x*\value x;
};
\end{tikzpicture}

\begin{tikzpicture}
\datavisualization [
scientific axes=clean,
x axis={length=3cm, ticks=few},
visualize as scatter
]
data [format=function] {
var x : interval [-1.5:1.5] samples 7;
func y = \value x*\value x;
};
\end{tikzpicture}

Section 78 provides more information on visualizers as well as reference lists.

\RRR{75-5 = Concept: Style Sheets and Legends }

\begin{tikzpicture}[baseline]
\datavisualization [ scientific axes=clean,
y axis=grid,
visualize as smooth line/.list={sin,cos,tan},
style sheet=strong colors,
style sheet=vary dashing,
sin={label in legend={text=$\sin x$}},
cos={label in legend={text=$\cos x$}},
tan={label in legend={text=$\tan x$}},
data/format=function ]
data [set=sin] {
var x : interval [-0.5*pi:4];
func y = sin(\value x r);
}
data [set=cos] {
var x : interval [-0.5*pi:4];
func y = cos(\value x r);
}
data [set=tan] {
var x : interval [-0.3*pi:.3*pi];
func y = tan(\value x r);
};
\end{tikzpicture}


Section 79 details style sheets and legends.

\RRR{75-6 = Usage}

\subsection{/pgf/data/read from file=filename} (no default, initially empty)

If you set the source attribute to a non-empty hfilenamei, the data will be read from this file. In
this case, no hinline datai may be present, not even empty curly braces should be provided.
%\datavisualization ...
data [read from file=file1.csv]
data [read from file=file2.csv];
The other way round, if read from file is empty, the data must directly follow as hinline datai.
%\datavisualization ...
data {
x, y
1, 2
2, 3
};

The second important key is format, which is used to specify the data format:

\subsection{/pgf/data/format}

Use this key to locally set the format used for parsing the data, see Section 76 for a list of predefined
formats.

\tikz
\datavisualization [school book axes, visualize as line]
data [/data point/x=1] {
y
1
2
}
data [/data point/x=2] {
y
2
0
.5
};

\BS{datavisualization} . . . data point[options] . . . ;

\tikz \datavisualization [school book axes, visualize as line]
data point [x=1, y=1] data point [x=1, y=2]
data point [x=2, y=2] data point [x=2, y=0.5];

/tikz/data visualization/data point=options

\tikzdatavisualizationset{
horizontal/.style={
data point={x=#1, y=1}, data point={x=#1, y=2}},
}
\tikz \datavisualization
[ school book axes, visualize as line,
horizontal=1,
horizontal=2 ];

\BS{datavisualization} . . . data group[options]\AC{name}+=\AC{data specifications} . . . ;


\tikz \datavisualization data group {points} = {
data {
x, y
0, 1
1, 2
2, 2
5, 1
2, 0
1, 1
}
};

\tikz \datavisualization [school book axes, visualize as line] data group {points};
\qquad
\tikz \datavisualization [scientific axes=clean, visualize as line] data group {points};


\BS{datavisualization} . . . scope[options]{data specification} . . . ;

%\datavisualization...
%scope [/data point/experiment=7]
%{
%data [read from file=experiment007-part1.csv]
%data [read from file=experiment007-part2.csv]
%data [read from file=experiment007-part3.csv]
%}
%scope [/data point/experiment=23, format=foo]
%{
%data [read from file=experiment023-part1.foo]
%data [read from file=experiment023-part2.foo]
%};


\BS{datavisualization} . . . info[options]{code} . . . ;

\begin{tikzpicture}[baseline]
\datavisualization [ school book axes, visualize as line ]
data [format=function] {
var x : interval [-0.1*pi:2];
func y = sin(\value x r);
}
info {
\draw [red] (visualization cs: x={(.5*pi)}, y=1) circle [radius=1pt]
node [above,font=\footnotesize] {extremal point};
};
\end{tikzpicture}

\subsection{Coordinate system visualization}

\BS{datavisualization} . . . info’[options]{code} . . . ;

\begin{tikzpicture}[baseline]
\datavisualization [ school book axes, visualize as line ]
data [format=function] {
var x : interval [-0.1*pi:2];
func y = sin(\value x r);
}
info' {
\fill [red] (visualization cs: x={(.5*pi)}, y=1) circle [radius=2mm];
};
\end{tikzpicture}


\subsection{Predefined node data visualization bounding box}
This rectangle node stores a bounding box of the data visualization that is currently being constructed.
This node can be useful inside info commands or when labels need to be added.

\subsection{Predefined node data bounding box}
This rectangle node is similar to data visualization bounding box, but it keeps track only of the actual
data. The spaces needed for grid lines, ticks, axis labels, tick labels, and other all other information
that is not part of the actual data is not part of this box.


\RRR{75-7 = Advanced: Executing User Code During a Data Visualization}

\RRR{75-8 = Advanced: Creating New Objects}


\section{76 Providing Data for a Data Visualization}


%
%\newpage
%
%\RRR{17-2-1} fini
%
%
%/tikz/node font=font commands
%
%\begin{tikzpicture}
%\draw[node font=\itshape] (1,0) -- +(1,1) node[above] {italic};
%\end{tikzpicture}
%
%\tikz \node [node font=\tiny, minimum height=3em, draw] {tiny};
%\tikz \node [node font=\small, minimum height=3em, draw] {small};
%
%
%
%
%/tikz/node align header=
%
%
%\RRR{17-4-4} OK
%
%
%\RRR{17-5} Positioning Nodes
%
%
%
%\RRR{17-5-3 }Advanced Placement Options
%
%
%  
%
%\subsubsection{title}
%
%\begin{tikzpicture}[every node/.style={draw}]
%\draw[help lines](0,0) grid (3,2);
%\draw (1,0) node{A}
%(2,0) node[rotate=90,scale=1.5] {B};
%\draw[rotate=30] (1,0) node{A}
%(2,0) node[rotate=90,scale=1.5] {B};
%\draw[rotate=60] (1,0) node[transform shape] {A}
%(2,0) node[transform shape,rotate=90,scale=1.5] {B};
%\end{tikzpicture}
%
%m:::::::::::::::%\begin{tikzpicture}
%%% Install a nonlinear transformation:
%%\pgfsetcurvilinearbeziercurve
%%{\pgfpoint{0mm}{20mm}}
%%{\pgfpoint{10mm}{20mm}}
%%{\pgfpoint{10mm}{10mm}}
%%{\pgfpoint{20mm}{10mm}}
%%\pgftransformnonlinear{\pgfpointcurvilinearbezierorthogonal\pgf@x\pgf@y}%
%%% Draw something:
%%\draw [help lines] (0,-30pt) grid [step=10pt] (80pt,30pt);
%%\foreach \x in {0,20,...,80}
%%\node [fill=red!20] at (\x pt, -20pt) {\x};
%%\foreach \x in {0,20,...,80}
%%\node [fill=blue!20, transform shape nonlinear] at (\x pt, 20pt) {\x};
%%\end{tikzpicture}
%
%
%\newpage
%
%
%
%
\SbSSCT{Coordonnées}{Coordinates}
\begin{center}
\RRR{13-2-1}
\end{center}


\SbSbSSCT{Système de coordonnées \og canvas \fg}{Canvas coordinates}

\noindent


\tikzset{every picture/.style=blue,very thick,inner sep=0pt}

\begin{tabular}{|c|c|} \hline 
\TFRGB{Explicite}{explicit}  & \TFRGB{Implicite}{implicit}
\\ \hline
\begin{tikzpicture}
\draw[help lines] (0,0) grid (3,2);
\fill (canvas cs:x=2cm,y=1.5cm) circle (2pt);
\end{tikzpicture}
&
\begin{tikzpicture}
\draw[help lines] (0,0) grid (3,2);
\fill (2,1.5) circle (2pt);
\end{tikzpicture}

\\ \hline  
 \BS{fill} (\RDD{canvas cs}:\blll{x=2cm,y=1.5cm}) circle (2pt);
& \BS{fill} {\color{blue}(2cm,1.5cm)} circle (2pt);
\\ \hline 
\end{tabular} 


\SbSbSSCT{Système de coordonnées polaire \og canvas \fg}{Polar coordinates}

\noindent


\begin{tabular}{|c|c|c|} \hline
\TFRGB{Explicite}{explicit}  & \TFRGB{Implicite}{implicit}
\\ \hline
\begin{tikzpicture}
\draw[help lines] (0,0) grid (3,2);
\draw [dotted](0,2) arc (90 :0 :2);
\draw [dotted](0,0) --(2,2);
\fill (canvas polar cs:angle=45,radius=2cm) circle (2pt);
\end{tikzpicture}
&
\begin{tikzpicture}
\draw[help lines] (0,0) grid (3,2);
\draw [dotted](0,2) arc (90 :0 :2);
\draw [dotted](0,0) --(2,2);
\fill (45:2cm) circle (2pt);
\end{tikzpicture}
\\ \hline 
\BS{fill} (\RDD{canvas polar cs}:\RDD{angle}=45,\RDD{radius}=2cm) circle (2pt);
&
\BS{fill} {\color{blue}(45:2cm)} circle (2pt);
\\ \hline 
\end{tabular} 

\bigskip
\begin{tabular}{|c|} \hline  
\begin{tikzpicture}
\draw[help lines] (0,0) grid (3,2);
\draw [dotted](0,2) arc (90 :0 :3 and 2);
\draw [dotted](0,0) --(3,2);
\fill (canvas polar cs:angle=45,x radius=3cm,y radius=2cm) circle (2pt);
\end{tikzpicture}
\\ \hline  
\BS{fill} (canvas polar cs:angle=45,\RDD{x radius}=3cm,\RDD{y radius}=2cm) circle (2pt);
\\ \hline 
\end{tabular}


\SbSbSSCT{Système de coordonnées  xyz}{xyz coordinates}

\noindent


\begin{tabular}{|c|c|c|} \hline 
\begin{tikzpicture}[->]
\draw (0,0) -- (xyz cs:x=1);
\draw[red] (0,0) -- (xyz cs:y=1);
\draw[magenta] (0,0) -- (xyz cs:z=1);
\end{tikzpicture}
&
\begin{tikzpicture}[->]
\draw (0,0) -- (1,0,0);
\draw[red]  (0,0) -- (0,1,0);
\draw[magenta]  (0,0) -- (0,0,1);
\end{tikzpicture}
\\ \hline 
\BS{draw} (0,0) - - (\RDD{xyz cs}:x=1); & \BS{draw}  (0,0) - - (1,0,0); \\
\BS{draw}[red]  (0,0) - - (\RDD{xyz cs}:y=1); &  \BS{draw}[red] (0,0) - - (0,1,0); \\
\BS{draw}[magenta]  (0,0) - - (\RDD{xyz cs}:z=1); &  \BS{draw}[magenta]   (0,0) - - (0,0,1); 
\\ \hline 

\end{tabular} 

 
\newpage

\SbSbSSCT{Coordinate system xyz polar}{Coordinate system xyz polar}

\noindent

\begin{tabular}{|c|c|c|} \hline
\TFRGB{Explicite}{explicit}  & \TFRGB{Implicite}{implicit}
\\ \hline
\begin{tikzpicture}
\draw[help lines] (0,0) grid (3,2);
\draw [dotted](0,2) arc (90 :0 :2);
\draw [dotted](0,0) --(2,2);
\fill (xyz polar cs:angle=45,radius=2) circle (2pt);
\end{tikzpicture}
&
\begin{tikzpicture}
\draw[help lines] (0,0) grid (3,2);
\draw [dotted](0,2) arc (90 :0 :2);
\draw [dotted](0,0) --(2,2);
\fill (45:2) circle (2pt);
\end{tikzpicture}
\\ \hline 
\BS{fill} (\RDD{xyz polar cs}:\RDD{angle}=45,\RDD{radius}=2) circle (2pt);
&
\BS{fill} {\color{blue}(45:2cm)} circle (2pt);
\\ \hline 
\end{tabular} 

\bigskip
\begin{tabular}{|c|} \hline  
\begin{tikzpicture}
\draw[help lines] (0,0) grid (3,2);
\draw [dotted](0,2) arc (90 :0 :3 and 2);
\draw [dotted](0,0) --(3,2);
\fill (xyz polar cs:angle=45,x radius=3,y radius=2) circle (2pt);
\end{tikzpicture}
\\ \hline  
\BS{fill} (xyz polar cs:angle=45,\RDD{x radius}=3,\RDD{y radius}=2) circle (2pt);
\\ \hline 
\end{tabular} 

\bigskip

\begin{tabular}{|c|c|c|} \hline
\multicolumn{2}{|c|}{\BS{begin}\AC{tikzpicture}{\color{red}[x=1.5cm,y=1cm]} }
\\ \hline
\begin{tikzpicture}[x=1.5cm,y=1cm]
\draw[help lines] (0,0) grid (3,2);
\draw [dotted](0,2) arc (90 :0 :2);
\draw [dotted](0,0) --(2,2);
\fill (xyz polar cs:angle=45,radius=2) circle (2pt);
\end{tikzpicture}
&
\begin{tikzpicture}[x=1.5cm,y=1cm]
\draw[help lines] (0,0) grid (3,2);
\draw [dotted](0,2) arc (90 :0 :2);
\draw [dotted](0,0) --(2,2);
\fill (45:2) circle (2pt);
\end{tikzpicture}
\\ \hline 
\BS{fill} (\RDD{xyz polar cs}:\RDD{angle}=45,\RDD{radius}=2) circle (2pt);
&
\BS{fill} {\color{blue}(45:2cm)} circle (2pt);
\\ \hline 
\end{tabular} 
\bigskip

\begin{tabular}{|c|c|c|} \hline
\multicolumn{2}{|c|}{\BS{begin}\AC{tikzpicture}{\color{red}[x=\AC{(0cm,1cm)},y=\AC{(-1cm,0cm)}]} }
\\ \hline
\begin{tikzpicture}[x={(0cm,1cm)},y={(-1cm,0cm)}]
\draw[help lines] (0,0) grid (3,2);
\draw [dotted](0,2) arc (90 :0 :2);
\draw [dotted](0,0) --(2,2);
\fill (xyz polar cs:angle=45,radius=2) circle (2pt);
\end{tikzpicture}
&
\begin{tikzpicture}[x={(0cm,1cm)},y={(-1cm,0cm)}]
\draw[help lines] (0,0) grid (3,2);
\draw [dotted](0,2) arc (90 :0 :2);
\draw [dotted](0,0) --(2,2);
\fill (45:2) circle (2pt);
\end{tikzpicture}
\\ \hline 
\BS{fill} (\RDD{xyz polar cs}:\RDD{angle}=45,\RDD{radius}=2) circle (2pt);
&
\BS{fill} {\color{blue}(45:2cm)} circle (2pt);
\\ \hline 
\end{tabular} 

\SbSbSSCT{Coordonnées barycentriques}{Barycentric coordinates}

\begin{center}
\RRR{13-2-2}
\end{center}

\begin{tabular}{|c|c|c|} \hline
\multicolumn{3}{|c|}{  \BS{node} [circle,fill=red!20] at (\RDD{barycentric cs}:A=0.6,B=0.3 ) \AC{X};   }\\ 
\hline
\begin{tikzpicture}[scale=.6]
\draw[help lines] (0,0) grid (4,4);
\node[circle,fill=green!20,] (A) at (0,0) {A};
\node[circle,fill=green!20,] (B) at (4,0) {B};
\node[circle,fill=red!20] at (barycentric cs:A=0.3,B=0.3 ) {X};
\end{tikzpicture}
&
\begin{tikzpicture}[scale=.6]
\draw[help lines] (0,0) grid (4,4);
\node[circle,fill=green!20,] (A) at (0,0) {A};
\node[circle,fill=green!20,] (B) at (4,0) {B};
\node[circle,fill=green!20,] (C) at (4,4) {C};
\node[circle,fill=red!20] at (barycentric cs:A=0.4,B=0.4 ,C=.4) {X};
\end{tikzpicture}
&
\begin{tikzpicture}[scale=.6]
\draw[help lines] (0,0) grid (4,4);
\node[circle,fill=green!20,] (A) at (0,0) {A};
\node[circle,fill=green!20,] (B) at (4,0) {B};
\node[circle,fill=green!20,] (C) at (1,4) {C};
\node[circle,fill=green!20,] (D) at (4,4) {D};
\node[circle,fill=red!20] at (barycentric cs:A=0.5,B=0.5,C=.5,D=.5 ) {X};
\end{tikzpicture}
\\ \hline
A=0.3,B=0.3 & A=0.4,B=0.4 ,C=.4 & A=0.5,B=0.5,C=.5,D=.5 
\\ \hline
\begin{tikzpicture}[scale=.6]
\draw[help lines] (0,0) grid (4,4);
\node[circle,fill=green!20,] (A) at (0,0) {A};
\node[circle,fill=green!20,] (B) at (4,0) {B};
\node[circle,fill=red!20] at (barycentric cs:A=0.6,B=0.3 ) {X};
\end{tikzpicture}
&
\begin{tikzpicture}[scale=.6]
\draw[help lines] (0,0) grid (4,4);
\node[circle,fill=green!20,] (A) at (0,0) {A};
\node[circle,fill=green!20,] (B) at (4,0) {B};
\node[circle,fill=green!20,] (C) at (4,4) {C};
\node[circle,fill=red!20] at (barycentric cs:A=0.2,B=0.4 ,C=.6) {X};
\end{tikzpicture}
&
\begin{tikzpicture}[scale=.6]
\draw[help lines] (0,0) grid (4,4);
\node[circle,fill=green!20,] (A) at (0,0) {A};
\node[circle,fill=green!20,] (B) at (4,0) {B};
\node[circle,fill=green!20,] (C) at (1,4) {C};
\node[circle,fill=green!20,] (D) at (4,4) {D};
\node[circle,fill=red!20] at (barycentric cs:A=0.2,B=0.4,C=.6,D=.8 ) {X};
\end{tikzpicture}
\\ \hline
A=0.6,B=0.3 & A=0.2,B=0.4 ,C=.6 & A=0.2,B=0.4,C=.6,D=.8
\\ \hline
\end{tabular}

\SbSbSSCT{Coordonnées nominatives : n\oe ud}{Named coordinates: nodes}

\begin{center}
\RRR{13-2-3}
\end{center}

\begin{tabular}{|c|c|} \hline  
\begin{tikzpicture}[blue,very thick,baseline=1cm]
\draw[help lines] (0,0) grid (3,3);
\coordinate (centre) at (1.5,1.5) ;
\coordinate (A) at (.5,.5) ;
\coordinate (B) at (2.5,2.5) ;
\fill (centre) circle (3pt);
\draw[red] (A) rectangle (B) ;
\end{tikzpicture}
&  
\parbox[c]{8cm}{
\BSS{coordinate} {\color{blue}(centre)} at(1.5,1.5) ; \\
\BSS{coordinate} {\color{blue}(A)} at (.5,.5) ;\\
\BSS{coordinate} {\color{blue}(B)} at  (2.5,2.5) ;\\
\\
\BS{fill} {\color{blue}(centre)} circle (3pt);\\
\BS{draw}[red] {\color{blue}(A)} rectangle {\color{blue}(B)} ;\\
}
\\ \hline 
\end{tabular} 


\TFRGB{voir aussi}{see also} page \pageref{noeuds}


\SbSbSSCT{Coordonnées relatives à un noeud}{Coordinates relative to a node}

\noindent

\begin{tabular}{|c|c|c|c|} \hline
\multicolumn{4}{|l|}{  \BS{node} [draw,fill=green!20,] (A) at (1,1) \AC{\BS{huge}  noeud}; }\\ 
\multicolumn{4}{|l|}{  \BS{fill}[red] (\RDD{node cs}:\RDD{name}=A,\RDD{anchor}=south) circle (3pt);   }\\ 
\hline

\begin{tikzpicture}
\draw[help lines] (0,0) grid (2,2);
\node[draw,fill=green!20,] (A) at (1,1) {\huge noeud};
\fill[red] (node cs:name=A,anchor=south) circle (3pt);
\end{tikzpicture}
&
\begin{tikzpicture}
\draw[help lines] (0,0) grid (2,2);
\node[draw,fill=green!20,] (A) at (1,1) {\huge noeud};
\fill[red] (node cs:name=A,anchor=west) circle (3pt);
\end{tikzpicture}
&
\begin{tikzpicture}
\draw[help lines] (0,0) grid (2,2);
\node[draw,fill=green!20,] (A) at (1,1) {\huge noeud};
\fill[red] (node cs:name=A,anchor=north) circle (3pt);
\end{tikzpicture}
&
\begin{tikzpicture}
\draw[help lines] (0,0) grid (2,2);
\node[draw,fill=green!20,] (A) at (1,1) {\huge noeud};
\fill[red] (node cs:name=A,anchor=east) circle (3pt);
\end{tikzpicture}
\\ \hline
name=A,anchor=south & name=A,anchor=west & name=A,anchor=north & name=A,anchor=east
\\ \hline
\end{tabular}

\bigskip

\begin{tabular}{|c|c|c|c|} \hline
\multicolumn{4}{|l|}{  \BS{node} [draw,fill=green!20,] \blll{(A)} at (1,1) \AC{\BS{huge}  noeud}; }\\ 
\multicolumn{4}{|l|}{  \BS{fill}[red] (\blll{A}.south) circle (3pt);   }\\ 
\hline

\begin{tikzpicture}
\draw[help lines] (0,0) grid (2,2);
\node[draw,fill=green!20,] (A) at (1,1) {\huge noeud};
\fill[red] (A.south) circle (3pt);
\end{tikzpicture}
&
\begin{tikzpicture}
\draw[help lines] (0,0) grid (2,2);
\node[draw,fill=green!20,] (A) at (1,1) {\huge noeud};
\fill[red] (A.west) circle (3pt);
\end{tikzpicture}
&
\begin{tikzpicture}
\draw[help lines] (0,0) grid (2,2);
\node[draw,fill=green!20,] (A) at (1,1) {\huge noeud};
\fill[red] (A.north) circle (3pt);
\end{tikzpicture}
&
\begin{tikzpicture}
\draw[help lines] (0,0) grid (2,2);
\node[draw,fill=green!20,] (A) at (1,1) {\huge noeud};
\fill[red] (A.east) circle (3pt);
\end{tikzpicture}
\\ \hline
A.south & A.west & A.north & A.east
\\ \hline
\end{tabular}



\bigskip
\begin{tabular}{|c|c|c|c|} \hline
\multicolumn{4}{|c|}{  \BS{fill}[red] (node cs:\RDD{name}=A,\RDD{angle}=0) circle (3pt);  }\\ 
\hline

\begin{tikzpicture}
\draw[help lines] (0,0) grid (2,2);
\node[draw,fill=green!20,] (A) at (1,1) {\huge noeud};
\fill[red] (node cs:name=A,angle=0) circle (3pt);
\end{tikzpicture}
&
\begin{tikzpicture}
\draw[help lines] (0,0) grid (2,2);
\node[draw,fill=green!20,] (A) at (1,1) {\huge noeud};
\fill[red] (node cs:name=A,angle=-30) circle (3pt);
\end{tikzpicture}
&
\begin{tikzpicture}
\draw[help lines] (0,0) grid (2,2);
\node[draw,fill=green!20,] (A) at (1,1) {\huge noeud};
\fill[red] (node cs:name=A,angle=-90) circle (3pt);
\end{tikzpicture}
&
\begin{tikzpicture}
\draw[help lines] (0,0) grid (2,2);
\node[draw,fill=green!20,] (A) at (1,1) {\huge noeud};
\fill[red] (node cs:name=A,angle=-150) circle (3pt);
\end{tikzpicture}
\\ \hline
name=A,angle=0 & name=A,angle=-30 & nname=A,angle=-90 & name=A,angle=-150
\\ \hline
\end{tabular}

\bigskip


\begin{tabular}{|c|c|c|c|} \hline
\multicolumn{4}{|c|}{  \BS{fill}[red] (A.0) circle (3pt);  }\\ 
\hline

\begin{tikzpicture}
\draw[help lines] (0,0) grid (2,2);
\node[draw,fill=green!20,] (A) at (1,1) {\huge noeud};
\fill[red] (A.0) circle (3pt);
\end{tikzpicture}
&
\begin{tikzpicture}
\draw[help lines] (0,0) grid (2,2);
\node[draw,fill=green!20,] (A) at (1,1) {\huge noeud};
\fill[red] (A.-30) circle (3pt);
\end{tikzpicture}
&
\begin{tikzpicture}
\draw[help lines] (0,0) grid (2,2);
\node[draw,fill=green!20,] (A) at (1,1) {\huge noeud};
\fill[red] (A.-90) circle (3pt);
\end{tikzpicture}
&
\begin{tikzpicture}
\draw[help lines] (0,0) grid (2,2);
\node[draw,fill=green!20,] (A) at (1,1) {\huge noeud};
\fill[red] (A.-150) circle (3pt);
\end{tikzpicture}
\\ \hline
A.0 & A.-30 & A.-90 & A.-150
\\ \hline
\end{tabular}

\TFRGB{voir aussi}{see also} page \pageref{nomnoeud}


\newpage

\SbSbSSCT{Coordonnées relatives à deux points}{Coordinates relative to two points}
\begin{center}
\RRR{13-3-1}
\end{center}

\begin{tabular}{|c|c|} \hline
\multicolumn{2}{|c|}{  \BS{node} [circle,fill=red!20] at (1,1 {\color{red}|-} 3,3) \AC{X}   }\\ 
\hline
\begin{tikzpicture}
\draw[help lines] (0,0) grid (4,4);
\node[circle,fill=green!20,] (A) at (1,1) {A};
\node[circle,fill=green!20,] (B) at (3,3) {B};
\node[circle,fill=red!20] at (1,1 |- 3,3) {X};
\end{tikzpicture}
&
\begin{tikzpicture}
\draw[help lines] (0,0) grid (4,4);
\node[circle,fill=green!20,] (A) at (1,1) {A};
\node[circle,fill=green!20,] (B) at (3,3) {B};
\node[circle,fill=red!20] at (1,1 -| 3,3) {X};
\end{tikzpicture}
\\ \hline
at (1,1 {\color{red}|-} 3,3)
&
at (1,1 {\color{red}-|} 3,3)
\\ \hline
\end{tabular}



\SbSbSSCT{Coordonnée relative à une intersection}{Coordinates relative to an intersection}
\begin{center}
\RRR{13-3-2}
\end{center}

 \maboite{\BS{usetikzlibrary}\AC{intersections}}
\label{lib-intersections}


\begin{tabular}{|c|c|c|c|} \hline 
\multicolumn{4}{|l|}{  \BS{draw} [\RDD{name path}=XXX] (2,1) circle  (1cm);   }\\ 
\multicolumn{4}{|l|}{  \BS{draw} [\RDD{name path}=YYY] (0.5,0.5) rectangle +(3,1);   }\\ 
\multicolumn{4}{|l|}{ \BS{fill} [red,\RDD{ name intersections}=\AC{of=xxx and YYY}]
(\RDD{intersection}-1) circle (2pt)   }\\ 
\hline 
\begin{tikzpicture}[scale=.8]
\draw [help lines] grid (4,2);
\draw [name path=XXX] (2,1) circle  (1cm);
\draw [name path=YYY] (0.5,0.5) rectangle +(3,1);
\fill [red, name intersections={of=XXX and YYY}]
(intersection-1) circle (2pt)  ;
\end{tikzpicture}
& 
\begin{tikzpicture}[scale=.8]
\draw [help lines] grid (4,2);
\draw [name path=XXX] (2,1) circle  (1cm);
\draw [name path=YYY] (0.5,0.5) rectangle +(3,1);
\fill [red, name intersections={of=XXX and YYY}] (intersection-2) circle (2pt) ;
\end{tikzpicture} 
&  
\begin{tikzpicture}[scale=.8]
\draw [help lines] grid (4,2);
\draw [name path=XXX] (2,1) circle  (1cm);
\draw [name path=YYY] (0.5,0.5) rectangle +(3,1);
\fill [red, name intersections={of=XXX and YYY}] (intersection-3) circle (2pt) ;
\end{tikzpicture}
&  
\begin{tikzpicture}[scale=.8]
\draw [help lines] grid (4,2);
\draw [name path=XXX] (2,1) circle  (1cm);
\draw [name path=YYY] (0.5,0.5) rectangle +(3,1);
\fill [red, name intersections={of=XXX and YYY}] (intersection-4) circle (2pt) ;
\end{tikzpicture}
\\ 
\hline intersection-1 & intersection-2 &intersection-3  & intersection-4 \\ 
\hline 
\end{tabular} 

\bigskip

\begin{tabular}{|c|} \hline  
\BS{fill} [red, name intersections=\AC{of=XXX and YYY}] \\
(intersection-1) circle (2pt) {\color{red} node[black,above right] \AC{point a}} ;
\\ \hline  
\begin{tikzpicture}
\draw [help lines] grid (4,2);
\draw [name path=XXX] (2,1) circle  (1cm);
\draw [name path=YYY] (0.5,0.5) rectangle +(3,1);
\fill [red, name intersections={of=XXX and YYY}]
(intersection-1) circle (2pt) node[black,above right] {point a} ;
\end{tikzpicture} 
\\ \hline 
\end{tabular} 

\bigskip

\begin{tabular}{|c|} \hline 
\BS{fill} [red, name intersections=\AC{of=XXX and YYY, \RDD{name}=ZZZ}]; \\
\BS{draw} [red] (ZZZ-1) - - (ZZZ-3); \BS{draw} [green] (ZZZ-2) - - (ZZZ-4);
\\ \hline  
\begin{tikzpicture}
\draw [help lines] grid (4,2);
\draw [name path=XXX] (2,1) circle  (1cm);
\draw [name path=YYY] (0.5,0.5) rectangle +(3,1);
\fill [red, name intersections={of=XXX and YYY, name=ZZZ}];
\draw [red] (ZZZ-1) -- (ZZZ-3);
\draw [green] (ZZZ-2) -- (ZZZ-4);
\end{tikzpicture}
\\ \hline 
\end{tabular} 

\bigskip
\begin{tabular}{|c|} \hline  
\BS{fill} [red, name intersections=\AC{of=XXX and YYY , \RDD{by}=\AC{a,b,c,d}}]; \\
\BS{draw} [red] (a) - - (c); \hspace{1cm} \BS{draw} [green] (b) - - (d);
\\ \hline   
\begin{tikzpicture}
\draw [help lines] grid (4,2);
\draw [name path=XXX] (2,1) circle  (1cm);
\draw [name path=YYY] (0.5,0.5) rectangle +(3,1);
\fill [red, name intersections={of=XXX and YYY, by={a,b,c,d}}];
\draw [red] (a) -- (c);
\draw [green] (b) -- (d);
\end{tikzpicture}
\\ \hline 
\end{tabular} 

\bigskip

\begin{tabular}{|c|} \hline  
\BS{fill} [name intersections=\AC{of=XXX and YYY, name=i, \RDD{total}=\BS{t}}] [red] \\
\BS{foreach} \BS{s} in \AC{1,...,\BS{t}} \AC{(i-\BS{s}) circle (2pt) node[black,above right] \AC{\BS{s}}}
\\ \hline  
\begin{tikzpicture}
\draw [help lines] grid (4,2);
\draw [name path=XXX] (2,1) circle  (1cm);
\draw [name path=YYY] (0.5,0.5) rectangle +(3,1);
\fill [name intersections={of=XXX and YYY , name=i, total=\t}]
[red]
\foreach \s in {1,...,\t}{(i-\s) circle (2pt) node[black,above right] {\s}};
\end{tikzpicture}
\\ \hline 
\end{tabular} 



\newpage

\SbSbSSCT{Position calculée avec le module  \og  pgfmath \fg}{Calculated positions with  \og  pgfmath \fg }

\begin{center}
\RRR{13-2-1}
\end{center}

\TFRGB{Ce module est chargé automatiquement avec le module Tikz}{Package automatically loaded with Tikz} 

\begin{tabular}{|c|} \hline 
\begin{tikzpicture}
\draw[help lines] (0,0) grid (4,2);
\fill [red] (canvas cs:x=2cm+1.5cm,y=1.5cm-1cm) circle (3pt);
\fill [blue] (2cm,1.5cm) circle (3pt);
\draw[dashed] (2,1.5) -| (3.5,.5);
\end{tikzpicture}
\\ \hline 
\emph{\TFRGB{Explicite}{explicit}} 
 : \BS{fill} [red] (\RDD{canvas cs}:x=2cm+1.5cm,y=1.5cm-1cm) circle (3pt);
 \\  \hline 
\emph{\TFRGB{Implicite}{implicit}} :  \BS{fill} [red] {\color{red}(2cm+1.5cm,1.5cm-1cm)} circle (3pt);
\\ \hline 
\end{tabular} 

\bigskip
\begin{tabular}{|c|c|c|} \hline 
\begin{tikzpicture}[baseline=0pt]
\draw[help lines] (0,0) grid (4,4);
 \draw[dashed] (2,2) circle (2);
\fill[red](2+ 2*cos 30,2+2*sin 30) circle (3pt);
\fill[magenta](2+ 2*cos{(120)},2+2*sin{(120)}) circle (3pt);
\end{tikzpicture}
&
\parbox[c]{8cm}{
 \BS{draw}[dashed] (2,2) circle (2);\\
 \smallskip
 \BS{fill} [red]{\color{red}(2+ 2*cos 30 , 2+2*sin 30)} circle (3pt);\\
  \smallskip
 \BS{fill}[magenta] {\color{red}(2+2*cos\AC{(120)} , 2+2*sin\AC{(120)})} circle (3pt); 
 }
\\ \hline 
\end{tabular} 

\SbSbSSCT{Position calculée avec \og library calc \fg}{Calculated positions with \og  calc  library calc \fg}

\begin{center}
\RRR{13-5}
\end{center}
\label{lib-calc}

 \maboite{\BS{usetikzlibrary}\AC{calc}}
 
\begin{tabular}{|c|c|} \hline  
\begin{tikzpicture}[baseline=0pt]
\draw [help lines] (0,0) grid (3,2);
\node (a) at (1,1) {A};
\fill [red] ($(a) + 2/3*(1cm,0)$) circle (2pt);
\fill [red] ($(a) + 4/3*(1cm,0)$) circle (2pt);
\end{tikzpicture}
&
\parbox{8cm}{
\BS{node} (a) at (1,1) \AC{A}; \\
\BS{fill} [red] {\color{red} (\$(a) + 2/3*(1cm,0)\$)} circle (2pt); \\
\BS{fill} [red] {\color{red}(\$(a) + 4/3*(1cm,0)\$)} circle (2pt); \\
}
\\ 
\hline 
\end{tabular} 

\SbSbSSCT{Tangentes avec \og library calc \fg}{Tangents with  \og calc library  \fg}

\begin{center}
\RRR{13-2-4}
\end{center}

\begin{tabular}{|c|c|} \hline 
\multicolumn{2}{|l|}{\BS{node}[fill=green!20] (a) at (3,1.5) \AC{A}; } \\
\multicolumn{2}{|l|}{\BS{fill}[red] (\RDD{tangent cs}:\RDD{node}=c,\RDD{point}=\AC{(A)},\RDD{solution}=1);  }\\ 
\hline
\begin{tikzpicture}
\draw[help lines] (0,0) grid (4,2);
\node[fill=green!20] (A) at (3,1.5) {A};
\node [circle,draw] (c) at (1,1) [minimum size=1.5cm] {$c$};
\draw[red,dashed] (A) - -(tangent cs:node=c,point={(A)},solution=1) ;
\draw[red,dashed] (1,1) - -(tangent cs:node=c,point={(3,1.5)},solution=1) ;
\fill[red] (tangent cs:node=c,point={(A)},solution=1) circle (3pt);
\end{tikzpicture}
&
\begin{tikzpicture}
\draw[help lines] (0,0) grid (4,2);
\node[fill=green!20] (A) at (3,1.5) {A};
\node [circle,draw] (c) at (1,1) [minimum size=1.5cm] {$c$};
\draw[red,dashed] (A) - -(tangent cs:node=c,point={(A)},solution=2) ;
\draw[red,dashed] (1,1) - -(tangent cs:node=c,point={(A)},solution=2) ;
\fill[red] (tangent cs:node=c,point={(A)},solution=2) circle (3pt);
\end{tikzpicture}
\\ \hline
\RDD{solution}=1 & \RDD{solution}=2
\\ \hline
\end{tabular} 

\newpage

\SbSbSSCT{Point à pourcentage donné }{Percentage position }

\begin{center}
\RRR{13-5-3}
\end{center}


\begin{tabular}{|c|c|} \hline  
\multicolumn{2}{|c|}{\BS{fill}[red] ({\color{red}\$(0,1)!.25!(4,1)\$}) circle (4pt); } \\  \hline  

\begin{tikzpicture}
\draw [help lines] (0,0) grid (4,2);
\draw [line width= 3pt] (0,1) -- (4,1);
\fill[red] ($(0,1)!.25!(4,1)$) circle (4pt);
\end{tikzpicture}
&  
\begin{tikzpicture}
\draw [help lines] (0,0) grid (4,2);
\draw [line width= 3pt] (0,1) -- (4,1);
\fill[red] ($(0,1)!.75!(4,1)$) circle (4pt);
\end{tikzpicture}
\\ \hline (0,1)!{\color{red}0.25}!(4,1) & (0,1)!{\color{red}0.75}!(4,1) \\ 
\hline 
\end{tabular} 

\bigskip

\begin{tabular}{|c|} \hline  
\begin{tikzpicture}
\draw [help lines] (0,0) grid (4,3);
\draw [line width=2pt ](0,2) -- (4,2);
\draw[red] ($(0,2)!.75!(4,2)$) -- (0,0);
\fill[red] ($(0,2)!.75!(4,2)!.66!(0,0)$) circle (4pt);
\end{tikzpicture}
\\ \hline 
\BS{fill}[red] (\${\color{blue}(0,2)!0.75!(4,2)}!{\color{red}0.66!(0,0)}\$) circle (2pt);
\\ \hline 
\end{tabular} 


\SbSbSSCT{Point à distance donnée}{Position at a given distance }

\begin{center}
\RRR{13-5-4}
\end{center}

\begin{tabular}{|c|c|} \hline  
\multicolumn{2}{|c|}{\BS{fill}[red] ({\color{red}\$(0,1)!1.5cm!(4,1)\$}) circle (4pt); } \\  \hline  

\begin{tikzpicture}
\draw [help lines] (0,0) grid (4,2);
\draw [line width= 2pt] (0,1) -- (4,1);
\fill[red] ($(0,1)!1.5cm!(4,1)$) circle (4pt);
\end{tikzpicture}
&  
\begin{tikzpicture}
\draw [help lines] (0,0) grid (4,2);
\draw [line width= 2pt] (0,1) -- (4,1);
\fill[red] ($(0,1)!3cm!(4,1)$) circle (4pt);
\end{tikzpicture}
\\ \hline (0,1)!{\color{red}1.5cm}!(4,1) & (0,1)!{\color{red}3cm}!(4,1) \\ 
\hline 
\end{tabular} 

\bigskip

\begin{tabular}{|c|} \hline  
\begin{tikzpicture}
\draw [help lines] (0,0) grid (4,4);
\coordinate (a) at (1,0);
\coordinate (b) at (4,1);
\draw [line width= 3pt] (0,0) -- (4,1);
\draw [line width= 2pt,red](2,.5) -- ($ (2,.5)!2cm!90:(4,1) $);
\end{tikzpicture}
\\ \hline
\BS{draw} (2,.05) - - (\$ (2,0.5)!{\color{red}2cm!90:(4,1)} \$);
\\ \hline 
\end{tabular} 

\newpage

\SbSbSSCT{Coordonnées relatives}{Relative coordinates}


\Par{Cartésienne}{Cartesian coordinates}

\begin{center}
\RRR{13-4-1}
\end{center}

\begin{tabular}{|c|c|c|} \hline  
\TFRGB{relative à l'origine}{relative to the origin}  & \TFRGB{relative à une position}{relative to a position}  &  \TFRGB{relative à la dernière position}{relative to the last position}   
\\ \hline  
 
\begin{tikzpicture}
\draw[help lines] (0,-1) grid (3,1); 
 \draw[blue,very thick] (0,0) -- (1,0) - - (2,1) - - (2,-1);
 \fill[red] (0,0) circle (4pt);
\end{tikzpicture}
&
\begin{tikzpicture} %[scale=.8]
\draw[help lines] (0,-1) grid (4,1);
 \draw[blue,very thick] (0,0) - - (1,0) -- +(2,1) -- +(2,-1) ; %–- +(2,-1) ;
 \fill[red] (1,0) circle (4pt);
\end{tikzpicture}
&
\begin{tikzpicture} %[scale=.8]
\draw[help lines] (0,-1) grid (5,1);  
 \draw[blue,very thick] (0,0) -- (1,0)  - - ++(2,1) - - ++(2,-1);
 \fill[red] (1,0) circle (4pt);
 \fill[red] (3,1) circle (4pt);
\end{tikzpicture}
\\ \hline 
\tikz \fill node[fill=green!20,inner sep=0pt]{(0,0)}; - - (1,0) &
 (0,0) - - \tikz \fill node[fill=green!20,inner sep=0pt]{(1,0)};  & (0,0) - - \tikz \fill node[fill=green!20,inner sep=0pt]{(1,0)}; \\
 - - (2,1) - - (2,-1)  &
   - - +(2,1) - - +(2,-1) & - - ++\tikz \fill node[fill=green!20,inner sep=0pt]{(2,1)}; - - ++(2,-1)
\\ \hline 
\end{tabular} 

\bigskip

\begin{tabular}{|c|c|c|} \hline  
\begin{tikzpicture} [scale=.5]
\draw[help lines] (0,-1) grid (6,6);
 \draw[red,dotted,line width=2pt] (0,0) rectangle (2,2) ;
  \draw[green,dotted,line width=2pt] (0,0) rectangle (3,3) ;  
 \draw[blue,line width=2pt] (0,0) rectangle (1,1)  rectangle (2,2) rectangle (3,3);

\end{tikzpicture}

&  
\begin{tikzpicture} [scale=.5]
\draw[help lines] (0,-1) grid (6,6); 
  \draw[green,dotted,line width=2pt] (1,1) rectangle (4,4) ;   
 \draw[blue,line width=2pt] (0,0) rectangle (1,1)  rectangle +(2,2) rectangle +(3,3);
    \fill[red] (1,1) circle (4pt);
\end{tikzpicture}
&  
\begin{tikzpicture} [scale=.5]
\draw[help lines] (0,-1) grid (6,6);  
 \draw[blue,line width=2pt] (0,0) rectangle (1,1)  rectangle ++(2,2) rectangle ++(3,3);
    \fill[red] (1,1) circle (4pt);
     \fill[green] (3,3) circle (4pt); 
\end{tikzpicture}
\\ 
\hline 
\BS{draw} (0,0) rectangle (1,1)   &
\BS{draw} (0,0) rectangle (1,1)   & 
\BS{draw} (0,0) rectangle (1,1)  \\
rectangle (2,2) rectangle (3,3);  &
rectangle +(2,2) rectangle +(3,3);  &
rectangle ++(2,2) rectangle ++(3,3); \\
\hline 
\end{tabular}


\Par{Polaire }{Polar} {}

\bigskip


\noindent

\begin{tabular}{|c|c|c|c|} \hline
\TFRGB{relative à l'origine}{relative to the origin}  & \TFRGB{relative à une position}{relative to a position}  &  \TFRGB{relative à la dernière position}{relative to the last position}   
\\ \hline    
\begin{tikzpicture} %[scale=.8] 
\draw[help lines] (0,-1) grid (3,1);
 \fill[red] (0:0) circle (4pt);
 \draw[blue,very thick] (0:0)-- (0:1) -- (30:2) -- (-30:2);
\end{tikzpicture}
&
\begin{tikzpicture} %[scale=.8] 
\draw[help lines] (0,-1) grid (4,1);
 \fill[red] (1,0) circle (4pt);
 \draw[blue,very thick] (0:0) -- (0:1) -- +(30:2) -- +(-30:2);
\end{tikzpicture}
&
\begin{tikzpicture} %[scale=.8] 
\draw[help lines] (0,-1) grid (5,1);
 \fill[red] (1,0) circle (4pt);
 \fill[red] (2.732,1) circle (4pt);
 \draw[blue,very thick] (0:0)-- (0:1) -- ++(30:2) -- ++(-30:2);
\end{tikzpicture}
\\ \hline
\tikz \fill node[fill=green!20,inner sep=0pt] {(0:0)}; - - (0:1)&
 (0:0) - - \tikz \fill node[fill=green!20,inner sep=0pt] {(0:1)}; & (0:0)- - \tikz \fill node[fill=green!20,inner sep=0pt] {(0:1)}; \\
 - - (30:2) - - (-30:2)  &  - -  +(30:2) - - +(-30:2) & - -  ++\tikz \fill node[fill=green!20,inner sep=0pt] {(30:2)}; - - ++(-30:2)
\\ \hline 
\end{tabular} 

%\subsubsection{coordonnée relative en polaire}
\Par{coordonnée relative en polaire}{Relative polar coordinate}

\begin{center}
\RRR{13-4-2}
\end{center}
\bigskip

\begin{tabular}{|c|c|} \hline 
\multicolumn{2}{|c|}{ \BS{draw}[blue,very thick] (0,0) -- (2,1) -- ([turn]-45:1cm);}
 \\ \hline
\begin{tikzpicture} %[scale=.8] 
\draw[help lines] (0,0) grid (4,2);
 \draw[dotted] (0,0) -- (4,2);
 \draw[blue,very thick] (0,0) -- (2,1) -- ([turn]-45:1cm);
\end{tikzpicture}
&  
\begin{tikzpicture} %[scale=.8] 
\draw[help lines] (0,0) grid (4,2);
 \draw[dotted] (0,0) -- (4,2);
 \draw[blue,very thick] (0,0) -- (2,1) -- ([turn]45:1cm);
\end{tikzpicture}
\\ \hline ([\RDD{turn}]-45:1cm) & ([\RDD{turn}]45:1cm) \\ 
\hline 
\end{tabular}

\bigskip

\begin{tabular}{|c|c|} \hline  
\begin{tikzpicture}  
\draw[help lines] (-1,0) grid (4,3);
\draw [line width=2pt] (4,0) arc (0 :120 :2)  -- ([turn]90:2cm) ;

\end{tikzpicture}
&  
\begin{tikzpicture} %[scale=.8] 
\draw[help lines] (0,0) grid (4,3);
\draw [line width=2pt]  (0,0) to [bend left] (2,2) --  ([turn]0:2cm);
\fill [red](2,2) circle (4pt);
\end{tikzpicture}
\\ \hline  
\BS{draw} (4,0) arc (0 :120 :2)  - - ([\RDD{turn}]90:2cm) ;
& \BS{draw}  (0,0) to [bend left] (2,2) - -  ([\RDD{turn}]0:2cm); \\

\hline 
\end{tabular} 


%\bigskip 
%
%
%\tikz [delta angle=30, radius=1cm]
%\draw (0,0) arc [start angle=0] -- ([turn]0:1cm)
%arc [start angle=30] -- ([turn]0:1cm)
%arc [start angle=60] -- ([turn]30:1cm);



\bigskip

\begin{tabular}{|c|c|c|} \hline  
\multicolumn{3}{|c|}{ \BS{draw}(1,2)
.. controls ([turn]0:2cm) .. ([turn]-90:2cm); }
\\ \hline
\begin{tikzpicture} %[scale=.8] 
\draw[help lines] (0,0) grid (4,4);
 \draw [line width=2pt] (1,2)
.. controls ([turn]0:2cm) .. ([turn]-90:2cm);
\end{tikzpicture}
&  
\begin{tikzpicture} %[scale=.8] 
\draw[help lines] (0,0) grid (4,4);
 \draw [line width=2pt] (1,2)
.. controls ([turn]30:2cm) .. ([turn]-90:2cm);
\end{tikzpicture}
&  
\begin{tikzpicture} %[scale=.8] 
\draw[help lines] (-2,0) grid (2,4);
 \draw [line width=2pt] (1,2)
.. controls ([turn]0:2cm) .. ([turn]90:2cm);

\end{tikzpicture}
\\ \hline ([turn]0:2cm) .. ([turn]-90:2cm) & ([turn]30:2cm) .. ([turn]-90:2cm) & ([turn]0:2cm) .. ([turn]90:2cm) \\ 
\hline 
\end{tabular} 


\tikzset{every picture/.style=blue,very thick,inner sep=.3333em}

%
%
%
%\newpage
%
%\SSCT{Les n\oe uds }{Nodes}
%
%\SbSSCT{Définition des  n\oe uds}{Creation of nodes}
\tikzset{blue}

\label{noeuds}
\noindent

\begin{tabular}{|c | c | c | c | c |} \hline
\multicolumn{5}{|c|}{  \BS{draw} (1,1) node[\RDD{fill}=red!20] \AC{};   }\\ 
\hline 
\tikz \draw (0,0) grid (2,2) (1,1) node[fill=red!20] {};
&
\tikz \draw (0,0) grid (2,2) (1,1) node[fill=red!20,draw] {}; 
&
\tikz \draw (0,0) grid (2,2) (1,1) node[circle,fill=red!20] {};
&
\tikz \draw (0,0) grid (2,2) (1,1) node[circle,fill=red!20,draw] {};
&
\tikz \draw (0,0) grid (2,2) (1,1) node[coordinate] {};
\\  \hline
\dft
&
node[\RDD{draw}] 
&
 node[\RDD{circle}]  
&
 node[\RDD{circle},\RDD{draw}]
 &
  node[\RDD{coordinate}]
 \\  \hline
\end{tabular}
\bigskip

\begin{tabular}{|c | c | c | c | } \hline
\multicolumn{4}{|c|}{ \BSS{node} \RDD{at} (1,1) [fill=red!20] \AC{};   }\\ 
\hline 
 \begin{tikzpicture}
\draw (0,0) grid (2,2) ; 
\node at (1,1) [fill=red!20] {};
 \end{tikzpicture}
&
 \begin{tikzpicture}
\draw (0,0) grid (2,2) ; 
\node at (1,1) [draw] {};
 \end{tikzpicture}
&
 \begin{tikzpicture}
\draw (0,0) grid (2,2) ; 
\node at (1,1) [fill=red!20,circle] {};
 \end{tikzpicture}
&
 \begin{tikzpicture}
\draw (0,0) grid (2,2) ; 
\node at (1,1) [circle,draw] {};
 \end{tikzpicture}

\\  \hline
[fill=red!20]
&
[\RDD{draw}] 
&
[\RDD{circle},fill=red!20]
 &
[\RDD{circle},draw] 
 \\  \hline
\end{tabular}
\bigskip

\TFRGB{Autres types de n\oe uds voir page}{Other type of nodes see page} \pageref{noeudboite}
\bigskip


\begin{tabular}{|c|c|} \hline 
\BS{draw} (0,0) node at (1,0) \AC{1} node at (2,0) \AC{2} & \BS{draw}(0,0) node foreach \BS{x} in \AC{1,2,...,5}\\ 
node at (3,0) \AC{3} node at (4,0) \AC{4} node at (5,0) \AC{5}; &  at (\BS{x},0) \AC{\BS{x}};\\ 
\hline 
\tikz \draw (0,0) node at (1,0) {1} node at (2,0) {2} node at (3,0) {3} node at (4,0) {4} node at (5,0) {5};
&
\tikz \draw (0,0) node foreach \x in {1,2,...,5} at (\x,0) {\x};
\\ \hline 
\end{tabular} 



\bigskip

\begin{tabular}{|c|} \hline 
\BS{draw}[\rouge{every node/.style=\AC{draw,red}}](0,0) node foreach \BS{x} in \AC{1,2,...,5} at (\BS{x},0) \AC{\BS{x}};
\\ \hline 
\rule[-3pt]{0pt}{.8cm}\tikz \draw[every node/.style={draw,red}] (0,0) node foreach \x in {1,2,...,5} at (\x,0) {\x};
\\ \hline 
\end{tabular} 

\bigskip

\begin{tabular}{|c|} \hline 
\BS{draw}[\rouge{every rectangle node/.style=\AC{draw,red}},\\
\rouge{every circle node/.style=\AC{draw,double}}]\\ (0,0) node at (1,0) \AC{1} node[circle] at (2,0) \AC{2} \\ node[circle] at (3,0) \AC{3} node at (4,0) \AC{4} node at (5,0) \AC{5};
\\ \hline 
\rule[-3pt]{0pt}{1cm} \tikz \draw[every rectangle node/.style={draw,red},
every circle node/.style={draw,double}] (0,0) node at (1,0) {1} node[circle] at (2,0) {2} node[circle] at (3,0) {3} node at (4,0) {4} node at (5,0) {5};
\\ \hline 
\end{tabular} 

\SbSSCT{Nom des  n\oe uds}{Node name}


\begin{tabular}{|c|c|c|}
\hline 
\multicolumn{3}{|c|}{} \\ 
\hline 
\begin{tikzpicture}
\node[name=A,fill=red] at (0,0) {};
\draw  (-1,-1) grid (1,1) ;
\draw (A) circle (.5) ;
\end{tikzpicture} 
&  
\begin{tikzpicture}
\node[name=A,alias=B,fill=red] at (0,0) {} ;
\draw  (-1,-1) grid (1,1) ;
\draw (B) circle (.5) ;
\end{tikzpicture}
& 
\begin{tikzpicture}
\node[fill=red] (C) at (0,0) {};
\draw  (-1,-1) grid  (1,1) ;
\draw (C) circle (.5);
\end{tikzpicture} \\ 
\hline 
\BS{node}[\RDD{name}=A] at (0,0) \AC{}  & \BS{node}[\RDD{name}=A,\RDD{alias}=B] at (0,0) \AC{}  & 
\BS{node}\rouge {(C)} at (0,0) \AC{} \\ 
\BS{draw} (A) circle (.5); & \BS{draw}  (B) circle (.5); &\BS{draw} (C) circle (.5);
\\ \hline 
\end{tabular} 
\newpage

\SbSSCT{Contenu des  n\oe uds}{Node contents}
\tikzset{blue}

\begin{center}
\RRR{17-2-1}
\end{center}

\begin{tabular}{|c|c|} \hline 
\BS{node} at (1,1) [fill=red!20]\rouge { \AC{XXX} };
&  
\BS{node} at (1,1) [fill=red!20,\RDD{node contents}=XXX] \AC{};
\\  \hline 
 \begin{tikzpicture}
\draw (0,0) grid (2,2) ; 
\node at (1,1) [fill=red!20] {XXX};
\end{tikzpicture}
&  
\begin{tikzpicture}
\draw (0,0) grid (2,2) ; 
\node at (1,1) [fill=red!20,node contents=XXX] {};
\end{tikzpicture} 
\\ \hline 
\end{tabular} 

\bigskip

\begin{tabular}{|c|c|} \hline 
\BS{node}[red] at (1,1) [fill=blue!20] \AC{XXX} ;
&  
\BS{node}[red] at (1,1) [fill=blue20,node contents=XXX] \AC{};
\\  \hline 
 \begin{tikzpicture}
\draw (0,0) grid (2,2) ; 
\node[red] at (1,1) [fill=blue!20] {XXX};
\end{tikzpicture}
&  
\begin{tikzpicture}
\draw (0,0) grid (2,2) ; 
\node[red] at (1,1) [fill=blue!20,node contents=XXX] {};
\end{tikzpicture} 
\\ \hline 
\end{tabular} 


\SbSSCT{Premier ou arrière plan}{Behind or in front}

\begin{tabular}{|c|c|} \hline 
\multicolumn{2}{|l|}{\BS{tikz} \BS{fill} [fill=blue!50, draw=blue, very thick]
(0,0) } \\ 
\multicolumn{2}{|l|}{node [\RDD{behind path}, fill=red!50] \AC{XXXXX} }  \\
\multicolumn{2}{|l|}{- - (1.5,0) - - (1.5,1) - - (0,1) ;}
\\ \hline 
\tikz \fill [fill=blue!50, draw=blue, very thick]
(0,0) node [behind path, fill=red!50] {XXXXX}
-- (1.5,0) 
-- (1.5,1) 
-- (0,1) ;
&  
\tikz \fill [fill=blue!50, draw=blue, very thick]
(0,0) node [in front of path, fill=red!50] {XXXXX}
-- (1.5,0) 
-- (1.5,1) 
-- (0,1) ;
\\ \hline 
\RDD{behind path}
&  
\RDD{in front of path}
\\ \hline 
\end{tabular}



\SbSSCT{Noms à préfixe ou suffixe}{Name prefix or name suffix}


\begin{tabular}{|c|c|}
\hline 
\begin{tikzpicture}[every node/.style={draw},baseline=0pt]
\draw[name prefix = top-] node (A) at (1,1) {A} node (B) at (2,1) {B} node (C) at (3,1) {C};
\draw[name prefix = bottom-] node (1) at (1,0) {1} node (2) at (2,0) {2} node(3) at  (3,0) {3};
\draw [red] (top-A) -- (bottom-3);
\end{tikzpicture} 
&
\parbox{12cm}{
\BS{draw}[\RDD{name prefix} = \blll{top-} ] node (A) at (1,1) \AC{A} node (B) at (2,1) \AC{B} node (C) at (3,1) \AC{C}; \\
\BS{draw}[\RDD{name prefix} = \blll{bottom-}] node (1) at (1,0) \AC{1} node (2) at (2,0) \AC{2} node(3) at  (3,0) \AC{3}; \\
\BS{draw} [red] (\blll{top-}A) -- (\blll{bottom-}3);}
\\ \hline
\begin{tikzpicture}[every node/.style={draw},baseline=0pt]
\draw[name suffix= -top] node (A) at (1,1) {A} node (B) at (2,1) {B} node (C) at (3,1) {C};
\draw[name suffix=  -bottom] node (1) at (1,0) {1} node (2) at (2,0) {2} node(3) at  (3,0) {3};
\draw [red] (A-top) -- (3-bottom);
\end{tikzpicture}
&
\parbox{12cm}{
\BS{draw}[\RDD{name suffix} = \blll{-top}] node (A) at (1,1) \AC{A} node (B) at (2,1) \AC{B} node (C) at (3,1) \AC{C}; \\
\BS{draw}[\RDD{name suffix} = \blll{-bottom}] node (1) at (1,0) \AC{1} node (2) at (2,0) \AC{2} node(3) at  (3,0) \AC{3}; \\
\BS{draw} [red] (A \blll{-top}) - - (3 \blll{-bottom});}
\\ \hline 

\end{tabular} 


\SbSSCT{Liaisons}{Links}
\label{liaisons}

\begin{tabular}{|c|c|c|} \hline 
\multicolumn{3}{|l|}{\BS{node}[draw] (A) at (0,0) \AC{A}; \hspace{.5cm} \BS{node}[draw] (B) at (1.5,1.5) \AC{B}; \hspace{.5cm} \BS{draw} (A) - - (B) } \\ \hline 
\begin{tikzpicture}[blue]
\node[draw] (A) at (0,0) {A};
\node[draw] (B) at (1.5,1.5) {B};
\draw (A) -- (B);
\end{tikzpicture}
&  
\begin{tikzpicture}[blue]
\node[draw] (A) at (0,0) {A};
\node[draw] (B) at (1.5,1.5) {B};
\draw (A) |- (B);
\end{tikzpicture}
&  
\begin{tikzpicture}[blue]
\node[draw] (A) at (0,0) {A};
\node[draw] (B) at (1.5,1.5) {B};
\draw (A) -| (B);
\end{tikzpicture}
\\ \hline  
(A){\color{red} - -} (B) & (A) {\color{red}|-} (B) &  (A) {\color{red}-|} (B)
\\ \hline 
\begin{tikzpicture}[blue]
\node[draw] (A) at (0,0) {A};
\node[draw] (B) at (1.5,1.5) {B};
\draw (A) to [bend right] (B);
\end{tikzpicture}
&  
\begin{tikzpicture}[blue]
\node[draw] (A) at (0,0) {A};
\node[draw] (B) at (1.5,1.5) {B};
\draw (A) to [bend left] (B);
\end{tikzpicture}
&  
\begin{tikzpicture}[blue]
\node[draw] (A) at (0,0) {A};
\node[draw] (B) at (1.5,1.5) {B};
\draw (A) to[bend left=0] (B);
\end{tikzpicture}
\\ \hline  
(A) to [\RDD{bend right}] (B) & (A) to [\RDD{bend left}] (B) &  (A) to[\RDD{bend left}=0] (B)
\\ \hline 
\begin{tikzpicture}[blue]
\node[draw] (A) at (0,0) {A};
\node[draw] (B) at (1.5,1.5) {B};
\draw (A) to[bend left=120]  (B);
\end{tikzpicture}
&  
\begin{tikzpicture}[blue]
\node[draw] (A) at (0,0) {A};
\node[draw] (B) at (1.5,1.5) {B};
\draw (A) to[bend left=45] (B);
\end{tikzpicture}
&  
\begin{tikzpicture}[blue]
\node[draw] (A) at (0,0) {A};
\node[draw] (B) at (1.5,1.5) {B};
\draw (A) to[bend left=90] (B);
\end{tikzpicture}
\\ \hline  
(A)  to[\RDD{bend left}=120]  (B) & (A) to[\RDD{bend left}=45] (B) &  (A) to[\RDD{bend left}=90] (B)
\\ \hline 
\begin{tikzpicture}[blue]
\node[draw] (A) at (0,0) {A};
\node[draw] (B) at (1.5,1.5) {B};
\draw (A)  to[out=90]  (B);
\end{tikzpicture}
&  
\begin{tikzpicture}[blue]
\node[draw] (A) at (0,0) {A};
\node[draw] (B) at (1.5,1.5) {B};
\draw (A) to[out=30] (B);
\end{tikzpicture}
&  
\begin{tikzpicture}[blue]
\node[draw] (A) at (0,0) {A};
\node[draw] (B) at (1.5,1.5) {B};
\draw (A)  to[in=-90]  (B);
\end{tikzpicture}
\\ \hline  
(A)  to[\RDD{out}=90] (B) & (A) to[\RDD{out}=30]  (B) &  (A)  to[\RDD{in}=-90]  (B)
\\ \hline  
\end{tabular} 

\bigskip
\begin{tabular}{|c|c|c|} \hline  
\multicolumn{2}{|c|}{ \BS{draw} (A) .. controls +(right:2cm) and +(down:2cm) .. (B);  }\\ 
\hline  
\begin{tikzpicture}[blue]
\node[draw] (A) at (0,0) {A};
\node[draw] (B) at (2,2) {B};
\draw  (A) .. controls +(right:2cm) and +(down:2cm) .. (B);
\end{tikzpicture}
&
\begin{tikzpicture}[blue]
\node[draw] (A) at (0,0) {A};
\node[draw] (B) at (2,2) {B};
\draw  (A) .. controls +(up:1cm) and +(left:1cm) .. (B);
\end{tikzpicture}
\\ \hline 
controls +(right:2cm) and +(down:2cm)  &
controls +(up:1cm) and +(left:1cm)
\\ \hline 
\begin{tikzpicture}[blue]
\node[draw] (A) at (0,0) {A};
\node[draw] (B) at (2,2) {B};
\draw  (A) .. controls +(right:1cm) and +(right:2cm) .. (B);
\end{tikzpicture}
&
\begin{tikzpicture}[blue]
\node[draw] (A) at (0,0) {A};
\node[draw] (B) at (2,2) {B};
\draw  (A) .. controls +(up:1cm) and +(right:2cm) .. (B);
\end{tikzpicture}
\\ \hline 
controls +(right:1cm) and +(right:2cm)  &
controls +(up:1cm) and +(right:2cm) 
\\ \hline 
\begin{tikzpicture}[blue]
\node[draw] (A) at (0,0) {A};
\node[draw] (B) at (2,2) {B};
\draw  (A) .. controls +(120:2cm) and +(200:1cm) .. (B);
\end{tikzpicture}
 &
 \begin{tikzpicture}[blue]
 \node[draw] (A) at (0,0) {A};
 \node[draw] (B) at (2,2) {B};T
 \draw  (A) .. controls +(120:2cm) and +(200:1cm) .. (A);
 \end{tikzpicture}
\\  \hline  
controls +(120:2cm) and +(200:1cm) & controls +(120:2cm) and +(200:1cm) 
\\ \hline 
\begin{tikzpicture}[blue]
\node[draw] (A) at (0,0) {A};
\node[draw] (B) at (2,2) {B};
\node[draw] (C) at (0,1) {C};
\node[draw] (D) at (3,0) {D};
\draw  (A) .. controls +(C) and +(D) .. (B);
\end{tikzpicture}
&
\begin{tikzpicture}[blue]
\node[draw] (A) at (0,0) {A};
\node[draw] (B) at (2,2) {B};
\node[draw] (C) at (0,1) {C};
\node[draw] (D) at (3,0) {D};
\draw (A) .. controls +(D)  .. (B);
\end{tikzpicture}
\\ \hline 
controls +(C) and +(D) &
controls +(D) 
\\ \hline 
\end{tabular} 
 \bigskip
 
\begin{tabular}{|c|c|c|} \hline 
\multicolumn{3}{|l|}{ \BS{node}[draw] (A) at (0,0) \AC{A}  }\\

\multicolumn{3}{|l|}{ \BS{node}[draw] (B) at (2,2) \AC{B} \RDD{edge}  [->] (A);  }\\
\multicolumn{3}{|c|}{\RRR{17-12-1}}  \\
\hline 
 \begin{tikzpicture}
 \node[draw] (A) at (0,0) {A};
 \node[draw] (B) at (2,2) {B} edge [->] (A);
 \end{tikzpicture}
 &
 \begin{tikzpicture}
 \node[draw] (A) at (0,0) {A};
 \node[draw] (B) at (2,2) {B} edge [red]  (A);
 \end{tikzpicture}
 &
 \begin{tikzpicture}
 \node[draw] (A) at (0,0) {A};
 \node[draw] (B) at (2,2) {B} edge [dashed] (A);
 \end{tikzpicture}
\\ \hline 
[->] & [red]  & [dashed]
\\ \hline 
\end{tabular}

\SbSSCT{\'Etiquettes sur les n\oe uds}{Node labels}

\begin{tabular}{|c|c|c|c|} \hline
\multicolumn{4}{|c|}{  \BS{fill}(0,0) circle (2pt) node[\RDD{above}] \AC{texte} ; \RRR{17-5-2}   }\\ 
\hline 
  
\begin{tikzpicture} \draw[help lines] (-1,-1) grid (1,1) ;\fill (0,0) circle (2pt) node[above] {texte};\end{tikzpicture}
& 
\begin{tikzpicture} \draw[help lines] (-1,-1) grid (1,1) ;\fill (0,0) circle (2pt) node[below] {texte};\end{tikzpicture}
 &  
\begin{tikzpicture} \draw[help lines] (-1,-1) grid (1,1);\fill (0,0) circle (2pt) node[left] {texte};\end{tikzpicture}
 &  
\begin{tikzpicture} \draw[help lines] (-1,-1) grid (1,1); \fill (0,0) circle (2pt) node[right] {texte};\end{tikzpicture}
 \\  \hline 
 [\RDD{above}] & [\RDD{below}] & [\RDD{left}] &  [\RDD{right}]
 \\ \hline 
 \begin{tikzpicture} \draw[help lines] (-1,-1) grid (1,1) ;\fill (0,0) circle (2pt) node[above left] {texte};\end{tikzpicture}
 & 
 \begin{tikzpicture} \draw[help lines] (-1,-1) grid (1,1) ;\fill (0,0) circle (2pt) node[below left] {texte};\end{tikzpicture}
  &  
 \begin{tikzpicture} \draw[help lines] (-1,-1) grid (1,1);\fill (0,0) circle (2pt) node[above right] {texte};\end{tikzpicture}
  &  
 \begin{tikzpicture} \draw[help lines] (-1,-1) grid (1,1); \fill (0,0) circle (2pt) node[below right] {texte};\end{tikzpicture}
  \\  \hline 
  [\RDD{above left}] & [\RDD{below left}] & [\RDD{above right}] &  [\RDD{below right}]
  \\ \hline 
 \begin{tikzpicture} \draw[help lines] (-1,-1) grid (1,1) ;\fill (0,0) circle (2pt) node[anchor=south] {texte};\end{tikzpicture}
 & 
 \begin{tikzpicture} \draw[help lines] (-1,-1) grid (1,1) ;\fill (0,0) circle (2pt) node[anchor=west] {texte};\end{tikzpicture}
  &  
 \begin{tikzpicture} \draw[help lines] (-1,-1) grid (1,1);\fill (0,0) circle (2pt) node[anchor=north] {texte};\end{tikzpicture}
  &  
 \begin{tikzpicture} \draw[help lines] (-1,-1) grid (1,1); \fill (0,0) circle (2pt) node[anchor=east] {texte};\end{tikzpicture}
  \\  \hline 
  [\RDD{anchor}=south] & [\RDD{anchor}=west] & [\RDD{anchor}=north] & [\RDD{anchor}=east]                                                                                                                                                               ]
  \\ \hline 
 \begin{tikzpicture} \draw[help lines] (-1,-1) grid (1,1) ;\fill (0,0) circle (2pt) node[anchor=south east] {texte};\end{tikzpicture}
 & 
\begin{tikzpicture} \draw[help lines] (-1,-1) grid (1,1) ;\fill (0,0) circle (2pt) node[anchor=south west] {texte};\end{tikzpicture}
&  
\begin{tikzpicture} \draw[help lines] (-1,-1) grid (1,1);\fill (0,0) circle (2pt) node[anchor=north west] {texte};\end{tikzpicture}
&  
\begin{tikzpicture} \draw[help lines] (-1,-1) grid (1,1); \fill (0,0) circle (2pt) node[anchor=east] {texte};\end{tikzpicture}
\\  \hline 
[\RDD{anchor}=south east] & [\RDD{anchor}=south west] & [\RDD{anchor}=north west] & [\RDD{anchor==north east                                                                                                                                                       }]
  \\ \hline 
\end{tabular} 


\bigskip
\begin{tabular}{|c|c|c|c|} \hline
\multicolumn{4}{|c|}{  \BS{fill}(0,0) circle (2pt) node[\RDD{above}=.3cm] \AC{texte} ; \RRR{17-5-2}  }\\ 
\hline 
  
\begin{tikzpicture} \draw[help lines] (-1,-1) grid (1,1) ;\fill (0,0) circle (2pt) node[above=.3cm] {texte};\end{tikzpicture}
& 
\begin{tikzpicture} \draw[help lines] (-1,-1) grid (1,1) ;\fill (0,0) circle (2pt) node[below=.3cm] {texte};\end{tikzpicture}
 &  
\begin{tikzpicture} \draw[help lines] (-1,-1) grid (1,1);\fill (0,0) circle (2pt) node[left=.3cm] {texte};\end{tikzpicture}
 &  
\begin{tikzpicture} \draw[help lines] (-1,-1) grid (1,1); \fill (0,0) circle (2pt) node[right=.3cm] {texte};\end{tikzpicture}
 \\  \hline 
 [\RDD{above}=.3cm] & [\RDD{below}=.3cm] & [\RDD{left}=.3cm] &  [\RDD{right}=.3cm]]
 \\ \hline 
\begin{tikzpicture} \draw[help lines] (-1,-1) grid (1,1) ;\fill (0,0) circle (2pt) node[above left=.3cm] {texte};\end{tikzpicture}
& 
\begin{tikzpicture} \draw[help lines] (-1,-1) grid (1,1) ;\fill (0,0) circle (2pt) node[below left=.3cm] {texte};\end{tikzpicture}
 &  
\begin{tikzpicture} \draw[help lines] (-1,-1) grid (1,1);\fill (0,0) circle (2pt) node[above right=.3cm] {texte};\end{tikzpicture}
 &  
\begin{tikzpicture} \draw[help lines] (-1,-1) grid (1,1); \fill (0,0) circle (2pt) node[below right=.3cm] {texte};\end{tikzpicture}
 \\  \hline 
 [\RDD{above left}=.3cm] & [\RDD{below left}=.3cm] & [\RDD{above right}=.3cm] &  [\RDD{below right}=.3cm]]
 \\ \hline 
 
 \end{tabular} 

 
 \newpage
\selectlanguage{french}
 
 \begin{tabular}{|c|c|c|c|c|} \hline
 \multicolumn{5}{|l|}{ \BSS{shorthandoff}\AC{:} \footnotemark[1]  } \\
 \multicolumn{5}{|l|}{  \BS{node} [draw,\RDD{label}=right:texte] \AC{}   }\\
 \multicolumn{5}{|l|}{ \BSS{shorthandon}\AC{:} } \\ 
 \hline 
     \shorthandoff{:} 
 \tikz \node [draw,label=right:texte] {};
 \shorthandon{:}
 &
  \shorthandoff{:}
 \tikz \node [draw,label=left:texte] {};
 \shorthandon{:}
 &
  \shorthandoff{:}
 \tikz \node [draw,label=above:texte] {};
 \shorthandon{:}
 &
  \shorthandoff{:}
 \tikz \node [draw,label=below:texte] {};
 \shorthandon{:}
 &
  \shorthandoff{:}
 \tikz \node [draw,label=45:texte] {};
    \shorthandon{:}
   \\ \hline
  label=right & label=left &  label=above & label=below & label=45
    \\ \hline 
 \end{tabular}
 \footnotetext[1]{\TFRGB{désactivation et ré-activation de \og : \fg  conflit entre les modules Tikz et Babel en français}{Only useful when the package babel is loaded with the frenchb option    }}
 
 \bigskip
  \begin{tabular}{|c|c|c|c|c|} \hline
  \BS{fill}(0,0) circle (2pt) node[below right=.3cm,draw,label=45:étiquette] \AC{texte} ;
      \\ \hline 
  
  \shorthandoff{:}
\begin{tikzpicture} \draw[help lines] (-1,-1) grid (2,1); \fill (0,0) circle (2pt) node[below right=.3cm,draw,label=45:étiquette] {texte};\end{tikzpicture}
 \shorthandon{:}
 
    \\ \hline 
 \end{tabular}
\bigskip

 \shorthandoff{:}

\SbSSCT{\'Etiquettes épinglées}{The Pin Option} 

\begin{center}
\RRR{17-10-3}
\end{center}
 
\begin{tabular}{|c|c|c|} \hline
\multicolumn{3}{|c|}{  \BSS{shorthandoff}\AC{:} \BS{node}[circle,draw,blue,\RDD{pin}=texte] \AC{} ;   \BSS{shorthandon}\AC{:}  \footnotemark[1] }\\ 
\hline
\begin{tikzpicture} 
\node [circle,draw,blue,pin=texte] {};
\end{tikzpicture}
&
\begin{tikzpicture} 
\node [circle,draw,blue,pin=60:texte] {};
\end{tikzpicture}
&
\begin{tikzpicture} 
\node [circle,draw,blue,pin=right:texte] {};
\end{tikzpicture}
 \\ \hline
[circle,pin=texte] &   [circle,pin=60:texte] & [circle,pin=right:texte]
 \\ \hline 
\end{tabular}  

\bigskip
\begin{tabular}{|c|c|c|} \hline
\multicolumn{3}{|c|}{  \BS{tikz}[\RDD{pin position}=60] \BS{node} [circle,pin=texte] \AC{} ;   }\\ 
\hline 
\tikz[pin position=60] \node [circle,draw,blue,pin=texte] {};
&
\tikz[pin distance=0 cm] \node [circle,draw,blue,pin=60:texte] {};
&
\tikz[pin distance=2 cm] \node [circle,draw,blue,pin=60:texte,pin distance=0cm] {};
  \\ \hline
  [\RDD{pin position}=60] & [\RDD{pin distance}=0 cm] & [\RDD{pin distance}=2 cm]
    \\ \hline
  \dft{ : above} & \multicolumn{2}{|c|}{ \dft{ : 3 ex}}
      \\ \hline
\end{tabular}  

\newpage

   \shorthandon{:} 
   
\selectlanguage{english}   

\SbSSCT{ N\oe uds  sur un chemin}{Nodes on a path}

\RRR{17-8}

\begin{tabular}{|c|c|c|} \hline
\multicolumn{3}{|c|}{  \BS{draw}(0,0) .. controls (1,2) and (2,-1) .. (4,0) node[\RDD{at end}] \AC{texte} ;   }\\ 
\hline 
\tikz \draw (0,0) .. controls (1,2) and (2,-1) .. (4,0) node[pos=0] {texte}; 
&
\tikz \draw (0,0) .. controls (1,2) and (2,-1) .. (4,0) node[pos=.33] {texte}; 
&
\tikz \draw (0,0) .. controls (1,2) and (2,-1) .. (4,0) node[at end] {texte}; 
  \\ \hline 
\RDD{pos}{\color{red}  =0} & \RDD{pos}{\color{red}  =.33} & \RDD{at end} (pos=1)
  \\ \hline 

\tikz \draw (0,0) .. controls (1,2) and (2,-1) .. (4,0) node[very near end] {texte}; 
&
\tikz \draw (0,0) .. controls (1,2) and (2,-1) .. (4,0) node[near end] {texte}; 
&
\tikz \draw (0,0) .. controls (1,2) and (2,-1) .. (4,0) node[midway] {texte}; 
  \\ \hline 
\RDD{very near end} (pos=0.875.) & \RDD{ near end} (pos=0.75) & \RDD{midway} (pos=0.5)
  \\ \hline 
  
\tikz \draw (0,0) .. controls (1,2) and (2,-1) .. (4,0) node[near start] {texte}; 
&
\tikz \draw (0,0) .. controls (1,2) and (2,-1) .. (4,0) node[very near start] {texte}; 
&
\tikz \draw (0,0) .. controls (1,2) and (2,-1) .. (4,0) node[at start] {texte};
\\ \hline 
\RDD{near start} (pos=0.25) & \RDD{very near start} (pos=0.125) & \RDD{at start} (pos=0)
  \\ \hline 
  
\end{tabular} 

\bigskip
\begin{tabular}{|c|c|c|} \hline
\multicolumn{3}{|c|}{  \BS{draw}(0,0) .. controls (1,2) and (2,1) .. (4,0) node[\RDD{sloped},midway] \AC{texte} ;   }\\ 
\hline 
\tikz \draw (0,0) .. controls (1,2) and (2,-1) .. (4,0) node[sloped,midway] {texte};
&
\tikz \draw (0,0) .. controls (1,2) and (2,-1) .. (4,0) node[above,midway] {texte};
&
\tikz \draw (0,0) .. controls (1,2) and (2,-1) .. (4,0) node[below,midway] {texte};
  \\ \hline
\RDD{sloped} & \RDD{above} &\RDD{below}
  \\ \hline
\end{tabular}
\bigskip

\begin{tabular}{|c|c|c|} \hline
\multicolumn{3}{|c|}{  \BS{draw}(0,0) .. controls (1,2) and (2,1) .. (5,0) node[\RDD{sloped},midway,allow upside down] \AC{texte} ;   }\\ 
\hline 
\tikz \draw (0,0) .. controls (1,2) and (2,-1) .. (4,0) node[sloped,midway,allow upside down] {texte};
&
\tikz \draw (0,0) .. controls (1,2) and (2,-1) .. (4,0) node[above,midway,allow upside down] {texte};
&
\tikz \draw (0,0) .. controls (1,2) and (2,-1) .. (4,0) node[below,midway,allow upside down] {texte};
  \\ \hline
\RDD{sloped} & \RDD{above} &\RDD{below}
  \\ \hline
\end{tabular}  


\begin{tabular}{|c|c|c|} \hline
\multicolumn{3}{|c|}{  \BS{draw}(A)  to [bend right]  node [\RDD{bend right}] \AC{texte} (B);   }\\ 
\hline 
\begin{tikzpicture} 
\node[draw] (A) at (0,0) {A};
\node[draw] (B) at (2,2) {B};
\draw (A) to [bend right] node [bend right] {texte} (B);
\end{tikzpicture}
&
\begin{tikzpicture} 
\node[draw] (A) at (0,0) {A};
\node[draw] (B) at (2,2) {B};
\draw (A) to [bend right] node [auto,bend right] {texte} (B);
\end{tikzpicture}
&
\begin{tikzpicture} 
\node[draw] (A) at (0,0) {A};
\node[draw] (B) at (2,2) {B};
\draw (A) to[bend right] node [auto,swap,bend right] {texte} (B);
\end{tikzpicture}
  \\ \hline
[bend right]  & [\RDD{auto},bend right] & [auto,\RDD{swap},bend right] 
  \\ \hline
\end{tabular}  

\SbSSCT{ N\oe uds  sur un \og edge\fg}{Nodes on an edge}

\begin{tabular}{|c|c|c|}\hline  
\multicolumn{3}{|c|}{  \BS{draw}(0,0) edge \rouge{["abc", ->]} (4,0);  }\\ 
\multicolumn{3}{|c|}{  \RRR{17-12-2} }\\ 
\hline 
\begin{tikzpicture}[blue] 
\useasboundingbox  (0,-.5) rectangle (4,.5); 
\draw (0,0) edge ["abc", ->] (4,0);
\end{tikzpicture}
&
\begin{tikzpicture}[blue] 
\useasboundingbox  (0,-.5) rectangle (4,.5); 
\draw (0,0) edge ["abc", near start] (4,0);
\end{tikzpicture}
&
\begin{tikzpicture}[blue] 
\useasboundingbox  (0,-.5) rectangle (4,.5); 
\draw (0,0) edge ["abc", style={auto=right}] (4,0);
\end{tikzpicture}
\\ \hline 
["abc", ->]
& 
["abc", near start] &  ["abc", style=\AC{auto=right}] 
\\ \hline  
\begin{tikzpicture}[blue] 
\useasboundingbox  (0,-.5) rectangle (4,.5); 
\draw (0,0) edge [font=\Large,"abc" ] (4,0);
\end{tikzpicture}
&
\begin{tikzpicture}[blue] 
\useasboundingbox  (0,-.5) rectangle (4,.5); 
\draw (0,0) edge ["abc" color=red ] (4,0);
\end{tikzpicture}
&
\begin{tikzpicture}[blue] 
\useasboundingbox  (0,-.5) rectangle (4,.5); 
 \draw (0,0) edge ["abc" '] (4,0);
\end{tikzpicture}
\\ \hline 
[font=\BS{Large},"abc" ] & ["abc" color=red ]
&["abc" ' ]
\\ \hline 

\begin{tikzpicture}[blue] 
\useasboundingbox  (0,-.5) rectangle (4,.75); 
\draw (0,0) edge ["abc" draw ] (4,0);
\end{tikzpicture}
&
\begin{tikzpicture}[blue] 
\useasboundingbox  (0,-.5) rectangle (4,.5); 
\draw (0,0) edge ["abc" inner sep=0pt ] (4,0);
\end{tikzpicture}
&
\begin{tikzpicture}[blue] 
\useasboundingbox  (0,-.5) rectangle (4,.5); 
\draw (0,0) edge ["abc" fill ,fill=yellow ] (4,0);
\end{tikzpicture}
\\ \hline
["abc" draw ]
&
["abc" inner sep=0pt ]
&
["abc" fill ,fill=yellow ]
\\ \hline
\end{tabular} 



\bigskip

\begin{tabular}{|c|} \hline  
\BS{draw}[every edge quotes/.style=\AC{fill=yellow}] (0,0) edge ["abc"] (4,0);
\\ \hline  
\begin{tikzpicture}[blue] 
\useasboundingbox  (0,-.5) rectangle (4,.5); 
 \draw[every edge quotes/.style={fill=yellow}] (0,0) edge ["abc"] (4,0);
\end{tikzpicture}
\\ \hline 
\end{tabular} 

%
%\newpage
%
%\subsection{Positionnement relatif de n\oe uds}
\label{lib-pos}

\maboite{\BS{usetikzlibrary}\AC{positioning}}


\begin{center}
\RRR{17-5-3}
\end{center}

\begin{tabular}{|c|c|c|}  \hline 
\multicolumn{2}{|c|}{\BS{node} (a) at (1,0) [above=.4cm+.6cm,draw] \AC{XXX};} &  \\ \hline 
\begin{tikzpicture}
\draw[help lines] (0,0) grid (3,2);
\node (a) at (1,0) [above=.4cm+.6cm,draw] {XXX};
\draw[->,blue,line width=2pt,dotted] (1,0) -- (a.south) node [midway,right,draw=none,fill=red!10] {.4cm+.6cm} ;
\end{tikzpicture} 
&
\begin{tikzpicture}
\draw[help lines] (0,0) grid (3,2);
\node (a) at (1,0) [above=.5+sin(60),draw] {XXX};
\draw[->,blue,line width=2pt,dotted] (1,0) -- (a.south) node [midway,right,draw=none,fill=red!10] {.5+sin(60)} ;
\end{tikzpicture}  
&
\begin{tikzpicture}
\draw[help lines] (0,0) grid (2,2);
\node (a) at (1,0) [above=1,draw] {XXX};
\draw[->,blue,line width=2pt,dotted] (1,0) -- (a.south) node [midway,right,draw=none,fill=red!10] {1} ;
\end{tikzpicture}  
\\ \hline 
above = \rouge{0.4cm+0.6cm} & above = \rouge{.5+sin(60)}  & above = \rouge{1} \\ 
\hline 
\end{tabular} 

\bigskip

\begin{tabular}{|c|c|} \hline 
\multicolumn{2}{|c|}{\BS{node} (a) at (1,0) [\rouge{above right=3cm and 2cm},draw] \AC{XXX};} \\  \hline 
\begin{tikzpicture}
\draw[help lines] (0,0) grid (5,5);
\node (a) at (1,1) [above right=3cm and 2cm,draw] {XXX};
\draw[->,blue,line width=2pt,dotted] (1,1) |- (a.south west);
\end{tikzpicture}
&  
\begin{tikzpicture}
\draw[help lines] (0,0) grid (5,5);

\node (b) at (1,4) [below right=3cm and 2cm,draw] {XXX};
\draw[->,blue,line width=2pt,dotted] (1,4) |- (b.north west);
\end{tikzpicture}
\\ \hline 
\rouge{above right=3cm and 2cm} & \rouge{below right=3cm and 2cm}
\\ \hline 
\end{tabular}  

\bigskip
 
\begin{tabular}{|c|c|}  \hline 
\begin{tikzpicture}[every node/.style=draw,baseline=1.5cm]
\draw[help lines] (0,0) grid (5,4);
\node (a) at (1,1) {node a};
\node (b) [above=2cm of a.north east] {XXX};
\draw[->,blue,line width=2pt,dotted] (a.north) -- (b.south) node [midway,right,draw=none,fill=red!10] {2cm of a.north east} ;
\end{tikzpicture}
&  
\parbox{8cm}{
\BS{node} (a) at (1,1) \AC{node a}; \\
\BS{node} (b) [\rouge{above=2cm of a.north east}] \AC{XXX};}
\\ \hline 
\end{tabular} 

\bigskip

\begin{tabular}{|c|c|}  \hline 
\begin{tikzpicture}[every node/.style=draw]
\draw[help lines] (0,0) grid (2,3);
\node (a) at (1,0) {node a};
\node (b) [above=1cm of a] {node b};
\node (c) [above=1cm of b] {node c};
\draw[->,blue,line width=2pt,dotted] (a.north) -- (b.south) node [midway,right,draw=none,fill=red!10] {1cm} ;
\draw[->,blue,line width=2pt,dotted] (b.north) -- (c.south) node [midway,right,draw=none,fill=red!10] {1cm} ;
\end{tikzpicture}
&  
\begin{tikzpicture}[every node/.style=draw]
\draw[help lines] (0,0) grid (2,3);
\node (a) at (1,0) {node a };
\node (b) [on grid,above=1cm of a] {node b};
\node (c) [on grid,above=1cm of b] {node c};
\draw[->,blue,line width=2pt,dotted] (a.center) -- (b.center) node [midway,right,draw=none,fill=red!10] {1cm} ;
\draw[->,blue,line width=2pt,dotted] (b.center) -- (c.center) node [midway,right,draw=none,fill=red!10] {1cm} ;
\end{tikzpicture}
\\  \hline 
\BS{node} (a) at (1,0) \AC{node a};  &\BS{node} (a) at (1,0) \AC{node a};   \\ 
\BS{node} (b) [above=1cm of a] \AC{node b};  &\BS{node} (b) [\RDD{on grid},above=1cm of a] \AC{node b};   \\ 
\BS{node} (c) [above=1cm of b] \AC{node c};  &\BS{node} (c) [\RDD{on grid},above=1cm of b] \AC{node c};   \\ 
\hline 
\end{tabular} 

\begin{tabular}{|c|c|} \hline 
\begin{tikzpicture}[every node/.style=draw,node distance=1cm,baseline = 1.5cm]
\draw[help lines] (0,0) grid (2,3);
\node (a1) at (1,0) {node a};
\node (b) [above=of a] {node b};
\node (c) [above=of b] {node c};
\draw[->,blue,line width=2pt,dotted] (a.north) -- (b.south) node [midway,right,draw=none,fill=red!10] {1cm} ;
\draw[->,blue,line width=2pt,dotted] (b.north) -- (c.south) node [midway,right,draw=none,fill=red!10] {1cm} ;
\end{tikzpicture}
 & 
 \parbox{12cm}{ 
\BS{begin}\AC{tikzpicture}[every node/.style=draw,\RDD{node distance}=1mm] \\
\BS{node} (a1) at (1,0) \AC{node a}; \\
\BS{node} (b) [above=of a] \AC{node b}; \\
\BS{node} (c) [above=of b] \AC{node c}; \\
\BS{end}\AC{tikzpicture}
} 
 \\ 
\hline 
\end{tabular} 

\bigskip

\begin{tabular}{|l|l|} \hline 
\begin{tikzpicture}[node distance=2cm]
\draw[help lines] (0,-1) grid (6,1);
\huge
\node[draw] (X) at (0,0) {X};
\node[draw] (a) [right=of X] {a};
\node[draw] (y) [right=of a] {y};
\draw[->,blue,line width=2pt,dotted] (X.east) -- (a.west) node [midway,draw=none,fill=red!10] {\small{2cm}} ;
\draw[->,blue,line width=2pt,dotted] (a.east) -- (y.west) node [midway,draw=none,fill=red!10] {\small{2cm}} ;
\end{tikzpicture}
&  
\begin{tikzpicture}[node distance=2cm]
\draw[help lines] (0,-1) grid (6,1);
\huge
\node[draw] (X) at (0,0) {X};
\node[draw] (a) [base right=of X] {a};
\node[draw] (y) [base right=of a] {y};
\draw[->,blue,line width=2pt,dotted] (X.base east) -- (a.base west) node [midway,draw=none,fill=red!10] {\small{2cm}} ;
\draw[->,blue,line width=2pt,dotted] (a.base east) -- (y.base west) node [midway,draw=none,fill=red!10] {\small{2cm}} ;
\end{tikzpicture}
\\ \hline 
\BS{node}[draw] (X) at (0,0) \AC{X};
&  
\BS{node}[draw] (X) at (0,0) \AC{X};
\\
\BS{node}[draw] (a) [right=of X] \AC{a};
&
\BS{node}[draw] (a) [base right=of X] \AC{a};
\\
\BS{node}[draw] (y) [right=of a] \AC{y};
&
\BS{node}[draw] (y) [base right=of a] \AC{y};
\\ \hline 
\end{tabular} 


%
%\newpage
%
%\SSCT{Mettre du texte  en valeur}{Text highlighting}
%
%\label{ndbt}

\tikzset{blue}


\SbSSCT{Dans un n\oe ud de Tikz}{In a TikZ node}
\label{noeudboite}

\begin{tabular}{|c | c | c | c |} \hline
\multicolumn{4}{|c|}{ \BS{tikz} \BS{draw} (0,0) grid (2,2) (1,1) node[ fill=red!20 ] \AC{texte};   }\\ 
\hline 
\tikz \draw (0,0) grid (2,2) (1,1) node[fill=red!20] {texte};
&
\tikz \draw (0,0) grid (2,2) (1,1) node[fill=red!20,draw] {texte}; 
&
\tikz \draw (0,0) grid (2,2) (1,1) node[circle,fill=red!20] {texte};
&
\tikz \draw (0,0) grid (2,2) (1,1) node[circle,fill=red!20,draw] {texte};
\\  \hline
node[fill=red!20] 
&
node[fill=red!20,\RDD{draw}] 
&
 node[fill=red!20,\RDD{circle}]  
&
 node[fill=red!20,\RDD{circle},\RDD{draw}]
 \\  \hline
\end{tabular}
\bigskip


\subsubsection{Options}
\begin{tabular}{|c | c | c | c |c |c |c |c |} \hline
\multicolumn{8}{|c|}{ \BS{tikz} \BS{draw} node[draw,\RDD{double},blue] \AC{texte};   }\\ 
\hline 

\tikz \draw  node[draw,double,blue] {texte};
&
\tikz \draw  node[draw,rounded corners,blue] {texte};
&
\tikz \draw  node[draw,ultra thick,blue] {texte};
&
\tikz \draw  node[draw,dashed,blue] {texte};
&
\tikz \draw  node[draw,red] {texte};
&
\tikz \draw  node[draw,rotate=45,blue] {texte};
&
\tikz \draw  node[draw,shading=radial,blue] {texte};
&
\tikz \draw  node[draw,blue,text=red] {texte};
\\ \hline
\RDD{double} & \RDD{rounded corners} &  ultra thick & dashed & red & rotate=45 & shading=radial & text=red 
\\ \hline
\end{tabular}
\bigskip


\begin{tabular}{|c | c | c | c |c |} \hline
\multicolumn{4}{|c|}{ \BS{tikz} \BS{draw}  node[draw,\RDD{inner sep}=0pt] \AC{texte}; \RRR{17-2-3}  }\\ 
\hline 
\tikz \draw  node[draw,inner sep=0pt,blue] {texte};
&
\tikz \draw node[draw,inner sep=1cm,blue] {texte};
&
\tikz \draw  node[draw,inner xsep=1cm,blue] {texte};
&
\tikz \draw  node[draw,inner ysep=1cm,blue] {texte};
\\ \hline
 \RDD{inner sep}=0pt & \RDD{inner sep}=1cm & \RDD{inner xsep}=1cm & \RDD{inner ysep}=1cm
\\ \hline
\multicolumn{4}{|c|}{ \dft{} : 0.3333em }\\ 
\hline 

\end{tabular}

\bigskip

\begin{tabular}{|c | c | c | c |} \hline
\multicolumn{4}{|l|}{ \BS{node} [fill=red!20,\RDD{outer sep}=1cm] (A) at (1,1) \AC{texte}; \RRR{17-2-3} } \\ 
\multicolumn{4}{|l|}{ \BS{fill} (node cs:name=A,anchor=east) circle (3pt);  }\\ 
\multicolumn{4}{|l|}{ \BS{fill} (node cs:name=A,anchor=south) circle (3pt);  }\\ 
\hline 
\begin{tikzpicture}
\draw[help lines] (0,0) grid (3,2);
\node[fill=red!20,outer sep=1cm] (A) at (1,1) {texte};
\fill[red] (node cs:name=A,anchor=east) circle (3pt);
\fill[red] (node cs:name=A,anchor=south) circle (3pt);
\end{tikzpicture}
&
\begin{tikzpicture}
\draw[help lines] (0,0) grid (3,2);
\node[fill=red!20,outer sep=0pt] (A) at (1,1) {texte};
\fill[red] (node cs:name=A,anchor=east) circle (3pt);
\fill[red] (node cs:name=A,anchor=south) circle (3pt);
\end{tikzpicture}
&
\begin{tikzpicture}
\draw[help lines] (0,0) grid (3,2);
\node[fill=red!20,outer xsep=1cm] (A) at (1,1){texte};
\fill[red] (node cs:name=A,anchor=east) circle (3pt);
\fill[red] (node cs:name=A,anchor=south) circle (3pt);
\end{tikzpicture}
&
\begin{tikzpicture}
\draw[help lines] (0,0) grid (3,2);
\node[fill=red!20,outer ysep=1cm] (A) at (1,1) {texte};
\fill[red] (node cs:name=A,anchor=east) circle (3pt);
\fill[red] (node cs:name=A,anchor=south) circle (3pt);
\end{tikzpicture}
\\ \hline
 \RDD{outer sep}=1cm & \RDD{outer sep}=0pt & \RDD{outer xsep}=1cm & \RDD{outer ysep}=1cm
\\ \hline
\multicolumn{4}{|c|}{ \dft{} : 0.5\BS{pgflinewidth} }\\ 
\hline 
\end{tabular}

\SbSbSSCT{Taille minimale des noeuds}{Minimum size}

\begin{tabular}{|c|c|} \hline  
\multicolumn{2}{|c|}{  \BS{draw}((0,0) node[fill=blue!20,\RDD{minimum height}=1.5cm,draw]  \AC{texte} ;  \RRR{17-2-3}  }\\ 
\hline 
\tikz \draw (0,0) node[fill=red!20,minimum height=1.5cm,draw] {texte};
&  
\tikz \draw (0,0) node[fill=red!20,minimum width=3cm,draw] {texte};

\\ \hline  

\RDD{minimum height}=1.5cm
&  
\RDD{minimum width}=3cm
\\ \hline  
\tikz \draw (0,0) node[fill=red!20,minimum size=1.5cm,draw] {texte};
&  
\tikz \draw (0,0) node[fill=red!20,minimum size=1.5cm,draw,circle] {texte};

\\ \hline 
\RDD{minimum size}=1.5cm,draw
&  
\RDD{minimum size}=1.5cm,circle

\\ \hline 
\end{tabular} 

\newpage

\SbSSCT{Dans un n\oe ud à formes géométriques}{Geometric Shapes nodes}

\label{lib-geom}
\label{formes}


 \maboite{\BS{usetikzlibrary}\AC{shapes.geometric}}
 
 
\begin{center}
\RRR{67-3}
\end{center}

\SbSbSSCT{Formes disponibles}{Available shapes}

\label{nd1}

\begin{tabular}{|c|c|c|c|} \hline  
\multicolumn{4}{|l|}{ 2 syntaxes :   }\\ 
\multicolumn{4}{|l|}{ \BS{tikz} \BS{node}[fill=green!20,\RDD{shape}=diamond,draw,blue] \AC{texte};   }\\ 
\multicolumn{4}{|l|}{ \BS{tikz} \BS{node}[fill=green!20,\RDD{diamond},draw] \AC{texte};   }\\ 
\hline 
\tikz  \node[fill=green!20,diamond,draw] {texte}; 
&  
\tikz  \node[fill=green!20,ellipse,draw] {texte};
&  
\tikz  \node[fill=green!20,trapezium, regular polygon sides=6,draw] {texte};
&
\tikz  \node[fill=green!20,semicircle,draw] {texte}; 
\\ \hline 
diamond & ellipse  & trapezium & semicircle
\\ \hline 
\tikz  \node[fill=green!20,star,draw] {texte};
&  
\tikz  \node[fill=green!20,regular polygon,draw] {texte};
&  
\tikz  \node[fill=green!20,isosceles triangle,draw] {texte};
&
\tikz  \node[fill=green!20,kite,draw] {texte};
\\ \hline 
star & regular polygon  & isosceles triangle & kite 
\\ \hline 
\tikz  \node[fill=green!20,dart,draw] {texte};
&
\tikz  \node[fill=green!20,circular sector,draw] {texte};
&
\tikz  \node[fill=green!20,cylinder,draw] {texte};
&

\\ \hline 
dart & circular sector & cylinder &
\\ \hline 
\end{tabular} 

\subsubsection{Options}

\begin{tabular}{|c|c|c|} \hline
\multicolumn{3}{|c|}{  \BS{node} [trapezium,draw,\RDD{trapezium left angle}=90,draw,blue] \AC{texte};   }\\ 
\hline
\begin{tikzpicture}
\node[trapezium,draw,red,dashed] {texte};
\node[trapezium,draw,trapezium left angle=90,draw,blue] {texte};
\end{tikzpicture}
& 
\begin{tikzpicture}
\node[trapezium,draw,red,dashed] {texte};
\node[trapezium,draw,trapezium right angle=90,draw,blue] {texte};
\end{tikzpicture} 
& 
\begin{tikzpicture}
\node[trapezium,draw,red,dashed] {texte};
\node[trapezium,draw,trapezium angle=120,draw,blue] {texte};
\end{tikzpicture} 
\\ \hline
\RDD{trapezium left angle}=90  & \RDD{trapezium right angle}=90  & \RDD{trapezium  angle}=120 \\ 
\hline 
\begin{tikzpicture}
\node[trapezium,draw,red,dashed] {texte};
\node[trapezium,draw,minimum height=1.5cm,trapezium stretches=true,draw,blue] {texte};
\end{tikzpicture}
& 
\begin{tikzpicture}
\node[trapezium,draw,red,dashed] {texte};
\node[trapezium,draw,minimum height=1.5cm,trapezium stretches=false,draw,blue] {texte};
\end{tikzpicture} 
& 
\begin{tikzpicture}
\node[trapezium,draw,red,dashed] {texte};
\node[trapezium,draw,minimum width=3cm,trapezium stretches =false,draw,blue] {texte};
\end{tikzpicture} 

\\ \hline
minimum height=1.5cm & minimum height=1.5cm & minimum width=1.5cm \\
\RDD{trapezium stretches}=true & \RDD{trapezium stretches}=false & \RDD{trapezium stretches}  \\ 
\hline

\end{tabular} 


\bigskip
\begin{tabular}{|c|c|c|} \hline
\multicolumn{3}{|c|}{ \BS{tikz} \BS{node} [fill=green!20,star,\RDD{star points}=6,draw] \AC{texte};   }\\ 
\hline
\begin{tikzpicture}
\node[star,draw,red,dashed] {texte};
\node[star,star points=7,draw,blue] {texte};
\end{tikzpicture}
&  
\begin{tikzpicture}
\node[star,draw,red,dashed] {texte};
\node[star,star point height = 2cm,draw,blue] {texte};
\end{tikzpicture} 
&  
\begin{tikzpicture}
\node[star,draw,red,dashed] {texte};
\node[star,star point ratio = 3,draw,blue] {texte};
\end{tikzpicture} 
\\ \hline  
\RDD{star points}=7 & \RDD{star point height} = 2cm & \RDD{star point ratio} = 3 \\ \hline
\dft{5} & \dft.5cm &  \dft{1.5}\\ 
\hline 
\end{tabular} 
\bigskip

\begin{tabular}{|c|c|c|} \hline
\multicolumn{3}{|c|}{  \BS{node} [isosceles triangle,\RDD{isosceles triangle apex angle}=90,draw,blue] \AC{texte};   }\\ 
\multicolumn{3}{|c|}{  \BS{node} [regular polygon, \RDD{regular polygon sides}=6,draw,blue] \AC{texte};   }\\ 
\hline
\begin{tikzpicture}
\node[isosceles triangle,draw,red,dashed] {texte};
 \node[isosceles triangle,isosceles triangle apex angle=90,draw,blue] {texte};
\end{tikzpicture} 
& 
\begin{tikzpicture}
\node[isosceles triangle,draw,red,dashed] {texte};
 \node[isosceles triangle,isosceles triangle stretches=true,draw,blue] {texte};
\end{tikzpicture}
&  
\begin{tikzpicture}
\node[regular polygon,draw,red,dashed] {texte};
\node[regular polygon, regular polygon sides=6,draw,blue] {texte};
\end{tikzpicture} 
\\ \hline  
\RDD{isosceles triangle apex angle}=90 & \RDD{isosceles triangle stretches} & \RDD{regular polygon sides}=6 \\ 
\hline 
\end{tabular} 
\bigskip

\begin{tabular}{|c|c|c|} \hline 
\multicolumn{3}{|c|}{  \BS{node} [kite,\RDD{kite upper vertex angle}=90,draw,blue] \AC{texte};   }\\ 
\hline 
\begin{tikzpicture}
\node[red,kite,draw,dashed] {texte} ;
 \node[kite,kite upper vertex angle=90,draw,blue] {texte};
\end{tikzpicture} 
&  
\begin{tikzpicture}
\node[red,kite,draw,dashed] {texte} ;
 \node[kite,kite lower vertex angle=90,draw,blue] {texte};
\end{tikzpicture} 
&  
\begin{tikzpicture}
\node[red,kite,draw,dashed] {texte} ;
\node[kite,kite vertex angles=90,draw,blue] {texte};
\end{tikzpicture} 
\\ \hline  
\RDD{kite upper vertex angle}=90 & \RDD{kite lower vertex angle}=90 &\RDD{kite vertex angles}=90
\\ \hline 
initially 120 & initially 60 &  \\ 
\hline 
\end{tabular} 

\bigskip

\begin{tabular}{|c|c|c|} \hline
\multicolumn{3}{|c|}{  \BS{node} [dart,\RDD{dart tip angle}=90,draw,blue] \AC{texte};   }\\ 
\hline 
\begin{tikzpicture}
\node[dart,draw,red,dashed] {texte};
\node[dart,dart tip angle=90,draw,blue] {texte};
\end{tikzpicture} 
&  
\begin{tikzpicture}
\node[dart,draw,red,dashed] {texte};
\node[dart,dart tail angle=90,draw,blue] {texte};
\end{tikzpicture} 
&  
\begin{tikzpicture}
\node[,circular sector,draw,red,dashed] {texte};
\node[circular sector,circular sector angle=90,draw,blue] {texte};
\end{tikzpicture} 
\\ \hline  
\RDD{dart tip angle}=90 & \RDD{dart tail angle}=90  & \RDD{circular sector angle}=90
\\ \hline  
initially 45 & initially 135 & initially 60  \\ 
\hline 
\end{tabular} 

\bigskip

\begin{tabular}{|c|c|} \hline  
\multicolumn{2}{|c|}{  \BS{node} [cylinder,\RDD{aspect=2},draw,blue] \AC{texte};   }\\ 
\hline
\tikz  \node[cylinder,aspect=2,draw,blue] {texte};
& 
 \tikz  \node[cylinder,aspect=4,draw,blue] {texte};
\\ \hline 
\RDD{aspect}=2 & \RDD{aspect}=4 
\\ \hline
\tikz  \node[cylinder,cylinder uses custom fill, cylinder end fill=yellow,draw,blue] {texte};
&  
\tikz  \node[cylinder,cylinder uses custom fill, cylinder body fill=yellow,draw,blue] {texte};
\\ \hline
\RDD{cylinder uses custom fill}, & \RDD{cylinder uses custom fill}, \\ 
\RDD{cylinder end fill}=yellow & \RDD{cylinder body fill}=yellow  \\ 
\hline 
\end{tabular} 

\bigskip

\begin{tabular}{|c|c|c|c|} \hline 
\multicolumn{4}{|c|}{  \BS{draw}(0,0) node[\RDD{shape aspect}=1,diamond,draw]  \AC{texte} ;   }
\\ \hline
 
\tikz \draw (0,0) node[shape aspect=1,diamond,draw,blue] {texte};
&  
\tikz \draw (0,-2) node[shape aspect=2,diamond,draw,blue] {texte};
&
\tikz \draw (0,0) node[shape aspect=3,diamond,draw,blue] {texte};
&
\tikz \draw (0,0) node[shape aspect=4,diamond,draw,blue] {texte};
\\ \hline  
\RDD{shape aspect}=1
&  
\RDD{shape aspect}=2
&
\RDD{shape aspect}=3
&
\RDD{shape aspect}=4
\\ \hline 
\end{tabular} 

\bigskip

\begin{tabular}{|c|} \hline 
\BS{draw} node[\rouge {shape border rotate}=30,shape=dart, draw, \rouge {shape border uses incircle}] \AC{texte};
\\ \hline 
\tikz[] \draw node[shape border rotate=30,shape=dart, draw, shape border uses incircle] {texte};
\\ \hline 
\end{tabular} 

\newpage

\SbSSCT{Dans un n\oe ud en forme de symboles}{Symbol Shapes nodes}

\label{lib-symb}

\maboite{\BS{usetikzlibrary}\AC{shapes.symbols}}

\begin{center}
\RRR{67-4}
\end{center}

\SbSbSSCT{Formes disponibles}{Available shapes}

\label{nd2}

\begin{tabular}{|c|c|c|} \hline  
\tikz  \node[fill=green!20,forbidden sign,draw] {texte};
&  
\tikz  \node[fill=green!20,magnifying glass,draw] {texte};
&  
\tikz  \node[fill=green!20,cloud,draw] {texte};
\\ \hline 
forbidden sign & magnifying glass & cloud
\\ \hline  
\tikz  \node[fill=green!20,starburst,draw] {texte};
&  
\tikz  \node[fill=green!20,signal,draw] {texte};

&  
\tikz  \node[fill=green!20,tape,draw] {texte};
\\ \hline 
starburst & signal & tape
\\ \hline 
\end{tabular} 
\bigskip

\subsubsection{Options}

\begin{tabular}{|c|c|c|} \hline  
\multicolumn{3}{|c|}{  \BS{node}[magnifying glass,\RDD{magnifying glass handle angle}=45,draw,blue]  \AC{texte} ;   }
\\ \hline
\tikz  \node[magnifying glass,magnifying glass handle angle=45,draw,blue] {texte};
&  
\tikz  \node[,magnifying glass,magnifying glass handle aspect=3,draw,blue] {texte};
& 
\tikz  \node[magnifying glass,line width=1ex,draw,blue] {texte};

\\ \hline  
\RDD{magnifying glass handle angle}=45 & \RDD{magnifying glass handle aspect}=3  & line width=1ex  
\\ \hline 
\dft{ : -45} & \dft{ : 1.5}& 
\\ \hline 
\end{tabular} 

\bigskip

\begin{tabular}{|c|c|c|c|} \hline 
\multicolumn{4}{|c|}{  \BS{node} [cloud,\RDD{cloud puffs}=5,draw,blue] \AC{texte};   }\\ 
\hline 
\begin{tikzpicture}
\node[cloud,draw,red,dashed] {texte};
\node[cloud,cloud puffs=5,draw,blue] {texte};
\end{tikzpicture} 
&  
\begin{tikzpicture}
\node[cloud,draw,red,dashed] {texte};
\node[cloud,cloud puff arc=270,draw,blue] {texte};
\end{tikzpicture} 
&  
\begin{tikzpicture}
\node[cloud,draw,red,dashed] {texte};
\node[cloud,cloud ignores aspect=true,draw,blue] {texte};
\end{tikzpicture} 
&
\begin{tikzpicture}
\node[cloud,draw,red,dashed] {texte};
\node[cloud,cloud ignores aspect=false,draw,blue] {texte};
\end{tikzpicture} 
\\ \hline  
\RDD{cloud puffs}=5 & \RDD{cloud puff arc}=270 & \RDD{cloud ignores aspect}=false & \RDD{cloud ignores aspect}=true  \\ 
\hline 
\dft :  10 & \dft :  135 &\multicolumn{2}{|c|}{ \dft :  true } \\ \hline
\end{tabular} 

\bigskip

\begin{tabular}{|c|c|c|c|} \hline 
\multicolumn{4}{|c|}{  \BS{node} [starburst,\RDD{starburst points}=5,draw,blue] \AC{texte};   }\\ 
\hline  
\tikz  \node[starburst,starburst points=5,draw,blue] {texte};
&  
\tikz  \node[starburst,starburst point height=1cm,draw,blue] {texte};
&  
\tikz  \node[starburst,random starburst=50,draw,blue] {texte};
&
\tikz  \node[,starburst,random starburst=0,draw,blue] {texte};
\\ \hline  
\RDD{starburst points}=5 & \RDD{starburst point height}=1cm & \RDD{random starburst}=50 & \RDD{random starburst}=0  \\ 
\hline 
\end{tabular} 

\bigskip


\begin{tabular}{|c|c|c|} \hline 
\multicolumn{3}{|c|}{  \BS{node} [signal,\RDD{signal pointer angle}=45,draw,blue] \AC{texte};   }\\ 
\hline 
\tikz  \node[signal,signal pointer angle=45,draw,blue] {texte};
&
\tikz  \node[signal,signal pointer angle=10,draw,blue] {texte};
&
\tikz  \node[signal,signal pointer angle=300,draw,blue] {texte};
\\ \hline 
\RDD{signal pointer angle}=45
&
signal pointer angle=10
&
signal pointer angle=300
\\ \hline 
\multicolumn{3}{|c|}{  \dft{ : signal pointer angle= 90}  }
\\  \hline 

\end{tabular} 
\bigskip

\begin{tabular}{|c|c|c|c|c|} \hline 
\multicolumn{4}{|c|}{  \BS{node} [signal,\RDD{signal to}=above,draw,blue] \AC{texte};   }
\\ \hline 
\tikz  \node[signal,signal to=above,draw,blue] {texte};
&  
\tikz  \node[signal,signal to=below,draw,blue] {texte};
&
\tikz  \node[signal,signal to=right,draw,blue] {texte};
&
\tikz  \node[signal,signal to=above,draw,blue] {texte};
\\ \hline  
  \RDD{signal to}=above  & \RDD{signal to}=below & \RDD{signal to}=right  & \RDD{signal to}=above \\ 
\hline 
\end{tabular} 
\bigskip

\begin{tabular}{|c|c|c|c|c|} \hline 
\multicolumn{4}{|c|}{ \BS{tikz} [signal to=nowhere] \BS{node} [signal,\RDD{signal from=above}=45,draw,blue] \AC{texte};   }\\ 
\hline 
\tikz [signal to=nowhere] \node[signal,signal from=above,draw,blue] {texte};
&  
\tikz [signal to=nowhere] \node[signal,signal from=below,draw,blue] {texte};
&
\tikz [signal to=nowhere] \node[signal,signal from=right,draw,blue] {texte};
&
\tikz [signal to=nowhere] \node[signal,signal from=above,draw,blue] {texte};
\\ \hline  
  \RDD{signal from}=above  & \RDD{signal from}=below & \RDD{signal from}=right  & \RDD{signal from}=above \\ 
\hline 
\end{tabular} 

\bigskip
\begin{tabular}{|c|c|c|c|} \hline
\multicolumn{2}{|c|}{ \tikz  \node[draw,signal, signal from=east , signal to=west,blue] at (0,0) {texte};}
&
\multicolumn{2}{|c|}{ \tikz  \node[draw,signal,signal from=south, signal to=north,blue] at (0,0) {texte};}
\\ \hline 
\multicolumn{2}{|c|}{ \RDD{signal from}=east , \RDD{signal to}=west}
&
\multicolumn{2}{|c|}{\RDD{signal from}=south, \RDD{signal to}=north}

\\ \hline 
\end{tabular}
\bigskip

\begin{tabular}{|c | c | c | c |} \hline
\multicolumn{3}{|c|}{ \BS{tikz} \BS{node}  [tape, draw,\RDD{tape bend top}=out and in] \AC{texte};   }\\ 
\hline  
\tikz \node [tape, draw,tape bend top=out and in,blue] {texte};
&
\tikz \node [tape, draw, tape bend bottom=out and in,blue] {texte};
&
\tikz \node [tape, draw, tape bend bottom=in and in,blue] {texte};
 \\  \hline
 \RDD{tape bend top}=out and in & \RDD{tape bend bottom}=out and in &  \RDD{tape bend bottom}=in and in 
  \\  \hline
 \tikz \node [tape, draw, tape bend top=none,blue] {texte};
 &
 \tikz \node [tape, draw,tape bend top=out and in,tape bend bottom=out and in,blue] {texte};
 &
  \tikz \node [tape, draw,tape bend top=in and out,tape bend bottom=in and out,blue] {texte};
  \\  \hline
 \RDD{tape bend top}=none & \RDD{tape bend bottom}=out and in 	&  \RDD{tape bend bottom}=in and out  \\
 					& \RDD{tape bend top}=out and in 		& \RDD{tape bend top}=in and out  \\
 					& & (\dft{} ) 
  \\  \hline 
\end{tabular}
\bigskip

\begin{tabular}{|c | c | c | c |} \hline
\BS{tikz} \BS{node} [tape, draw, \RDD{tape bend height}=1cm,blue] \AC{texte}; 
  \\  \hline 
\tikz \node [tape, draw, tape bend height=1cm,blue] {texte};

  \\  \hline 
\dft{ : tape bend height = 5pt}
  \\  \hline 
\end{tabular}

\newpage

\SbSSCT{Dans un n\oe ud en forme de flèche}{Arrow Shapes nodes}

\label{lib-arr}

\maboite{\BS{usetikzlibrary}\AC{shapes.arrows}}

\begin{center}
\RRR{67-5}
\end{center}

\SbSbSSCT{Formes disponibles}{Available shapes}
\label{nd3}

\begin{tabular}{|c|c|c|} \hline  
\tikz \node[fill=green!20,single arrow,draw] {texte};
&  
\tikz  \node[fill=green!20,double arrow,draw] {texte};
&  
\tikz  \node[fill=green!20,arrow box,draw] {texte};
\\ \hline 
single arrow & double arrow & arrow box \\ 
\hline 
\end{tabular} 

\subsubsection{Options}

\begin{tabular}{|c|c|c|c|c|} \hline  
 \multicolumn{5}{|c|}{  \BS{node}[single arrow,draw,\RDD{single arrow tip angle}=45] \AC{texte};   }\\ 
  \multicolumn{5}{|c|}{  \BS{node}[single arrow,draw,\RDD{single arrow head extend}=.75cm] \AC{texte};   }\\
 \hline
\begin{tikzpicture}
 \node[single arrow,draw,red,dashed,text=black] {texte};
 \node[single arrow,draw,single arrow tip angle=45,blue] {texte};
\end{tikzpicture}
&
\begin{tikzpicture}
 \node[single arrow,draw,red,dashed,text=black] {texte};
\node[single arrow,draw,single arrow tip angle=120,blue] {texte};
\end{tikzpicture}
&
\begin{tikzpicture}
 \node[single arrow,draw,red,dashed,text=black] {texte};
 \node[single arrow,draw,single arrow head extend=.75cm,blue] {texte};
\end{tikzpicture}
&
\begin{tikzpicture}
 \node[single arrow,draw,red,dashed,text=black] {texte};
 \node[single arrow,draw,single arrow head extend=0cm,blue] {texte};
 \end{tikzpicture}
 &
 \begin{tikzpicture}
  \node[single arrow,draw,red,dashed,text=black] {texte};
  \node[single arrow,draw,single arrow head extend=-1mm,blue] {texte};
 \end{tikzpicture}

\\ \hline
angle=45 & angle=120 & extend=.75cm] & extend=0cm & extend=-1mm
\\ \hline 
\multicolumn{2}{|c|}{  \dft : single arrow tip angle= 90   }
&
\multicolumn{3}{|c|}{  \dft : single arrow head extend=0.5cm   }
\\ \hline 
\end{tabular} 
\bigskip


\begin{tabular}{|c|c|c|c|} \hline
 \multicolumn{4}{|c|}{  \BS{node}[minimum size=2cm,single arrow,draw,\RDD{single arrow head indent}=1cm,blue] \AC{texte};   }\\ 
 \hline   
\begin{tikzpicture}
 \node[minimum size=2cm,single arrow,draw,red,dashed,text=black] {texte};
\node[minimum size=2cm,single arrow,draw,single arrow head indent=1cm,blue] {texte};
\end{tikzpicture}
&
\begin{tikzpicture}
 \node[minimum size=2cm,single arrow,draw,red,dashed,text=black] {texte};
  \node[minimum size=2cm,single arrow,draw,single arrow head indent=10pt,blue] {texte};
  \end{tikzpicture}
&
\begin{tikzpicture}
 \node[minimum size=2cm,single arrow,draw,red,dashed,text=black] {texte};
  \node[minimum size=2cm,single arrow,draw,single arrow head indent=1ex,blue] {texte};
  \end{tikzpicture}
  &
  \begin{tikzpicture}
   \node[minimum size=2cm,single arrow,draw,red,dashed,text=black] {texte};
    \node[minimum size=2cm,single arrow,draw,single arrow head indent=-1ex,blue] {texte};
    \end{tikzpicture}
\\ \hline
indent=1cm & indent=10pt & indent=1ex & indent=-1ex
\\ \hline 
\end{tabular}
\bigskip

 



\begin{tabular}{|c|c|c|c|c|} \hline
 \multicolumn{5}{|c|}{  \BS{node}[minimum size=2cm,double arrow,draw,\RDD{double arrow tip angle}=45] \AC{texte};   }\\ 
  \multicolumn{5}{|c|}{  \BS{node}[minimum size=2cm,double arrow,draw,\RDD{double arrow head extend}=1ex] \AC{texte};   }\\
   \multicolumn{5}{|c|}{  \BS{node}[minimum size=2cm,double arrow,draw,\RDD{double arrow head indent}=1ex] \AC{texte};   }\\ 
 \hline  
\begin{tikzpicture}
\node[minimum size=2cm,double arrow,draw,red,dashed,text=black] {texte};
\node[minimum size=2cm,double arrow,draw,double arrow tip angle=45,blue] {texte};
\end{tikzpicture}
&
\begin{tikzpicture}
\node[minimum size=2cm,double arrow,draw,red,dashed,text=black] {texte};
\node[minimum size=2cm,double arrow,draw,double arrow tip angle=120,blue] {texte};
\end{tikzpicture}
&
\begin{tikzpicture}
 \node[minimum size=2cm,double arrow,draw,red,dashed,text=black] {texte};
 \node[minimum size=2cm,double arrow,draw,double arrow head extend=1ex,blue] {texte};
   \end{tikzpicture}
&
\begin{tikzpicture}
 \node[minimum size=2cm,double arrow,draw,red,dashed,text=black] {texte};
  \node[minimum size=2cm,double arrow,draw,double arrow head extend=0,blue] {texte};
    \end{tikzpicture}
&
\begin{tikzpicture}
 \node[minimum size=2cm,double arrow,draw,red,dashed,text=black] {texte};
  \node[,minimum size=2cm,double arrow,draw,double arrow head indent=1ex,blue] {texte};
    \end{tikzpicture}
\\ \hline 
angle=45 & angle=120 & extend=1ex & extend=0 & indent=1ex
\\ \hline
\end{tabular}

\bigskip

\begin{tabular}{|c|c|c|c|c|} \hline
\multicolumn{4}{|c|}{ \BS{node} [arrow box, draw, \RDD{arrow box arrows}=\AC{north:.25cm}] \AC{texte}; }\\ 
\hline 
\begin{tikzpicture}
\node[arrow box, draw,red,text=white,dashed] {texte};
\node[arrow box, draw, arrow box arrows={north:.25cm},blue] {texte};
\end{tikzpicture}
& 
\begin{tikzpicture}
\node[arrow box, draw,red,text=white,dashed] {texte};
\node[arrow box, draw, arrow box arrows={west:.25cm},blue] {texte};
\end{tikzpicture}
 &
 \begin{tikzpicture}
 \node[arrow box, draw,red,text=white,dashed] {texte};
 \node[arrow box, draw, arrow box arrows={south:.25cm},blue] {texte};
 \end{tikzpicture}
&
 \begin{tikzpicture}
 \node[arrow box, draw,red,text=white,dashed] {texte};
 \node[arrow box, draw, arrow box arrows={east:.25cm},blue] {texte};
 \end{tikzpicture}   
 \\ \hline
\AC{north:.25cm} & \AC{west:.25cm} & \AC{south:.25cm}& \AC{east:.25cm} 
\\ \hline
\multicolumn{4}{|c|}{  \dft{} : 0.5 cm}
 \\ \hline 
 \end{tabular}
 
 
 \bigskip
 
 \begin{tabular}{|c|c|} \hline
 \multicolumn{2}{|c|}{ \BS{node} [arrow box, draw, \RDD{arrow box tip angle}=45] \AC{texte}; }\\ 
 \hline 
  \begin{tikzpicture}
  \node[arrow box, draw,red,text=white,dashed] {texte};
  \node[arrow box, draw, arrow box tip angle=45,blue] {texte};
  \end{tikzpicture} 
  &
    \begin{tikzpicture}
   \node[arrow box, draw,red,text=white,dashed] {texte};
   \node[arrow box, draw, arrow box head extend=.25cm,blue] {texte};
   \end{tikzpicture}
\\ \hline  
\RDD{arrow box tip angle}=45 & \RDD{arrow box head extend}=.25cm
\\ \hline 
\dft : 90  & \dft : 0.125cm 
\\ \hline 
   \begin{tikzpicture}
   \node[arrow box, draw,red,text=white,dashed] {texte};
   \node[arrow box, draw, arrow box head indent=.25cm,blue] {texte};
   \end{tikzpicture} 
 &
    \begin{tikzpicture}
    \node[arrow box, draw,red,text=white,dashed] {texte};
    \node[arrow box, draw,arrow box shaft width=.25cm,blue] {texte};
    \end{tikzpicture} 
 \\ \hline 
\RDD{arrow box head indent}=.25cm  &  \RDD{arrow box shaft width}=.25cm
 \\ \hline  
 \dft{ : 0cm } &  \dft{ : 0.125cm }
 \\ \hline  
 \end{tabular}

\newpage

\SbSSCT{Dans un n\oe ud en forme de bulle}{Callout Shapes nodes}
\label{lib-call}

 \maboite{\BS{usetikzlibrary}\AC{shapes.callouts}}
 
\begin{center}
\RRR{67-7}
\end{center}

\SbSbSSCT{Formes disponibles}{Available shapes}

\begin{tabular}{|c|c|c|} \hline 
\tikz  \node[fill=green!20,ellipse callout,draw] {texte};
 &  
 \tikz  \node[fill=green!20,rectangle callout,draw] {texte};
  &  
  \tikz  \node[fill=green!20,cloud callout,draw] {texte};
 \\ \hline
 ellipse callout  &  rectangle callout  & cloud callout \\ 
\hline 
\end{tabular} 

\subsubsection{Options}


\begin{tabular}{|c | c | c | c |} \hline
\multicolumn{4}{|c|}{  \BS{node} [rectangle callout,draw,\RDD{callout absolute pointer}={(0,1)}] at (2,1) \AC{texte};   }\\ 
\hline 
\begin{tikzpicture} 
\draw [help lines] grid(3,3);
\node [rectangle callout,draw,blue, callout relative pointer={(0,1)}] at (2,1) {texte};
\end{tikzpicture}
&
\begin{tikzpicture} 
\draw [help lines] grid(3,3);
\node [ellipse callout,draw, callout relative pointer={(0,1)},blue] at (2,1) {texte};
\end{tikzpicture}
&
\begin{tikzpicture} 
\draw [help lines] grid(3,3);
\node [rectangle callout,draw,blue,callout absolute pointer={(0,1)}] at (2,1) {texte};
\end{tikzpicture}
&
\begin{tikzpicture} 
\draw [help lines] grid(3,3);
\node [ellipse callout,draw, callout absolute pointer={(0,1)},blue] at (2,1) {texte};
\end{tikzpicture}
 \\  \hline
\multicolumn{2}{|c|}{ \RDD{callout relative pointer}=\AC{(0,1)} } & 
\multicolumn{2}{|c|}{  \RDD{callout absolute pointer}=\AC{(0,1)} }
 \\  \hline 
 \begin{tikzpicture} 
 \draw [help lines] grid(3,3);
 \node [rectangle callout,draw, callout relative pointer={(0,1)},callout pointer shorten=.5cm,blue] at (2,1) {texte};
 \end{tikzpicture}
 &
  \begin{tikzpicture} 
  \draw [help lines] grid(3,3);
  \node [ellipse callout,draw, callout relative pointer={(0,1)},callout pointer shorten=.5cm,blue] at (2,1) {texte};
  \end{tikzpicture}
  &
 \begin{tikzpicture} 
 \draw [help lines] grid(3,3);
 \node [rectangle callout,draw, callout absolute pointer={(0,1)},callout pointer shorten=.5cm,blue] at (2,1) {texte};
 \end{tikzpicture}
  &
  \begin{tikzpicture} 
  \draw [help lines] grid(3,3);
  \node [ellipse callout,draw, callout absolute pointer={(0,1)},callout pointer shorten=.5cm,blue] at (2,1) {texte};
  \end{tikzpicture}
  \\  \hline
\multicolumn{4}{|c|}{ \RDD{callout pointer shorten}=.5cm} 
  \\  \hline 
\end{tabular}


\bigskip

\begin{tabular}{|c | c | c | c |} \hline
\multicolumn{3}{|c|}{  \BS{node} [ellipse callout,draw,\RDD{callout pointer arc}=1] at (0,1.5) \AC{texte};   }\\ 
\hline
\begin{tikzpicture}
\node[ellipse callout,draw, callout pointer arc=1,blue] at (0,1.5) {texte};
\end{tikzpicture}
&
\begin{tikzpicture}
\node[ellipse callout,draw, callout pointer arc=30,blue] at (0,1.5) {texte};
\end{tikzpicture}
 &
\begin{tikzpicture}
\node[ellipse callout,draw, callout pointer arc=90,blue] at (0,1.5) {texte};
\end{tikzpicture}
  \\  \hline 
   callout pointer arc=1 & callout pointer arc=30 & callout pointer arc=90
  \\  \hline  
  \multicolumn{3}{|c|}{  \dft{ : callout pointer arc=15}}
 \\  \hline  
 \end{tabular}

\bigskip

\begin{tabular}{|c | c | c | c |} \hline
\multicolumn{3}{|c|}{  \BS{node}[draw,cloud callout, aspect=2.5] \AC{texte};   }\\ 
\hline 
 \begin{tikzpicture}
  \node[draw,cloud callout, dashed,red,text=black] {texte};
 \node[draw,cloud callout, cloud puffs=5,blue] {texte};
 \end{tikzpicture}
&
 \begin{tikzpicture}
 \node[draw,cloud callout, dashed,red,text=black] {texte};
 \node[draw,cloud callout, aspect=2.5,blue] {texte};
 \end{tikzpicture}
&
  \begin{tikzpicture}
  \node[draw,cloud callout, dashed,red,text=black] {texte};
  \node[draw,cloud callout,cloud puff arc=120,blue] {texte};
  \end{tikzpicture}
   \\  \hline 
cloud puffs=5 & aspect=2.5 &  cloud puff arc=120
\\  \hline 
 \end{tabular}

\bigskip

\begin{tabular}{|c | c | c | c |c |} \hline
\multicolumn{3}{|c|}{  \BS{node} [draw,cloud callout,\RDD{callout pointer start size}=.1] \AC{texte};   }\\ 
\hline 
  \begin{tikzpicture}
  \node[draw,cloud callout, dashed,red,text=black] {texte};
  \node[draw,cloud callout,callout pointer start size=.1,blue] {texte};
  \end{tikzpicture}
&
  \begin{tikzpicture}
  \node[draw,cloud callout, dashed,red,text=black] {texte};
  \node[draw,cloud callout,callout pointer start size=.8cm,blue] {texte};
  \end{tikzpicture}
&
  \begin{tikzpicture}
  \node[draw,cloud callout, dashed,red,text=black] {texte};
 \node[draw,cloud callout,callout pointer start size=1cm and 0.1cm,blue] {texte};
  \end{tikzpicture}
\\  \hline 
\RDD{callout pointer start size}=.1 &start size=.8cm & start size=20pt and 1pt
\\  \hline 
\multicolumn{3}{|c|}{  \dft{} : callout pointer start size =.2 of callout  }
\\ 
\hline 
  \begin{tikzpicture}
  \node[draw,cloud callout, dashed,red,text=black] {texte};
  \node[draw,cloud callout,callout pointer end size=5,blue] {texte};
  \end{tikzpicture}
&
  \begin{tikzpicture}
  \node[draw,cloud callout, dashed,red,text=black] {texte};
  \node[draw,cloud callout,callout pointer end size=.8cm,blue] {texte};
  \end{tikzpicture}
&
    \begin{tikzpicture}
    \node[draw,cloud callout, dashed,red,text=black] {texte};
    \node[draw,cloud callout,callout pointer segments=3,blue] {texte};
    \end{tikzpicture}
\\  \hline 
\RDD{callout pointer end size}=.5 & \RDD{callout pointer end size}=.8cm & \RDD{callout pointer segments}=3
\\  \hline 
\multicolumn{2}{|c|}{  \dft{} : callout pointer start size = .1 of callout  }
& \dft{} : segments=2
\\  \hline  

 \end{tabular}

\newpage


\SbSSCT{Dans un n\oe ud en diverses formes  diverses}{Miscellaneous Shapes nodes}

\label{lib-misc}


 \maboite{\BS{usetikzlibrary}\AC{shapes.misc}}
 
\begin{center}
\RRR{67-8}
\end{center}

\SbSbSSCT{Formes disponibles}{Available shapes}

\begin{tabular}{|c|c|c|c|} \hline  
\tikz  \node[fill=green!20,cross out,draw] {texte};
&  
\tikz  \node[fill=green!20,strike out,draw] {texte};
&  
\tikz  \node[fill=green!20,rounded rectangle,draw] {texte};
&  
\tikz  \node[fill=green!20,chamfered rectangle,draw] {texte};
\\ \hline  
cross out & strike out & rounded rectangle & chamfered rectangle \\ 
\hline 
\end{tabular} 


\subsubsection{Options}

\paragraph{Options \TFRGB{pour}{for} \og rounded rectangle \fg} :


\begin{tabular}{|c|c|c|c|c|} \hline
\multicolumn{5}{|c|}{  \BS{node} [draw, rounded rectangle,\RDD{rounded rectangle arc length}=270] \AC{texte};   }\\ 

\hline 

\tikz \node[draw, rounded rectangle,rounded rectangle arc length=270,blue] {texte}; 
&
\tikz \node[draw, rounded rectangle,rounded rectangle arc length=180,blue]  {texte}; 
&
\tikz \node[draw, rounded rectangle,rounded rectangle arc length=120,blue] {texte}; 
&
\tikz \node[draw, rounded rectangle,rounded rectangle arc length=90,blue]  {texte}; 
&
\tikz \node[draw, rounded rectangle,rounded rectangle arc length=45,blue] {texte}; 
 \\ \hline 
270 & 180 & 120 & 90& 45 
\\ \hline 


\end{tabular} 

\bigskip


\begin{tabular}{|c|c|c|c|} \hline 
\multicolumn{4}{|c|}{  \BS{node} [draw, rounded rectangle,\RDD{rounded rectangle west arc}=concave] \AC{texte};   }\\ 
\multicolumn{4}{|c|}{  \BS{node} [draw, rounded rectangle,\RDD{rounded rectangle left arc}=concave] \AC{texte};   }\\ 
\hline 
\tikz \node[draw, rounded rectangle,rounded rectangle west arc=concave,blue] {texte}; 
&
\tikz \node[draw, rounded rectangle,rounded rectangle left arc=concave,blue] {texte}; 
&
\tikz \node[draw, rounded rectangle,rounded rectangle west arc=convex,blue] {texte}; 
&
\tikz \node[draw, rounded rectangle,rounded rectangle left arc=none,blue] {texte};
 \\\hline 
concave & convex & none 
 \\\hline 
\end{tabular} 

\bigskip

\begin{tabular}{|c|c|c|c|} \hline 
\multicolumn{3}{|c|}{  \BS{node} [draw, rounded rectangle,\RDD{rounded rectangle east arc}=concave] \AC{texte};   }\\ 
\multicolumn{3}{|c|}{  \BS{node} [draw, rounded rectangle,\RDD{rounded rectangle right arc}=concave] \AC{texte};   }\\ 

\hline 
\tikz \node[draw, rounded rectangle,rounded rectangle east arc=concave,blue] {texte}; 
&
\tikz \node[draw, rounded rectangle,rounded rectangle  east arc=convex,blue] {texte}; 
&
\tikz \node[draw, rounded rectangle,rounded rectangle right arc=none,blue] {texte};
 \\\hline 
concave & convex & none 
 \\\hline 
\end{tabular} 

\paragraph{Options  \TFRGB{pour}{for} \og chamfered rectangle \fg} :


\begin{tabular}{|c|c|c|c|} \hline 
\multicolumn{4}{|c|}{  \BS{node} [draw, chamfered rectangle,\RDD{chamfered rectangle angle}=30] \AC{texte};   }\\ 
\hline 
\tikz \node[draw, chamfered rectangle,chamfered rectangle angle=10,blue] {texte}; 
&
\tikz \node[draw, chamfered rectangle,chamfered rectangle angle=30,blue] {texte}; 
&
\tikz \node[draw,chamfered rectangle,chamfered rectangle angle=60,blue] {texte};
&
\tikz \node[draw,chamfered rectangle,chamfered rectangle angle=80,blue] {texte};
 \\ \hline 
10 & 30 & 60 & 80
\\ \hline 
\multicolumn{4}{|c|}{  \dft :  45 }
  \\\hline  

\end{tabular}

\bigskip

\begin{tabular}{|c|c|c|c|c|} \hline 
\multicolumn{5}{|c|}{  \BS{node} [draw, chamfered rectangle,\RDD{chamfered rectangle xsep}=10pt] \AC{texte};   }\\ 
\hline 
\tikz \node[draw, chamfered rectangle,chamfered rectangle xsep=0pt,blue] {texte}; 
&
\tikz \node[draw, chamfered rectangle,chamfered rectangle xsep=5pt,blue] {texte}; 
&
\tikz \node[draw, chamfered rectangle,chamfered rectangle xsep=10pt,blue] {texte}; 
&
\tikz \node[draw,chamfered rectangle,chamfered rectangle xsep=-10pt,blue] {texte};
&
\tikz \node[draw,chamfered rectangle,chamfered rectangle xsep=2cm,blue] {texte};
 \\\hline 
  xsep=0pt & xsep=5pt & xsep=10pt & xsep=-10pt  & xsep=2cm
  \\\hline  
\multicolumn{5}{|c|}{  \dft :  0.666ex }
  \\\hline   
\end{tabular}

\bigskip

\begin{tabular}{|c|c|c|c|c|} \hline 
\multicolumn{5}{|c|}{  \BS{node} [draw, chamfered rectangle,\RDD{chamfered rectangle ysep}=10pt] \AC{texte};   }\\ 
\hline 
\tikz \node[draw, chamfered rectangle,chamfered rectangle ysep=0pt,blue] {texte}; 
&
\tikz \node[draw, chamfered rectangle,chamfered rectangle ysep=5pt,blue] {texte}; 
&
\tikz \node[draw,chamfered rectangle,chamfered rectangle ysep=10pt,blue] {texte};
&
\tikz \node[draw,chamfered rectangle,chamfered rectangle ysep=-10pt,blue] {texte};
&
\tikz \node[draw,chamfered rectangle,chamfered rectangle ysep=1cm,blue] {texte};
 \\ \hline 
 ysep=0pt & ysep=5pt & ysep=10pt & ysep=-10pt & ysep=1cm
 \\\hline  
\end{tabular}

\bigskip

\begin{tabular}{|c|c|c|c|c|} \hline 
\multicolumn{5}{|c|}{  \BS{node} [draw, chamfered rectangle,\RDD{chamfered rectangle ysep}=10pt] \AC{texte};   }\\ 
\hline 
\tikz \node[draw, chamfered rectangle,chamfered rectangle sep=0pt,blue] {texte}; 
&
\tikz \node[draw, chamfered rectangle,chamfered rectangle sep=5pt,blue] {texte}; 
&
\tikz \node[draw, chamfered rectangle,chamfered rectangle sep=10pt,blue] {texte}; 

&
\tikz \node[draw, chamfered rectangle,chamfered rectangle sep=-10pt,blue] {texte}; 
&
\tikz \node[draw,chamfered rectangle,chamfered rectangle sep=1cm,blue] {texte};
 \\\hline 
 sep=0pt & sep=5pt & sep=10pt& sep=-10pt & sep=1cm
 \\\hline  
\end{tabular}

\bigskip

\begin{tabular}{|c|c|c|c|} \hline 
\multicolumn{3}{|c|}{  \BS{node} [draw, chamfered rectangle,\RDD{chamfered rectangle corners}=north west] \AC{texte};   }\\ 
\hline
\tikz \node[draw, chamfered rectangle,chamfered rectangle corners=north west,blue] {texte}; 
&
\tikz \node[draw, chamfered rectangle,chamfered rectangle corners={north east, south east},blue] {texte}; 
&
\tikz \node[draw,chamfered rectangle,chamfered rectangle corners={north east, south west},blue] {texte};
 \\ \hline 
 north west & \AC{north east, south east}  & \AC{north east, south west}
 \\ \hline 
\end{tabular}

\newpage

\SbSSCT{N\oe uds à plusieurs parties}{Shapes with Multiple Text Parts}

\label{lib-mult}


 \maboite{\BS{usetikzlibrary}\AC{shapes.multipart}}

\begin{center}
\RRR{67-6}
\end{center}



\begin{tabular}{|c|c|c|c|} \hline 
\multicolumn{4}{|c|}{  \BS{node} [\RDD{circle split},draw,fill=green!20]\AC{haut  \BSS{nodepart}\AC{lower} bas };   }\\ 
\hline 
 
\tikz  \node [circle split,draw,blue,fill=green!20] {haut  \nodepart{lower} bas }; 

&  
\tikz  \node [circle solidus,draw,blue,fill=green!20]{haut  \nodepart{lower} bas };
&  
\tikz  \node [ellipse split,draw,blue,fill=green!20]{texte haut  \nodepart{lower} texte bas };
& 
\tikz  \node [rectangle split,draw,blue,fill=green!20]{haut  \nodepart{lower} bas}; 

\\ \hline 
\RDD{circle split} & \RDD{circle solidus} & \RDD{ellipse split} & \RDD{rectangle split} \\ 
\hline 
\end{tabular} 

 \bigskip
 
 \begin{tabular}{|c|c|}  \hline  
 \begin{tikzpicture} [baseline=0pt]
 \node[rectangle split,rectangle split parts=5,draw,blue,fill=green!20] at(0,0)
 {texte 1
 \nodepart{second}
 texte 2
 \nodepart{four}
 texte 3};
 \end{tikzpicture}
&
\parbox[c]{10cm}{
 \BS{node}[rectangle split,\RDD{rectangle split parts}=5,\\
 draw] \\
 \AC{texte 1 \\
 \BSS{nodepart}\AC{second} texte 2 \\
 \BSS{nodepart}\AC{four} texte 3}; \\
 \\
\dft : rectangle split parts=4 }
 \\  \hline 
 \end{tabular} 
 
\bigskip

\begin{tabular}{|c|}\hline  
\BS{node} [rectangle split,rectangle split parts=3,\RDD{rectangle split horizontal},draw,blue] \\
\AC{texte1\BSS{nodepart}\AC{two}texte2\BSS{nodepart}\AC{three}texte3};
\\ \hline  
\tikz \node [rectangle split,rectangle split parts=3, rectangle split horizontal,draw,blue]
{texte 1\nodepart{two}texte 2\nodepart{three}texte 3}; 
\\ \hline 
\end{tabular} 
 
 \bigskip
 
% % % <<<<<<<<<<<<<<<<< A Voir rectangle split allocate boxes= >>>>>>>>>>>>>>>>>>>>>>>>>>>>>>>>

% \begin{tikzpicture} [baseline=0pt]%[every text node part/.style={text centered}]
% \node[rectangle split,draw,rectangle split parts=5,fill=green!20,rectangle split allocate boxes=3] at(0,0)
% {texte 1  \nodepart{second}  texte 2  \nodepart{four}  texte 3};
% \end{tikzpicture}
% 
 
\bigskip
 \begin{tabular}{|c|c|}  \hline  
\begin{tikzpicture}[baseline=0pt]
\node[rectangle split, rectangle split parts=3, draw,blue, text width=2.75cm]
{texte 1
\nodepart{two}
texte 2a \\
texte 2b \\
texte 2c
\nodepart{three}
texte 3a \\
texte 3b};
\end{tikzpicture}
&
\parbox{8cm}{
 \BS{node}[rectangle split,\RDD{rectangle split parts}=5, draw] \\
 \AC{texte 1 \\
 \BSS{nodepart}\AC{second} texte 2a  \BS{}\BS{}texte 2b  \BS{}\BS{}  texte 2c \\
 \BSS{nodepart}\AC{three} texte 3a \BS{}\BS{} texte 3b }; \\
}
 \\  \hline 
 \end{tabular} 
\bigskip


 \begin{tabular}{|c|c|}  \hline  
 \multicolumn{2}{|c|}{  \BS{node}[rectangle split, draw,blue,minimum size = 2cm,\RDD{rectangle split draw splits}= true] } \\
  \multicolumn{2}{|c|}{ 
  \AC{texte 1 \BS{nodepart}\AC{two} texte 2 \BS{nodepart}\AC{three} texte 3 \BS{nodepart}\AC{four} texte 4};   }\\ 
 \hline 
\tikz \node[rectangle split, draw,blue,minimum size = 2cm,rectangle split draw splits= true] {texte 1 \nodepart{two} texte 2 \nodepart{three} texte 3 \nodepart{four} texte 4};
&
\tikz \node[rectangle split, draw,blue,minimum size = 2cm,rectangle split draw splits= false] {texte 1 \nodepart{two} texte 2 \nodepart{three} texte 3 \nodepart{four} texte 4};
 \\ \hline
 \RDD{rectangle split draw splits}= true & \RDD{rectangle split draw splits}= false \\
 \dft &
 \\ \hline 
 \end{tabular}
 
\bigskip

 \begin{tabular}{|c|c|}  \hline  
\multicolumn{2}{|c|}{  
\BS{node} [rectangle split,rectangle split parts=3,draw,\RDD{rectangle split ignore empty parts}=false] }\\
 \multicolumn{2}{|c|}{ \AC{texte 1 \BS{nodepart}\AC{second} \BS{nodepart}\AC{third}texte 3};} 
\\ \hline  
\begin{tikzpicture} 
\node[rectangle split,rectangle split parts=3,draw,blue,rectangle split ignore empty parts=false] {texte 1 \nodepart{second} \nodepart{third}texte 3};
\end{tikzpicture}
&
\begin{tikzpicture}
\node[rectangle split,rectangle split parts=3,draw,blue,rectangle split ignore empty parts] 
{texte 1 \nodepart{second} \nodepart{third}texte 3};
\end{tikzpicture}
 \\  \hline 
\RDD{rectangle split ignore empty parts}=false & \RDD{rectangle split ignore empty parts}=true 
\\ \hline
 \end{tabular}
 
\bigskip

 \begin{tabular}{|c|c|}  \hline  
\multicolumn{2}{|c|}{  
\BS{node} [rectangle split,rectangle split parts=3,draw,\RDD{rectangle split empty part depth}=1cm] }\\
 \multicolumn{2}{|c|}{ \AC{texte 1 \BS{nodepart}\AC{second} \BS{nodepart}\AC{third}texte 3};} 
\\ \hline 
\begin{tikzpicture} 
\node[rectangle split,rectangle split parts=3,draw,blue,rectangle split empty part depth=1cm] {texte 1 \nodepart{second} \nodepart{third}texte 3};
\end{tikzpicture}
&
\begin{tikzpicture} 
\node[rectangle split,rectangle split parts=3,draw,blue,text depth=1cm] {texte 1 \nodepart{second} \nodepart{third}texte 3};
\end{tikzpicture}
\\ \hline 
\RDD{rectangle split empty part depth}=1cm & \RDD{text depth}=1cm
\\ \hline
\dft : 0ex & \dft : 0ex
\\ \hline 
\begin{tikzpicture}
\node[rectangle split,rectangle split parts=3,draw,blue,rectangle split empty part  height=1cm] 
{texte 1 \nodepart{second} \nodepart{third}texte 3};
\end{tikzpicture}
&
\begin{tikzpicture}
\node[rectangle split,rectangle split parts=3,draw,blue,text height=1cm] 
{texte 1 \nodepart{second} \nodepart{third}texte 3};
\end{tikzpicture}
\\  \hline 
\RDD{rectangle split empty part height}=1cm & \RDD{text height}=1cm
\\ \hline
\dft : 1ex & \dft : 1ex
\\ \hline 
 \end{tabular}
 
\bigskip



 \begin{tabular}{|c|c|}  \hline 
 \multicolumn{2}{|c|}{ 
 \BS{node} [rectangle split,rectangle split parts=3,draw,\RDD{rectangle split empty part width}=1cm]   \AC{};  } 
 \\ \hline 
\begin{tikzpicture} 
\node[rectangle split,rectangle split parts=3,draw,blue,rectangle split empty part width=2cm]{};
\end{tikzpicture}

&
\begin{tikzpicture} 
\node[rectangle split,rectangle split parts=3,draw,blue]{}; 
\end{tikzpicture}
\\  \hline 
 \RDD{rectangle split empty part width}=2cm  &  \dft : 1ex
\\ \hline
 \end{tabular} 
 
 \bigskip



% % % % <<<<<<<<<< A voir   /pgf/rectangle split use custom fill= (default true) <<<<<<<<<<<<<<<<<<<<<<<<<<<<
 


 \begin{tabular}{|c|c|}  \hline 
 \tikz[baseline=0pt] \node[rectangle split, draw,blue,minimum size = 2cm,rectangle split part align={center, left,right}] {texte 1 \nodepart{two} texte 2 \nodepart{three} texte 3 \nodepart{four} texte 4};
&
\parbox{8cm}{
\BS{node}[rectangle split, draw,blue,minimum size = 2cm,\\
\RDD{rectangle split part align}=\AC{center, left,right}]\\
 \AC{texte 1 \BS{nodepart}\AC{two} texte 2  \\
 \BS{nodepart}\AC{three} texte 3  \BS{nodepart}\AC{four} texte 4};
}
\\ \hline
 \tikz[baseline=0pt] \node[rectangle split, draw,blue,minimum size = 2cm, rectangle split horizontal,rectangle split part align={center,base, top,bottom}] {texte 1 \nodepart{two} texte 2 \nodepart{three} texte 3 \nodepart{four} texte 4};
 &
 \parbox{8cm}{
 \BS{node}[rectangle split, draw,blue,minimum size = 2cm,\\
 rectangle split horizontal,\\
 \RDD{rectangle split part align}=\AC{center,base, top,bottom}]\\
  \AC{texte 1 \BS{nodepart}\AC{two} texte 2  \\
  \BS{nodepart}\AC{three} texte 3  \BS{nodepart}\AC{four} texte 4};
 }
 \\ \hline
 \end{tabular}
 
\bigskip


 \begin{tabular}{|c|c|}  \hline  
\tikz[baseline=0pt] \node[rectangle split, draw,blue, minimum width=1cm,rectangle split part fill={red, green,cyan}]{};
&
\parbox{12cm}{
\BS{node}[rectangle split, draw,blue, minimum width=1cm,\\
 \RDD{rectangle split part fill}=\AC{red, green,cyan}]\AC{};}
\\ \hline
\end{tabular} 

\newpage

\SbSSCT{Mise en forme du texte}{Text attributes}

\subsubsection{Position}

\begin{center}
\RRR{17-4-3}
\end{center}

\begin{tabular}{|c|c|c|c|} \hline  
\multicolumn{4}{|l|}{ \BS{tikz} \BS{draw} (0,0) node[fill=blue!10,\RDD{text width}=2cm,\RDD{text justified}]   }\\ 

\multicolumn{4}{|l|}{ \AC{Ceci est une démonstration d'un texte  sur une largeur de 2cm};  }\\ 
\hline 
\tikz \draw (0,0) node[fill=blue!10,text width=2cm]
{Ceci est une démonstration d'un texte  sur une largeur de 2cm.};
&  
\tikz \draw (0,0) node[fill=blue!10,text width=2cm,text justified]
{Ceci est une démonstration d'un texte  sur une largeur de 2cm};
&  
\tikz \draw (0,0) node[fill=blue!10,text width=2cm,text centered]
{Ceci est une démonstration d'un texte  sur une largeur de 2cm .};
&  
\tikz \draw (0,0) node[fill=blue!10,text width=2cm,text ragged]
{Ceci est une démonstration d'un texte  sur une largeur de 2cm .};
\\  \hline  
\TFRGB{sans}{without} option & \RDD{text justified} & \RDD{text centered }& \RDD{text ragged}   
\\ \hline  
\tikz \draw (0,0) node[fill=blue!10,text width=2cm,text badly ragged]
{Ceci est une démonstration d'un texte  sur une largeur de 2cm.};
&  
\tikz \draw (0,0) node[fill=blue!10,text width=2cm,text badly centered]
{Ceci est une démonstration d'un texte  sur une largeur de 2cm .};
&
\tikz \draw (0,0) node[fill=blue!10,text width=2cm,align=center]
{Ceci est une démonstration d'un texte  sur une largeur de 2cm .};
&
\tikz \draw (0,0) node[fill=blue!10,text width=2cm,align=flush center]
{Ceci est une démonstration d'un texte  sur une largeur de 2cm .};
\\  \hline 
\RDD{text badly ragged} &  \RDD{text badly centered} &  \RDD{align}=center & \RDD{align}=flush center 
\\  \hline 
\tikz \draw (0,0) node[fill=blue!10,text width=2cm,align=justify]
{Ceci est une démonstration d'un texte  sur une largeur de 2cm .};
&
\tikz \draw (0,0) node[fill=blue!10,text width=2cm,align=flush right]
{Ceci est une démonstration d'un texte  sur une largeur de 2cm .};
&
\tikz \draw (0,0) node[fill=blue!10,text width=2cm,align=right]
{Ceci est une démonstration d'un texte  sur une largeur de 2cm .};
&
\tikz \draw (0,0) node[fill=blue!10,text width=2cm,align=flush left]
{Ceci est une démonstration d'un texte  sur une largeur de 2cm .};
\\ \hline 
\RDD{align}=justify & \RDD{align}=flush right &  \RDD{align}=right & \RDD{align}=flush left
\\ \hline 

\end{tabular} 
\bigskip

\begin{tabular}{|c|c|} \hline 
\tikz[baseline=0cm] \node [draw] {
\begin{tabular}{|c|c|} \hline
AAA & BBB \\ \hline
CCC & DDD \\ \hline
\end{tabular}
};
& 
\parbox{8cm}{
\BS{tikz} \BS{node} [draw] \AC{
\BS{begin}\AC{tabular}\AC{|c|c|} \BS{hline} \\
AAA \& BBB \BS{}\BS{} \BS{hline} \\
CCC \& DDD \BS{}\BS{} \BS{hline} \\
\BS{end}\AC{tabular}
};}
\\ \hline 
\end{tabular} 

\bigskip


\begin{tabular}{|c|c|c|}  \hline 
\multicolumn{3}{|c|}{\BS{tikz}[align=left] \BS{node}[draw] \AC{AAA \rouge{ \BS{}\BS{} } BBBBBBBB \rouge{ \BS{}\BS{} } CC};} \\ \hline
\tikz[align=left] \node[draw] {AAA\\BBBBBBBB\\CC};
&  
\tikz[align=center] \node[draw] {AAA\\BBBBBBBB\\CC};
&
\tikz[align=right] \node[draw] {AAA\\BBBBBBBB\\CC};
\\ \hline
[align=left]  & [align=center] &[align=right] 
\\ \hline
\end{tabular} 


\bigskip

\begin{tabular}{|c|c|} \hline 
\multicolumn{2}{|c|}{\BS{tikz}[align=left] \BS{node}[draw] \AC{AAA  \BS{}\BS{} \rouge{[1cm] } BBBBBBBB };} 
\\ \hline 
\rule[-1cm]{0pt}{1,5cm} \tikz[align=left] \node[draw] {AAA\\[1cm]BBBBBBBB\\}; 
& 
\tikz[align=left] \node[draw] {AAA\\[-1cm]BBBBBBBB\\}; 
\\ \hline 
\rouge{ [1cm] } & \rouge{[ -1cm] }
\\ \hline 
\end{tabular} 

\SbSbSSCT{Couleur et fontes }{Colors and Fonts}

\begin{tabular}{|c|c|c|c|c|c|} \hline  
\tikz \draw (0,0) node[text= red]{Texte.};
&
\tikz \draw (0,0) node[font=\itshape]{Texte.};
&
\tikz \draw (0,0) node[font=\slshape]{Texte.};
&
\tikz \draw (0,0) node[font=\scshape]{Texte.};
&
\tikz \draw (0,0) node[font=\upshape]{Texte.};
&
\tikz \draw (0,0) node[font=\bfseries]{Texte.};
\\ \hline 



[text= red] & [font=\BS{itshape}]  & [font=\BS{slshape}] & [font=\BS{scshape}] & [font=\BS{upshape}] & [font=\BS{bfseries}]
\\ \hline 
\end{tabular} 



\bigskip
 
\SbSbSSCT{Taille des fontes}{Font Sizes}

\begin{tabular}{|c|c|c|c|c|c|c|}\hline
\multicolumn{7}{|c|}{ \BS{tikz} \BS{draw} (0,0) node[\RDD{font}=\BS{tiny}]\AC{Texte.}   }
\\  \hline
\tikz \draw (0,0) node[font=\tiny]{Texte.};
&
\tikz \draw (0,0) node[font=\footnotesize]{Texte.};
&
\tikz \draw (0,0) node[font=\small]{Texte.};
&
\tikz \draw (0,0) node[font=\large]{Texte.};
&
\tikz \draw (0,0) node[font=\Large]{Texte.};
&
\tikz \draw (0,0) node[font=\huge]{Texte.};
&
\tikz \draw (0,0) node[font=\Huge]{Texte.};
\\ \hline \BS{tiny} & \BS{footnotesize}  & \BS{small} & \BS{large} & \BS{Large} & \BS{huge} & \BS{Huge} \\ 
\hline 
\end{tabular} 

\bigskip
\begin{center}
\RRR{17-4-4}
\end{center}

\begin{tabular}{|c|c|c|} \hline  
\tikz \draw (0,0) node[fill=blue!10,text height=1cm,draw]{Texte.};
&  
\tikz \draw (0,0) node[fill=blue!10,text depth=1cm,draw]{Texte.};
&  
\tikz \draw (0,0) node[fill=blue!10,text depth=0.5cm,,text height=.5cm,draw]{Texte.};
\\ \hline  
\RDD{text height}=1cm
&  
\RDD{text depth}=1cm
&
\RDD{text height}=0.5cm, \RDD{text depth}=0.5cm
\\ \hline 
\end{tabular} 

\newpage

\SbSSCT{Positions prédéfinies  sur un n\oe ud}{Positions on a node}
\label{nomnoeud}

\SbSbSSCT{pour l'ensemble des n\oe uds}{For all types of node}
\begin{center}
\RRR{17-5-1}
\end{center}

\begin{tabular}{|c|c|c|c|} \hline  
\begin{tikzpicture}
\node[rectangle,draw,minimum size=3cm] (A) at (1,1) {\Huge texte};
\fill[red] (node cs:name=A,anchor=north west) circle (3pt);
\end{tikzpicture}
&
\begin{tikzpicture}
\node[rectangle,draw,minimum size=3cm] (A) at (1,1) {\Huge texte};
\fill[red] (node cs:name=A,anchor=north) circle (3pt);
\end{tikzpicture}
&
\begin{tikzpicture}
\node[rectangle,draw,minimum size=3cm] (A) at (1,1) {\Huge texte};
\fill[red] (node cs:name=A,anchor=north east) circle (3pt);
\end{tikzpicture}
&
\begin{tikzpicture}
\node[rectangle,draw,minimum size=3cm] (A) at (1,1) {\Huge texte};
\fill[red] (node cs:name=A,anchor=text) circle (3pt);
\end{tikzpicture}
\\ \hline 
north west & north & north east & text
\\ \hline 

\begin{tikzpicture}
\node[rectangle,draw,minimum size=3cm] (A) at (1,1) {\Huge texte};
\fill[red] (node cs:name=A,anchor= west) circle (3pt);
\end{tikzpicture}
&
\begin{tikzpicture}
\node[rectangle,draw,minimum size=3cm] (A) at (1,1) {\Huge texte};
\fill[red] (node cs:name=A,anchor=mid  west) circle (3pt);
\end{tikzpicture}
&
\begin{tikzpicture}
\node[rectangle,draw,minimum size=3cm] (A) at (1,1) {\Huge texte};
\fill[red] (node cs:name=A,anchor= base west) circle (3pt);
\end{tikzpicture}
&
\begin{tikzpicture}
\node[rectangle,draw,minimum size=3cm] (A) at (1,1) {\Huge texte};
\fill[red] (node cs:name=A,anchor= base) circle (3pt);
\end{tikzpicture}
\\ \hline 
west & mid west & base west &  base
\\ \hline
 
\begin{tikzpicture}
\node[rectangle,draw,minimum size=3cm] (A) at (1,1) {\Huge texte};
\fill[red] (node cs:name=A,anchor=east) circle (3pt);
\end{tikzpicture}
&
\begin{tikzpicture}
\node[rectangle,draw,minimum size=3cm] (A) at (1,1) {\Huge texte};
\fill[red] (node cs:name=A,anchor=mid east) circle (3pt);
\end{tikzpicture}
&
\begin{tikzpicture}
\node[rectangle,draw,minimum size=3cm] (A) at (1,1) {\Huge texte};
\fill[red] (node cs:name=A,anchor=base east) circle (3pt);
\end{tikzpicture}
&
\begin{tikzpicture}
\node[rectangle,draw,minimum size=3cm] (A) at (1,1) {\Huge texte};
\fill[red] (node cs:name=A,anchor= mid) circle (3pt);
\end{tikzpicture}
\\ \hline 
east & mid esat & base east & mid
\\ \hline 

\begin{tikzpicture}
\node[rectangle,draw,minimum size=3cm] (A) at (1,1) {\Huge texte};
\fill[red] (node cs:name=A,anchor= south east) circle (3pt);
\end{tikzpicture}
&
\begin{tikzpicture}
\node[rectangle,draw,minimum size=3cm] (A) at (1,1) {\Huge texte};
\fill[red] (node cs:name=A,anchor= south) circle (3pt);
\end{tikzpicture}
&
\begin{tikzpicture}                                       
\node[rectangle,draw,minimum size=3cm] (A) at (1,1) {\Huge texte};
\fill[red] (node cs:name=A,anchor= south west) circle (3pt);
\end{tikzpicture}
&
\begin{tikzpicture}
\node[rectangle,draw,minimum size=3cm] (A) at (1,1) {\Huge texte};
\fill[red] (node cs:name=A,anchor=center ) circle (3pt);
\end{tikzpicture}
\\ \hline 
south east & south & south west & center
\\ \hline
 
\begin{tikzpicture}
\node[rectangle,draw,minimum size=3cm] (A) at (1,1) {\Huge texte};
\fill[red] (node cs:name=A,anchor=0) circle (3pt);
\end{tikzpicture}
&
\begin{tikzpicture}
\node[rectangle,draw,minimum size=3cm] (A) at (1,1) {\Huge texte};
\fill[red] (node cs:name=A,anchor=120) circle (3pt);
\end{tikzpicture}
&
\begin{tikzpicture}
\node[rectangle,draw,minimum size=3cm] (A) at (1,1) {\Huge texte};
\fill[red] (node cs:name=A,anchor=-60) circle (3pt);
\end{tikzpicture}
&


\\ \hline 
0 & 120 & -60 &  
\\ \hline 
\end{tabular}
 
\newpage 

\SbSbSSCT{spécifique à un n\oe ud}{Specific to a node}

\TFRGB{Consultez }{see} \RRR{67 }


\begin{tabular}{|c|c|} \hline 
shape=circle & shape=diamond
\\  \hline 
\begin{tikzpicture}[]
\node[circle,draw,minimum size=3.5cm] (A) at (1,1) {\Huge XXX};
\foreach \anchor/\placement in
{north west/above left, north/above, north east/above right,
west/above left, center/above, east/right,
mid west/left, mid/below right, mid east/right,
base west/below left, base/below, base east/below right,
south west/below left, south/below, south east/below right,
text/below, 20/right, 120/above}
\fill[blue,pin position=\placement] (node cs:name=A,anchor= \anchor) circle (2pt) node[blue,pin=\scriptsize{ \anchor} ] {} ;
\end{tikzpicture}
&
\begin{tikzpicture}[]
\node[diamond,draw,minimum size=3.5cm] (A) at (1,1) {\Huge XXX};
\foreach \anchor/\placement in
{north west/above left, north/above, north east/above right,
west/left, center/above, east/right,
mid/10,
base/below,
south west/below left, south/below, south east/below right,
text/left, 10/right, 120/above}
\fill[blue,pin position=\placement] (node cs:name=A,anchor= \anchor) circle (2pt) node[blue,pin=\scriptsize{ \anchor} ] {} ;
\end{tikzpicture}
\\ \hline 
\end{tabular} 

\bigskip

\begin{tabular}{|c|} \hline 
shape=ellipse
\\  \hline 
\begin{tikzpicture}[]
\node[ellipse,draw,minimum size=3.5cm] (A) at (1,1) {\Huge XXXXXXX};
\foreach \anchor/\placement in
{north west/above left, north/above, north east/above right, west/left, center/above, east/right,
mid west/left, mid/-75, mid east/right,
base west/200, base/-105, base east/-20,
south west/below left, south/below, south east/below right,
text/-75, 10/right, 130/above}
\fill[blue,pin position=\placement] (node cs:name=A,anchor= \anchor) circle (2pt) node[blue,pin=\scriptsize{ \anchor} ] {} ;
\end{tikzpicture}
\\ \hline 
\end{tabular}

\bigskip

\begin{tabular}{|c|} \hline 
shape=trapezium
\\  \hline 
\begin{tikzpicture}[]
\node[ trapezium,draw,minimum size=3cm] (A) at (1,1) {\Huge XXX};
\foreach \anchor/\placement in
{center/120, text/below, mid/-45, base/below, mid west/left, base west/-175, mid east/right, base east/-25,
west/175, east/above, north/-75, south/-60,
north west/above, north east/above,
south west/-150, south east/-30, 150/above}
\fill[blue,pin position=\placement] (node cs:name=A,anchor= \anchor) circle (2pt) node[blue,pin=\scriptsize{ \anchor} ] {} ;

\foreach \anchor/\placement in
{bottom left corner/below, top right corner/right,
top left corner/left, bottom right corner/below,
bottom side/-120, left side/left, right side/right, top side/above}
\fill[red,pin position=\placement] (node cs:name=A,anchor= \anchor) circle (2pt) node[blue,pin=\scriptsize{ \anchor} ] {} ;
\end{tikzpicture}
\\ \hline 
\end{tabular}

\bigskip

\begin{tabular}{|c|} \hline 
shape=semicircle,shape border rotate=0
\\  \hline 
\begin{tikzpicture}[]
\node[ semicircle,shape border rotate=0,draw,minimum size=3cm] (A) at (1,1) {\Huge XXX};
\foreach \anchor/\placement in
{center/above, base/-160, mid/-40, text/left, base west/-120, base east/-60, mid west/left, mid east/right, north/below, south/-75, east/60, west/120, north west/above left, north east/above right, south west/-140, south east/-60, 30/right}
\fill[blue,pin position=\placement] (node cs:name=A,anchor= \anchor) circle (2pt) node[blue,pin=\scriptsize{ \anchor} ] {} ;
\foreach \anchor/\placement in
{apex/above, arc start/-60, arc end/-120, chord center/-100}
\fill[red,pin position=\placement] (node cs:name=A,anchor= \anchor) circle (2pt) node[blue,pin=\scriptsize{ \anchor} ] {} ;
\end{tikzpicture}
\\ \hline 
\end{tabular}

\bigskip

\begin{tabular}{|c|} \hline 
shape=regular polygon
\\  \hline 
\begin{tikzpicture}[]
\node[ regular polygon,draw,minimum size=3cm] (A) at (1,1) {\Huge XXX};
\foreach \anchor/\placement in
{ center/97, text/97 , mid/-30, base/below, 75/above,
west/left, east/right, north/-87, south/-60,
north east/right, south east/right, north west/left, south west/left}
\fill[blue,pin position=\placement] (node cs:name=A,anchor= \anchor) circle (2pt) node[blue,pin=\scriptsize{ \anchor} ] {} ;

\foreach \anchor/\placement in
{corner 1/above, corner 2/left, corner 3/left, corner 4/right, corner 5/right,
side 1/above, side 2/left, side 3/-120, side 4/right, side 5/above}
\fill[red,pin position=\placement] (node cs:name=A,anchor= \anchor) circle (2pt) node[blue,pin=\scriptsize{ \anchor} ] {} ;
\end{tikzpicture}
\\ \hline 
\end{tabular}


\bigskip

\begin{tabular}{|c|} \hline 
shape=star
\\  \hline 
\begin{tikzpicture}[]
\node[  shape=star, star points=5, star point ratio=1.65,draw,minimum size=3cm] (A) at (1,1) {\Huge XXX};
\foreach \anchor/\placement in
{center/above, 
text/below, 
mid/-30, 
base/-80, 
75/above,
west/left, 
east/right, 
north/below, 
south/94,
north east/right, 
south east/right, 
north west/left, 
south west/left}
\fill[blue,pin position=\placement] (node cs:name=A,anchor= \anchor) circle (2pt) node[blue,pin=\scriptsize{ \anchor} ] {} ;

\foreach \anchor/\placement in
{inner point 1/above left, 
inner point 2/left, 
inner point 3/below, 
inner point 4/right,
inner point 5/above right, 
outer point 1/above, 
outer point 2/left, 
outer point 3/left,
outer point 4/right, 
outer point 5/right}
\fill[red,pin position=\placement] (node cs:name=A,anchor= \anchor) circle (2pt) node[blue,pin=\scriptsize{ \anchor} ] {} ;
\end{tikzpicture}
\\ \hline 
\end{tabular}



\bigskip

\begin{tabular}{|c|c|} \hline 
shape= isosceles triangle & shape= kite
\\  \hline 
\begin{tikzpicture}[]
\node[ shape=isosceles triangle,draw,minimum size=3cm] (A) at (1,1) {\Huge XXX};
\foreach \anchor/\placement in
{center/above,text/above,150/left,mid/-5, mid west/left, mid east/right,base/-120, base west/-150, 
base east/below right ,west/left, east/right,  north/above, north west/left, north east/above right,
south /-120 , south east/below right}
\fill[blue,pin position=\placement] (node cs:name=A,anchor= \anchor) circle (2pt) node[blue,pin=\scriptsize{ \anchor} ] {} ;
\foreach \anchor/\placement in
{apex/above, left corner/left, right corner/left,left side/above, right side/below, lower side/160}
\fill[red,pin position=\placement] (node cs:name=A,anchor= \anchor) circle (2pt) node[blue,pin=\scriptsize{ \anchor} ] {} ;
\end{tikzpicture}
&
\begin{tikzpicture}[]
\node[ shape=kite,draw,minimum size=3cm] (A) at (1,1) {\Huge XXX};
\foreach \anchor/\placement in
{center/above, text/85, mid/-85, base/-95,mid west/left, base west/-160, 
mid east/right, base east/-20,west/left, east/right, north/80, south/below left,north west/above left, north east/above right,south west/left, south east/right, 
110/above left}
\fill[blue,pin position=\placement] (node cs:name=A,anchor= \anchor) circle (2pt) node[blue,pin=\scriptsize{ \anchor} ] {} ;
\foreach \anchor/\placement in
{upper vertex/110, 
left vertex/left, 
lower vertex/below right,
right vertex/right, 
upper left side/left, 
upper right side/right,
lower left side/left, 
lower right side/below right}
\fill[red,pin position=\placement] (node cs:name=A,anchor= \anchor) circle (2pt) node[blue,pin=\scriptsize{ \anchor} ] {} ;
\end{tikzpicture}
\\ \hline 
\end{tabular}


\bigskip

\begin{tabular}{|c|c|} \hline 
shape= dart & shape= circular sector
\\  \hline 
\begin{tikzpicture}[]
\node[shape=dart, shape border rotate=90,,draw,minimum size=3cm] (A) at (1,1) {\Huge XXX};
\foreach \anchor/\placement in
{west/left  , east/above right , north/below,south/left,
north west/left, north east/right, south west/below, south east/below,110/above left}
\fill[blue,pin position=\placement] (node cs:name=A,anchor= \anchor) circle (2pt) node[blue,pin=\scriptsize{ \anchor} ] {} ;
\foreach \anchor/\placement in
{tip/above, tail center/right, right tail/below,
left tail/below, right tail/below, left side/above left, right side/above right}
\fill[red,pin position=\placement] (node cs:name=A,anchor= \anchor) circle (2pt) node[blue,pin=\scriptsize{ \anchor} ] {} ;
\end{tikzpicture}
&
\begin{tikzpicture}[]
\node[shape=circular sector,draw,minimum size=3cm] (A) at (1,1) {\Huge XXX};
\foreach \anchor/\placement in
{west/170  , east/right , north/above , south/below, north west/left, north east/above, south west/left, south east/below, 120/left}
\fill[blue,pin position=\placement] (node cs:name=A,anchor= \anchor) circle (2pt) node[blue,pin=\scriptsize{ \anchor} ] {} ;
\foreach \anchor/\placement in
{sector center/above, arc start/above, arc end/below, arc center/190}
\fill[red,pin position=\placement] (node cs:name=A,anchor= \anchor) circle (2pt) node[blue,pin=\scriptsize{ \anchor} ] {} ;
\end{tikzpicture}
\\ \hline 
\end{tabular}



\bigskip

\begin{tabular}{|c|c|} \hline 
shape=cylinder & shape=cloud
\\  \hline 
\begin{tikzpicture}[]
\node[shape=cylinder,draw,minimum size=3cm] (A) at (1,1) {\Huge XXX};
\foreach \anchor/\placement in
{west/170  , east/-10 , north/above , south/below, north west/left, north east/above, south west/left, south east/below, 120/left}
\fill[blue,pin position=\placement] (node cs:name=A,anchor= \anchor) circle (2pt) node[blue,pin=\scriptsize{ \anchor} ] {} ;
\foreach \anchor/\placement in
{before top/10 , top/10, after top/below right, before bottom/below left, bottom/190, after bottom/above left}
\fill[red,pin position=\placement] (node cs:name=A,anchor= \anchor) circle (2pt) node[blue,pin=\scriptsize{ \anchor} ] {} ;
\end{tikzpicture}
&
\begin{tikzpicture}[]
\node[shape=cloud,draw,minimum size=3cm] (A) at (1,1) {\Huge XXX};
\foreach \anchor/\placement in
{west/west  , east/east , north/below , south/below left, north west/left, north east/above right, south west/left, south east/right, 110/above}
\fill[blue,pin position=\placement] (node cs:name=A,anchor= \anchor) circle (2pt) node[blue,pin=\scriptsize{ \anchor} ] {} ;
\foreach \anchor/\placement in
{puff 1/above, puff 2/above left , puff 3/left, puff 4/left,
puff 5/below left, puff 6/below right, puff 7/below right, puff 8/right,
puff 9/right, puff 10/above}
\fill[red,pin position=\placement] (node cs:name=A,anchor= \anchor) circle (2pt) node[blue,pin=\scriptsize{ \anchor} ] {} ;
\end{tikzpicture}

\\ \hline 
\end{tabular}

\bigskip

\begin{tabular}{|c|} \hline 
shape=starburst
\\  \hline 
\begin{tikzpicture}[]
\node[shape=starburst, starburst points=9, starburst point height=2cm,draw,minimum size=3cm] (A) at (1,1) {\Huge XXX};
\foreach \anchor/\placement in
{west/west  , east/east , north/70 , south/above, north west/below , north east/below, south west/below left, south east/-85, 30/above right}
\fill[blue,pin position=\placement] (node cs:name=A,anchor= \anchor) circle (2pt) node[blue,pin=\scriptsize{ \anchor} ] {} ;
\foreach \anchor/\placement in
{outer point 1/105, outer point 2/above left , 
outer point 3/left, outer point 4/left, 
outer point 5/below, outer point 6/below, 
outer point 7/below, outer point 8/right, 
outer point 9/above,
inner point 1/93, inner point 2/160, 
inner point 3/190, inner point 4/below left, 
inner point 5/below, inner point 6/-85,
inner point 7/-30, inner point 8/above right, 
inner point 9/above}
\fill[red,pin position=\placement] (node cs:name=A,anchor= \anchor) circle (2pt) node[blue,pin=\scriptsize{ \anchor} ] {} ;
\end{tikzpicture}
\\ \hline 
\end{tabular}


\bigskip

\begin{tabular}{|c|} \hline 
shape=signal
\\  \hline 
\begin{tikzpicture}[]
\node[signal,signal from=west,draw,minimum size=3.5cm] (A) at (1,1) {\Huge XXX};
\foreach \anchor/\placement in
{north west/above left, north/above, north east/above right,
west/left, center/above, east/right,
mid west/left, mid/below right, mid east/right,
base west/-160, base/below, base east/below right,
south west/below left, south/below, south east/below right,
text/below, 20/right, 120/above}
\fill[blue,pin position=\placement] (node cs:name=A,anchor= \anchor) circle (2pt) node[blue,pin=\scriptsize{ \anchor} ] {} ;
\end{tikzpicture}
\\ \hline 
\end{tabular}


\bigskip

\begin{tabular}{|c|} \hline 
shape=tape
\\  \hline 
\begin{tikzpicture}[]
\node[tape, tape bend height=1cm,draw,minimum size=3.5cm] (A) at (1,1) {\Huge XXX};
\foreach \anchor/\placement in
{north west/above left, north/above, north east/above right,
west/left, center/above, east/right,
mid west/left, mid/below right, mid east/right,
base west/-160, base/110, base east/below right,
south west/below left, south/below, south east/below right,
text/110, 20/right, 120/above}
\fill[blue,pin position=\placement] (node cs:name=A,anchor= \anchor) circle (2pt) node[blue,pin=\scriptsize{ \anchor} ] {} ;
\end{tikzpicture}
\\ \hline 
\end{tabular}

\begin{tabular}{|c|} \hline 
shape=magnetic tape
\\  \hline 
\begin{tikzpicture}[]
\node[shape=magnetic tape,draw,minimum size=3cm] (A) at (1,1) {\Huge XXX};
\foreach \anchor/\placement in
{west/west  , east/east , north/above , south/below, north west/above left , north east/above right, south west/left, south east/below, 30/above right}
\fill[blue,pin position=\placement] (node cs:name=A,anchor= \anchor) circle (2pt) node[blue,pin=\scriptsize{ \anchor} ] {} ;
\foreach \anchor/\placement in
{tail east/right, tail south east/below right, tail north east/above right}
\fill[red,pin position=\placement] (node cs:name=A,anchor= \anchor) circle (2pt) node[blue,pin=\scriptsize{ \anchor} ] {} ;
\end{tikzpicture}
\\ \hline 
\end{tabular}



\bigskip

\begin{tabular}{|c|} \hline 
shape=single arrow
\\  \hline 
\begin{tikzpicture}[]
\node[shape=single arrow,draw,minimum size=3cm] (A) at (1,1) {\Huge XXXXXX};
\foreach \anchor/\placement in
{west/170  , east/below right , north/above , south/below, north west/above left, north east/above right, south west/below left, south east/below right, 30/east}
\fill[blue,pin position=\placement] (node cs:name=A,anchor= \anchor) circle (2pt) node[blue,pin=\scriptsize{ \anchor} ] {} ;
\foreach \anchor/\placement in
{tip/above right, before tip/above, after tip/below, before head/190 , after head/170, after tail/left, before tail/left, tail/190}
\fill[red,pin position=\placement] (node cs:name=A,anchor= \anchor) circle (2pt) node[blue,pin=\scriptsize{ \anchor} ] {} ;
\end{tikzpicture}
\\ \hline 
\end{tabular}


\bigskip

\begin{tabular}{|c|} \hline 
shape=double arrow
\\  \hline 
\begin{tikzpicture}[]
\node[shape=double arrow, double arrow head extend=1.5cm,,draw,minimum size=3cm] (A) at (1,1) {\Huge XXXXXXXXX};
\foreach \anchor/\placement in
{west/170  , east/-10 , north/above , south/below, north west/above left, north east/above right, south west/below left, south east/below right, 35/above right}
\fill[blue,pin position=\placement] (node cs:name=A,anchor= \anchor) circle (2pt) node[blue,pin=\scriptsize{ \anchor} ] {} ;
\foreach \anchor/\placement in
{before head 1/above right, before tip 1/above, 
tip 1/10, after tip 1/below, 
after head 1/below right, before head 2/below left, 
before tip 2/below left, tip 2/190, 
after tip 2/above left, after head 2/above left}
\fill[red,pin position=\placement] (node cs:name=A,anchor= \anchor) circle (2pt) node[blue,pin=\scriptsize{ \anchor} ] {} ;
\end{tikzpicture}
\\ \hline 
\end{tabular}


\bigskip

\begin{tabular}{|c|} \hline 
shape=arrow box
\\  \hline 
\begin{tikzpicture}[]
\node[shape=arrow box,draw,minimum size=3cm,arrow box arrows={north:2cm from border, south, east:2cm from border, west},arrow box shaft width=1cm,arrow box head extend=0.25cm] (A) at (1,1) {\Huge XXXXXXXXX};
\foreach \anchor/\placement in
{west/right  , east/left , north/below , south/above, north west/left, north east/right, south west/left, south east/right}
\fill[blue,pin position=\placement] (node cs:name=A,anchor= \anchor) circle (2pt) node[blue,pin=\scriptsize{ \anchor} ] {} ;
\foreach \anchor/\placement in
{north arrow tip/above,
south arrow tip/below, 
east arrow tip/right, 
west arrow tip/left,
before north arrow/above left, 
before north arrow head/110, 
before north arrow tip/left,
after north arrow tip/right, 
after north arrow head/70, 
after north arrow/above right,
before south arrow/below right, 
before south arrow head/-70, 
before south arrow tip/right,
after south arrow tip/left, 
after south arrow head/-110, 
after south arrow/below left,
before east arrow/above right, 
before east arrow head/right, 
before east arrow tip/right,
after east arrow tip/right, 
after east arrow head/right, 
after east arrow/below right,
before west arrow/below left, 
before west arrow head/left, 
before west arrow tip/left,
after west arrow tip/west, 
after west arrow head/left, 
after west arrow/above left}
\fill[red,pin position=\placement] (node cs:name=A,anchor= \anchor) circle (2pt) node[blue,pin=\scriptsize{ \anchor} ] {} ;
\end{tikzpicture}
\\ \hline 
\end{tabular}


\bigskip

\begin{tabular}{|c|} \hline 
shape=circle split
\\  \hline 
\begin{tikzpicture}[]
\node[shape=circle split,draw,minimum size=3.5cm](A) at (1,1) {XXX\nodepart{lower}YYY}  ;
\foreach \anchor/\placement in
{north west/above left, north/above, north east/above right,
west/left, center/above, east/right,
mid west/left, mid/below right, mid east/right,
base west/-160, base/110, base east/below right,
south west/below left, south/below, south east/below right,
text/110, 20/right, 120/above}
\fill[blue,pin position=\placement] (node cs:name=A,anchor= \anchor) circle (2pt) node[blue,pin=\scriptsize{ \anchor} ] {} ;
\foreach \anchor/\placement in
{text/left, lower/left}
\fill[red,pin position=\placement] (node cs:name=A,anchor= \anchor) circle (2pt) node[blue,pin=\scriptsize{ \anchor} ] {} ;
\end{tikzpicture}
\\ \hline 
\end{tabular}

\begin{tabular}{|c|} \hline 
shape=circle solidus
\\  \hline 
\begin{tikzpicture}[]
\node[shape=circle solidus,draw,minimum size=3.5cm](A) at (1,1) {XXX\nodepart{lower}YYY}  ;
\foreach \anchor/\placement in
{north west/above left, north/above, north east/above right,
west/left, center/above, east/right,
mid west/left, mid/below right, mid east/right,
base west/-160, base/110, base east/below right,
south west/below left, south/below, south east/below right,
text/110, 20/right, 120/above}
\fill[blue,pin position=\placement] (node cs:name=A,anchor= \anchor) circle (2pt) node[blue,pin=\scriptsize{ \anchor} ] {} ;
\foreach \anchor/\placement in
{text/left, lower/left}
\fill[red,pin position=\placement] (node cs:name=A,anchor= \anchor) circle (2pt) node[blue,pin=\scriptsize{ \anchor} ] {} ;
\end{tikzpicture}
\\ \hline 
\end{tabular}


\bigskip

\begin{tabular}{|c|} \hline 
shape=ellipse split
\\  \hline 
\begin{tikzpicture}[]
\node[shape=ellipse split,draw,minimum size=3.5cm](A) at (1,1) {XXX\nodepart{lower}YYY}  ;
\foreach \anchor/\placement in
{north west/above left, north/above, north east/above right,
west/left, center/above, east/right,
mid west/left, mid/below right, mid east/right,
base west/-160, base/110, base east/below right,
south west/below left, south/below, south east/below right,
text/110, 20/right, 120/above}
\fill[blue,pin position=\placement] (node cs:name=A,anchor= \anchor) circle (2pt) node[blue,pin=\scriptsize{ \anchor} ] {} ;
;
\end{tikzpicture}
\\ \hline 
\end{tabular}

\bigskip

\begin{tabular}{|c|} \hline 
shape=rectangle split
\\  \hline 
\begin{tikzpicture}[]
\node[name=s,shape=rectangle split, rectangle split parts=4,draw,inner ysep=0.75cm](A) at (1,1)
{\nodepart{text}XXXXXXXXXXXXXX\nodepart{two}YYY
\nodepart{three}ZZZ\nodepart{four}four};
\foreach \anchor/\placement in
{north/above, south/below, east/10, west/170,
north west/above, north east/above, south west/below, south east/below,
center/145, 20/right, mid/30, base/-145}
\fill[blue,pin position=\placement] (node cs:name=A,anchor= \anchor) circle (2pt) node[blue,pin=\scriptsize{ \anchor} ] {} ;
\foreach \anchor/\placement in
{text split/10, text split east/0, text split west/180,two split/30, two split east/right, two split west/left,
three split/30, three split east/east, three split west/west,text/-170, text east/east, text west/west,
two/left, two east/east, two west/west,
three/left, three east/east, three west/west,
four/west, four east/east, four west/west
}
\fill[red,pin position=\placement] (node cs:name=A,anchor= \anchor) circle (2pt) node[blue,pin=\scriptsize{ \anchor} ] {} ;
\end{tikzpicture}
\\ \hline 
\end{tabular}

\bigskip

\begin{tabular}{|c|} \hline 
shape=rectangle callout
\\  \hline 
\begin{tikzpicture}[]
\node[shape=rectangle callout, callout relative pointer={(1.5cm,-.5cm)},draw,
callout pointer width=2cm, inner xsep=1cm, inner ysep=.5cm] (A) at (1,1) {\Huge XXXXXXX};
\foreach \anchor/\placement in
{west/west  , east/east , north/above , south/below, north west/west , north east/right, south west/left, south east/right, 25/right}
\fill[blue,pin position=\placement] (node cs:name=A,anchor= \anchor) circle (2pt) node[blue,pin=\scriptsize{ \anchor} ] {} ;
\foreach \anchor/\placement in
{pointer/right}
\fill[red,pin position=\placement] (node cs:name=A,anchor= \anchor) circle (2pt) node[blue,pin=\scriptsize{ \anchor} ] {} ;
\end{tikzpicture}
\\ \hline 
\end{tabular}

\bigskip

\begin{tabular}{|c|} \hline 
shape=ellipse callout
\\  \hline 
\begin{tikzpicture}[]
\node[shape=ellipse callout,draw] (A) at (1,1) {\Huge XXXXXX};
\foreach \anchor/\placement in
{west/west  , east/right , north/above, south/below, north west/above left, north east/above right, south west/below left, south east/below right}
\fill[blue,pin position=\placement] (node cs:name=A,anchor= \anchor) circle (2pt) node[blue,pin=\scriptsize{ \anchor} ] {} ;
\foreach \anchor/\placement in
{pointer/below right}
\fill[red,pin position=\placement] (node cs:name=A,anchor= \anchor) circle (2pt) node[blue,pin=\scriptsize{ \anchor} ] {} ;
\end{tikzpicture}
\\ \hline 
\end{tabular}


\bigskip

\begin{tabular}{|c|} \hline 
shape=cloud callout
\\  \hline 
\begin{tikzpicture}[]
\node[shape=cloud callout,draw,aspect=1.5] (A) at (1,1) {\Huge XXXXXX};
\foreach \anchor/\placement in
{west/west  , east/right , north/below , south/above, north west/above left, north east/above right, south west/below left, south east/below right}
\fill[blue,pin position=\placement] (node cs:name=A,anchor= \anchor) circle (2pt) node[blue,pin=\scriptsize{ \anchor} ] {} ;
\foreach \anchor/\placement in
{puff 1/above, puff 2/above, puff 3/left, puff 4/left,
puff 5/below left, puff 6/below, puff 7/below right, puff 8/right,
puff 9/right, puff 10/above,pointer/below right}
\fill[red,pin position=\placement] (node cs:name=A,anchor= \anchor) circle (2pt) node[blue,pin=\scriptsize{ \anchor} ] {} ;
\end{tikzpicture}
\\ \hline 
\end{tabular}


\bigskip

\begin{tabular}{|c|} \hline 
shape=cross out
\\  \hline 
\begin{tikzpicture}[]
\node[shape=cross out,draw,minimum size=3cm] (A) at (1,1) {\Huge XXXXXXXXXX};
\foreach \anchor/\placement in
{west/west  , east/right , north/above , south/below, north west/above left, north east/above right, south west/below left, south east/below right}
\fill[blue,pin position=\placement] (node cs:name=A,anchor= \anchor) circle (2pt) node[blue,pin=\scriptsize{ \anchor} ] {} ;
\end{tikzpicture}
\\ \hline 
\end{tabular}

\bigskip

\begin{tabular}{|c|} \hline 
shape=rounded rectangle
\\  \hline 
\begin{tikzpicture}[]
\node[shape=rounded rectangle,draw,minimum size=3cm] (A) at (1,1) {\Huge XXXXXXXXXX};
\foreach \anchor/\placement in
{west/west  , east/right , north/above , south/below, north west/above left, north east/above right, south west/below left, south east/below right}
\fill[blue,pin position=\placement] (node cs:name=A,anchor= \anchor) circle (2pt) node[blue,pin=\scriptsize{ \anchor} ] {} ;

\end{tikzpicture}
\\ \hline 
\end{tabular}


\bigskip

\begin{tabular}{|c|} \hline 
shape=chamfered rectangle
\\  \hline 
\begin{tikzpicture}[]
\node[shape=chamfered rectangle,draw,minimum size=3cm, chamfered rectangle sep=.5cm,] (A) at (1,1) {\Huge XXXXXX};
\foreach \anchor/\placement in
{west/west  , east/right , north/above , south/below, north west/above left, north east/above right, south west/below left, south east/below right}
\fill[blue,pin position=\placement] (node cs:name=A,anchor= \anchor) circle (2pt) node[blue,pin=\scriptsize{ \anchor} ] {} ;
\foreach \anchor/\placement in
{before north east/above right, after north east/above right, before south east/below right,after south east/below right, before north west/above left, after north west/above left, before south west/below left,after south west/below left}
\fill[red,pin position=\placement] (node cs:name=A,anchor= \anchor) circle (2pt) node[blue,pin=\scriptsize{ \anchor} ] {} ;
\end{tikzpicture}
\\ \hline 
\end{tabular}

\normalsize


 

%
%%\begin{tikzpicture}
%%[spy using outlines={circle, magnification=4, size=2cm, connect spies}]
%%\draw [help lines] (0,0) grid (3,2);
%%\draw [decoration=Koch curve type 1]
%%decorate { decorate{ decorate{ decorate{ (0,0) -- (2,0) }}}};
%%\spy [red] on (1.6,0.3)
%%in node [left] at (3.5,-1.25);
%%\spy [blue, size=1cm] on (1,1)
%%in node [right] at (0,-1.25);
%%\end{tikzpicture}
%%
%%\begin{tikzpicture}
%%[spy using outlines={circle, magnification=4, size=2cm, connect spies}]
%%\draw [help lines] (0,0) grid (3,2);
%%\shadedraw[shading=Mandelbrot set ] (0,0) rectangle (2,2) ;
%%\spy [red] on (1,0.4)
%%in node [left] at (3.5,-1.25);
%%\spy [blue, size=1cm] on (.5,.8)
%%in node [right] at (0,-1.25);
%%\end{tikzpicture}
%
%
%%
 \maboite{\BS{usetikzlibrary}\AC{turtle}}
\label{lib-turtle}


\begin{center}
\RRR{ 73 }
\end{center}

\begin{tabular}{|c|c|c|c|} \hline 
\multicolumn{4}{|c|}{  \BS{draw} [blue,line width=3pt,turtle={home,forward}];} \\  \hline 
\begin{tikzpicture}
\draw[help lines] (-1.5,-2) grid (1.5,2) ; 
\draw [blue,line width=3pt,turtle={home,forward}];
\end{tikzpicture}
&  
\begin{tikzpicture}
\draw[help lines] (-1.5,-2) grid (1.5,2) ;  
\draw [blue,line width=3pt,turtle={home,forward=1.5cm}];
\end{tikzpicture}
&  
\begin{tikzpicture}
\draw[help lines] (-1.5,-2) grid (1.5,2) ;  
\draw [blue,line width=3pt,turtle={home,fd}];
\end{tikzpicture}
&  
\begin{tikzpicture}
\draw[help lines] (-1.5,-2) grid (1.5,2) ; 
\draw [blue,line width=3pt,turtle={home,fd=1.5cm}];
\end{tikzpicture}
\\ \hline 
turtle=\AC{home,forward}  & turtle=\AC{home,forward=1.5cm} & turtle=\AC{home,fd} & 
turtle=\AC{home,fd=1.5cm} \\ 
\hline 
\end{tabular} 

\bigskip


\begin{tabular}{|c|c|c|c|}
\hline 
\multicolumn{4}{|c|}{  \BS{draw} [blue,line width=3pt,turtle={home,left,fd];}} \\  \hline  
\hline 
\begin{tikzpicture}
\draw (-1,-1) grid (1,1) ; 
\draw [blue,line width=3pt,turtle={home,left,fd}];
\end{tikzpicture} 
&  
\begin{tikzpicture}
\draw (-1,-1) grid (1,1) ; 
\draw [blue,line width=3pt,turtle={home,left=45,fd}];
\end{tikzpicture}
&  
\begin{tikzpicture}
\draw (-1,-1) grid (1,1) ; 
\draw [blue,line width=3pt,turtle={home,lt,fd}];
\end{tikzpicture}
&  
\begin{tikzpicture}
\draw (-1,-1) grid (1,1) ; 
\draw [blue,line width=3pt,turtle={home,lt=45,fd}];
\end{tikzpicture}
\\ \hline
turtle=\AC{home,left,fd}  & turtle=\AC{home,left=45,fd} & turtle=\AC{home,lt,fd} & 
turtle=\AC{home,lt=45,fd} \\ 
\hline 
\end{tabular} 

\bigskip

\begin{tabular}{|c|c|c|c|}
\hline 
\multicolumn{4}{|c|}{  \BS{draw} [blue,line width=3pt,turtle={home,right,fd];}} \\  \hline  
\hline 
\begin{tikzpicture}
\draw (-1,-1) grid (1,1) ; 
\draw [blue,line width=3pt,turtle={home,right,fd}];
\end{tikzpicture} 
&  
\begin{tikzpicture}
\draw (-1,-1) grid (1,1) ; 
\draw [blue,line width=3pt,turtle={home,right=45,fd}];
\end{tikzpicture}
&  
\begin{tikzpicture}
\draw (-1,-1) grid (1,1) ; 
\draw [blue,line width=3pt,turtle={home,rt,fd}];
\end{tikzpicture}
&  
\begin{tikzpicture}
\draw (-1,-1) grid (1,1) ; 
\draw [blue,line width=3pt,turtle={home,rt=45,fd}];
\end{tikzpicture}
\\ \hline
turtle=\AC{home,right,fd}  & turtle=\AC{home,right=45,fd} & turtle=\AC{home,rt,fd} & 
turtle=\AC{home,rt=45,fd} \\ 
\hline 
\end{tabular} 

\bigskip

\begin{tabular}{|c|c|} \hline 
\tikz[blue,line width=3pt]
\draw [->,turtle={home,rt,fd,fd,lt,fd,lt,fd}];
&  
\tikz[blue,line width=3pt]
\draw [->,turtle/distance=2cm,turtle={home,rt,fd,fd,lt,fd,lt,fd}];
\\ \hline 
[->,turtle={home,rt,fd,fd,lt,fd,lt,fd}] & [->,turtle/distance=2cm,turtle={home,rt,fd,fd,lt,fd,lt,fd}] 
\\ \hline 
\end{tabular} 

\bigskip


\begin{tabular}{|c|} \hline 
\begin{tikzpicture}[turtle/distance=2cm]
\draw[help lines] (-1.5,-1) grid (6,3) ; 
\draw [blue,line width=3pt,dotted,turtle={home,forward,right,forward},fd];
\draw [red,line width=3pt,turtle={how/.style={bend left},home,fd,rt,fd,fd}] ;
\end{tikzpicture}
\\  \hline 
[red,turtle=\AC{\rouge{how/.style}=\AC{bend left},home,fd,rt,fd,fd}]
\\ \hline 
\end{tabular} 

\bigskip

\begin{tabular}{|c|c|}  \hline 
\tikz
\filldraw [turtle/distance=2cm,thick,blue,fill=red!20]
[turtle=home]
\foreach \i in {1,...,5}
{
[turtle={forward,right=144}]
}; 
& 
 
\parbox[b]{10cm}{
\BS{filldraw}[turtle/distance=2cm,thick,blue,fill=red!20] \\
$[$ turtle=home $]$ \\
\BS{foreach} \BS{i} in \AC{1,...,5} \\
{
[ turtle=\AC{forward,right=144} ]
};
}
\\ \hline  
\end{tabular} 

\bigskip



\begin{tabular}{|c|c|}  \hline 
\tikz \draw [thick,blue]
[turtle=home]
\foreach \i in {1,...,25}
{
[turtle={forward=\i/5,right=120}]
};
& 
 
\parbox[b]{10cm}{
\BS{draw}[thick,blue] \\
$[$ turtle=home $]$ \\
\BS{foreach} \BS{i} in \AC{1,...,25} \\
{
[turtle=\AC{forward=\BS{i}/5,right=120} ]
}; \\
\vspace{1cm}
}
\\ \hline  
\end{tabular}




%% 
%%
%% 
%%\newpage 
%%
%%\SbSSCT{Matrice de n\oe uds}{Matrices and Alignment}
%%
%%
%%
\label{matrix}
\begin{center}
\RRR{20}
\end{center}

\begin{tabular}{|c|c|} \hline  
\begin{tikzpicture}[baseline=1cm]
\draw[help lines] (0,0) grid (4,2);
\node [matrix,fill=red!20,draw=blue,very thick] (my matrix) at (2,1)
{
\draw (0,0) circle (4mm); & \node[rotate=45] {Hello}; \\
\draw (0.2,0) circle (2mm); & \fill[red] (0,0) circle (3mm); \\
};
\end{tikzpicture}
& 
\parbox{10cm}{
\BS{node} [\RDD{matrix},fill=red!10,draw=blue,very thick] at (2,1) \\
\{ \\
\BS{draw} (0,0) circle (4mm); \& \BS{node} [rotate=45] {Hello}; \BS{}\BS{} \\
\BS{draw}  (0.2,0) circle (2mm); \& \BS{fill}[red] (0,0) circle (3mm); \BS{}\BS{} \\
\}; \\
}
\\ \hline 
\end{tabular} 

\bigskip

\begin{tabular}{|c|c|} \hline  
\begin{tikzpicture}[baseline=0pt]
\matrix [fill=red!20,draw=blue,very thick] 
{
\draw (0,0) circle (4mm); & \node[rotate=45] {Hello}; \\
\draw (0.2,0) circle (2mm); & \fill[red] (0,0) circle (3mm); \\
};
\end{tikzpicture}
&  
\parbox{10cm}{
\BSS{matrix} [fill=red!10,draw=blue,very thick] \\
\{ \\
\BS{draw} (0,0) circle (4mm); \& \BS{node} [rotate=45] {Hello}; \BS{}\BS{} \\
\BS{draw}  (0.2,0) circle (2mm); \& \BS{fill}[red] (0,0) circle (3mm); \BS{}\BS{} \\
\}; \\
}
\\ \hline 
\end{tabular} 


\SbSbSSCT{Alignement des cellules}{Cell Pictures}


\begin{center}
\RRR{20-3}
\end{center}

\begin{tabular}{|c|c|c|} \hline  
\begin{tikzpicture}
[every node/.style={draw=black,font=\huge}]
\matrix [draw=red]
{
\node {a}; \fill[blue] (0,0) circle (2pt); &
\node {X}; \fill[blue] (0,0) circle (2pt); &
\node {g}; \fill[blue] (0,0) circle (2pt); \\
};
\end{tikzpicture}
&  
\begin{tikzpicture}
[every node/.style={draw=black,anchor=base,font=\huge}]
\matrix [draw=red]
{
\node {a}; \fill[blue] (0,0) circle (2pt); &
\node {X}; \fill[blue] (0,0) circle (2pt); &
\node {g}; \fill[blue] (0,0) circle (2pt); \\
};
\end{tikzpicture}
&  
\begin{tikzpicture}[every node/.style={draw=black}]
\matrix [draw=red,anchor=north,font=\huge]
{
\node {a}; \fill[blue] (0,0) circle (2pt); &
\node {X}; \fill[blue] (0,0) circle (2pt); &
\node {g}; \fill[blue] (0,0) circle (2pt); \\
};
\end{tikzpicture}
\\ \hline  
 & anchor=base &  anchor=north \\ \hline 
\end{tabular} 

\bigskip
\begin{tabular}{|c|c|c|} \hline  
\begin{tikzpicture}
[every node/.style={draw=black,font=\huge}]
\matrix [draw=red]
{

\node[left]  {X}; \fill[blue] (0,0) circle (2pt);  \\
};
\end{tikzpicture}
&  
\begin{tikzpicture}
[every node/.style={draw=black,anchor=base,font=\huge}]
\matrix [draw=red]
{
\node {a}; \fill[blue] (0,0) circle (2pt); ²\\
\node[right] {X}; \fill[blue] (0,0) circle (2pt);  \\
\node {g}; \fill[blue] (0,0) circle (2pt); \\
};
\end{tikzpicture}
&  
\begin{tikzpicture}[every node/.style={draw=black}]
\matrix [draw=red,anchor=north,font=\huge]
{
\node {a}; \fill[blue] (0,0) circle (2pt); &
\node[right] {X}; \fill[blue] (0,0) circle (2pt); &
\node {g}; \fill[blue] (0,0) circle (2pt); \\
};
\end{tikzpicture}
\\ \hline  
 & anchor=base &  anchor=north \\ \hline 
\end{tabular} 

\bigskip

\begin{tabular}{|c|c|} \hline  
\begin{tikzpicture}[baseline=0pt]
\matrix [draw=red,nodes=draw]
{
\node[left] {A}; \fill[blue] (0,0) circle (2pt); \\
\node {B}; \fill[blue] (0,0) circle (2pt); \\
\node[right] {C}; \fill[blue] (0,0) circle (2pt); \\
};
\end{tikzpicture}
&  
\parbox{12cm}{
\BS{matrix} [draw=red,nodes=draw]
\AC{\\
\BS{node}\rouge{[left]} {A}; \BS{fill}[blue] (0,0) circle (2pt); \BS{} \BS{} \\
\BS{node} {B}; \BS{fill}[blue] (0,0) circle (2pt);\BS{} \BS{} \\
\BS{node}\rouge{[right]} {C}; \BS{fill}[blue] (0,0) circle (2pt); \BS{} \BS{}\\
}; \\
}

\\ \hline 
\end{tabular} 

\bigskip

\begin{tabular}{|c|c|} \hline  
\multicolumn{2}{|c|}{\BS{matrix} [draw,\RDD{column  sep}=1cm,nodes=draw]} 
\\ \hline 
\begin{tikzpicture}
\matrix [draw,column sep=1cm,nodes=draw]
{
\node(a) {123}; & \node (b) {1}; & \node {1}; \\
\node {12}; & \node {12}; & \node {1}; \\
\node(c) {1}; & \node (d) {123}; & \node {1}; \\
};
\draw [red,thick] (a.east) -- (a.east |- c)
(d.west) -- (d.west |- b);
\draw [<->,red,thick] (a.east) -- (d.west |- b)
node [above,midway] {1cm};
\end{tikzpicture}
&  
\begin{tikzpicture}
\matrix [draw,column sep={1cm,between origins},nodes=draw]
{
\node(a) {123}; & \node (b) {1}; & \node {1}; \\
\node {12}; & \node {12}; & \node {1}; \\
\node {1}; & \node {123}; & \node {1}; \\
};
\draw [<->,red,thick] (a.center) -- (b.center) node [above,midway] {1cm};
\end{tikzpicture}
\\ \hline \RDD{column sep}=1cm & column sep=\AC{1cm,\RDD{between origins} } 
\\ \hline 
\end{tabular} 

\bigskip

\begin{tabular}{|c|c|} \hline
\multicolumn{2}{|c|}{\BS{matrix} [draw,\RDD{row sep}=1cm,nodes=draw]} 
\\ \hline 
\begin{tikzpicture}
\matrix [draw,row sep=1cm,nodes=draw]
{
\node (a) {123}; & \node {1}; & \node {1}; \\
\node (b) {12}; & \node {12}; & \node {1}; \\
\node {1}; & \node {123}; & \node {1}; \\
};
\draw [<->,red,thick] (a.south) -- (b.north) node [right,midway] {1cm};
\end{tikzpicture}
&
\begin{tikzpicture}
\matrix [draw,row sep={1cm,between origins},nodes=draw]
{
\node (a) {123}; & \node {1}; & \node {1}; \\
\node (b) {12}; & \node {12}; & \node {1}; \\
\node {1}; & \node {123}; & \node {1}; \\
};
\draw [<->,red,thick] (a.center) -- (b.center) node [right,midway] {1cm};
\end{tikzpicture}
\\  \hline 
\RDD{row sep}=1cm  & row sep=\AC{1cm,\RDD{between origins} } 
\\ \hline 


\end{tabular} 




\bigskip

\begin{tabular}{|c|c|} \hline  
\multicolumn{2}{|c|}{\BS{matrix} [ \rouge{row sep=5mm},draw,nodes=draw]} \\
\multicolumn{2}{|c|}{ \{ \BS{node} \AC{1}; \& \BS{node} \AC{2}; \& \BS{node} \AC{3}; \BS{}\BS{}  } \\
\multicolumn{2}{|c|}{ \BS{node} \AC{4} ; \& \BS{node}  \AC{5}; \& \BS{node}  \AC{6};  \BS{}\BS{} \rouge{[1cm]} } \\
\multicolumn{2}{|c|}{ \BS{node} \AC{7}; \& \BS{node}\AC{8}; \& \BS{node}\AC{9}; \BS{}\BS{} \}  } 
\\ \hline  
\begin{tikzpicture}
\matrix [row sep=5mm,draw,nodes=draw]
{
\node {1}; & \node {2};& \node {3}; \\
\node(a) {4} ; & \node {5}; & \node {6};\\[1cm]
\node(b) {7}; &\node {8}; & \node {9}; \\
};
\draw [<->,red,thick] (a.center) -- (b.center) node [right,midway] {1,5cm};
\end{tikzpicture}
&  
\begin{tikzpicture}
\matrix [row sep=5mm,draw,nodes=draw]
{
\node {1}; & \node {2};& \node {3}; \\
\node(a) {4} ; & \node {5}; & \node {6};\\[10mm,between origins]
\node(b) {7}; &\node {8}; & \node {9}; \\
};
\draw [<->,red,thick] (a.center) -- (b.center) node [right,midway] {1,5cm};
\end{tikzpicture}
\\ \hline 
\rouge{[1cm]} & \rouge{[1cm,between origins]}
\\ \hline 
\end{tabular} 

\bigskip

\begin{tabular}{|c|c|} \hline  
\multicolumn{2}{|c|}{\BS{matrix} [ \rouge{column sep=5mm},draw,nodes=draw]} \\
\multicolumn{2}{|c|}{ \{ \BS{node} \AC{1}; \& \BS{node} \AC{2}; \& \BS{node} \AC{3}; \BS{}\BS{}  } \\
\multicolumn{2}{|c|}{ \BS{node} \AC{4} ; \& \BS{node}  \AC{5}; \& \rouge{[1cm]}\BS{node}  \AC{6};  \BS{}\BS{}  } \\
\multicolumn{2}{|c|}{ \BS{node} \AC{7}; \& \BS{node}\AC{8}; \& \BS{node}\AC{9}; \BS{}\BS{} \}  } 
\\ \hline  

\begin{tikzpicture}
\matrix [draw,nodes=draw,column sep=5mm]
{
\node {1}; & \node(a) {2}; &[1cm] \node(b) {3}; \\
\node {4}; & \node{5}; & \node {6}; \\
\node {7}; & \node{8}; & \node {9}; \\
};
\draw [<->,red,thick] (a.east) -- (b.west) node [above,midway] {15mm};
\end{tikzpicture}
&  
\begin{tikzpicture}
\matrix [draw,nodes=draw,column sep=5mm]
{
\node {1}; &[2mm] \node(a){2}; &[1cm,between origins] \node(b){3}; \\
\node {4}; & \node {5}; & \node {6}; \\
\node {7}; & \node {8}; & \node {9}; \\
};
\draw [<->,red,thick] (a.center) -- (b.center) node [above,midway] {15mm};
\end{tikzpicture}
\\ \hline  
\rouge{[1cm]}
&  
\rouge{[1cm,between origins]}
\\ \hline 
\end{tabular} 




\bigskip

\begin{tikzpicture}
\matrix [draw,nodes=draw,column sep={1cm,between origins}]
{
\node (a) {8}; & \node (b) {1}; &[between borders] \node (c) {6}; \\
\node {3}; & \node {5}; & \node {7}; \\
\node {4}; & \node {9}; & \node {2}; \\
};
\draw [<->,red,thick] (a.center) -- (b.center) node [above,midway] {10mm};
\draw [<->,red,thick] (b.east) -- (c.west) node [above,midway] {1cm};
\end{tikzpicture}



\SbSbSSCT{Format des cellules}{Cell Styles and Options}

\noindent 

\begin{tabular}{|c|} \hline  
\BS{matrix} [nodes=draw,nodes=\AC{\rouge{fill}=blue!10\rouge{,minimum size}=1cm}]
\\ \hline  
\begin{tikzpicture}
\matrix [nodes=draw,nodes={fill=blue!10,minimum size=1cm}]
{
\node {1}; & \node{2}; & \node {3}; \\
\node {4}; & \node{5}; & \node {6}; \\
\node {7}; & \node{8}; & \node {9}; \\
};
\end{tikzpicture}
\\ \hline 
\end{tabular} 


\bigskip 


\begin{tabular}{|c|c|c|} \hline 
\multicolumn{3}{|c|}{\BS{matrix}[\rouge{row 2/.style}=\AC{red}]}
 \\ \hline 
\begin{tikzpicture}
\matrix[row 2/.style={red}]
{
\node {8}; & \node{1}; & \node {6}; \\
\node {3}; & \node{5}; & \node {7}; \\
\node {4}; & \node{9}; & \node {2}; \\
};
\end{tikzpicture}
&  
\begin{tikzpicture}
\matrix[column 2/.style={red}]
{
\node {8}; & \node{1}; & \node {6}; \\
\node {3}; & \node{5}; & \node {7}; \\
\node {4}; & \node{9}; & \node {2}; \\
};
\end{tikzpicture}
&  
\begin{tikzpicture}
\matrix[row 2 column 2/.style={red}]
{
\node {8}; & \node{1}; & \node {6}; \\
\node {3}; & \node{5}; & \node {7}; \\
\node {4}; & \node{9}; & \node {2}; \\
};
\end{tikzpicture}
\\ \hline 
row 2/.style=\AC{red} & column 2/.style=\AC{red}  & row 2 column 2/.style=\AC{red}\\ 
\hline 
\end{tabular} 

\bigskip 

\begin{tabular}{|c|c|c|} \hline 
\multicolumn{3}{|c|}{\BS{matrix}[column 1/.style=\AC{anchor=west}]}
 \\ \hline 
\begin{tikzpicture}
\matrix[column 1/.style={anchor=west}]
{
\node {12345};  & \node {67890}; \\
\node {123}; & \node{67};  \\
\node {1}; & \node{6}; & \\
};
\end{tikzpicture}
&  
\begin{tikzpicture}
\matrix[column 1/.style={anchor=east}]
{
\node {12345};  & \node {67890}; \\
\node {123}; & \node{67};  \\
\node {1}; & \node{6}; & \\
};
\end{tikzpicture}
&  
\begin{tikzpicture}
\matrix[column 1/.style={anchor=base}]
{
\node {12345};  & \node {67890}; \\
\node {123}; & \node{67};  \\
\node {1}; & \node{6}; & \\
};
\end{tikzpicture}
\\  \hline  
[\rouge{column 1/.style}={anchor=west}]& [\rouge{column 1/.style}={anchor=east}] & [\rouge{column 1/.style}={anchor=base}]\\ 
\hline 
\end{tabular} 

\bigskip

\begin{tabular}{|c|c|c|c|} \hline
\multicolumn{4}{|c|}{\BS{matrix}[matrix of nodes,\RDD{every odd column}/.style={red}]}
 \\ \hline 
\begin{tikzpicture}
\matrix [matrix of nodes,every odd column/.style={red}]
{
a & b & c & d \\
e & f & g & h \\
i & j & k & l \\
};
\end{tikzpicture}
&  
\begin{tikzpicture}
\matrix [matrix of nodes,every even column/.style={red}]
{
a & b & c & d \\
e & f & g & h \\
i & j & k & l \\
};
\end{tikzpicture}
&  
\begin{tikzpicture}
\matrix [matrix of nodes,every odd row/.style={red}]
{
a & b & c & d \\
e & f & g & h \\
i & j & k & l \\
};
\end{tikzpicture}
&  
\begin{tikzpicture}
\matrix [matrix of nodes,every even row/.style={red}]
{
a & b & c & d \\
e & f & g & h \\
i & j & k & l \\
};
\end{tikzpicture}
\\ 
\hline 
\RDD{every odd column} & \RDD{every even column} & \RDD{every odd row}  & \RDD{every even row} \\ 
\hline 
\end{tabular} 


\bigskip


\begin{tabular}{|c|} \hline  
\BS{matrix} [draw,matrix of nodes,\rouge{execute at begin cell}=\AC{(}]
\\ \hline  
\begin{tikzpicture}
\matrix [draw,matrix of nodes,execute at begin cell={(}]
{
1 & 2 &   \\
4 &   & 6 \\
  &   & 9 \\
};
\end{tikzpicture}
\\ \hline 
\end{tabular} 

\bigskip

\begin{tabular}{|c|} \hline  
\BS{tikz} 
[matrix of nodes/.style=\AC{
execute at begin cell=\BS{node}\BS{bgroup} , \\
\rouge{execute at end cell}=\$m\wedge 2\$\BS{egroup}; 
}] \\
\BS{matrix} [draw,matrix of nodes
]
\\ \hline  
\tikz 
[matrix of nodes/.style={
execute at begin cell=\node\bgroup ,
execute at end cell=$m^2$\egroup;
}]
\matrix [draw,matrix of nodes
]
{1 & 2 &  \\
4 &   & 6 \\
  & 8 & 9 \\
};
\\ \hline 
\end{tabular}

\bigskip

\begin{tabular}{|c|} \hline 

 \BS{matrix} [raw,matrix of nodes, \rouge {execute at empty cell}=\BS{node}\AC{- -}; ]
\\ \hline 
 
\begin{tikzpicture}
\matrix [draw,matrix of nodes,execute at empty cell=\node{--};]
{
1 & 2 & \\
4 & & 6 \\
& & 9 \\
};
\end{tikzpicture}
\\ \hline  
\end{tabular} 


\newpage
\SbSbSSCT{Points d'ancrage}{Anchoring a Matrix}

\begin{center}
\RRR{20-4}
\end{center}

\begin{tabular}{|c|c|c|} \hline 
\multicolumn{3}{|c|}{
\BS{matrix} [draw=red,nodes=draw,\RDD{matrix anchor}=east](XXX) at (1,1) }
\\ \hline  
\begin{tikzpicture}
\draw[help lines] (0,0) grid (3,3);
\matrix [draw=red,nodes=draw,matrix anchor=west](XXX) at (1,1)
{
\node {123}; \\ 
\node {12}; \\
\node {1}; \\
};
\fill[red](XXX.west) circle (3pt);
\end{tikzpicture}
&  
\begin{tikzpicture}
\draw[help lines] (0,0) grid (3,3);
\matrix [draw=red,nodes=draw,matrix anchor=east](XXX) at (1,1)
{
\node {123}; \\ 
\node {12}; \\
\node {1}; \\
};
\fill[red] (XXX.east) circle (3pt);
\end{tikzpicture}
&  
\begin{tikzpicture}
\draw[help lines] (0,0) grid (3,3);
\matrix [draw=red,nodes=draw,matrix anchor=south](XXX) at (1,1)
{
\node {123}; \\ 
\node {12}; \\
\node {1}; \\
};
\fill[red](XXX.south) circle (3pt);
\end{tikzpicture}

\\  \hline 
matrix anchor=west & matrix anchor=east & matrix anchor=south 
\\ \hline 
\end{tabular} 

\bigskip 
\begin{tabular}{|c|c|c|c|} \hline 
\multicolumn{2}{|c|}{\BS{matrix} [draw=red,nodes=draw,\rouge{anchor=west}] }
\\ \hline  
\begin{tikzpicture}
\matrix [draw=red,nodes=draw,anchor=west] 
{
\node {123}; & \node {abc}; \\ 
\node {12}; & \node {ab}; \\
\node {1}; & \node {a}; \\
};
\end{tikzpicture}
&  
\begin{tikzpicture}
\matrix [draw=red,nodes=draw,anchor=east] 
{
\node {123};& \node {abc}; \\ 
\node {12};  &\node {ab};\\
\node {1};  & \node {a}; \\
};
\end{tikzpicture}

\\ \hline  
anchor=west & anchor=east  \\ 
\hline 
\end{tabular} 

\bigskip 


\begin{tabular}{|c|c|}\hline  
\begin{tikzpicture}[baseline=1cm]
\draw[help lines] (0,0) grid (4,3);
\matrix[draw=red,nodes=draw ,matrix anchor=inner node.south,anchor=base, row sep=5mm, column sep=5mm] at (2,1)
{
\node {a}; & \node {b}; & \node {c}; & \node {d}; \\
\node {a}; & \node {b}; & \node(inner node){c}; & \node {d}; \\
\node {a}; & \node {b}; & \node {c}; & \node {d}; \\
};
\fill[red] (inner node.south) circle (3pt);
\end{tikzpicture}
&  
\parbox{10.5cm}{
\BS{matrix}[draw=red,nodes=draw, \\ 
\RDD{ matrix anchor}=\blll{inner node}.south, anchor=base, \\
  row sep=5mm,column sep=5mm] at (2,1) \\
\{ \\
\BS{node} \AC{a}; \& \BS{node} \AC{b}; \& \BS{node} \AC{c}; \& \BS{node} \AC{d};  \BS{}\BS{} \\
\BS{node} \AC{a}; \& \BS{node} \AC{b}; \& \BS{node}(\blll{inner node})\AC{c}; \& \BS{node} \AC{d};  \BS{}\BS{} \\
\BS{node}\AC{a}; \& \BS{node} \AC{b}; \& \BS{node}\AC{c}; \& \BS{node} \AC{d}; \BS{}\BS{}  \\
\};
}
\\ \hline 
\end{tabular} 


\SbSbSSCT{Changement du séparateur}{Considerations Concerning Active Characters}

\begin{center}
\RRR{20-5}
\end{center}

\begin{tabular}{|c|c|} \hline  
\tikz[baseline=0pt]
\matrix [ampersand replacement=\|]
{
\draw (0,0) circle (4mm); \| \node[rotate=10] {Hello}; \\
\draw (0.2,0) circle (2mm); \| \fill[red] (0,0) circle (3mm); \\
};
& 
\parbox{12cm}{ 
\BS{tikz}
\BS{matrix} [\RDD{ampersand replacement}=\blll{\BS{|}} ] \\
\{ \\
\BS{draw} (0,0) circle (4mm); \blll{\BS{|} }  \BS{node}[rotate=10] \AC{Hello}; \BS{}\BS{} \\
\BS{draw} (0.2,0) circle (2mm);  \blll{\BS{|} }   \BS{fill}[red] (0,0) circle (3mm); \BS{}\BS{} \\
\}; \\
}
\\ \hline 
\end{tabular} 


\SbSSCT{Matrice de n\oe uds (compléments) }{Matrix Library}

 \maboite{\BS{usetikzlibrary}\AC{matrix}}
\label{lib-matrix}


\begin{center}
\RRR{57-1}
\end{center}

\begin{tabular}{|c|c|} \hline  
\begin{tikzpicture}[baseline=0pt]
\matrix (XXX) [matrix of nodes]
{
1 & 2 & 3 \\
4 & 5 & 6 \\
7 & 8 & 9 \\
};
\end{tikzpicture}
& 
\parbox{10cm}{ 
\BS{begin}\AC{tikzpicture} \\
\BSS{matrix}  [matrix of nodes]\\
\{ \\
1 \hspace{3mm} \& \hspace{3mm}  2 \hspace{3mm} \& \hspace{3mm} 3 \hspace{3mm} \BS{}\BS{}   \\
4 \hspace{3mm} \& \hspace{3mm}  5 \hspace{3mm} \& \hspace{3mm} 6 \hspace{3mm} \BS{}\BS{}  \\
7 \hspace{3mm} \& \hspace{3mm}  8 \hspace{3mm} \& \hspace{3mm} 9 \hspace{3mm} \BS{}\BS{} \\
\}; \\
\BS{end}\AC{tikzpicture}
}
\\ \hline  
\end{tabular} 

\bigskip

\begin{tabular}{|c|c|} \hline  
\begin{tikzpicture}[baseline=0pt]
\matrix (XXX) [matrix of nodes,column sep=.5cm,row sep=.5cm,every node/.style=draw]
{
1 & 2 & 3 \\
4 & 5 & 6 \\
7 & 8 & 9 \\
};
\draw[thick,red,->] (XXX-1-1) -- (XXX-2-3);
\end{tikzpicture}
& 
\parbox{10cm}{ 
\BS{begin}\AC{tikzpicture} \\
\BSS{matrix} \blll{(XXX)} [matrix of nodes,column sep=.5cm,row sep=.5cm,every node/.style=draw]\\
\{ \\
1 \hspace{3mm} \& \hspace{3mm} 2 \hspace{3mm} \& \hspace{3mm} 3 \hspace{3mm} \BS{}\BS{}   \\
4 \hspace{3mm} \& \hspace{3mm} 5 \hspace{3mm} \& \hspace{3mm} 6 \hspace{3mm} \BS{}\BS{}  \\
7 \hspace{3mm} \& \hspace{3mm} 8 \hspace{3mm} \& \hspace{3mm} 9 \hspace{3mm} \BS{}\BS{} \\
\}; \\
\BS{draw}[thick,red,->] \blll{(XXX-1-1)} - - \blll{(XXX-2-3)} ; \\
\BS{end}\AC{tikzpicture}
}
\\ \hline  
\end{tabular} 

\bigskip


\begin{tabular}{|c|c|} \hline  
\begin{tikzpicture}
\matrix [matrix of nodes,column sep=.5cm,row sep=.5cm,every node/.style=draw]
{
8 & 1 & 6 \\
3 & 5 & |[red]| 7 \\
4 & 9 & 2 \\
};
\end{tikzpicture}
&  
\begin{tikzpicture}
\matrix [matrix of nodes]
{
1 & \& &  2 & \& &  3 				& \BS{}\BS{} \\
4 & \& & 5 	& \& & \rouge{ $|[$red$]|$} 6 & \BS{}\BS{} \\
7 & \& & 8 	& \& & 9 				& \BS{}\BS{} \\
};
\end{tikzpicture}
\\ \hline 
\end{tabular}  


\bigskip

\begin{tabular}{|c|c|} \hline 
\begin{tikzpicture}[baseline=-1cm] 
\matrix [matrix of nodes,column sep=.5cm,row sep=.5cm,every node/.style=draw]
{
AAA 			& |[circle]| BBB \\
CCC & |(d) [isosceles triangle]| DDD \\
| [ellipse]| EEE &  FFF \\
};
\end{tikzpicture}
& 
\begin{tikzpicture}
\matrix [matrix of nodes]
{
AAA & \& & \rouge{ $|[$circle$]|$} BBB &  \BS{}\BS{} \\
CCC & \& &\rouge{ $|[$isosceles triangle$]|$} DDD 	&  \BS{}\BS{} \\
\rouge{ $|[$ellipse$]|$} EEE & \& & FFF & \BS{}\BS{} \\
};
\end{tikzpicture}
\\ \hline 
\end{tabular} 


\bigskip

\begin{tabular}{|c|c|} \hline 
\begin{tikzpicture}[baseline=-2cm] 
\matrix [matrix of nodes,column sep=.5cm,row sep=.5cm,every node/.style=draw]
{
|(a)| AAA 	& |(b)| BBB \\
|(c)| CCC 	& |(d)| DDD \\
|(e)| EEE 	& |(f)| FFF \\
};
\draw (a) -- (d);
\draw (d) -- (f);
\end{tikzpicture}
&  
\begin{tikzpicture}
\node at (0,1.5) [text width=10cm]
{\BS{matrix} [matrix of nodes,column sep=.5cm,row sep=.5cm,every node/.style=draw] \\
\{ 
};
\matrix [matrix of nodes]
{
\rouge{ $|$(a)$|$} AAA & \& & \rouge{ $|$(b)$|$} BBB &  \BS{}\BS{} \\
\rouge{ $|$(c)$|$} CCC & \& & \rouge{ $|$(d)$|$} DDD 	&  \BS{}\BS{} \\
\rouge{ $|$(e)$|$} EEE & \& & \rouge{ $|$(f)$|$} FFF & \BS{}\BS{} \\
};

\node at (0,-1.2) [text width=10cm]
{  \}; \\ 
\BS{draw} (a) - - (d); \\ \BS{draw} (d) - - (f);
};
\end{tikzpicture}
\\ \hline 
\end{tabular} 

\bigskip


\begin{tabular}{|c|c|} \hline  
\begin{tikzpicture}
\matrix [matrix of nodes]
{
1 &[1cm] 2 &[5mm] |[red]| 3 \\
4 & 5 &  6 \\
7 & 8 & 9 \\
};
\end{tikzpicture}
&
\begin{tikzpicture}
\matrix [matrix of nodes]
{
1 & \& & \rouge{\lbrack 1cm \rbrack} 2 & \& &\rouge{\lbrack 5mm \rbrack} |[red]| 3 & \BS{}\BS{} \\
4 & \& & 5 & \& & 6 & \BS{}\BS{} \\
7 & \& & 8 & \& & 9 & \BS{}\BS{} \\
};
\end{tikzpicture}

\\ \hline 
\end{tabular} 



\bigskip

\begin{tabular}{|c|c|} \hline  
\begin{tikzpicture}[baseline=0pt]
\matrix [matrix of math nodes]
{
A_1 & A_2 & A_3 \\
a_4 & a_5 &  a_6 \\
a^7 & a^8 & a^9 \\
};
\end{tikzpicture}
&  
\parbox{8cm}{ 
\BSS{matrix}  [\rouge{ matrix of math nodes}]\\
\{ \\
A\_1 \hspace{2mm}  \& \hspace{2mm}  A\_2 \hspace{2mm}  \& \hspace{2mm}  A\_3 \hspace{2mm}   \BS{}\BS{}   \\
a\_4 \hspace{2mm}  \& \hspace{2mm}  a\_5 \hspace{2mm}  \&  \hspace{2mm}  a\_6 \hspace{2mm}  \BS{}\BS{}  \\
a\land 7 \hspace{2mm}  \& \hspace{2mm}  a\land 8 \hspace{2mm}  \& \hspace{2mm}  a\land 9 \hspace{2mm}   \BS{}\BS{} \\
\}; 
}
\\ \hline  
\end{tabular} 

\bigskip

\begin{tabular}{|c|c|} \hline  
\begin{tikzpicture}[baseline=0pt]
\matrix [matrix of math nodes,nodes={circle,draw}]
{
a_1 & & a_3 \\
a_4 & & a_6 \\
a_7 & a_8 & \\
};
\end{tikzpicture}
&  
\parbox{10cm}{ 
\BSS{matrix}  [matrix of math nodes,\rouge{nodes={circle,draw}}]\\
\{ \\
A\_1 \hspace{2mm}  \& \hspace{12mm}  \& \hspace{2mm}  A\_3 \hspace{2mm}   \BS{}\BS{}   \\
a\_4 \hspace{2mm}  \& \hspace{12mm}  \& \hspace{2mm}   a\_6 \hspace{2mm}  \BS{}\BS{}  \\
a\_ 7 \hspace{2mm}  \& \hspace{2mm}  a\_ 8 \hspace{2mm}  \& \hspace{12mm}    \BS{}\BS{} \\
\}; 
}
\\ \hline 
\end{tabular} 

\bigskip

\begin{tabular}{|c|c|} \hline  
\begin{tikzpicture}[baseline=0pt]
\matrix [matrix of math nodes,nodes={circle,draw},nodes in empty cells]
{
a_1 & & a_3 \\
a_4 & & a_6 \\
a_7 & a_8 & \\
};
\end{tikzpicture}
&  
\parbox{10cm}{ 
\BSS{matrix}  [matrix of math nodes,nodes={circle,draw} ,\rouge{nodes in empty cells}]\\
\{ \\
A\_1 \hspace{2mm}  \& \hspace{12mm}  \& \hspace{2mm}  A\_3 \hspace{2mm}   \BS{}\BS{}   \\
a\_4 \hspace{2mm}  \& \hspace{12mm}  \& \hspace{2mm}   a\_6 \hspace{2mm}  \BS{}\BS{}  \\
a\_ 7 \hspace{2mm}  \& \hspace{2mm}  a\_ 8 \hspace{2mm}  \& \hspace{12mm}    \BS{}\BS{} \\
\}; 
}
\\ \hline 
\end{tabular} 

\SbSbSSCT{Texte dans les n\oe uds}{Characters in Matrices of Nodes}

\begin{center}
\RRR{57-2}
\end{center}


\begin{tabular}{|c|c|} \hline  
\begin{tikzpicture}[baseline=0pt]
\matrix [matrix of nodes,nodes={text width=2cm,draw}]
{
aaa  & bbb \\ 
ccc \\
eee & fff\\
};
\end{tikzpicture}
&  
\parbox{10cm}{ 
\BSS{matrix}  [matrix of nodes,\rouge{nodes=\AC{text width=2cm,draw}} ]\\
\{ \\
aaa \&  bbb \BS{}\BS{}  \\
ccc \BS{}\BS{}  \\
eee \& fff \BS{}\BS{}  \\
\}; 
}
\\ \hline 
\end{tabular} 

\bigskip

\begin{tabular}{|c|c|}  \hline  
\begin{tikzpicture}[baseline=0cm]
\matrix [matrix of nodes,nodes={text width=2cm,draw}]
{
1 & {aaa \\ bbb \\ ccc } \\
2 & ddd \\
};
\end{tikzpicture}
&  
\parbox{10cm}{ 
\BSS{matrix}  [matrix of nodes,nodes=\AC{text width=2cm,draw} ]\\
\{ \\
1 \& \& \rouge { \AC{aaa \BS{}\BS{} bbb \BS{}\BS{} ccc } } \BS{}\BS{}   \\
2 \& \& ddd \BS{}\BS{}  \\
\}; 
}
\\ \hline 
\end{tabular} 

\bigskip

\SbSbSSCT{Délimiteurs}{Delimiters}


\begin{center}
\RRR{57-3}
\end{center}

\bigskip

\begin{tabular}{|c|c|c|c|} \hline 
\multicolumn{4}{|c|}{\BS{matrix} [matrix of math nodes,\RDD{left delimiter}=( ]}
\\ \hline  
\begin{tikzpicture}
\matrix [matrix of math nodes,left delimiter=( ]
{
a_1 & a_2 & a_3 \\
a_4 & a_5 & a_6 \\
a_7 & a_8 & a_9 \\
};
\end{tikzpicture}
&  
\begin{tikzpicture}
\matrix [matrix of math nodes,right delimiter=\}]
{
a_1 & a_2 & a_3 \\
a_4 & a_5 & a_6 \\
a_7 & a_8 & a_9 \\
};
\end{tikzpicture}
&
\begin{tikzpicture}
\matrix [matrix of math nodes,above delimiter=\| ]
{
a_1 & a_2 & a_3 \\
a_4 & a_5 & a_6 \\
a_7 & a_8 & a_9 \\
};
\end{tikzpicture}
&
\begin{tikzpicture}
\matrix [matrix of math nodes,below delimiter=\rmoustache ]
{
a_1 & a_2 & a_3 \\
a_4 & a_5 & a_6 \\
a_7 & a_8 & a_9 \\
};
\end{tikzpicture}


\\  \hline 
\RDD{left delimiter}=(  & \RDD{right delimiter}=\BS{\}} & \RDD{above delimiter}=\BS{|} & \RDD{below delimiter}=\BS{rmoustache}
\\  \hline
\end{tabular} 

\bigskip
\begin{tabular}{|c|} \hline  
\BS{tikz}
\BS{node} [fill=red!20,text width=2cm,\rouge{left delimiter}=\BS{\{} ] \\
\AC{Ceci est une démonstration d'un texte  sur une largeur de 2cm.};
\\ \hline  
\tikz
\node [fill=red!20,text width=2cm,left delimiter=\{]
{Ceci est une démonstration d'un texte  sur une largeur de 2cm.};
\\ \hline 
\end{tabular} 



%% 
%%\newpage 
%%
%%\SbSSCT{Chaine de n\oe uds}{Chains of nodes}
%%
%%\SbSbSSCT{Création d'une chaine de n\oe euds}{Starting and Continuing a Chain}

 \maboite{\BS{usetikzlibrary}\AC{chains}}
\label{lib-chains}


\begin{center}
\RRR{46-2}
\end{center}

\bigskip

\begin{tabular}{|l|} \hline  
\BS{begin}\AC{tikzpicture}[\RDD{start chain}] \\
\BS{node} [\RDD{on chain}] \AC{A};\\
\BS{node}  [\RDD{on chain}] \AC{B};\\
\BS{node}  [\RDD{on chain}] \AC{C};\\
\BS{end}\AC{tikzpicture} \\ \hline  
\begin{tikzpicture}[start chain]
\node [on chain] {A};
\node [on chain] {B};
\node [on chain] {C};
\end{tikzpicture}
\\ \hline 
\end{tabular} 

\bigskip

\begin{tabular}{|c|}  \hline  
\BS{begin}\AC{tikzpicture}[start chain, \RDD{node distance}= 0.5 cm] 
\\ \hline  
\begin{tikzpicture}[start chain, node distance= .5 cm]
\node [on chain] {A};
\node [on chain] {B};
\node [on chain] {C};
\end{tikzpicture}
\\ \hline 
\end{tabular} 

\bigskip

\begin{tabular}{|c|}  \hline 
\BS{begin}\AC{tikzpicture}[start chain=\rouge {going below} ]
\\   \hline 
\begin{tikzpicture}[start chain=going below]
\node [on chain] {A};
\node [on chain] {B};
\node [on chain] {C};
\end{tikzpicture}
\\   \hline 
\end{tabular} 

\bigskip

\begin{tabular}{|c|}  \hline 
\BS{begin}\AC{tikzpicture}[start chain=\rouge {going left} ] 
\\   \hline 
\rule[0cm]{0pt}{.7cm}  
\begin{tikzpicture}[start chain=going left]
\node [on chain] {A};
\node [on chain] {B};
\node [on chain] {C};
\end{tikzpicture} 
\\ \hline 
\end{tabular} 


\bigskip

\begin{tabular}{|c|}  \hline  
\BS{begin}\AC{tikzpicture}[start chain, \rouge{every node/.style=draw} ] 
\\ \hline 
\rule[0cm]{0pt}{.7cm}  
\begin{tikzpicture}[start chain, every node/.style=draw]
\node [on chain] {A};
\node [on chain] {B};
\node [on chain] {C};
\end{tikzpicture}
\\ \hline 
\end{tabular} 

\bigskip

\begin{tabular}{|c|c|}\hline
\begin{tikzpicture}[start chain=1 going right,
start chain=2 going left]
\node [draw,on chain=1] {A};
\node [draw,on chain=1] {B};
\node [draw,on chain=1] {C};
\node [draw,on chain=2] at (3,1) {0};
\node [draw,on chain=2] {1};
\node [draw,on chain=2] {2};
\node [draw,on chain=1] {D};
\end{tikzpicture} 
 &  
\parbox{10cm}{
\BS{begin}\AC{tikzpicture}[\rouge{start chain=1} going right , \\
\blll{start chain=2} going left] \\
\BS{node} [draw,\rouge{on chain=1}] \AC{A}; \\
\BS{node} [draw,\rouge{on chain=1}] \AC{B}; \\
\BS{node}[draw,\rouge{on chain=1}] \AC{C}; \\
\BS{node} [draw,\blll{on chain=2}] at (3,1) \AC{0}; \\
\BS{node} [draw,\blll{on chain=2}] \AC{1}; \\
\BS{node} [draw,\blll{on chain=2}] \AC{2}; \\
\BS{node}[draw,\rouge{on chain=1}] \AC{D}; \\
\BS{end}\AC{tikzpicture}} 
\\ \hline 
\end{tabular} 

\bigskip


\begin{tabular}{|c|c|} \hline  
\rule[-2cm]{0pt}{4cm} 
\begin{tikzpicture}[start chain=going right,baseline=-1.5cm]
\node [draw,on chain] {A};
\node [draw,on chain] {B};
\node [draw,continue chain=going below,on chain] {C};
\node [draw,on chain] {D};
\node [draw,continue chain=going right,on chain] {E};
\end{tikzpicture}
&  
\parbox{11cm}{
\BS{begin}\AC{tikzpicture}[start chain going right]
\BS{node} [draw,on chain] \AC{A}; \\
\BS{node} [draw,on chain] \AC{B}; \\
\BS{node} [draw,\RDD{continue chain}=going below,on chain] \AC{C}; \\
\BS{node}[draw,on chain] \AC{D}; \\
\BS{node} [draw,\RDD{continue chain}=going right,on chain] \AC{E}; \\
\BS{end}\AC{tikzpicture}} 
\\ \hline 
\end{tabular} 

\bigskip

\begin{tabular}{|c|c|}  \hline 
\begin{tikzpicture}[every node/.style=draw,baseline=-1.5cm]
{ [start chain=1]
\node [on chain] {A};
\node [on chain] {B};
\node [on chain] {C};
}
{ [start chain=2 going below]
\node [on chain=2] at (0.5,-.5) {0};
\node [on chain=2] {1};
\node [on chain=2] {2};
}
{ [continue chain=1]
\node [on chain] {D};
}
\end{tikzpicture}
&  
\parbox{10cm}{
\BS{begin}\AC{tikzpicture}[start chain going right] \\
\{ [\RDD{start chain}=1] \\
\BS{node} [draw,on chain] \AC{A}; \\
\BS{node} [draw,on chain] \AC{B}; \\
\BS{node} [draw,on chain] \AC{C}; \\
\} \\
\{ [\RDD{start chain}=2] \\
\BS{node}[draw,on chain=2] \AC{0}; \\
\BS{node}[draw,on chain=2] \AC{1}; \\
\BS{node}[draw,on chain=2] \AC{2}; \\
\} \\
\{ [\RDD{continue chain}=1] \\
\BS{node} [draw,on chain] \AC{D}; \\
\} \\
\BS{end}\AC{tikzpicture}} 
\\  \hline 
\end{tabular} 

\bigskip

\SbSbSSCT{N\oe uds sur la chaine}{Nodes on a Chain}

\begin{center}
\RRR{46-3} 
\end{center}

\bigskip

\begin{tabular}{|c|c|} \hline 
 \begin{tikzpicture}[start chain=XXX placed  {at=(\tikzchaincount*-30+90:1.5)},baseline=0pt]
 \foreach \i in {1,...,12}
 \node [on chain] {\i};
 \draw (0,0) -- (XXX-10);
 \draw (0,0) -- (XXX-2);
 \end{tikzpicture}
&
\parbox{11cm}{
\BS{begin}\AC{tikzpicture}[start chain=\blll{XXX} \RDD{placed} \\ \AC{at=(\BSS{tikzchaincount}*-30+90:1.5)}] \\
 \BS{foreach} \BS{i} in \AC{1,...,12} \\
\BS{node} [on chain] \AC{\BS{i}}; \\
\BS{draw }(0,0) -- \blll{(XXX-10)}; \\
\BS{draw }(0,0) -- \blll{(XXX-2)}; \\
\BS{end}\AC{tikzpicture}} 
\\ \hline 
\end{tabular} 

\bigskip


\begin{tabular}{|c|c|}  \hline 
\begin{tikzpicture}[start chain,baseline=-1cm]
\node [draw,on chain] {A};
\node [draw,on chain] {B};
\node [draw,on chain=going below] {C};
\node [draw,on chain] {D};
\node [draw,on chain] {E};
\end{tikzpicture}
&  
\parbox{11cm}{
\BS{begin}\AC{tikzpicture}[start chain] \\
\BS{node} [draw,on chain] \AC{A}; \\
\BS{node} [draw,on chain] \AC{B}; \\
\BS{node} [draw,on chain=\rouge{going below}] \AC{C}; \\
\BS{node} [draw,on chain] \AC{D}; \\
\BS{node} [draw,on chain] \AC{E}; \\
\BS{end}\AC{tikzpicture}} 
\\  \hline 
\end{tabular} 


\bigskip

\begin{tabular}{|c|c|} \hline 
\begin{tikzpicture}[start chain=going {at=(\tikzchainprevious),shift=(30:1)},baseline=1cm]
\node [draw,on chain] {A};
\node [draw,on chain] {B};
\node [draw,on chain] {C};
\node [draw,on chain] {D};
\end{tikzpicture}  
&  
\parbox{11cm}{
\BS{begin}\AC{tikzpicture}[start chain=going \\ \AC{at=(\BSS{tikzchainprevious},shift=(30:1)}] \\
\BS{node} [draw,on chain] \AC{A}; \\
\BS{node} [draw,on chain] \AC{B}; \\
\BS{node} [draw,on chain] \AC{C}; \\
\BS{node} [draw,on chain] \AC{D}; \\
\BS{end}\AC{tikzpicture}} 
\\ \hline 
\end{tabular} 

\bigskip

\begin{tabular}{|c|c|} \hline 
\begin{tikzpicture}[baseline=1cm]
\node[draw,red] (A) at (0,2) {A};
{ [start chain]
\node [draw,on chain] {B};
\node [draw,on chain] {C};
\chainin (A) [join];
\node [draw,on chain] {D};
\node [draw,on chain] {E};
}
\end{tikzpicture}
&  
\parbox{11cm}{
\BS{begin}\AC{tikzpicture} \\
\BS{node} [draw,red] (A) at (0,2)  \AC{A}; \\
\{ [start chain] \\
\BS{node} [draw,on chain] \AC{B}; \\
\BS{node} [draw,on chain] \AC{C}; \\
\BSS{chainin} (A) [join]; \\
\BS{node} [draw,on chain] \AC{D}; \\
\BS{node} [draw,on chain] \AC{E}; \\
\} \\
\BS{end}\AC{tikzpicture}} 
\\  \hline 
\end{tabular} 



\bigskip

\begin{tabular}{|c|c|} \hline 
\begin{tikzpicture}[baseline=-1cm]
\matrix [matrix of nodes,column sep=1cm,row sep=1cm,every node/.style=draw]
{
|(a) | A 	& |(b) |  B 	& |(c) | C \\
|(d) | D 	& |(e) | E 		& |(f) | F \\
};
{ [start chain,every on chain/.style={join=by ->}]
\chainin (a);
\chainin (b);
\chainin (d);
\chainin (c);
\chainin (f);
\chainin (e);
}
\end{tikzpicture}
&  
\parbox{11cm}{
\BS{begin}\AC{tikzpicture} \\
\BS{matrix} [matrix of nodes,column sep=5mm,row sep=5mm] ,every node/.style=draw \\
\{ \\
|(a) | A 	\& |(b) |  B 	\& |(c) | C \BS{}\BS{} \\
|(d) | D 	\& |(e) | E 	\& |(f) | F \BS{}\BS{} \\
\}; \\
\{ [start chain,every on chain/.style=\AC{join=by ->}] \\
\BSS{chainin} (a);
\BSS{chainin}(b);
\BSS{chainin}(d); \\
\BSS{chainin} (c);
\BSS{chainin}(f);
\BSS{chainin}(e);
\}
\BS{end}\AC{tikzpicture}
} 
\\ \hline 
\end{tabular} 

\bigskip

\SbSbSSCT{Jonction de n\oe uds}{Joining Nodes on a Chain}

\begin{center}
\RRR{46-4}
\end{center} 

\bigskip

\begin{tabular}{|c|c|} \hline 
\begin{tikzpicture}[start chain]
\node [draw,on chain] {A};
\node [draw,on chain,join] {B};
\node [draw,on chain] {C};
\node [draw,on chain,join] {D};
\end{tikzpicture}
&  
\parbox{11cm}{
\BS{begin}\AC{tikzpicture}[start chain] \\
\BS{node} [draw,on chain] \AC{A}; \\
\BS{node} [draw,on chain,\RDD{join}] \AC{B}; \\
\BS{node} [draw,on chain] \AC{C}; \\
\BS{node} [draw,on chain,\RDD{join}] \AC{D}; \\
\BS{end}\AC{tikzpicture}} 
\\ \hline 
\end{tabular} 

\bigskip

\begin{tabular}{|c|c|} \hline 
\begin{tikzpicture}[start chain, every on chain/.style=join, every join/.style=->]
\node [draw,on chain] {A};
\node [draw,on chain] {B};
\node [draw,on chain] {C};
\node [draw,on chain] {D};
\end{tikzpicture}
&  
\parbox{11cm}{
\BS{begin}\AC{tikzpicture}[start chain, \RDD{every on chain}/.style=join, \\ \RDD{every join}/.style=->] \\
\BS{node} [draw,on chain] \AC{A}; \\
\BS{node} [draw,on chain,\RDD{join}] \AC{B}; \\
\BS{node} [draw,on chain] \AC{C}; \\
\BS{node} [draw,on chain,\RDD{join}] \AC{D}; \\
\BS{end}\AC{tikzpicture}} 
\\ \hline 
\end{tabular} 

\bigskip

\begin{tabular}{|c|c|}  \hline 
\begin{tikzpicture}[start chain,baseline=-1cm]
\node [draw,on chain] {A};
\node [draw,on chain] {B};
\node [draw,on chain] {C};
\node [draw,on chain=going below,join=with chain-2 ] {D};
\end{tikzpicture} 
&  
\parbox{11cm}{
\BS{begin}\AC{tikzpicture}[start chain] \\
\BS{node} [draw,on chain] \AC{A}; \\
\BS{node} [draw,on chain] \AC{B}; \\
\BS{node} [draw,on chain] \AC{C}; \\
\BS{node} [draw,on chain=going below,\rouge{join=with chain-2} ] \AC{D}; \\
\BS{end}\AC{tikzpicture}} 
\\ \hline 
\begin{tikzpicture}[start chain,baseline=-1cm]
\node [draw,on chain] {A};
\node [draw,on chain] {B};
\node [draw,on chain] {C};
\node [draw,on chain=going below,join=with chain-1 by {blue,<-}] {D};
\end{tikzpicture}
&
\parbox{12cm}{
\BS{begin}\AC{tikzpicture}[start chain] \\
\BS{node} [draw,on chain] \AC{A}; \\
\BS{node} [draw,on chain] \AC{B}; \\
\BS{node} [draw,on chain] \AC{C}; \\
\BS{node} [draw,on chain=going below,join=with chain-1 \rouge{ by \AC{blue,<-}} ] \AC{D}; \\
\BS{end}\AC{tikzpicture}} 
\\ \hline 
\end{tabular} 



\bigskip

\SbSbSSCT{Branches}{Branches}

\begin{center}
\RRR{46-5}
\end{center} 


\bigskip

\begin{tabular}{|c|c|}  \hline 
\begin{tikzpicture} [baseline=-2cm]
{ [start chain=XXX]
\node [draw,on chain] {A};
\node [draw,on chain] {B};
{ [start branch=YYY going below]
\node [draw,on chain] {1};
\node [draw,on chain] {2};
\node [draw,on chain] {3};
}
\node [draw,on chain,join=with XXX/YYY-end,join=with XXX/YYY-2 ] {C};
}
\end{tikzpicture}
&  
\parbox{12cm}{
\BS{begin}\AC{tikzpicture}\\
\{ [start chain=\blll{XXX}] \\
\BS{node} [draw,on chain] \AC{A}; \\
\BS{node} [draw,on chain] \AC{B}; \\
\{ [\RDD{start branch}=\blll{YYY} going below] \\
\BS{node} [draw,on chain] \AC{1}; \\
\BS{node} [draw,on chain] \AC{2}; \\
\BS{node} [draw,on chain] \AC{3}; \\
\} \\
\BS{node} [ draw,on chain,join=with \blll{XXX/YYY}\rouge{-end}, \\ join=with \blll{XXX/YYY}\rouge{-2}]  \AC{C}; \\
\} \\
\BS{end}\AC{tikzpicture}   } 

\\ \hline 
\end{tabular} 

\bigskip

\begin{tabular}{|c|} \hline 
\BS{begin}\AC{tikzpicture}[ \RDD{node distance}=.2cm and 3cm]
\\ \hline 
\begin{tikzpicture}[ node distance=.2cm and 3cm]
{ [start chain=XXX]
\node [on chain] {A};
\node [on chain] {B};
{ [start branch=YYY going below]
\node [on chain] {1};
\node [on chain] {2};
\node [on chain] {3};
}
\node [on chain,join=with XXX/YYY-end] {C};
}
\end{tikzpicture}
\\ \hline 
\end{tabular} 

\bigskip

\begin{tabular}{|c|c|} \hline 
\begin{tikzpicture}[ node distance=2mm and 1cm,baseline=-2cm]
{ [start chain=XXX]
\node [draw,on chain] {A};
\node [draw,on chain] {B};
{ [start branch=YYY going below]
\node [draw,on chain] {1};
\node [draw,on chain] {2};
\node [draw,on chain] {3};
}
\node [draw,on chain,join=with XXX/YYY-end] {C};
{
[continue branch=YYY]
\node [draw,on chain] {4};
\node [draw,on chain] {5};
}
}
\end{tikzpicture}
&  
\parbox{12cm}{
\BS{begin}\AC{tikzpicture}[ node distance=2mm and 1cm]\\
\{ [start chain=\blll{XXX}] \\
\BS{node} [draw,on chain] \AC{A}; \\
\BS{node} [draw,on chain] \AC{B}; \\
\{ [start branch=\blll{YYY} going below] \\
\BS{node} [draw,on chain] \AC{1}; \\
\BS{node} [draw,on chain] \AC{2}; \\
\BS{node} [draw,on chain] \AC{3}; \} \\
\BS{node}  [draw,on chain,join=with \blll{XXX/YYY}-end]  \AC{C}; \\
\{ [\RDD{continue branch}=\blll{YYY}]\\
\BS{node} [on chain] \AC{4}; \\
\BS{node} [on chain] \AC{5}; \} \\
\} \\
\BS{end}\AC{tikzpicture}   } 
\\ \hline 
\end{tabular} 


\bigskip

\begin{tabular}{|c|c|} \hline 
\begin{tikzpicture}[node distance=2mm and 1cm, every node/.style=draw,baseline=-1cm]
{ [start chain]
\node [on chain] {1};
\node [on chain] {2};
{ [start branch=XXX going below] }
\node [on chain] {3};
{ [start branch=YYY going above] }
\node [on chain] {4};
{ [continue branch=XXX]
\node [on chain] {a};
\node [on chain] {b};
}{
[continue branch=YYY]
\node [on chain] {A};
\node [on chain] {B};
}
}
\end{tikzpicture}
&  
\parbox{12cm}{
\BS{begin}\AC{tikzpicture}[node distance=2mm and 1cm, every node/.style=draw]\\
\{ [start chain] \\
\BS{node} [on chain] \AC{1};  \\
\BS{node} [on chain] \AC{2}; \\
\{ [\RDD{start branch}=\blll{XXX} going below] \} \\
\BS{node} [on chain] \AC{3}; \\
\{ [\RDD{start branch}=\blll{YYY} going above] \} \\
\BS{node} [on chain] \AC{4}; \\
\{ [\RDD{continue branch}=\blll{XXX} ] \\
\BS{node} [on chain] \AC{a}; \\
\BS{node} [on chain] \AC{b};\} \\
\{ [\RDD{continue branch}=\blll{YYY} ] \\
\BS{node} [on chain] \AC{A}; \\
\BS{node} [on chain] \AC{B}; \}  }
\\ \hline 
\end{tabular} 



%%
%%
%%\bigskip
%
%%\begin{tikzpicture}[baseline=-1cm] 
%%\matrix [matrix of nodes,column sep=.5cm,row sep=.5cm,every node/.style=draw]
%%{
%%|(a) [red]|AAA 			& |(b) [circle]| BBB \\
%%|(c)| CCC 		& |(d) [isosceles triangle]| DDD \\
%%|(e) [ellipse]| EEE & |(f)| FFFF \\
%%};
%%\draw (a) -- (f);
%%\end{tikzpicture}
%
%
%\newpage
%
%
%
%%\section{symboles}
%%
%%
\SbSSCT{Coordonnées}{Coordinates}
\begin{center}
\RRR{13-2-1}
\end{center}


\SbSbSSCT{Système de coordonnées \og canvas \fg}{Canvas coordinates}

\noindent


\tikzset{every picture/.style=blue,very thick,inner sep=0pt}

\begin{tabular}{|c|c|} \hline 
\TFRGB{Explicite}{explicit}  & \TFRGB{Implicite}{implicit}
\\ \hline
\begin{tikzpicture}
\draw[help lines] (0,0) grid (3,2);
\fill (canvas cs:x=2cm,y=1.5cm) circle (2pt);
\end{tikzpicture}
&
\begin{tikzpicture}
\draw[help lines] (0,0) grid (3,2);
\fill (2,1.5) circle (2pt);
\end{tikzpicture}

\\ \hline  
 \BS{fill} (\RDD{canvas cs}:\blll{x=2cm,y=1.5cm}) circle (2pt);
& \BS{fill} {\color{blue}(2cm,1.5cm)} circle (2pt);
\\ \hline 
\end{tabular} 


\SbSbSSCT{Système de coordonnées polaire \og canvas \fg}{Polar coordinates}

\noindent


\begin{tabular}{|c|c|c|} \hline
\TFRGB{Explicite}{explicit}  & \TFRGB{Implicite}{implicit}
\\ \hline
\begin{tikzpicture}
\draw[help lines] (0,0) grid (3,2);
\draw [dotted](0,2) arc (90 :0 :2);
\draw [dotted](0,0) --(2,2);
\fill (canvas polar cs:angle=45,radius=2cm) circle (2pt);
\end{tikzpicture}
&
\begin{tikzpicture}
\draw[help lines] (0,0) grid (3,2);
\draw [dotted](0,2) arc (90 :0 :2);
\draw [dotted](0,0) --(2,2);
\fill (45:2cm) circle (2pt);
\end{tikzpicture}
\\ \hline 
\BS{fill} (\RDD{canvas polar cs}:\RDD{angle}=45,\RDD{radius}=2cm) circle (2pt);
&
\BS{fill} {\color{blue}(45:2cm)} circle (2pt);
\\ \hline 
\end{tabular} 

\bigskip
\begin{tabular}{|c|} \hline  
\begin{tikzpicture}
\draw[help lines] (0,0) grid (3,2);
\draw [dotted](0,2) arc (90 :0 :3 and 2);
\draw [dotted](0,0) --(3,2);
\fill (canvas polar cs:angle=45,x radius=3cm,y radius=2cm) circle (2pt);
\end{tikzpicture}
\\ \hline  
\BS{fill} (canvas polar cs:angle=45,\RDD{x radius}=3cm,\RDD{y radius}=2cm) circle (2pt);
\\ \hline 
\end{tabular}


\SbSbSSCT{Système de coordonnées  xyz}{xyz coordinates}

\noindent


\begin{tabular}{|c|c|c|} \hline 
\begin{tikzpicture}[->]
\draw (0,0) -- (xyz cs:x=1);
\draw[red] (0,0) -- (xyz cs:y=1);
\draw[magenta] (0,0) -- (xyz cs:z=1);
\end{tikzpicture}
&
\begin{tikzpicture}[->]
\draw (0,0) -- (1,0,0);
\draw[red]  (0,0) -- (0,1,0);
\draw[magenta]  (0,0) -- (0,0,1);
\end{tikzpicture}
\\ \hline 
\BS{draw} (0,0) - - (\RDD{xyz cs}:x=1); & \BS{draw}  (0,0) - - (1,0,0); \\
\BS{draw}[red]  (0,0) - - (\RDD{xyz cs}:y=1); &  \BS{draw}[red] (0,0) - - (0,1,0); \\
\BS{draw}[magenta]  (0,0) - - (\RDD{xyz cs}:z=1); &  \BS{draw}[magenta]   (0,0) - - (0,0,1); 
\\ \hline 

\end{tabular} 

 
\newpage

\SbSbSSCT{Coordinate system xyz polar}{Coordinate system xyz polar}

\noindent

\begin{tabular}{|c|c|c|} \hline
\TFRGB{Explicite}{explicit}  & \TFRGB{Implicite}{implicit}
\\ \hline
\begin{tikzpicture}
\draw[help lines] (0,0) grid (3,2);
\draw [dotted](0,2) arc (90 :0 :2);
\draw [dotted](0,0) --(2,2);
\fill (xyz polar cs:angle=45,radius=2) circle (2pt);
\end{tikzpicture}
&
\begin{tikzpicture}
\draw[help lines] (0,0) grid (3,2);
\draw [dotted](0,2) arc (90 :0 :2);
\draw [dotted](0,0) --(2,2);
\fill (45:2) circle (2pt);
\end{tikzpicture}
\\ \hline 
\BS{fill} (\RDD{xyz polar cs}:\RDD{angle}=45,\RDD{radius}=2) circle (2pt);
&
\BS{fill} {\color{blue}(45:2cm)} circle (2pt);
\\ \hline 
\end{tabular} 

\bigskip
\begin{tabular}{|c|} \hline  
\begin{tikzpicture}
\draw[help lines] (0,0) grid (3,2);
\draw [dotted](0,2) arc (90 :0 :3 and 2);
\draw [dotted](0,0) --(3,2);
\fill (xyz polar cs:angle=45,x radius=3,y radius=2) circle (2pt);
\end{tikzpicture}
\\ \hline  
\BS{fill} (xyz polar cs:angle=45,\RDD{x radius}=3,\RDD{y radius}=2) circle (2pt);
\\ \hline 
\end{tabular} 

\bigskip

\begin{tabular}{|c|c|c|} \hline
\multicolumn{2}{|c|}{\BS{begin}\AC{tikzpicture}{\color{red}[x=1.5cm,y=1cm]} }
\\ \hline
\begin{tikzpicture}[x=1.5cm,y=1cm]
\draw[help lines] (0,0) grid (3,2);
\draw [dotted](0,2) arc (90 :0 :2);
\draw [dotted](0,0) --(2,2);
\fill (xyz polar cs:angle=45,radius=2) circle (2pt);
\end{tikzpicture}
&
\begin{tikzpicture}[x=1.5cm,y=1cm]
\draw[help lines] (0,0) grid (3,2);
\draw [dotted](0,2) arc (90 :0 :2);
\draw [dotted](0,0) --(2,2);
\fill (45:2) circle (2pt);
\end{tikzpicture}
\\ \hline 
\BS{fill} (\RDD{xyz polar cs}:\RDD{angle}=45,\RDD{radius}=2) circle (2pt);
&
\BS{fill} {\color{blue}(45:2cm)} circle (2pt);
\\ \hline 
\end{tabular} 
\bigskip

\begin{tabular}{|c|c|c|} \hline
\multicolumn{2}{|c|}{\BS{begin}\AC{tikzpicture}{\color{red}[x=\AC{(0cm,1cm)},y=\AC{(-1cm,0cm)}]} }
\\ \hline
\begin{tikzpicture}[x={(0cm,1cm)},y={(-1cm,0cm)}]
\draw[help lines] (0,0) grid (3,2);
\draw [dotted](0,2) arc (90 :0 :2);
\draw [dotted](0,0) --(2,2);
\fill (xyz polar cs:angle=45,radius=2) circle (2pt);
\end{tikzpicture}
&
\begin{tikzpicture}[x={(0cm,1cm)},y={(-1cm,0cm)}]
\draw[help lines] (0,0) grid (3,2);
\draw [dotted](0,2) arc (90 :0 :2);
\draw [dotted](0,0) --(2,2);
\fill (45:2) circle (2pt);
\end{tikzpicture}
\\ \hline 
\BS{fill} (\RDD{xyz polar cs}:\RDD{angle}=45,\RDD{radius}=2) circle (2pt);
&
\BS{fill} {\color{blue}(45:2cm)} circle (2pt);
\\ \hline 
\end{tabular} 

\SbSbSSCT{Coordonnées barycentriques}{Barycentric coordinates}

\begin{center}
\RRR{13-2-2}
\end{center}

\begin{tabular}{|c|c|c|} \hline
\multicolumn{3}{|c|}{  \BS{node} [circle,fill=red!20] at (\RDD{barycentric cs}:A=0.6,B=0.3 ) \AC{X};   }\\ 
\hline
\begin{tikzpicture}[scale=.6]
\draw[help lines] (0,0) grid (4,4);
\node[circle,fill=green!20,] (A) at (0,0) {A};
\node[circle,fill=green!20,] (B) at (4,0) {B};
\node[circle,fill=red!20] at (barycentric cs:A=0.3,B=0.3 ) {X};
\end{tikzpicture}
&
\begin{tikzpicture}[scale=.6]
\draw[help lines] (0,0) grid (4,4);
\node[circle,fill=green!20,] (A) at (0,0) {A};
\node[circle,fill=green!20,] (B) at (4,0) {B};
\node[circle,fill=green!20,] (C) at (4,4) {C};
\node[circle,fill=red!20] at (barycentric cs:A=0.4,B=0.4 ,C=.4) {X};
\end{tikzpicture}
&
\begin{tikzpicture}[scale=.6]
\draw[help lines] (0,0) grid (4,4);
\node[circle,fill=green!20,] (A) at (0,0) {A};
\node[circle,fill=green!20,] (B) at (4,0) {B};
\node[circle,fill=green!20,] (C) at (1,4) {C};
\node[circle,fill=green!20,] (D) at (4,4) {D};
\node[circle,fill=red!20] at (barycentric cs:A=0.5,B=0.5,C=.5,D=.5 ) {X};
\end{tikzpicture}
\\ \hline
A=0.3,B=0.3 & A=0.4,B=0.4 ,C=.4 & A=0.5,B=0.5,C=.5,D=.5 
\\ \hline
\begin{tikzpicture}[scale=.6]
\draw[help lines] (0,0) grid (4,4);
\node[circle,fill=green!20,] (A) at (0,0) {A};
\node[circle,fill=green!20,] (B) at (4,0) {B};
\node[circle,fill=red!20] at (barycentric cs:A=0.6,B=0.3 ) {X};
\end{tikzpicture}
&
\begin{tikzpicture}[scale=.6]
\draw[help lines] (0,0) grid (4,4);
\node[circle,fill=green!20,] (A) at (0,0) {A};
\node[circle,fill=green!20,] (B) at (4,0) {B};
\node[circle,fill=green!20,] (C) at (4,4) {C};
\node[circle,fill=red!20] at (barycentric cs:A=0.2,B=0.4 ,C=.6) {X};
\end{tikzpicture}
&
\begin{tikzpicture}[scale=.6]
\draw[help lines] (0,0) grid (4,4);
\node[circle,fill=green!20,] (A) at (0,0) {A};
\node[circle,fill=green!20,] (B) at (4,0) {B};
\node[circle,fill=green!20,] (C) at (1,4) {C};
\node[circle,fill=green!20,] (D) at (4,4) {D};
\node[circle,fill=red!20] at (barycentric cs:A=0.2,B=0.4,C=.6,D=.8 ) {X};
\end{tikzpicture}
\\ \hline
A=0.6,B=0.3 & A=0.2,B=0.4 ,C=.6 & A=0.2,B=0.4,C=.6,D=.8
\\ \hline
\end{tabular}

\SbSbSSCT{Coordonnées nominatives : n\oe ud}{Named coordinates: nodes}

\begin{center}
\RRR{13-2-3}
\end{center}

\begin{tabular}{|c|c|} \hline  
\begin{tikzpicture}[blue,very thick,baseline=1cm]
\draw[help lines] (0,0) grid (3,3);
\coordinate (centre) at (1.5,1.5) ;
\coordinate (A) at (.5,.5) ;
\coordinate (B) at (2.5,2.5) ;
\fill (centre) circle (3pt);
\draw[red] (A) rectangle (B) ;
\end{tikzpicture}
&  
\parbox[c]{8cm}{
\BSS{coordinate} {\color{blue}(centre)} at(1.5,1.5) ; \\
\BSS{coordinate} {\color{blue}(A)} at (.5,.5) ;\\
\BSS{coordinate} {\color{blue}(B)} at  (2.5,2.5) ;\\
\\
\BS{fill} {\color{blue}(centre)} circle (3pt);\\
\BS{draw}[red] {\color{blue}(A)} rectangle {\color{blue}(B)} ;\\
}
\\ \hline 
\end{tabular} 


\TFRGB{voir aussi}{see also} page \pageref{noeuds}


\SbSbSSCT{Coordonnées relatives à un noeud}{Coordinates relative to a node}

\noindent

\begin{tabular}{|c|c|c|c|} \hline
\multicolumn{4}{|l|}{  \BS{node} [draw,fill=green!20,] (A) at (1,1) \AC{\BS{huge}  noeud}; }\\ 
\multicolumn{4}{|l|}{  \BS{fill}[red] (\RDD{node cs}:\RDD{name}=A,\RDD{anchor}=south) circle (3pt);   }\\ 
\hline

\begin{tikzpicture}
\draw[help lines] (0,0) grid (2,2);
\node[draw,fill=green!20,] (A) at (1,1) {\huge noeud};
\fill[red] (node cs:name=A,anchor=south) circle (3pt);
\end{tikzpicture}
&
\begin{tikzpicture}
\draw[help lines] (0,0) grid (2,2);
\node[draw,fill=green!20,] (A) at (1,1) {\huge noeud};
\fill[red] (node cs:name=A,anchor=west) circle (3pt);
\end{tikzpicture}
&
\begin{tikzpicture}
\draw[help lines] (0,0) grid (2,2);
\node[draw,fill=green!20,] (A) at (1,1) {\huge noeud};
\fill[red] (node cs:name=A,anchor=north) circle (3pt);
\end{tikzpicture}
&
\begin{tikzpicture}
\draw[help lines] (0,0) grid (2,2);
\node[draw,fill=green!20,] (A) at (1,1) {\huge noeud};
\fill[red] (node cs:name=A,anchor=east) circle (3pt);
\end{tikzpicture}
\\ \hline
name=A,anchor=south & name=A,anchor=west & name=A,anchor=north & name=A,anchor=east
\\ \hline
\end{tabular}

\bigskip

\begin{tabular}{|c|c|c|c|} \hline
\multicolumn{4}{|l|}{  \BS{node} [draw,fill=green!20,] \blll{(A)} at (1,1) \AC{\BS{huge}  noeud}; }\\ 
\multicolumn{4}{|l|}{  \BS{fill}[red] (\blll{A}.south) circle (3pt);   }\\ 
\hline

\begin{tikzpicture}
\draw[help lines] (0,0) grid (2,2);
\node[draw,fill=green!20,] (A) at (1,1) {\huge noeud};
\fill[red] (A.south) circle (3pt);
\end{tikzpicture}
&
\begin{tikzpicture}
\draw[help lines] (0,0) grid (2,2);
\node[draw,fill=green!20,] (A) at (1,1) {\huge noeud};
\fill[red] (A.west) circle (3pt);
\end{tikzpicture}
&
\begin{tikzpicture}
\draw[help lines] (0,0) grid (2,2);
\node[draw,fill=green!20,] (A) at (1,1) {\huge noeud};
\fill[red] (A.north) circle (3pt);
\end{tikzpicture}
&
\begin{tikzpicture}
\draw[help lines] (0,0) grid (2,2);
\node[draw,fill=green!20,] (A) at (1,1) {\huge noeud};
\fill[red] (A.east) circle (3pt);
\end{tikzpicture}
\\ \hline
A.south & A.west & A.north & A.east
\\ \hline
\end{tabular}



\bigskip
\begin{tabular}{|c|c|c|c|} \hline
\multicolumn{4}{|c|}{  \BS{fill}[red] (node cs:\RDD{name}=A,\RDD{angle}=0) circle (3pt);  }\\ 
\hline

\begin{tikzpicture}
\draw[help lines] (0,0) grid (2,2);
\node[draw,fill=green!20,] (A) at (1,1) {\huge noeud};
\fill[red] (node cs:name=A,angle=0) circle (3pt);
\end{tikzpicture}
&
\begin{tikzpicture}
\draw[help lines] (0,0) grid (2,2);
\node[draw,fill=green!20,] (A) at (1,1) {\huge noeud};
\fill[red] (node cs:name=A,angle=-30) circle (3pt);
\end{tikzpicture}
&
\begin{tikzpicture}
\draw[help lines] (0,0) grid (2,2);
\node[draw,fill=green!20,] (A) at (1,1) {\huge noeud};
\fill[red] (node cs:name=A,angle=-90) circle (3pt);
\end{tikzpicture}
&
\begin{tikzpicture}
\draw[help lines] (0,0) grid (2,2);
\node[draw,fill=green!20,] (A) at (1,1) {\huge noeud};
\fill[red] (node cs:name=A,angle=-150) circle (3pt);
\end{tikzpicture}
\\ \hline
name=A,angle=0 & name=A,angle=-30 & nname=A,angle=-90 & name=A,angle=-150
\\ \hline
\end{tabular}

\bigskip


\begin{tabular}{|c|c|c|c|} \hline
\multicolumn{4}{|c|}{  \BS{fill}[red] (A.0) circle (3pt);  }\\ 
\hline

\begin{tikzpicture}
\draw[help lines] (0,0) grid (2,2);
\node[draw,fill=green!20,] (A) at (1,1) {\huge noeud};
\fill[red] (A.0) circle (3pt);
\end{tikzpicture}
&
\begin{tikzpicture}
\draw[help lines] (0,0) grid (2,2);
\node[draw,fill=green!20,] (A) at (1,1) {\huge noeud};
\fill[red] (A.-30) circle (3pt);
\end{tikzpicture}
&
\begin{tikzpicture}
\draw[help lines] (0,0) grid (2,2);
\node[draw,fill=green!20,] (A) at (1,1) {\huge noeud};
\fill[red] (A.-90) circle (3pt);
\end{tikzpicture}
&
\begin{tikzpicture}
\draw[help lines] (0,0) grid (2,2);
\node[draw,fill=green!20,] (A) at (1,1) {\huge noeud};
\fill[red] (A.-150) circle (3pt);
\end{tikzpicture}
\\ \hline
A.0 & A.-30 & A.-90 & A.-150
\\ \hline
\end{tabular}

\TFRGB{voir aussi}{see also} page \pageref{nomnoeud}


\newpage

\SbSbSSCT{Coordonnées relatives à deux points}{Coordinates relative to two points}
\begin{center}
\RRR{13-3-1}
\end{center}

\begin{tabular}{|c|c|} \hline
\multicolumn{2}{|c|}{  \BS{node} [circle,fill=red!20] at (1,1 {\color{red}|-} 3,3) \AC{X}   }\\ 
\hline
\begin{tikzpicture}
\draw[help lines] (0,0) grid (4,4);
\node[circle,fill=green!20,] (A) at (1,1) {A};
\node[circle,fill=green!20,] (B) at (3,3) {B};
\node[circle,fill=red!20] at (1,1 |- 3,3) {X};
\end{tikzpicture}
&
\begin{tikzpicture}
\draw[help lines] (0,0) grid (4,4);
\node[circle,fill=green!20,] (A) at (1,1) {A};
\node[circle,fill=green!20,] (B) at (3,3) {B};
\node[circle,fill=red!20] at (1,1 -| 3,3) {X};
\end{tikzpicture}
\\ \hline
at (1,1 {\color{red}|-} 3,3)
&
at (1,1 {\color{red}-|} 3,3)
\\ \hline
\end{tabular}



\SbSbSSCT{Coordonnée relative à une intersection}{Coordinates relative to an intersection}
\begin{center}
\RRR{13-3-2}
\end{center}

 \maboite{\BS{usetikzlibrary}\AC{intersections}}
\label{lib-intersections}


\begin{tabular}{|c|c|c|c|} \hline 
\multicolumn{4}{|l|}{  \BS{draw} [\RDD{name path}=XXX] (2,1) circle  (1cm);   }\\ 
\multicolumn{4}{|l|}{  \BS{draw} [\RDD{name path}=YYY] (0.5,0.5) rectangle +(3,1);   }\\ 
\multicolumn{4}{|l|}{ \BS{fill} [red,\RDD{ name intersections}=\AC{of=xxx and YYY}]
(\RDD{intersection}-1) circle (2pt)   }\\ 
\hline 
\begin{tikzpicture}[scale=.8]
\draw [help lines] grid (4,2);
\draw [name path=XXX] (2,1) circle  (1cm);
\draw [name path=YYY] (0.5,0.5) rectangle +(3,1);
\fill [red, name intersections={of=XXX and YYY}]
(intersection-1) circle (2pt)  ;
\end{tikzpicture}
& 
\begin{tikzpicture}[scale=.8]
\draw [help lines] grid (4,2);
\draw [name path=XXX] (2,1) circle  (1cm);
\draw [name path=YYY] (0.5,0.5) rectangle +(3,1);
\fill [red, name intersections={of=XXX and YYY}] (intersection-2) circle (2pt) ;
\end{tikzpicture} 
&  
\begin{tikzpicture}[scale=.8]
\draw [help lines] grid (4,2);
\draw [name path=XXX] (2,1) circle  (1cm);
\draw [name path=YYY] (0.5,0.5) rectangle +(3,1);
\fill [red, name intersections={of=XXX and YYY}] (intersection-3) circle (2pt) ;
\end{tikzpicture}
&  
\begin{tikzpicture}[scale=.8]
\draw [help lines] grid (4,2);
\draw [name path=XXX] (2,1) circle  (1cm);
\draw [name path=YYY] (0.5,0.5) rectangle +(3,1);
\fill [red, name intersections={of=XXX and YYY}] (intersection-4) circle (2pt) ;
\end{tikzpicture}
\\ 
\hline intersection-1 & intersection-2 &intersection-3  & intersection-4 \\ 
\hline 
\end{tabular} 

\bigskip

\begin{tabular}{|c|} \hline  
\BS{fill} [red, name intersections=\AC{of=XXX and YYY}] \\
(intersection-1) circle (2pt) {\color{red} node[black,above right] \AC{point a}} ;
\\ \hline  
\begin{tikzpicture}
\draw [help lines] grid (4,2);
\draw [name path=XXX] (2,1) circle  (1cm);
\draw [name path=YYY] (0.5,0.5) rectangle +(3,1);
\fill [red, name intersections={of=XXX and YYY}]
(intersection-1) circle (2pt) node[black,above right] {point a} ;
\end{tikzpicture} 
\\ \hline 
\end{tabular} 

\bigskip

\begin{tabular}{|c|} \hline 
\BS{fill} [red, name intersections=\AC{of=XXX and YYY, \RDD{name}=ZZZ}]; \\
\BS{draw} [red] (ZZZ-1) - - (ZZZ-3); \BS{draw} [green] (ZZZ-2) - - (ZZZ-4);
\\ \hline  
\begin{tikzpicture}
\draw [help lines] grid (4,2);
\draw [name path=XXX] (2,1) circle  (1cm);
\draw [name path=YYY] (0.5,0.5) rectangle +(3,1);
\fill [red, name intersections={of=XXX and YYY, name=ZZZ}];
\draw [red] (ZZZ-1) -- (ZZZ-3);
\draw [green] (ZZZ-2) -- (ZZZ-4);
\end{tikzpicture}
\\ \hline 
\end{tabular} 

\bigskip
\begin{tabular}{|c|} \hline  
\BS{fill} [red, name intersections=\AC{of=XXX and YYY , \RDD{by}=\AC{a,b,c,d}}]; \\
\BS{draw} [red] (a) - - (c); \hspace{1cm} \BS{draw} [green] (b) - - (d);
\\ \hline   
\begin{tikzpicture}
\draw [help lines] grid (4,2);
\draw [name path=XXX] (2,1) circle  (1cm);
\draw [name path=YYY] (0.5,0.5) rectangle +(3,1);
\fill [red, name intersections={of=XXX and YYY, by={a,b,c,d}}];
\draw [red] (a) -- (c);
\draw [green] (b) -- (d);
\end{tikzpicture}
\\ \hline 
\end{tabular} 

\bigskip

\begin{tabular}{|c|} \hline  
\BS{fill} [name intersections=\AC{of=XXX and YYY, name=i, \RDD{total}=\BS{t}}] [red] \\
\BS{foreach} \BS{s} in \AC{1,...,\BS{t}} \AC{(i-\BS{s}) circle (2pt) node[black,above right] \AC{\BS{s}}}
\\ \hline  
\begin{tikzpicture}
\draw [help lines] grid (4,2);
\draw [name path=XXX] (2,1) circle  (1cm);
\draw [name path=YYY] (0.5,0.5) rectangle +(3,1);
\fill [name intersections={of=XXX and YYY , name=i, total=\t}]
[red]
\foreach \s in {1,...,\t}{(i-\s) circle (2pt) node[black,above right] {\s}};
\end{tikzpicture}
\\ \hline 
\end{tabular} 



\newpage

\SbSbSSCT{Position calculée avec le module  \og  pgfmath \fg}{Calculated positions with  \og  pgfmath \fg }

\begin{center}
\RRR{13-2-1}
\end{center}

\TFRGB{Ce module est chargé automatiquement avec le module Tikz}{Package automatically loaded with Tikz} 

\begin{tabular}{|c|} \hline 
\begin{tikzpicture}
\draw[help lines] (0,0) grid (4,2);
\fill [red] (canvas cs:x=2cm+1.5cm,y=1.5cm-1cm) circle (3pt);
\fill [blue] (2cm,1.5cm) circle (3pt);
\draw[dashed] (2,1.5) -| (3.5,.5);
\end{tikzpicture}
\\ \hline 
\emph{\TFRGB{Explicite}{explicit}} 
 : \BS{fill} [red] (\RDD{canvas cs}:x=2cm+1.5cm,y=1.5cm-1cm) circle (3pt);
 \\  \hline 
\emph{\TFRGB{Implicite}{implicit}} :  \BS{fill} [red] {\color{red}(2cm+1.5cm,1.5cm-1cm)} circle (3pt);
\\ \hline 
\end{tabular} 

\bigskip
\begin{tabular}{|c|c|c|} \hline 
\begin{tikzpicture}[baseline=0pt]
\draw[help lines] (0,0) grid (4,4);
 \draw[dashed] (2,2) circle (2);
\fill[red](2+ 2*cos 30,2+2*sin 30) circle (3pt);
\fill[magenta](2+ 2*cos{(120)},2+2*sin{(120)}) circle (3pt);
\end{tikzpicture}
&
\parbox[c]{8cm}{
 \BS{draw}[dashed] (2,2) circle (2);\\
 \smallskip
 \BS{fill} [red]{\color{red}(2+ 2*cos 30 , 2+2*sin 30)} circle (3pt);\\
  \smallskip
 \BS{fill}[magenta] {\color{red}(2+2*cos\AC{(120)} , 2+2*sin\AC{(120)})} circle (3pt); 
 }
\\ \hline 
\end{tabular} 

\SbSbSSCT{Position calculée avec \og library calc \fg}{Calculated positions with \og  calc  library calc \fg}

\begin{center}
\RRR{13-5}
\end{center}
\label{lib-calc}

 \maboite{\BS{usetikzlibrary}\AC{calc}}
 
\begin{tabular}{|c|c|} \hline  
\begin{tikzpicture}[baseline=0pt]
\draw [help lines] (0,0) grid (3,2);
\node (a) at (1,1) {A};
\fill [red] ($(a) + 2/3*(1cm,0)$) circle (2pt);
\fill [red] ($(a) + 4/3*(1cm,0)$) circle (2pt);
\end{tikzpicture}
&
\parbox{8cm}{
\BS{node} (a) at (1,1) \AC{A}; \\
\BS{fill} [red] {\color{red} (\$(a) + 2/3*(1cm,0)\$)} circle (2pt); \\
\BS{fill} [red] {\color{red}(\$(a) + 4/3*(1cm,0)\$)} circle (2pt); \\
}
\\ 
\hline 
\end{tabular} 

\SbSbSSCT{Tangentes avec \og library calc \fg}{Tangents with  \og calc library  \fg}

\begin{center}
\RRR{13-2-4}
\end{center}

\begin{tabular}{|c|c|} \hline 
\multicolumn{2}{|l|}{\BS{node}[fill=green!20] (a) at (3,1.5) \AC{A}; } \\
\multicolumn{2}{|l|}{\BS{fill}[red] (\RDD{tangent cs}:\RDD{node}=c,\RDD{point}=\AC{(A)},\RDD{solution}=1);  }\\ 
\hline
\begin{tikzpicture}
\draw[help lines] (0,0) grid (4,2);
\node[fill=green!20] (A) at (3,1.5) {A};
\node [circle,draw] (c) at (1,1) [minimum size=1.5cm] {$c$};
\draw[red,dashed] (A) - -(tangent cs:node=c,point={(A)},solution=1) ;
\draw[red,dashed] (1,1) - -(tangent cs:node=c,point={(3,1.5)},solution=1) ;
\fill[red] (tangent cs:node=c,point={(A)},solution=1) circle (3pt);
\end{tikzpicture}
&
\begin{tikzpicture}
\draw[help lines] (0,0) grid (4,2);
\node[fill=green!20] (A) at (3,1.5) {A};
\node [circle,draw] (c) at (1,1) [minimum size=1.5cm] {$c$};
\draw[red,dashed] (A) - -(tangent cs:node=c,point={(A)},solution=2) ;
\draw[red,dashed] (1,1) - -(tangent cs:node=c,point={(A)},solution=2) ;
\fill[red] (tangent cs:node=c,point={(A)},solution=2) circle (3pt);
\end{tikzpicture}
\\ \hline
\RDD{solution}=1 & \RDD{solution}=2
\\ \hline
\end{tabular} 

\newpage

\SbSbSSCT{Point à pourcentage donné }{Percentage position }

\begin{center}
\RRR{13-5-3}
\end{center}


\begin{tabular}{|c|c|} \hline  
\multicolumn{2}{|c|}{\BS{fill}[red] ({\color{red}\$(0,1)!.25!(4,1)\$}) circle (4pt); } \\  \hline  

\begin{tikzpicture}
\draw [help lines] (0,0) grid (4,2);
\draw [line width= 3pt] (0,1) -- (4,1);
\fill[red] ($(0,1)!.25!(4,1)$) circle (4pt);
\end{tikzpicture}
&  
\begin{tikzpicture}
\draw [help lines] (0,0) grid (4,2);
\draw [line width= 3pt] (0,1) -- (4,1);
\fill[red] ($(0,1)!.75!(4,1)$) circle (4pt);
\end{tikzpicture}
\\ \hline (0,1)!{\color{red}0.25}!(4,1) & (0,1)!{\color{red}0.75}!(4,1) \\ 
\hline 
\end{tabular} 

\bigskip

\begin{tabular}{|c|} \hline  
\begin{tikzpicture}
\draw [help lines] (0,0) grid (4,3);
\draw [line width=2pt ](0,2) -- (4,2);
\draw[red] ($(0,2)!.75!(4,2)$) -- (0,0);
\fill[red] ($(0,2)!.75!(4,2)!.66!(0,0)$) circle (4pt);
\end{tikzpicture}
\\ \hline 
\BS{fill}[red] (\${\color{blue}(0,2)!0.75!(4,2)}!{\color{red}0.66!(0,0)}\$) circle (2pt);
\\ \hline 
\end{tabular} 


\SbSbSSCT{Point à distance donnée}{Position at a given distance }

\begin{center}
\RRR{13-5-4}
\end{center}

\begin{tabular}{|c|c|} \hline  
\multicolumn{2}{|c|}{\BS{fill}[red] ({\color{red}\$(0,1)!1.5cm!(4,1)\$}) circle (4pt); } \\  \hline  

\begin{tikzpicture}
\draw [help lines] (0,0) grid (4,2);
\draw [line width= 2pt] (0,1) -- (4,1);
\fill[red] ($(0,1)!1.5cm!(4,1)$) circle (4pt);
\end{tikzpicture}
&  
\begin{tikzpicture}
\draw [help lines] (0,0) grid (4,2);
\draw [line width= 2pt] (0,1) -- (4,1);
\fill[red] ($(0,1)!3cm!(4,1)$) circle (4pt);
\end{tikzpicture}
\\ \hline (0,1)!{\color{red}1.5cm}!(4,1) & (0,1)!{\color{red}3cm}!(4,1) \\ 
\hline 
\end{tabular} 

\bigskip

\begin{tabular}{|c|} \hline  
\begin{tikzpicture}
\draw [help lines] (0,0) grid (4,4);
\coordinate (a) at (1,0);
\coordinate (b) at (4,1);
\draw [line width= 3pt] (0,0) -- (4,1);
\draw [line width= 2pt,red](2,.5) -- ($ (2,.5)!2cm!90:(4,1) $);
\end{tikzpicture}
\\ \hline
\BS{draw} (2,.05) - - (\$ (2,0.5)!{\color{red}2cm!90:(4,1)} \$);
\\ \hline 
\end{tabular} 

\newpage

\SbSbSSCT{Coordonnées relatives}{Relative coordinates}


\Par{Cartésienne}{Cartesian coordinates}

\begin{center}
\RRR{13-4-1}
\end{center}

\begin{tabular}{|c|c|c|} \hline  
\TFRGB{relative à l'origine}{relative to the origin}  & \TFRGB{relative à une position}{relative to a position}  &  \TFRGB{relative à la dernière position}{relative to the last position}   
\\ \hline  
 
\begin{tikzpicture}
\draw[help lines] (0,-1) grid (3,1); 
 \draw[blue,very thick] (0,0) -- (1,0) - - (2,1) - - (2,-1);
 \fill[red] (0,0) circle (4pt);
\end{tikzpicture}
&
\begin{tikzpicture} %[scale=.8]
\draw[help lines] (0,-1) grid (4,1);
 \draw[blue,very thick] (0,0) - - (1,0) -- +(2,1) -- +(2,-1) ; %–- +(2,-1) ;
 \fill[red] (1,0) circle (4pt);
\end{tikzpicture}
&
\begin{tikzpicture} %[scale=.8]
\draw[help lines] (0,-1) grid (5,1);  
 \draw[blue,very thick] (0,0) -- (1,0)  - - ++(2,1) - - ++(2,-1);
 \fill[red] (1,0) circle (4pt);
 \fill[red] (3,1) circle (4pt);
\end{tikzpicture}
\\ \hline 
\tikz \fill node[fill=green!20,inner sep=0pt]{(0,0)}; - - (1,0) &
 (0,0) - - \tikz \fill node[fill=green!20,inner sep=0pt]{(1,0)};  & (0,0) - - \tikz \fill node[fill=green!20,inner sep=0pt]{(1,0)}; \\
 - - (2,1) - - (2,-1)  &
   - - +(2,1) - - +(2,-1) & - - ++\tikz \fill node[fill=green!20,inner sep=0pt]{(2,1)}; - - ++(2,-1)
\\ \hline 
\end{tabular} 

\bigskip

\begin{tabular}{|c|c|c|} \hline  
\begin{tikzpicture} [scale=.5]
\draw[help lines] (0,-1) grid (6,6);
 \draw[red,dotted,line width=2pt] (0,0) rectangle (2,2) ;
  \draw[green,dotted,line width=2pt] (0,0) rectangle (3,3) ;  
 \draw[blue,line width=2pt] (0,0) rectangle (1,1)  rectangle (2,2) rectangle (3,3);

\end{tikzpicture}

&  
\begin{tikzpicture} [scale=.5]
\draw[help lines] (0,-1) grid (6,6); 
  \draw[green,dotted,line width=2pt] (1,1) rectangle (4,4) ;   
 \draw[blue,line width=2pt] (0,0) rectangle (1,1)  rectangle +(2,2) rectangle +(3,3);
    \fill[red] (1,1) circle (4pt);
\end{tikzpicture}
&  
\begin{tikzpicture} [scale=.5]
\draw[help lines] (0,-1) grid (6,6);  
 \draw[blue,line width=2pt] (0,0) rectangle (1,1)  rectangle ++(2,2) rectangle ++(3,3);
    \fill[red] (1,1) circle (4pt);
     \fill[green] (3,3) circle (4pt); 
\end{tikzpicture}
\\ 
\hline 
\BS{draw} (0,0) rectangle (1,1)   &
\BS{draw} (0,0) rectangle (1,1)   & 
\BS{draw} (0,0) rectangle (1,1)  \\
rectangle (2,2) rectangle (3,3);  &
rectangle +(2,2) rectangle +(3,3);  &
rectangle ++(2,2) rectangle ++(3,3); \\
\hline 
\end{tabular}


\Par{Polaire }{Polar} {}

\bigskip


\noindent

\begin{tabular}{|c|c|c|c|} \hline
\TFRGB{relative à l'origine}{relative to the origin}  & \TFRGB{relative à une position}{relative to a position}  &  \TFRGB{relative à la dernière position}{relative to the last position}   
\\ \hline    
\begin{tikzpicture} %[scale=.8] 
\draw[help lines] (0,-1) grid (3,1);
 \fill[red] (0:0) circle (4pt);
 \draw[blue,very thick] (0:0)-- (0:1) -- (30:2) -- (-30:2);
\end{tikzpicture}
&
\begin{tikzpicture} %[scale=.8] 
\draw[help lines] (0,-1) grid (4,1);
 \fill[red] (1,0) circle (4pt);
 \draw[blue,very thick] (0:0) -- (0:1) -- +(30:2) -- +(-30:2);
\end{tikzpicture}
&
\begin{tikzpicture} %[scale=.8] 
\draw[help lines] (0,-1) grid (5,1);
 \fill[red] (1,0) circle (4pt);
 \fill[red] (2.732,1) circle (4pt);
 \draw[blue,very thick] (0:0)-- (0:1) -- ++(30:2) -- ++(-30:2);
\end{tikzpicture}
\\ \hline
\tikz \fill node[fill=green!20,inner sep=0pt] {(0:0)}; - - (0:1)&
 (0:0) - - \tikz \fill node[fill=green!20,inner sep=0pt] {(0:1)}; & (0:0)- - \tikz \fill node[fill=green!20,inner sep=0pt] {(0:1)}; \\
 - - (30:2) - - (-30:2)  &  - -  +(30:2) - - +(-30:2) & - -  ++\tikz \fill node[fill=green!20,inner sep=0pt] {(30:2)}; - - ++(-30:2)
\\ \hline 
\end{tabular} 

%\subsubsection{coordonnée relative en polaire}
\Par{coordonnée relative en polaire}{Relative polar coordinate}

\begin{center}
\RRR{13-4-2}
\end{center}
\bigskip

\begin{tabular}{|c|c|} \hline 
\multicolumn{2}{|c|}{ \BS{draw}[blue,very thick] (0,0) -- (2,1) -- ([turn]-45:1cm);}
 \\ \hline
\begin{tikzpicture} %[scale=.8] 
\draw[help lines] (0,0) grid (4,2);
 \draw[dotted] (0,0) -- (4,2);
 \draw[blue,very thick] (0,0) -- (2,1) -- ([turn]-45:1cm);
\end{tikzpicture}
&  
\begin{tikzpicture} %[scale=.8] 
\draw[help lines] (0,0) grid (4,2);
 \draw[dotted] (0,0) -- (4,2);
 \draw[blue,very thick] (0,0) -- (2,1) -- ([turn]45:1cm);
\end{tikzpicture}
\\ \hline ([\RDD{turn}]-45:1cm) & ([\RDD{turn}]45:1cm) \\ 
\hline 
\end{tabular}

\bigskip

\begin{tabular}{|c|c|} \hline  
\begin{tikzpicture}  
\draw[help lines] (-1,0) grid (4,3);
\draw [line width=2pt] (4,0) arc (0 :120 :2)  -- ([turn]90:2cm) ;

\end{tikzpicture}
&  
\begin{tikzpicture} %[scale=.8] 
\draw[help lines] (0,0) grid (4,3);
\draw [line width=2pt]  (0,0) to [bend left] (2,2) --  ([turn]0:2cm);
\fill [red](2,2) circle (4pt);
\end{tikzpicture}
\\ \hline  
\BS{draw} (4,0) arc (0 :120 :2)  - - ([\RDD{turn}]90:2cm) ;
& \BS{draw}  (0,0) to [bend left] (2,2) - -  ([\RDD{turn}]0:2cm); \\

\hline 
\end{tabular} 


%\bigskip 
%
%
%\tikz [delta angle=30, radius=1cm]
%\draw (0,0) arc [start angle=0] -- ([turn]0:1cm)
%arc [start angle=30] -- ([turn]0:1cm)
%arc [start angle=60] -- ([turn]30:1cm);



\bigskip

\begin{tabular}{|c|c|c|} \hline  
\multicolumn{3}{|c|}{ \BS{draw}(1,2)
.. controls ([turn]0:2cm) .. ([turn]-90:2cm); }
\\ \hline
\begin{tikzpicture} %[scale=.8] 
\draw[help lines] (0,0) grid (4,4);
 \draw [line width=2pt] (1,2)
.. controls ([turn]0:2cm) .. ([turn]-90:2cm);
\end{tikzpicture}
&  
\begin{tikzpicture} %[scale=.8] 
\draw[help lines] (0,0) grid (4,4);
 \draw [line width=2pt] (1,2)
.. controls ([turn]30:2cm) .. ([turn]-90:2cm);
\end{tikzpicture}
&  
\begin{tikzpicture} %[scale=.8] 
\draw[help lines] (-2,0) grid (2,4);
 \draw [line width=2pt] (1,2)
.. controls ([turn]0:2cm) .. ([turn]90:2cm);

\end{tikzpicture}
\\ \hline ([turn]0:2cm) .. ([turn]-90:2cm) & ([turn]30:2cm) .. ([turn]-90:2cm) & ([turn]0:2cm) .. ([turn]90:2cm) \\ 
\hline 
\end{tabular} 


\tikzset{every picture/.style=blue,very thick,inner sep=.3333em}

%
%%\newpage
%%
%%\SbSSCT{Le peuple TikZ}{Tikzpeople}

\label{people}

 \maboite{\BS{usepackage}\AC{tikzpeople} \cite {tikzpeople} \footnote{ conflit \BS{usetikzlibrary}\AC{patterns} page \pageref{lib-patterns} : placer cette commande en premier} }

\bigskip
\begin{tabular}{|c|c|}\hline  
\BS{tikz} \BS{node}[\RDD{alice}] at (0,0) {};  &  \tikz \node[alice] at (0,0) {};\\ 
\hline 
\end{tabular} 
 
\SbSbSSCT{Personages disponibles}{available characters}

\noindent



\begin{tabular}{|c|c|c|c|c|c|c|}\hline 
\multicolumn{7}{|c|}{ \BS{tikz} \BS{node}[\RDD{alice},minimum size=1.5cm] at (0,0) {};  }
\\ \hline  
\tikz \node[alice,minimum size=1.5cm] at (0,0) {}; &  
\tikz \node[bob,minimum size=1.5cm] at (0,0) {}; &  
\tikz \node[bride,minimum size=1.5cm] at (0,0) {}; &  
\tikz \node[builder,minimum size=1.5cm] at (0,0) {}; &  
\tikz \node[businessman,minimum size=1.5cm] at (0,0) {}; &  
\tikz \node[charlie,minimum size=1.5cm] at (0,0) {}; &  
\tikz \node[chef,minimum size=1.5cm] at (0,0) {}; 
\\  \hline  
\RDD{alice} & \RDD{bob} & \RDD{bride} & \RDD{builder} & \RDD{businessman} & \RDD{charlie}  & \RDD{chef} 
\\  \hline  
\tikz \node[conductor,minimum size=1.5cm] at (0,0) {}; &  
\tikz \node[cowboy,minimum size=1.5cm] at (0,0) {}; &  
\tikz \node[criminal,minimum size=1.5cm] at (0,0) {}; &  
\tikz \node[dave,minimum size=1.5cm] at (0,0) {}; &  
\tikz \node[graduate,minimum size=1.5cm] at (0,0) {}; &  
\tikz \node[groom,minimum size=1.5cm] at (0,0) {}; & 
\tikz \node[guard,minimum size=1.5cm] at (0,0) {}; 
\\ \hline  
\RDD{conductor} & \RDD{cowboy} & \RDD{criminal} & \RDD{dave} & \RDD{graduate} & \RDD{groom} & \RDD{guard} 
\\ \hline  
\tikz \node[jester,minimum size=1.5cm] at (0,0) {}; &  
%\tikz \node[judge,minimum size=1.5cm] at (0,0) {}; 
&  
\tikz \node[mexican,minimum size=1.5cm] at (0,0) {}; &  
\tikz \node[nun,minimum size=1.5cm] at (0,0) {}; &  
\tikz \node[nurse,minimum size=1.5cm] at (0,0) {}; &  
\tikz \node[physician,minimum size=1.5cm] at (0,0) {}; &  
\tikz \node[pilot,minimum size=1.5cm] at (0,0) {};\\ \hline  
\RDD{jester} &  \RDD{judge} &  \RDD{mexican}  & \RDD{nun} &  \RDD{nurse} & \RDD{physician} 
&  \RDD{pilot}
\\ \hline  

\tikz \node[police,minimum size=1.5cm] at (0,0) {}; &  
\tikz \node[priest,minimum size=1.5cm] at (0,0) {}; &  
\tikz \node[sailor,minimum size=1.5cm] at (0,0) {}; &  
\tikz \node[santa,minimum size=1.5cm] at (0,0) {}; &  
\tikz \node[surgeon,minimum size=1.5cm] at (0,0) {};&  &  \\ 
\hline \RDD{police} & \RDD{priest}  & \RDD{sailor} & \RDD{santa} & \RDD{surgeon} &  &  \\ 
\hline 
\end{tabular} 

\subsubsection{Options}

\noindent

\begin{tabular}{|c|c|c|c|c|}\hline
\multicolumn{5}{|c|}{ \BS{tikz} \BS{node}[businessman,\RDD{evil},minimum size=1.5cm] at (0,0) {};  }
\\ \hline  
\tikz \node[businessman,evil,minimum size=1.5cm] at (0,0) {}; &  
\tikz \node[businessman,female,minimum size=1.5cm] at (0,0) {}; &  
\tikz \node[businessman,good,minimum size=1.5cm] at (0,0) {}; &  
\tikz \node[businessman,mirrored,minimum size=1.5cm] at (0,0) {}; &  
\tikz \node[businessman,monitor,minimum size=1.5cm] at (0,0) {};  
\\  \hline
\RDD{evil} & \RDD{female} & \RDD{good} & \RDD{mirrored} & \RDD{monitor}

\\  \hline 
\end{tabular}

\SbSbSSCT{Point d'ancrage spécifique}{Anchor specific}

\noindent

\begin{tabular}{|c|c|} \hline  
\begin{tikzpicture}[baseline=0pt,blue]
 \node[name=a,shape=bob,minimum
size=1.5cm] {};
 \node at (1.25,.5) [ellipse callout, draw,
callout absolute pointer={(a.mouth)},
font=\tiny] {Hey!};
\end{tikzpicture}
&  
\parbox{12cm}{
\BS{begin}\AC{tikzpicture}[blue] \\
\BS{node}[name=a,shape=bob,minimum size=1.5cm] \AC{};\\
\BS{node} at (1.25,.5) [ellipse callout, draw,
callout absolute pointer\AC{(a.\RDD{mouth})},
font=\BS{tiny}] {Hey!};\\
\BS{end}\AC{tikzpicture} \\
}
\\ \hline 
\end{tabular} 


\SbSbSSCT{Couleurs }{Colors}

\noindent


\begin{tabular}{|c|c|c|c|}\hline
\multicolumn{4}{|c|}{ \BS{tikz} \BS{node}[\blll{alice},\RDD{hair}=red,minimum size=1.5cm] at (0,0) {};  }
\\ \hline    
\tikz \node[alice,hair=red,minimum size=1.5cm] at (0,0) {}; &  
\tikz \node[alice,skin=red,minimum size=1.5cm] at (0,0) {}; &  
\tikz \node[alice,shirt=red,minimum size=1.5cm] at (0,0) {}; &  
\tikz \node[alice,undershirt=red,minimum size=1.5cm] at (0,0) {};  
\\  \hline
\RDD{hair}=red & \RDD{skin}=red & \RDD{shirt}=red & \RDD{details}=red 

\\  \hline 
\end{tabular}

\bigskip
\begin{tabular}{|c|c|c|c|}\hline 
\multicolumn{4}{|c|}{ \BS{tikz} \BS{node}[\blll{bob},\RDD{hair}=red,minimum size=1.5cm] at (0,0) {};  }
\\ \hline   
\tikz \node[bob,hair=red,minimum size=1.5cm] at (0,0) {}; &  
\tikz \node[bob,skin=red,minimum size=1.5cm] at (0,0) {}; &  
\tikz \node[bob,shirt=red,minimum size=1.5cm] at (0,0) {}; &  
\tikz \node[bob,details=red,minimum size=1.5cm] at (0,0) {};  
\\  \hline
\RDD{hair}=red & \RDD{skin}=red & \RDD{shirt}=red & \RDD{details}=red 
\\  \hline 
\end{tabular}

\bigskip
\begin{tabular}{|c|c|c|c|c|}\hline 
\multicolumn{5}{|c|}{ \BS{tikz} \BS{node}[\blll{bride},\RDD{hair}=red,minimum size=1.5cm] at (0,0) {};  }
\\ \hline    
\tikz \node[bride,hair=red,minimum size=1.5cm] at (0,0) {}; &  
\tikz \node[bride,skin=red,minimum size=1.5cm] at (0,0) {}; &  
\tikz \node[bride,shirt=red,minimum size=1.5cm] at (0,0) {}; &  
\tikz \node[bride,pearls=red,minimum size=1.5cm] at (0,0) {}; &
\tikz \node[bride,veil=red,minimum size=1.5cm] at (0,0) {};  
\\  \hline
\RDD{hair}=red & \RDD{skin}=red & \RDD{shirt}=red & \RDD{pearls}=red & \RDD{veil}=red 

\\  \hline 
\end{tabular}

\bigskip
\begin{tabular}{|c|c|c|c|c|}\hline
\multicolumn{5}{|c|}{ \BS{tikz} \BS{node}[\blll{builder},\RDD{hair}=red,minimum size=1.5cm] at (0,0) {};  }
\\ \hline    
\tikz \node[builder,hair=red,minimum size=1.5cm] at (0,0) {}; &  
\tikz \node[builder,skin=red,minimum size=1.5cm] at (0,0) {}; &  
\tikz \node[builder,shirt=red,minimum size=1.5cm] at (0,0) {}; &  
\tikz \node[builder,trousers=red,minimum size=1.5cm] at (0,0) {}; &
\tikz \node[builder,hat=red,minimum size=1.5cm] at (0,0) {};  
\\  \hline
\RDD{hair}=red & \RDD{skin}=red & \RDD{shirt}=red & \RDD{trousers}=red & \RDD{hat}=red   
\\  \hline 
\end{tabular}

\bigskip
\begin{tabular}{|c|c|c|c|c|c|}\hline
\multicolumn{6}{|c|}{ \BS{tikz} \BS{node}[\blll{businessman},\RDD{hair}=red,minimum size=1.5cm] at (0,0) {};  }
\\ \hline    
\tikz \node[businessman,hair=red,minimum size=1.5cm] at (0,0) {}; &  
\tikz \node[businessman,skin=red,minimum size=1.5cm] at (0,0) {}; &  
\tikz \node[businessman,shirt=red,minimum size=1.5cm] at (0,0) {}; &  
\tikz \node[businessman,tie=red,minimum size=1.5cm] at (0,0) {}; &
\tikz \node[businessman,undershirt=red,minimum size=1.5cm] at (0,0) {};  &
\tikz \node[businessman,monogram=red,minimum size=1.5cm] at (0,0) {};
\\  \hline
\RDD{hair}=red & \RDD{skin}=red & \RDD{shirt}=red & \RDD{tie}=red & \RDD{undershirt}=red & \RDD{monogram}=red 
\\  \hline 
\end{tabular}


\bigskip
\begin{tabular}{|c|c|c|c|}\hline
\multicolumn{4}{|c|}{ \BS{tikz} \BS{node}[\blll{charlie},\RDD{hair}=red,minimum size=1.5cm] at (0,0) {};  }
\\ \hline    
\tikz \node[charlie,hair=red,minimum size=1.5cm] at (0,0) {}; &  
\tikz \node[charlie,skin=red,minimum size=1.5cm] at (0,0) {}; &  
\tikz \node[charlie,shirt=red,minimum size=1.5cm] at (0,0) {}; &  
\tikz \node[charlie,buttons=red,minimum size=1.5cm] at (0,0) {}; 

\\  \hline
\RDD{hair}=red & \RDD{skin}=red & \RDD{shirt}=red & \RDD{buttons}=red 
\\  \hline 
\end{tabular}


\bigskip
\begin{tabular}{|c|c|c|c|c|}\hline
\multicolumn{5}{|c|}{ \BS{tikz} \BS{node}[\blll{chef},\RDD{hair}=red,minimum size=1.5cm] at (0,0) {};  }
\\ \hline 
\tikz \node[chef,hair=red,minimum size=1.5cm] at (0,0) {}; &  
\tikz \node[chef,skin=red,minimum size=1.5cm] at (0,0) {}; &  
\tikz \node[chef,shirt=red,minimum size=1.5cm] at (0,0) {}; &  
\tikz \node[chef,hat=red,minimum size=1.5cm] at (0,0) {}; &
\tikz \node[chef,details=red,minimum size=1.5cm] at (0,0) {};  
\\  \hline
\RDD{hair}=red & \RDD{skin}=red & \RDD{shirt}=red & \RDD{hat}=red & \RDD{details}=red 
\\  \hline 
\end{tabular}


\bigskip
\begin{tabular}{|c|c|c|c|c|}\hline
\multicolumn{5}{|c|}{ \BS{tikz} \BS{node}[\blll{conductor},\RDD{hair}=red,minimum size=1.5cm] at (0,0) {};  }
\\ \hline 
\tikz \node[conductor,hair=red,minimum size=1.5cm] at (0,0) {}; &  
\tikz \node[conductor,skin=red,minimum size=1.5cm] at (0,0) {}; &  
\tikz \node[conductor,shirt=red,minimum size=1.5cm] at (0,0) {}; &  \tikz \node[conductor,hat=red,minimum size=1.5cm] at (0,0) {}; &
\tikz \node[conductor,hatshield=red,minimum size=1.5cm] at (0,0) {};  
\\  \hline
\RDD{hair}=red & \RDD{skin}=red & \RDD{shirt}=red & \RDD{hat}=red & \RDD{hatshield}=red  
\\  \hline 
\tikz \node[conductor,undershirt=red,minimum size=1.5cm] at (0,0) {}; &  
\tikz \node[conductor,tie=red,minimum size=1.5cm] at (0,0) {}; &  
\tikz \node[conductor,hatbadge=red,minimum size=1.5cm] at (0,0) {}; &
\tikz \node[conductor,badge=red,minimum size=1.5cm] at (0,0) {}; &
\\  \hline 
\RDD{undershirt}=red &  \RDD{shirt}=red & \RDD{hatbadge}=red & \RDD{badge}=red &
\\  \hline 
\end{tabular}

\bigskip
\begin{tabular}{|c|c|c|c|}\hline
\multicolumn{4}{|c|}{ \BS{tikz} \BS{node}[\blll{cowboy},\RDD{hair}=red,minimum size=1.5cm] at (0,0) {};  }
\\ \hline 
\tikz \node[cowboy,hair=red,minimum size=1.5cm] at (0,0) {}; &  
\tikz \node[cowboy,skin=red,minimum size=1.5cm] at (0,0) {}; &  
\tikz \node[cowboy,shirt=green,minimum size=1.5cm] at (0,0) {}; &  
\tikz \node[cowboy,hat=red,minimum size=1.5cm] at (0,0) {};   
\\  \hline
\RDD{hair}=red & \RDD{skin}=red & \RDD{shirt}=green & \RDD{hat}=red 
\\  \hline 
\tikz \node[cowboy,patches=red,minimum size=1.5cm] at (0,0) {}; &  
\tikz \node[cowboy,tie=green ,minimum size=1.5cm] at (0,0) {}; &  
\tikz \node[cowboy,stitching=red,minimum size=1.5cm] at (0,0) {}; &
\tikz \node[cowboy,vest=red,minimum size=1.5cm] at (0,0) {}; 
\\  \hline 
\RDD{patches}=red &  \RDD{tie}=green & \RDD{stitching}=red & \RDD{vest}=red
\\  \hline 
\end{tabular}


\bigskip
\begin{tabular}{|c|c|c|c|}\hline 
\multicolumn{4}{|c|}{ \BS{tikz} \BS{node}[\blll{criminal},\RDD{hat}=red,minimum size=1.5cm] at (0,0) {};  }
\\ \hline 
\tikz \node[criminal,hat=red,minimum size=1.5cm] at (0,0) {}; &  
\tikz \node[criminal,skin=red,minimum size=1.5cm] at (0,0) {}; &  
\tikz \node[criminal,shirt=red,minimum size=1.5cm] at (0,0) {}; &  
\tikz \node[criminal,details=red,minimum size=1.5cm] at (0,0) {}; 
\\  \hline
\RDD{hat}=red & \RDD{skin}=red & \RDD{shirt}=red & \RDD{details}=red 
\\  \hline 
\end{tabular}

\bigskip
\begin{tabular}{|c|c|c|c|c|}\hline
\multicolumn{5}{|c|}{ \BS{tikz} \BS{node}[\blll{dave},\RDD{hair}=red,minimum size=1.5cm] at (0,0) {};  }
\\ \hline 
\tikz \node[dave,hair=red,minimum size=1.5cm] at (0,0) {}; &  
\tikz \node[dave,skin=red,minimum size=1.5cm] at (0,0) {}; &  
\tikz \node[dave,shirt=red,minimum size=1.5cm] at (0,0) {}; &  
\tikz \node[dave,undershirt=green,minimum size=1.5cm] at (0,0) {}; &
\tikz \node[dave,tie=green,minimum size=1.5cm] at (0,0) {};
\\  \hline
\RDD{hair}=red & \RDD{skin}=red & \RDD{shirt}=red & \RDD{undershirt}=green & \RDD{tie}=green
\\  \hline 
\end{tabular}

\bigskip
\begin{tabular}{|c|c|c|c|c|c|}\hline
\multicolumn{6}{|c|}{ \BS{tikz} \BS{node}[\blll{graduate},\RDD{hair}=red,minimum size=1.5cm] at (0,0) {};  }
\\ \hline 
\tikz \node[graduate,hair=red,minimum size=1.5cm] at (0,0) {}; &  
\tikz \node[graduate,skin=red,minimum size=1.5cm] at (0,0) {}; &  
\tikz \node[graduate,shirt=red,minimum size=1.5cm] at (0,0) {}; &  
\tikz \node[graduate,undershirt=red,minimum size=1.5cm] at (0,0) {}; &
\tikz \node[graduate,stripes=red,minimum size=1.5cm] at (0,0) {};
&
\tikz \node[graduate,hat=red,minimum size=1.5cm] at (0,0) {};
\\  \hline
\RDD{hair}=red & \RDD{skin}=red & \RDD{shirt}=red & \RDD{undershirt}=red & \RDD{stripes}=red & \RDD{hat}=red
\\  \hline 
\end{tabular}

\bigskip
\begin{tabular}{|c|c|c|c|c|c|}\hline
\multicolumn{6}{|c|}{ \BS{tikz} \BS{node}[\blll{groom},\RDD{hair}=red,minimum size=1.5cm] at (0,0) {};  }
\\ \hline
\tikz \node[groom,hair=red,minimum size=1.5cm] at (0,0) {}; &  
\tikz \node[groom,skin=red,minimum size=1.5cm] at (0,0) {}; &  
\tikz \node[groom,shirt=red,minimum size=1.5cm] at (0,0) {}; &  
\tikz \node[groom,undershirt=green,minimum size=1.5cm] at (0,0) {}; &
\tikz \node[groom,tie=green,minimum size=1.5cm] at (0,0) {}; &
\tikz \node[groom,hat=red,minimum size=1.5cm] at (0,0) {};
\\  \hline
\RDD{hair}=red & \RDD{skin}=red & \RDD{shirt}=red & \RDD{undershirt}=green & \RDD{tie}=green & \RDD{hat}=red
\\  \hline 
\end{tabular}


\bigskip
\begin{tabular}{|c|c|c|c|c|c|}\hline 
\multicolumn{6}{|c|}{ \BS{tikz} \BS{node}[\blll{guard},\RDD{hat}=red,minimum size=1.5cm] at (0,0) {};  }
\\ \hline
\tikz\node[guard,hat=red,minimum size=1.5cm] at (0,0) {}; &  
\tikz \node[guard,skin=red,minimum size=1.5cm] at (0,0) {}; &  
\tikz \node[guard,shirt=red,minimum size=1.5cm] at (0,0) {}; &  
\tikz \node[guard,collar=red,minimum size=1.5cm] at (0,0) {}; &
\tikz \node[guard,lining=red,minimum size=1.5cm] at (0,0) {}; &
\tikz \node[guard,details=red,minimum size=1.5cm] at (0,0) {};
\\  \hline
\RDD{hat}=red & \RDD{skin}=red & \RDD{shirt}=red & \RDD{collar}=red & \RDD{lining}=red & \RDD{details}=red
\\  \hline 
\end{tabular}


\bigskip
\begin{tabular}{|c|c|c|c|c|c|}\hline
\multicolumn{6}{|c|}{ \BS{tikz} \BS{node}[\blll{jester},\RDD{hat}=red,minimum size=1.5cm] at (0,0) {};  }
\\ \hline
\tikz \node[jester,hair=red,minimum size=1.5cm] at (0,0) {}; &  
\tikz \node[jester,skin=red,minimum size=1.5cm] at (0,0) {}; &  
\tikz \node[jester,shirt=yellow,minimum size=1.5cm] at (0,0) {}; &  
\tikz \node[jester,hat=red,minimum size=1.5cm] at (0,0) {}; &
%\tikz \node[jester,pattern=yellow,minimum size=1.5cm] at (0,0) {};
&
\tikz \node[jester,details=blue,minimum size=1.5cm] at (0,0) {};
\\  \hline
\RDD{hair}=red & \RDD{skin}=red & \RDD{shirt}=yellow & \RDD{hat}=red & \RDD{pattern}=yellow \footnote{voir confit} & \RDD{details}=blue
\\  \hline 
\end{tabular}

\bigskip
\begin{tabular}{|c|c|c|c|c|}\hline
\multicolumn{5}{|c|}{ \BS{tikz} \BS{node}[\blll{judge},\RDD{hair}=red,minimum size=1.5cm] at (0,0) {};  }
\\ \hline
%\tikz \node[judge,hair=red,minimum size=1.5cm] at (0,0) {}; 
&  
%\tikz \node[judge,skin=red,minimum size=1.5cm] at (0,0) {}; &  
%\tikz \node[judge,shirt=red,minimum size=1.5cm] at (0,0) {}; &  
%\tikz \node[judge,undershirt=red,minimum size=1.5cm] at (0,0) {}; &
%\tikz \node[judge,hairshadow=red,minimum size=1.5cm] at (0,0) {};

\\  \hline
\RDD{hair}=red & \RDD{skin}=red & \RDD{shirt}=red & \RDD{undershirt}=red & \RDD{hairshadow}=red 
\\  \hline 
\end{tabular}

\bigskip
\begin{tabular}{|c|c|c|c|c|c|c|}\hline
\multicolumn{7}{|c|}{ \BS{tikz} \BS{node}[\blll{mexican},\RDD{hair}=red,minimum size=1.5cm] at (0,0) {};  }
\\ \hline
\tikz \node[mexican,hair=red,minimum size=1.5cm] at (0,0) {}; &  
\tikz \node[mexican,skin=red,minimum size=1.5cm] at (0,0) {}; &  
\tikz \node[mexican,shirt=red,minimum size=1.5cm] at (0,0) {}; &  
\tikz \node[mexican,hat=green,minimum size=1.5cm] at (0,0) {}; &
\tikz \node[mexican,ringtop=red,minimum size=1.5cm] at (0,0) {};
&
\tikz \node[mexican,ringmid=red,minimum size=1.5cm] at (0,0) {};
&
\tikz \node[mexican,ringbot=yellow,minimum size=1.5cm] at (0,0) {};
\\  \hline
\RDD{hair}=red & \RDD{skin}=red & \RDD{shirt}=red & \RDD{hat}=green & \RDD{ringtop}=red &\RDD{ringmid}=red & \RDD{ringbot}=yellow
\\  \hline 
\end{tabular}


\bigskip

\begin{tabular}{|c|c|c|}\hline 
\multicolumn{3}{|c|}{ \BS{tikz} \BS{node}[\blll{nun},\RDD{plaid}=red,minimum size=1.5cm] at (0,0) {};  }
\\ \hline
\tikz \node[nun,plaid=red,minimum size=1.5cm] at (0,0) {}; &  
\tikz \node[nun,skin=red,minimum size=1.5cm] at (0,0) {}; &  
\tikz \node[nun,shirt=red,minimum size=1.5cm] at (0,0) {}; 
\\  \hline
\RDD{plaid}=red & \RDD{skin}=red & \RDD{shirt}=red 
\\  \hline 
\end{tabular}


\bigskip
\begin{tabular}{|c|c|c|c|c|c|c|}\hline 
\multicolumn{7}{|c|}{ \BS{tikz} \BS{node}[\blll{nurse},\RDD{hair}=red,minimum size=1.5cm] at (0,0) {};  }
\\ \hline
\tikz \node[nurse,hair=red,minimum size=1.5cm] at (0,0) {}; &  
\tikz \node[nurse,skin=red,minimum size=1.5cm] at (0,0) {}; &  
\tikz \node[nurse,shirt=red,minimum size=1.5cm] at (0,0) {}; &  
\tikz \node[nurse,badgeclip=green,minimum size=1.5cm] at (0,0) {}; &
\tikz \node[nurse,redcross=green,minimum size=1.5cm] at (0,0) {};
&
\tikz \node[nurse,badge=red,minimum size=1.5cm] at (0,0) {};
&
\tikz \node[nurse,badgename=red,minimum size=1.5cm] at (0,0) {};
\\  \hline
\RDD{hair}=red & \RDD{skin}=red & \RDD{shirt}=red & \RDD{badgeclip}=green & \RDD{redcross}=green & \RDD{badge}=red &  \RDD{badgename}=red
\\  \hline 
\end{tabular}


\bigskip
\begin{tabular}{|c|c|c|c|c|c|}\hline
\multicolumn{6}{|c|}{ \BS{tikz} \BS{node}[\blll{physician},\RDD{hair}=red,minimum size=1.5cm] at (0,0) {};  }
\\ \hline
\tikz \node[physician,hair=red,minimum size=1.5cm] at (0,0) {}; &  
\tikz \node[physician,skin=red,minimum size=1.5cm] at (0,0) {}; &  
\tikz \node[physician,shirt=red,minimum size=1.5cm] at (0,0) {}; &  
\tikz \node[physician,hat=red,minimum size=1.5cm] at (0,0) {}; &
\tikz \node[physician,stethoscope=red,minimum size=1.5cm] at (0,0) {};
&
\tikz \node[physician,tube=red,minimum size=1.5cm] at (0,0) {};
\\  \hline
\RDD{hair}=red & \RDD{skin}=red & \RDD{shirt}=red & \RDD{hat}=red & \RDD{stethoscope}=red &   \RDD{tube}=red
\\  \hline 
\end{tabular}


\bigskip
\begin{tabular}{|c|c|c|c|c|c|c|}\hline
\multicolumn{7}{|c|}{ \BS{tikz} \BS{node}[\blll{pilot},\RDD{hat}=red,minimum size=1.5cm] at (0,0) {};  }
\\ \hline
\tikz \node[pilot,hat=red,minimum size=1.5cm] at (0,0) {}; &  
\tikz \node[pilot,skin=red,minimum size=1.5cm] at (0,0) {}; &  
\tikz \node[pilot,shirt=red,minimum size=1.5cm] at (0,0) {}; &  
\tikz \node[pilot,undershirt=red,minimum size=1.5cm] at (0,0) {}; &
\tikz \node[pilot,visor=red,minimum size=1.5cm] at (0,0) {}; &
\tikz \node[pilot,straps=red,minimum size=1.5cm] at (0,0) {}; &
%\tikz \node[pilot,decoration=red,minimum size=1.5cm] at (0,0) {};
\\  \hline
\RDD{hat}=red & \RDD{skin}=red & \RDD{shirt}=red & \RDD{undershirt}=red & \RDD{visor}=red & \RDD{straps}=red & \RDD{decoration}=red
\\  \hline 
\end{tabular}


\bigskip
\begin{tabular}{|c|c|c|c|}\hline
\multicolumn{4}{|c|}{ \BS{tikz} \BS{node}[\blll{police},\RDD{hair}=red,minimum size=1.5cm] at (0,0) {};  }
\\ \hline 
\tikz \node[police,hair=red,minimum size=1.5cm] at (0,0) {}; &  
\tikz \node[police,skin=red,minimum size=1.5cm] at (0,0) {}; &  
\tikz \node[police,shirt=red,minimum size=1.5cm] at (0,0) {}; &  
\tikz \node[police,hat=red,minimum size=1.5cm] at (0,0) {};   
\\  \hline
\RDD{hair}=red & \RDD{skin}=red & \RDD{shirt}=red & \RDD{hat}=red
\\  \hline
\tikz \node[police,badge=red,minimum size=1.5cm] at (0,0) {}; &  
\tikz \node[police,hatbadge=red ,minimum size=1.5cm] at (0,0) {}; &  
\tikz \node[police,hatshield=red,minimum size=1.5cm] at (0,0) {}; &
\tikz \node[police,undershirt=red,minimum size=1.5cm] at (0,0) {}; 
\\  \hline 
\RDD{badge}=red &  \RDD{hatbadge}=red & \RDD{hatshield}=red & \RDD{undershirt}=red
\\  \hline 
\end{tabular}



\bigskip
\begin{tabular}{|c|c|c|c|c|c|}\hline
\multicolumn{6}{|c|}{ \BS{tikz} \BS{node}[\blll{priest},\RDD{hair}=red,minimum size=1.5cm] at (0,0) {};  }
\\ \hline
\tikz \node[priest,hair=red,minimum size=1.5cm] at (0,0) {}; &  
\tikz \node[priest,skin=red,minimum size=1.5cm] at (0,0) {}; &  
\tikz \node[priest,shirt=red,minimum size=1.5cm] at (0,0) {}; &  
\tikz \node[priest,hat=red,minimum size=1.5cm] at (0,0) {}; &
\tikz \node[priest,collar=red,minimum size=1.5cm] at (0,0) {}; &
\tikz \node[priest,cross=red,minimum size=1.5cm] at (0,0) {};
\\  \hline
\RDD{hair}=red & \RDD{skin}=red & \RDD{shirt}=red & \RDD{hat}=red & \RDD{collar}=red &   \RDD{cross}=red
\\  \hline 
\end{tabular}


\bigskip
\begin{tabular}{|c|c|c|c|c|c|c|}\hline 
\multicolumn{7}{|c|}{ \BS{tikz} \BS{node}[\blll{sailor},\RDD{hair}=red,minimum size=1.5cm] at (0,0) {};  }
\\ \hline
\tikz \node[sailor,hair=red,minimum size=1.5cm] at (0,0) {}; &  
\tikz \node[sailor,skin=red,minimum size=1.5cm] at (0,0) {}; &  
\tikz \node[sailor,shirt=red,minimum size=1.5cm] at (0,0) {}; &  
\tikz \node[sailor,hat=red,minimum size=1.5cm] at (0,0) {}; &
\tikz \node[sailor,undershirt=red,minimum size=1.5cm] at (0,0) {};
&
\tikz \node[sailor,stripes=red,minimum size=1.5cm] at (0,0) {};
&
\tikz \node[sailor,details=red,minimum size=1.5cm] at (0,0) {};
\\  \hline
\RDD{hair}=red & \RDD{skin}=red & \RDD{shirt}=red & \RDD{hat}=red & \RDD{undershirt}=red & \RDD{stripes}=red &  \RDD{details}=red
\\  \hline 
\end{tabular}


\bigskip
\begin{tabular}{|c|c|c|c|c|}\hline 
\multicolumn{5}{|c|}{ \BS{tikz} \BS{node}[\blll{santa},\RDD{hat}=green,minimum size=1.5cm] at (0,0) {};  }
\\ \hline
\tikz \node[santa,hat=green,minimum size=1.5cm] at (0,0) {}; &  
\tikz \node[santa,skin=green,minimum size=1.5cm] at (0,0) {}; &  
\tikz \node[santa,shirt=green,minimum size=1.5cm] at (0,0) {}; &  
\tikz \node[santa,beard=green,minimum size=1.5cm] at (0,0) {}; &
\tikz \node[santa,details=green,minimum size=1.5cm] at (0,0) {};

\\  \hline
\RDD{hat}=green & \RDD{skin}=green& \RDD{shirt}=green & \RDD{beard}=green & \RDD{details}=green 
\\  \hline 
\end{tabular}


\bigskip
\begin{tabular}{|c|c|c|c|c|}\hline 
\multicolumn{5}{|c|}{ \BS{tikz} \BS{node}[\blll{surgeon},\RDD{hat}=red,minimum size=1.5cm] at (0,0) {};  }
\\ \hline
\tikz \node[surgeon,hat=red,minimum size=1.5cm] at (0,0) {}; &  
\tikz \node[surgeon,skin=red,minimum size=1.5cm] at (0,0) {}; &  
\tikz \node[surgeon,shirt=red,minimum size=1.5cm] at (0,0) {}; &  
\tikz \node[surgeon,hair=red,minimum size=1.5cm] at (0,0) {}; &
\tikz \node[surgeon,mask=red,minimum size=1.5cm] at (0,0) {};

\\  \hline
\RDD{hat}=red & \RDD{skin}=red & \RDD{shirt}=red & \RDD{hair}=red & \RDD{mask}=red 
\\  \hline 
\end{tabular}


%
%%\newpage
%%
%%\SbSSCT{Canards}{Ducks}
%
%%\label{ducks}

 \maboite{\BS{usepackage}\AC{tikzducks} \cite {tikzducks}}


\begin{center}
\begin{tabular}{|c|}\hline  
\BS{tikz} \BSS{duck} ;
\\ \hline  
\tikz \duck ;
\\ \hline 
\end{tabular} 
\end{center}


\subsubsection{Options}

\noindent

\begin{tabular}{|c|c|c|c|c|} \hline 
\multicolumn{4}{|c|}{\BS{tikz} \BS{duck}[\RDD{body}=red] ;} 
\\ \hline
\tikz \duck[body=red] ;
&  
\tikz \duck[head=red] ;
&
\tikz \duck[bill=red] ;
  &
  \tikz \duck[eye=red] ;
    \\ 
\hline  
[\RDD{body}=red] & [\RDD{head}=red] & [\RDD{bill}=red] & [\RDD{eye}=red] \\ 
\hline 
\end{tabular} 
\begin{tabular}{|c|} \hline
\BS{tikz}  \BS{duck}[\RDD{grumpy}] ;
\\ \hline   
\tikz  \duck[grumpy] ;
\\ \hline  

\end{tabular} 



\bigskip

\begin{tabular}{|c|c|c|c|c|c|} \hline  
\tikz \duck[longhair] ;
&  
\tikz \duck[shorthair] ;
&
\tikz \duck[crazyhair] ;
&
\tikz \duck[recedinghair] ;
&
\tikz \duck[mohican] ;
&
\tikz \duck[mullet] ;
\\ \hline  
[\RDD{longhair}] & [\RDD{shorthair}] & [\RDD{crazyhair}] & [\RDD{recedinghair}] &  [\RDD{mohican}] &  [\RDD{mullet}]\\ 
\hline
\tikz \duck[longhair=red] ;
&  
\tikz \duck[shorthair=red] ;
&
\tikz \duck[crazyhair=red] ;
  &
  \tikz \duck[recedinghair=red] ;
  &
  \tikz \duck[mohican=red] ;
  &
  \tikz \duck[mullet=red] ;
    \\ 
\hline  
[longhair=red] & [shorthair=red] & [crazyhair=red] & [recedinghair=red] &  [mohican=red] &  [mullet=red]  \\ 
\hline 
\end{tabular}

\bigskip





\begin{tabular}{|c|c|c|c|c|} \hline  
\tikz \duck[eyebrow] ;
&  
\tikz \duck[eyebrow=red] ;
&
\tikz \duck[beard] ;
  &
  \tikz \duck[beard=red] ;
    \\ 
\hline  
[\RDD{eyebrow}] & [eyebrow=red] & [\RDD{beard}] & [beard=red] \\ 
\hline
 
\end{tabular}

\bigskip

\begin{tabular}{|c|c|c|c|c|} \hline  
\tikz \duck[tshirt] ;
&  
\tikz \duck[tie] ;
&
\tikz \duck[jacket] ;
&
\tikz \duck[cape] ;
&
\tikz \duck[tshirt,tie ,jacket ,cape] ;
\\ \hline
[\RDD{tshirt}] & [\RDD{tie}] & [\RDD{jacket}] & [\RDD{cape}]& [tshirt,tie ,jacket ,cape]
\\ \hline
\dft{white} & \dft{blue} & \dft{blue} & \dft{red}&
\\ \hline
\tikz \duck[tshirt=red] ;
&  
\tikz \duck[tie=red] ;
&
\tikz \duck[jacket=red] ;
&
\tikz \duck[cape=blue] ;
&

\\ \hline
[tshirt=red] & [tie=red] & [jacket=red] & [cape=blue]& 
\\ \hline
\end{tabular}

\bigskip

\begin{tabular}{|c|c|c|c|c|} \hline 
\tikz \duck[water];
&
\tikz \duck[alien];
&
\tikz \duck[hat];
&
\tikz \duck[tophat];
&
\tikz \duck[cap];
\\ \hline
[\RDD{water}] & [\RDD{alien}] & [\RDD{hat}]& [\RDD{tophat}] & [\RDD{cap}]
\\ \hline
\tikz \duck[santa];
&
\tikz \duck[graduate];
&
\tikz \duck[graduate,tassel];
&
\tikz \duck[beret];
&
\tikz \duck[peakedcap];
\\ \hline
[\RDD{santa}] & [\RDD{graduate}] & [graduate,\RDD{tassel}] & [\RDD{beret}] & [\RDD{peakedcap}]
\\ \hline
\tikz \duck[crown];
&
\tikz \duck[queencrown];
&
\tikz \duck[kingcrown];
&
\tikz \duck[sheep];
&
\tikz \duck[horsetail];
\\ \hline
[\RDD{crown}] &[\RDD{queencrown}]&[\RDD{kingcrown}] & [\RDD{sheep}] &[\RDD{horsetail}]
\\ \hline

\tikz \duck[crozier];
&
\tikz \duck[unicorn];
&

\tikz \duck[bunny];
&
\tikz \duck[bunny=red,inear=blue];
&
\tikz \duck[witch];
\\ \hline
 [\RDD{crozier}] & [\RDD{unicorn}] &[\RDD{bunny}] & [bunny=red,\RDD{inear}=blue] & [\RDD{witch}]
\\ \hline
\tikz \duck[magicwand];
&
\tikz \duck[magichat];
&
\tikz \duck[magichat=teal,
magicstars=blue!30!cyan,
magicwand];
&
\tikz \duck[glasses];
&
\tikz \duck[sunglasses];
\\ \hline
[\RDD{magicwand}] & [\RDD{magichat}] & \parbox{3cm}{[magichat,\\ \RDD{magicstars}]} & [\RDD{glasses}] & [\RDD{sunglasses}]
\\ \hline
\end{tabular}  


\begin{tabular}{|c|c|c|c|c|} \hline

\tikz \duck[squareglasses];
&
\tikz \duck[signpost=42];
&
\tikz \duck[signpost=XXX,signcolour=green];
&
\tikz \duck[signpost=XXX,signback=green];
&
\tikz \duck[speech={XXX}];
\\ \hline
[\RDD{squareglasses}] & [\RDD{signpost}=42] & \parbox{3cm}{[signpost=XXX,\\ \RDD{signcolour}=green]} & \parbox{3cm}{[signpost=XXX, \\ \RDD{signback}=green]} & [\RDD{speech}=\AC{XXX}] 
\\ \hline 
\end{tabular}


\begin{tabular}{|c|c|c|c|c|} \hline 
\tikz \duck[speech={XXX},bubblecolour=green];
&
\tikz \duck[think={XXX}];
&
\tikz \duck[think=XXX,bubblecolour=green];
&
\tikz \duck[book={XXX}];
\\ \hline
\parbox{3cm}{[speech={XXX},\\ \RDD{bubblecolour}=green]} & [\RDD{think}=\AC{XXX}] & 
\parbox{3cm}{[think={XXX},\\ \RDD{bubblecolour}=green]}
&[\RDD{book}=\AC{XXX}] 
\\ \hline
%\end{tabular}  
%
%\begin{tabular}{|c|c|c|c|c|} \hline
\tikz \duck[book=XXX,bookcolour=green];
&
\tikz \duck[book=\scalebox{0.5}{XXX}];
&
\multicolumn{2}{|c|}{\tikz 
\duck[signpost=\scalebox{0.4}{
\parbox{2cm}{
\centering XXX \\ XXXXX}}]
;}
\\ \hline

\parbox{3cm}{[book={XXX},\\ \RDD{bookcolour}=green]} 
&
\parbox{3.5cm}{\BS{tikz} \BS{duck}[book=\\ \BS{scalebox}\AC{0.5}\AC{XXX}]; }
&
\multicolumn{2}{|c|}{ \parbox{7cm}{ \BS{tikz} 
\BS{duck}[signpost=\BS{scalebox}\AC{0.4}\AC{ \\
\BS{parbox}\AC{2cm}{  
\BS{centering} XXX \ XXXXX}}]}
;}
\\ \hline
\end{tabular}


\begin{tabular}{|c|c|c|c|c|} \hline 
\tikz \duck[cricket];
&
\tikz \duck[hockey];
&
\tikz \duck[football];
&
\tikz \duck[lightsaber];
&
\tikz \duck[torch];
\\ \hline
[\RDD{cricket}]& [\RDD{hockey}] & [\RDD{football}] & [\RDD{lightsaber}] & [\RDD{torch}]
\\ \hline
\tikz \duck[prison];
&
\tikz \duck[necklace];
&
\tikz \duck[icecream];
&
\tikz \duck[icecream,flavoura=green];
&
\tikz \duck[icecream,flavourb=green];
\\ \hline
[\RDD{prison}] & [\RDD{necklace}] & [\RDD{icecream}] &
\parbox{3cm}{[icecream,\\ \RDD{flavoura}=green]} & 
\parbox{3cm}{[icecream,\\ \RDD{flavourb}=green]}
\\ \hline
\tikz \duck[icecream,flavourc=green];
&
\tikz \duck[chef];
&
\tikz \duck[rollingpin];
&
\tikz \duck[cake];
&
\tikz \duck[pizza];
\\ \hline
\parbox{3cm}{[icecream,\\ \RDD{flavourc}=green]} &
 [\RDD{chef}] & [\RDD{rollingpin}] & [\RDD{cake}] & [\RDD{pizza}]
\\ \hline
\tikz \duck[baguette];
&
\tikz \duck[milkshake];
&
\tikz \duck[wine];
&
\tikz \duck[mask];
&
\tikz \duck[buttons];
\\ \hline
  [\RDD{baguette}] & [\RDD{milkshake}] & [\RDD{wine}] & [\RDD{mask}] & [\RDD{buttons}]
\\ \hline
\end{tabular}

\begin{tabular}{|c|c|c|c|c|} \hline 
\tikz \duck[basket];
&
\tikz \duck[easter];
&
\tikz \duck[easter,egga=red];
&
\tikz \duck[easter,eggb=red];
&
\tikz \duck[easter,eggc=red];
\\ \hline
[\RDD{basket}]& [\RDD{easter}] & [easter,\RDD{egga}=red] & [easter,\RDD{eggb}=red] & [easter,\RDD{eggc}=red]
\\ \hline

\end{tabular}

\bigskip

\begin{tabular}{|c|c|c|c|} \hline 
\multicolumn{4}{|c|}{\BS{tikz} \BS{duck} \BS{path}[preaction=\AC{fill,green},pattern=dots, pattern  color=red]  \BSS{duckpathbody} ;} 
\\ \hline
\tikz \duck
\path[preaction={fill,green},pattern=dots, pattern  color=red]  \duckpathbody;
&
\tikz   \duck
\path[preaction={fill,green},pattern=dots, pattern  color=red]  \duckpathgrumpybill;
&
\tikz   \duck
\path[preaction={fill,green},pattern=dots, pattern  color=red]  \duckpathbill;
&
\tikz   \duck
\path[preaction={fill,green},pattern=dots, pattern  color=red]  \duckpathtshirt;
\\ \hline 
\BSS{duckpathbody } & \BSS{duckpathgrumpybill} & \BSS{duckpathbill} & \BSS{duckpathtshirt} 
\\ \hline
\tikz   \duck
\path[preaction={fill,green},pattern=dots, pattern  color=red]  \duckpathjacket;
&
\tikz   \duck
\path[preaction={fill,green},pattern=dots, pattern  color=red]  \duckpathcape ;
&
\tikz   \duck
\path[preaction={fill,green},pattern=dots, pattern  color=red]  \duckpathshorthair ;
&
\tikz   \duck
\path[preaction={fill,green},pattern=dots, pattern  color=red]  \duckpathlonghair;
\\ \hline 
\BSS{duckpathjacket} & \BSS{duckpathcape}  & \BSS{duckpathshorthair} & \BSS{duckpathlonghair} 
\\ \hline 
\tikz   \duck
\path[preaction={fill,green},pattern=dots, pattern  color=red]  \duckpathcrazyhair;
&
\tikz   \duck
\path[preaction={fill,green},pattern=dots, pattern  color=red]  \duckpathrecedinghair;
&
\tikz   \duck
\path[preaction={fill,green},pattern=dots, pattern  color=red]  \duckpathcrown ;
&
\tikz   \duck
\path[preaction={fill,green},pattern=dots, pattern  color=red]   \duckpathmohican ;
\\ \hline 
\BSS{duckpathcrazyhair} & \BSS{duckpathrecedinghair} & \BSS{duckpathcrown} & \BSS{duckpathmohican}
\\ \hline 
\tikz   \duck
\path[preaction={fill,green},pattern=dots, pattern  color=red]   \duckpathmullet;
&
\tikz   \duck
\path[preaction={fill,green},pattern=dots, pattern  color=red]   \duckpathqueencrown ;
&
\tikz   \duck
\path[preaction={fill,green},pattern=dots, pattern  color=red]   \duckpathkingcrown ;
&
\tikz   \duck
\path[preaction={fill,green},pattern=dots, pattern  color=red]  \duckpathdarthvader ;
\\ \hline 
\BSS{duckpathmullet} & \BSS{duckpathqueencrown} & \BSS{duckpathkingcrown} & \BSS{duckpathdarthvader}
\\ \hline 
\tikz   \duck
\path[preaction={fill,green},pattern=dots, pattern  color=red]  \duckpathhorsetail ;
& & & 
\\ \hline 
\BSS{duckpathhorsetail}& & & 
\\ \hline 
\end{tabular}


\SbSbSSCT{Canards aléatoires}{Random ducks}

\noindent

\begin{tabular}{|c|}\hline  
\BS{tikz} \BSS{randuck} ; \BS{tikz} \BSS{randuck} ; \BS{tikz} \BSS{randuck} ; \BS{tikz} \BSS{randuck} ; \BS{tikz} \BSS{randuck} ; 
\\ \hline    
\tikz \randuck ; \tikz \randuck  ; \tikz  \randuck  ; \tikz  \randuck  ; \tikz  \randuck ;
\\ \hline 
\end{tabular} 

\bigskip

\begin{tabular}{|c|} \hline  
\BS{tikz} \BSS{shuffleducks} \BS{duck}[\BSS{randomhead}] ;
\\ \hline  
\tikz \shuffleducks \duck[\randomhead] ; \tikz \shuffleducks \duck[\randomhead] ; \tikz \shuffleducks \duck[\randomhead] ; \tikz \shuffleducks \duck[\randomhead] ;
\tikz \shuffleducks \duck[\randomhead] ;
\\ \hline 
\end{tabular} 

\bigskip

\begin{tabular}{|c|} \hline  
\BS{tikz} \BSS{shuffleducks} \BS{duck}[\BSS{randomaccessories}] ;
\\ \hline  
\tikz \shuffleducks \duck[\randomaccessories] ; \tikz \shuffleducks \duck[\randomaccessories] ; \tikz \shuffleducks \duck[\randomaccessories] ; \tikz \shuffleducks \duck[\randomaccessories] ; \tikz \shuffleducks \duck[\randomaccessories] ; 
\\ \hline 
\end{tabular} 


\SbSbSSCT{Coordonnées}{Coordinates}

\noindent


\begin{tabular}{|c|c|c|} \hline 
\multicolumn{3}{|c|}{\BS{tikz} \BS{duck} \BS{fill}[red] (wing) circle (3pt);}
\\ \hline  
\tikz \duck \fill[red] (wing) circle (3pt);
&
\tikz \duck \fill[red] (head) circle (3pt);
&
\tikz \duck \fill[red] (bill) circle (3pt);
\\ \hline wing &  head & bill \\ 
\hline 
\end{tabular} 

\bigskip

\begin{tabular}{|c|} \hline
\BS{tikz} \BS{duck}[\RDD{name}=XXX] \\ \BS{begin}\AC{scope} [xshift=4cm] \BS{duck}[\RDD{name}=YYY] 
\BS{end}\AC{scope} \\ \BS{draw}[red] (XXX-wing) - - (YYY-bill) ;
\\ \hline
\tikz \duck[name=XXX] \begin{scope} [xshift=4cm] 
\duck[name=YYY]
\end{scope}  \draw[red] (XXX-wing) -- (YYY-bill);
\\ \hline
\end{tabular} 


\SbSbSSCT{Rayures}{Stripes}

\noindent

\begin{tabular}{|c|c|} \hline  
\tikz \duck \stripes ;
&  
\tikz \duck[stripes] ;
\\ \hline  
\BS{tikz} \BS{duck} \BSS{stripes} ; 
&  
\BS{tikz} \BS{}duck[\RDD{stripes}] ;
\\ \hline 
\end{tabular} 

\bigskip

\begin{tabular}{|c|c|} \hline  
\tikz \duck[rollingpin] \stripes ;
&  
\tikz  \duck[rollingpin,stripes] ;
\\ \hline  
\BS{tikz} \BS{duck}[rollingpin] \BS{stripes} ;
&  
\BS{tikz}  \BS{duck}[rollingpin,stripes] ;
\\ \hline 
\end{tabular} 



\bigskip

\begin{tabular}{|c|c|c|c|}\hline 
\multicolumn{4}{|c|}{\BS{tikz} \BS[duck] \BS{stripes}[\RDD{color}=red];}
\\ \hline  
\tikz \duck \stripes[color=red];
& 
\tikz \duck \stripes[distance=.5]; 
&  
\tikz \duck \stripes[width=.05];
&
\tikz \duck \stripes[height=1];  
\\ \hline 
[\RDD{color}=red] & [\RDD{distance}=.5] & [\RDD{width}=.05] & [\RDD{height}=1] 
\\ \hline 
\dft{black} & \dft{0.3}  & \dft{0.15}  &  \dft{2.7} \\ 
\hline  

\tikz \duck \stripes[rotate=45]; & \tikz \duck \stripes[initialx=1]; & \tikz \duck \stripes[initialy=1]; &  
\\ \hline 
[\RDD{rotate=}45] & [\RDD{initialx}=1] & [\RDD{initialy}=1] &
\\ \hline 
\dft{-10} & \dft{0.1} & \dft{-0.3} & 
\\ \hline 
\end{tabular} 


\bigskip

\begin{tabular}{|c|c|c|} \hline
\multicolumn{3}{|c|}{\BS{tikz} \BS[duck] \BS{stripes}[\RDD{emblem}=XXX];}
\\ \hline 
\tikz \duck \stripes[emblem={XXX}];
&  
\tikz \duck \stripes[emblem={\includegraphics[width=6mm]{LogoIUT}}];
&  
\tikz \duck \stripes[emblem={\DFR}];
\\ \hline  
[emblem=XXX]
& \parbox{5cm}{ [emblem=\{\BS{includegraphics}  [width=6mm]\AC{LogoIUT} \} ] }  
&  [emblem=\AC{\BS{DFR}}  ] 
\\ \hline 
& &  \BS{DFR} : \TFRGB{voir}{see} page \pageref{DFR} 
\\ \hline   
\end{tabular} 

\bigskip

\begin{tabular}{|c|c|c|} \hline
\tikz \duck[stripes={ \stripes \stripes[rotate=45]}] ;
\\ \hline 
\BS{tikz}
\BS{duck}[stripes=\AC{
\BS{stripes}
\BS{stripes}[rotate=45] } ]
;

\\ \hline  
\end{tabular}



%%
%%\newpage
%
%
%\newpage
%
%\section{A Voir}
%
%
%%\RRR{75-2 = Concept: Data Points and Data Formats}

\begin{tikzpicture}
\datavisualization [school book axes, visualize as smooth line]
data {
x, y
-1.5, 2.25
-1, 1
-.5, .25
0, 0
.5, .25
1, 1
1.5, 2.25
};
\end{tikzpicture}


\begin{tikzpicture}
\datavisualization [school book axes, visualize as smooth line]
data [format=function] {
var x : interval [-1.5:1.5] samples 7;
func y = \value x*\value x;
};
\end{tikzpicture}


\begin{tikzpicture}
\datavisualization [school book axes, visualize as smooth line]
data [format=function] {
var x : interval [-1.5:1.5] samples 3;
func y = \value x*\value x;
};
\end{tikzpicture}

Section 76 gives an in-depth coverage of the available data formats and explains how new data formats
can be defined.


\RRR{75-3 = Concept: Axes, Ticks, and Grids}


\begin{tikzpicture}
\datavisualization [
scientific axes,
x axis={length=3cm, ticks=few},
visualize as smooth line
]
data [format=function] {
var x : interval [-1.5:1.5] samples 20;
func y = \value x*\value x;
};
\end{tikzpicture}

\begin{tikzpicture}
\datavisualization [
scientific axes=clean,
x axis={length=3cm, ticks=few},
all axes={grid},
visualize as smooth line
]
data [format=function] {
var x : interval [-1.5:1.5] samples 7;
func y = \value x*\value x;
};
\end{tikzpicture}

Section 77 explains in more detail how axes, ticks, and grid lines can be chosen and configured.


\RRR {75-4 = Concept: Visualizers}

\begin{tikzpicture}
\datavisualization [
scientific axes=clean,
x axis={length=3cm, ticks=few},
visualize as smooth line
]
data [format=function] {
var x : interval [-1.5:1.5] samples 7;
func y = \value x*\value x;
};
\end{tikzpicture}

\begin{tikzpicture}
\datavisualization [
scientific axes=clean,
x axis={length=3cm, ticks=few},
visualize as scatter
]
data [format=function] {
var x : interval [-1.5:1.5] samples 7;
func y = \value x*\value x;
};
\end{tikzpicture}

Section 78 provides more information on visualizers as well as reference lists.

\RRR{75-5 = Concept: Style Sheets and Legends }

\begin{tikzpicture}[baseline]
\datavisualization [ scientific axes=clean,
y axis=grid,
visualize as smooth line/.list={sin,cos,tan},
style sheet=strong colors,
style sheet=vary dashing,
sin={label in legend={text=$\sin x$}},
cos={label in legend={text=$\cos x$}},
tan={label in legend={text=$\tan x$}},
data/format=function ]
data [set=sin] {
var x : interval [-0.5*pi:4];
func y = sin(\value x r);
}
data [set=cos] {
var x : interval [-0.5*pi:4];
func y = cos(\value x r);
}
data [set=tan] {
var x : interval [-0.3*pi:.3*pi];
func y = tan(\value x r);
};
\end{tikzpicture}


Section 79 details style sheets and legends.

\RRR{75-6 = Usage}

\subsection{/pgf/data/read from file=filename} (no default, initially empty)

If you set the source attribute to a non-empty hfilenamei, the data will be read from this file. In
this case, no hinline datai may be present, not even empty curly braces should be provided.
%\datavisualization ...
data [read from file=file1.csv]
data [read from file=file2.csv];
The other way round, if read from file is empty, the data must directly follow as hinline datai.
%\datavisualization ...
data {
x, y
1, 2
2, 3
};

The second important key is format, which is used to specify the data format:

\subsection{/pgf/data/format}

Use this key to locally set the format used for parsing the data, see Section 76 for a list of predefined
formats.

\tikz
\datavisualization [school book axes, visualize as line]
data [/data point/x=1] {
y
1
2
}
data [/data point/x=2] {
y
2
0
.5
};

\BS{datavisualization} . . . data point[options] . . . ;

\tikz \datavisualization [school book axes, visualize as line]
data point [x=1, y=1] data point [x=1, y=2]
data point [x=2, y=2] data point [x=2, y=0.5];

/tikz/data visualization/data point=options

\tikzdatavisualizationset{
horizontal/.style={
data point={x=#1, y=1}, data point={x=#1, y=2}},
}
\tikz \datavisualization
[ school book axes, visualize as line,
horizontal=1,
horizontal=2 ];

\BS{datavisualization} . . . data group[options]\AC{name}+=\AC{data specifications} . . . ;


\tikz \datavisualization data group {points} = {
data {
x, y
0, 1
1, 2
2, 2
5, 1
2, 0
1, 1
}
};

\tikz \datavisualization [school book axes, visualize as line] data group {points};
\qquad
\tikz \datavisualization [scientific axes=clean, visualize as line] data group {points};


\BS{datavisualization} . . . scope[options]{data specification} . . . ;

%\datavisualization...
%scope [/data point/experiment=7]
%{
%data [read from file=experiment007-part1.csv]
%data [read from file=experiment007-part2.csv]
%data [read from file=experiment007-part3.csv]
%}
%scope [/data point/experiment=23, format=foo]
%{
%data [read from file=experiment023-part1.foo]
%data [read from file=experiment023-part2.foo]
%};


\BS{datavisualization} . . . info[options]{code} . . . ;

\begin{tikzpicture}[baseline]
\datavisualization [ school book axes, visualize as line ]
data [format=function] {
var x : interval [-0.1*pi:2];
func y = sin(\value x r);
}
info {
\draw [red] (visualization cs: x={(.5*pi)}, y=1) circle [radius=1pt]
node [above,font=\footnotesize] {extremal point};
};
\end{tikzpicture}

\subsection{Coordinate system visualization}

\BS{datavisualization} . . . info’[options]{code} . . . ;

\begin{tikzpicture}[baseline]
\datavisualization [ school book axes, visualize as line ]
data [format=function] {
var x : interval [-0.1*pi:2];
func y = sin(\value x r);
}
info' {
\fill [red] (visualization cs: x={(.5*pi)}, y=1) circle [radius=2mm];
};
\end{tikzpicture}


\subsection{Predefined node data visualization bounding box}
This rectangle node stores a bounding box of the data visualization that is currently being constructed.
This node can be useful inside info commands or when labels need to be added.

\subsection{Predefined node data bounding box}
This rectangle node is similar to data visualization bounding box, but it keeps track only of the actual
data. The spaces needed for grid lines, ticks, axis labels, tick labels, and other all other information
that is not part of the actual data is not part of this box.


\RRR{75-7 = Advanced: Executing User Code During a Data Visualization}

\RRR{75-8 = Advanced: Creating New Objects}


\section{76 Providing Data for a Data Visualization}

%
%\begin{tikzpicture}
%\draw[help lines] (0,0) grid (2,2);
%\node[draw,fill=green!20,] (A) at (1,1) {\huge noeud};
%\fill[red] (A.-30) circle (3pt);
%\end{tikzpicture}
%
%
%
%\begin{tikzpicture}
%\draw[help lines] (0,0) grid (3,2);
%\draw (0,0) -- (1,1);
%\draw[red] (0,0) -- ([xshift=3pt] 1,1);
%\draw (1,0) -- +(30:2cm);
%\draw[red] (1,0) -- +([shift=(135:5pt)] 30:2cm);
%\end{tikzpicture}
%
%
%
%\section{Problèmes a Voir}
%
%%\tikz \node[jester,pattern=yellow,minimum size=1.5cm] at (0,0) {};
%
%\newpage
%
%
%
\begin{tabular}{|c|c|l c|}\hline 
\multicolumn{4}{|c|}{ \textbf{\TFRGB{module de base TikZ}{Basic TikZ package} : } }
\\ \hline

\TFRGB{nom}{name} & \TFRGB{A insérer dans le préambule}{Load package}& documentation \footnotemark[1] 	& \\  \hline 
tikz & \BS{usepackage}\AC{tikz}  	& pgfmanual.pdf			& \DGB \\

\hline 
\end{tabular} 

\bigskip

\begin{tabular}{|c|c|l c|}\hline 
\multicolumn{4}{|c|}{ \textbf{\TFRGB{Autres modules}{Other packages}} }
\\ \hline
\TFRGB{nom}{name} & \TFRGB{voir page}{see page} & documentation  \footnotemark[2] 	& \\  \hline 
animate 	& \pageref{anim} 	& animate.pdf 			& \DGB \\
tikz-optics 	& \pageref{optics} 	& tikz-optics.pdf 			& \DFR \\
pgfplots 	& \pageref{pgfplots} & pgfplots.pdf 		& \DGB \\
tikzpeople  & \pageref{people} 	& tikzpeople.pdf 		& \DGB \\
tikzducks  & \pageref{ducks} 	& tikzducks-doc.pdf 		& \DGB \\
tikzsymbols  & \pageref{symbol} 	& tikzsymbols.pdf 		& \DGB \\
tkz-tab  	& \pageref{tabl} 	& tkz-tab-screen.pdf 	& \DFR \\
\hline 
\end{tabular} 
\bigskip



\begin{tabular}{|l|c|l|}\hline 
\multicolumn{3}{|c|}{ \textbf{\TFRGB{Compléments optionnels}{Optional library} (documentation : pgfmanual.pdf)} }
\\ \hline
\TFRGB{nom}{name} 				& \TFRGB{voir page}{see page}						& \TFRGB{A insérer dans le préambule}{Load package}\\ \hline 
angles & \pageref{lib-angles} &  \BS{usetikzlibrary}\AC{angles}\\
arrows.meta	& \pageref{lib-arrows.meta}	&  \BS{usetikzlibrary}\AC{arrows.meta}\\
bending				& \pageref{lib-bending}			&  \BS{usetikzlibrary}\AC{bending}
\\
backgrounds			& \pageref{lib-bkgd} 			&  \BS{usetikzlibrary}\AC{backgrounds}
\\
calc				& \pageref{lib-calc}			&  \BS{usetikzlibrary}\AC{calc}
\\
chains			& \pageref{lib-chains} 			& \BS{usetikzlibrary}\AC{chains} 
\\
circuits.ee.IEC				& \pageref{lib-ee}			&  \BS{usetikzlibrary}\AC{circuits.ee.IEC}
\\
circuits.logic.IEC	& \pageref{lib-gate}			&  \BS{usetikzlibrary}\AC{circuits.logic.IEC}
\\ 
circuits.logic.US	& \pageref{lib-gate}			&  \BS{usetikzlibrary}\AC{circuits.logic.US}
\\ 
circuits.logic.CDH	& \pageref{lib-gate}			&  \BS{usetikzlibrary}\AC{circuits.logic.CDH}
\\ 
fit & \pageref{lib-fit} 	& \BS{usetikzlibrary}\AC{fit} 
\\
decorations.footprints & \pageref{lib-footprints} 	& \BS{usetikzlibrary}\AC{decorations.footprints} 
\\
decorations.fractals & \pageref{lib-fractals} 		& \BS{usetikzlibrary}\AC{decorations.fractals} 
\\
decorations.markings & \pageref{lib-mark} 			& \BS{usetikzlibrary}\AC{decorations.markings} 
\\
decorations.pathmorphing  & \pageref{lib-morph}		& \BS{usetikzlibrary}\AC{decorations.pathmorphing}
\\
decorations.pathreplacing & \pageref{lib-replac}	& \BS{usetikzlibrary}\AC{decorations.pathreplacing} 
\\
decorations.shapes & \pageref{lib-shapes} 			& \BS{usetikzlibrary}\AC{decorations.shapes} 
\\
decorations.text & \pageref{lib-text} 				& \BS{usetikzlibrary}\AC{decorations.text} 
\\
fadings 			& \pageref{lib-fadings}			&  \BS{usetikzlibrary}\AC{fadings }
\\
intersections		& \pageref{lib-intersections}	&  \BS{usetikzlibrary}\AC{intersections}
\\
matrix			& \pageref{lib-matrix} 			& \BS{usetikzlibrary}\AC{matrix} 
\\
patterns			& \pageref{lib-patterns}		&  \BS{usetikzlibrary}\AC{patterns}
\\
plotmarks			& \pageref{plotmarks} 			&  \BS{usetikzlibrary}\AC{plotmarks}
\\
positioning			& \pageref{lib-pos} 			&  \BS{usetikzlibrary}\AC{positioning}
\\ 
scopes				& \pageref{lib-scopes}			&  \BS{usetikzlibrary}\AC{scopes}
\\
shadings			& \pageref{lib-shadings}		&  \BS{usetikzlibrary}\AC{shadings}
\\
shapes.arrows		& \pageref{lib-arr}				&\BS{usetikzlibrary}\AC{shapes.arrows} 
\\shapes.callouts		& \pageref{lib-call}			& \BS{usetikzlibrary}\AC{shapes.callouts} 
\\
shapes.geometric	& \pageref{lib-geom} 			& \BS{usetikzlibrary}\AC{shapes.geometric}
\\

shapes.misc			& \pageref{lib-misc} 			& \BS{usetikzlibrary}\AC{shapes.misc} 
\\
shapes.multipart	& \pageref{lib-mult} 			& \BS{usetikzlibrary}\AC{shapes.multipart} 
\\
shapes.symbols		& \pageref{lib-symb}			& \BS{usetikzlibrary}\AC{shapes.symbols} 
\\
through				& \pageref{lib-through}			&  \BS{usetikzlibrary}\AC{through}
\\ 
trees				& \pageref{lib-trees}
\BS{usetikzlibrary}\AC{trees}
\\ 
through				& \pageref{lib-turtle}			&  \BS{usetikzlibrary}\AC{turtle}
\\ 
\hline
 \end{tabular} 

\TFRGB{ 
\footnotetext[1]{voir dans le répertoire :  \BS{texlive}\BS{2016}\BS{tesmf-dist}\BS{doc}\BS{generic}\BS{pgf}}
\footnotetext[2]{chercher  dans le répertoire  :  \BS{texlive}\BS{2016}\BS{tesmf-dist}\BS{doc}\BS{latex}} }
{ 
\footnotetext[1]{look in repertory :  \BS{texlive}\BS{2016}\BS{tesmf-dist}\BS{doc}\BS{generic}\BS{pgf}}
\footnotetext[2]{search in repertory :  \BS{texlive}\BS{2016}\BS{tesmf-dist}\BS{doc}\BS{latex}} }

\bigskip



\begin{tabular}{|l|c|}\hline
\multicolumn{2}{|c|}{ \TFRGB{dans une prochaine mise à jour}{In a a future update } }
\\ \hline
automata			& \RRR{41} \\
babel				& \RRR{42} \\
calendar			& \RRR{45} \\
%chains				& \RRR{46} \\ 
%circuits.ee		& \RRR{47-4} \\ 
 
circular graph drawing library 				& \RRR{32} \\
curvilinear library 						& \RRR{103-4-7} \\
datavisualization library					& \RRR{75} \\
datavisualization.formats.functions library & \RRR{76-4} \\
datavisualization.polar library 			& \RRR{80}  \\
 er 										& \RRR{49}  \\
examples graph drawing library 				& \RRR{35-8} \\ 
external 									& \RRR{50}  \\  
%fit 										& \RRR{52} \\ 
fixedpointarithmetic 						& \RRR{53} \\ 
folding 									& \RRR{59} \\
force graph drawing library 				& \RRR{31}  \\
fpu											& \RRR{54}  \\
graph.standard library 						& \RRR{19-10}\\
graphdrawing library 						& \RRR{27} \\
graphs library 								& \RRR{19} \\ 
layered graph drawing library 				& \RRR{30}  \\
lindenmayersystems							& \RRR{55} \\  
mindmap										& \RRR{58} \\ 
petri										& \RRR{61}  \\ 
phylogenetics graph drawing library 		& \RRR{33} \\
plothandlers								& \RRR{62}  \\  
profiler									& \RRR{64}   \\ 
quotes library 								& \RRR{17-10-4} \\
routing graph drawing library 				& \RRR{34} \\
shadows										& \RRR{66}   \\ 
 
spy											&  \RRR{68} \\ 
svg.path									&  \RRR{69} \\ 
%through										&  \RRR{71} \\ 
topaths										&  \RRR{70} \\ 
trees graph drawing library					& 
\\ \hline
\end{tabular}  


%\newpage
%%
%% \tableofcontents
%\renewcommand{\bibname}{Sources}
%
\label{sources}
%\input{bib}

\newpage

\begin{thebibliography}{99}
\bibitem{pgfmanual} pgfmanual.pdf  	\hspace{1cm}	version 3.0.1a \hspace{1cm} 	1161 pages 	\hspace{1cm}	\DGB
\bibitem{pgfplots} pgfplots.pdf 	\hspace{1cm}	version 1.80 \hspace{1cm} 	439 pages 	\hspace{1cm}	\DGB
\bibitem{tikstab} tkz-tab-screen.pdf 	\hspace{1cm}	version 1.1c \hspace{1cm} 	83 pages 	\hspace{1cm}	\DFR
\bibitem{tikzpeople} tikzpeople.pdf 	\hspace{1cm}	 \hspace{1cm} 	19 pages 	\hspace{1cm}	\DGB

\bibitem{tikzducks} tikzducks-doc.pdf 	\hspace{1cm}	version 0.6  \hspace{1cm} 	28 pages 	\hspace{1cm}	\DGB

\bibitem{tikzsymbols} tikzsymbols.pdf 	\hspace{1cm}	version sept 2017  \hspace{1cm} 	15 pages 	\hspace{1cm}	\DGB

\bibitem{animate} animate.pdf 	\hspace{1cm}	 \hspace{1cm} 	26 pages 	\hspace{1cm}	\DGB

\bibitem{optics} tikz-optics.pdf	\hspace{1cm}	version 0.2.2  \hspace{1cm} 	39 pages 	\hspace{1cm}	\DFR
\end{thebibliography}


%
%\newpage 
%
%%\printindex 
  
\end{document}:'
\end{verbatim}

Primero se define la orden \verb|\blancoynegro| y después se lee el \filenomo{} real.  Poniendo \verb|\blancoynegro| a falso se producirá la versión en color del documento.

\subsection{Su propio paquete}

Si define muchos nuevos entornos y órdenes, el preámbulo de su documento se hará muy largo.  En situaciones así es buena idea crear un paquete \LaTeX{} que contenga todas sus definiciones de órdenes y entornos.  Puede usar después la orden \ci{usepackage} para cargar el paquete en su documento actual o en otros similares.

\begin{figure}[!htbp]
\begin{lined}{\textwidth}
\begin{verbatim}
% Paquete Demo de Tobias Oetiker
\ProvidesPackage{demopack}
\newcommand{\intc}{La introducción no-tan-corta
                   a \LaTeXe}
\newcommand{\txsit}[1]{La introducción \emph{#1}-corta
                       a \LaTeXe}
\newenvironment{king}{\begin{quote}}{\end{quote}}
\end{verbatim}
\end{lined}
\caption{Paquete de ejemplo.} \label{package}
\end{figure}

Escribir un paquete básicamente consiste en copiar el contenido del preámbulo de su documento en un \filenomo{} separado con un nombre que termine en \texttt{.sty}.  Hay una orden especial,
\begin{lscommand}
\ci{ProvidesPackage}\verb|{|\emph{nombre paquete}\verb|}|
\end{lscommand}
para usar justo al principio de su \filenomo{} de paquete. \verb|\ProvidesPackage| dice a  \LaTeX{} el nombre del paquete y le permite emitir un mensaje de error notable cuando intente incluir el paquete dos veces.  La figura~\ref{package} muestra un pequeño paquete de ejemplo que contiene órdenes definidas en ejemplos anteriores.

\section{\Fontsnomo{} y tamaños}

\subsection{Órdenes que cambian la \fontnomo{}}
\index{\fontnomo{}}\index{tamaño de la \fontnomo{}} \LaTeX{} escoge la \fontnomo{} y el tamaño de \fontnomo{} apropiados basándose en la estructura lógica del documento (secciones, notas al pie, \ldots).   En algunos casos, quizá desee cambiar \fontsnomo{} y tamaños a mano.  Para hacerlo, puede usar las órdenes listadas en los cuadros~\ref{fonts} y~\ref{sizes}.  El tamaño real de cada \fontnomo{} es una cuestión de diseño y depende de la clase de documento y de sus opciones.  El cuadro~\ref{tab:pointsizes} muestra los tamaños absolutos en puntos para estas órdenes según se implementan en las clases de documentos normales.

\begin{example}
{\small Pequeña \textbf{negrita}
 del África tropical,}
{\Large grande y \textit{cursi}va
 eres tú ya.}
\end{example}

Una característica importante de \LaTeXe{} es que los atributos de \fontnomo{} son independientes.  Esto significa que puede poner órdenes para cambiar el tamaño o incluso la \fontnomo{}, y todavía se mantendrán los atributos de negrita o cursiva establecidos anteriormente.

En \emph{modo mates} puede usar las \emph{órdenes} de cambio de \fontnomo{} para salir temporalmente del \emph{modo mates} e introducir texto normal.  Si quiere cambiar a otra \fontnomo{} para composición de mates necesita otro conjunto especial de órdenes; véase el cuadro~\ref{mathfonts}.

\begin{table}[!bp]
\caption{\Fontsnomo{}.} \label{fonts}
\begin{lined}{12cm}
%
% Alan suggested not to tell about the other form of the command
% eg \verb|\sffamily| or \verb|\bfseries|. This seems a good thing to me.
%
\begin{tabular}{@{}rl@{\qquad}rl@{}}
\fni{textrm}\verb|{...}|        &      \textrm{\wi{rematada}}&
\fni{textsf}\verb|{...}|        &      \textsf{\wi{palo seco}}\\
\fni{texttt}\verb|{...}|        &      \texttt{de máquina}\\[6pt]
\fni{textmd}\verb|{...}|        &      \textmd{peso medio}&
\fni{textbf}\verb|{...}|        &      \textbf{\wi{negrita}}\\[6pt]
\fni{textup}\verb|{...}|        &       \textup{\wi{recta}}&
\fni{textit}\verb|{...}|        &       \textit{\wi{cursiva}}\\
\fni{textsl}\verb|{...}|        &       \textsl{\wi{oblicua}}&
\fni{textsc}\verb|{...}|        &       \textsc{\wi{Versalitas}}\\[6pt]
\ci{emph}\verb|{...}|          &            \emph{destacada} &
\fni{textnormal}\verb|{...}|    &    \textnormal{por omisión}
\end{tabular}

\bigskip
\end{lined}
\end{table}


\begin{table}[!bp]
\index{font size}
\caption{Tamaños de \fontnomo{}.} \label{sizes}
\begin{lined}{12cm}
\begin{tabular}{@{}ll}
\fni{tiny}      & \tiny        \fontnomo{} minúscula \\
\fni{scriptsize}   & \scriptsize  \fontnomo{} muy pequeña\\
\fni{footnotesize} & \footnotesize  bastante pequeña \\
\fni{small}        &  \small          \fontnomo{} pequeña \\
\fni{normalsize}   &  \normalsize  \fontnomo{} normal \\
\fni{large}        &  \large       \fontnomo{} grande
\end{tabular}%
\qquad\begin{tabular}{ll@{}}
\fni{Large}        &  \Large       más grande \\[5pt]
\fni{LARGE}        &  \LARGE       muy grande \\[5pt]
\fni{huge}         &  \huge        enorme \\[5pt]
\fni{Huge}         &  \Huge        la más
\end{tabular}

\bigskip
\end{lined}
\end{table}

\begin{table}[!tbp]
\caption{Tamaños absolutos en puntos para las clases normales.}\label{tab:pointsizes}
\label{tab:sizes}
\begin{lined}{12cm}
\begin{tabular}{lrrr}
\multicolumn{1}{c}{tamaño} &
\multicolumn{1}{c}{10pt (por omisión) } &
           \multicolumn{1}{c}{opción 11pt}  &
           \multicolumn{1}{c}{opción 12pt}\\
\verb|\tiny|       & 5pt  & 6pt & 6pt\\
\verb|\scriptsize| & 7pt  & 8pt & 8pt\\
\verb|\footnotesize| & 8pt & 9pt & 10pt \\
\verb|\small|        & 9pt & 10pt & 11pt \\
\verb|\normalsize| & 10pt & 11pt & 12pt \\
\verb|\large|      & 12pt & 12pt & 14pt \\
\verb|\Large|      & 14pt & 14pt & 17pt \\
\verb|\LARGE|      & 17pt & 17pt & 20pt\\
\verb|\huge|       & 20pt & 20pt & 25pt\\
\verb|\Huge|       & 25pt & 25pt & 25pt\\
\end{tabular}

\bigskip
\end{lined}
\end{table}


\begin{table}[!bp]
\caption{\Fontsnomo{} para mates.} \label{mathfonts}
\begin{lined}{0.7\textwidth}
\begin{tabular}{@{}ll@{}}
\fni{mathrm}\verb|{...}|&     $\mathrm{Fundici\acute{o}n\ Rematada}$\\
\fni{mathbf}\verb|{...}|&     $\mathbf{Fundici\acute{o}n\ Negrita}$\\
\fni{mathsf}\verb|{...}|&     $\mathsf{Fundici\acute{o}n\ Palo\ Seco}$\\
\fni{mathtt}\verb|{...}|&     $\mathtt{Fundici\acute{o}n\ De\
  M\acute{a}quina}$\\
\fni{mathit}\verb|{...}|&     $\mathit{Fundici\acute{o}n\ Cursiva}$\\
\fni{mathcal}\verb|{...}|&    $\mathcal{FUNDICI\acute{O}N\ CALIGR\acute{A}FICA}$\\
\fni{mathnormal}\verb|{...}|& $\mathnormal{Fundici\acute{o}n\ Normal}$\\
\end{tabular}

%\begin{tabular}{@{}lll@{}}
%\textit{Command}&\textit{Example}&    \textit{Output}\\[6pt]
%\fni{mathcal}\verb|{...}|&    \verb|$\mathcal{B}=c$|&     $\mathcal{B}=c$\\
%\fni{mathscr}\verb|{...}|&    \verb|$\mathscr{B}=c$|&     $\mathscr{B}=c$\\
%\fni{mathrm}\verb|{...}|&     \verb|$\mathrm{K}_2$|&      $\mathrm{K}_2$\\
%\fni{mathbf}\verb|{...}|&     \verb|$\sum x=\mathbf{v}$|& $\sum x=\mathbf{v}$\\
%\fni{mathsf}\verb|{...}|&     \verb|$\mathsf{G\times R}$|&        $\mathsf{G\times R}$\\
%\fni{mathtt}\verb|{...}|&     \verb|$\mathtt{L}(b,c)$|&   $\mathtt{L}(b,c)$\\
%\fni{mathnormal}\verb|{...}|& \verb|$\mathnormal{R_{19}}\neq R_{19}$|&
%$\mathnormal{R_{19}}\neq R_{19}$\\
%\fni{mathit}\verb|{...}|&     \verb|$\mathit{ffi}\neq ffi$|& $\mathit{ffi}\neq ffi$
%\end{tabular}

\bigskip
\end{lined}
\end{table}

En relación a las órdenes de tamaño de \fontnomo{}, las \wi{llaves} representan un papel significativo.  Se usan para construir \emph{grupos}.  Los grupos limitan el alcance de la mayoría de las órdenes de \LaTeX{}.\index{grupos}\index{agrupar}

\begin{example}
Adora los {\LARGE grandes y
{\small pequeños} placeres}. 
\end{example}
 
Las órdenes de tamaño de \fontnomo{} también cambian el espaciado entre renglones, pero sólo si el párrafo termina dentro del ámbito de la orden de tamaño de \fontnomo{}.  La llave de cierre \verb|}| debería por tanto no llegar demasiado pronto.  Fíjese en la posición de la orden \ci{par} en los siguientes dos ejemplos.\footnote{\texttt{\bs{}par} equivale a un renglón en blanco.}


\begin{example}
{\Large ¡No lea esto!
 No es verdad.
 ¡Puede creerme!\par}
\end{example}

\begin{example}
{\Large Tampoco esto es verdad.
Mas recuerde qué mendaz soy.}\par
\end{example}

Si quiere activar una orden de cambio de tamaño para un párrafo entero de texto o incluso más, puede usar la sintaxis de entorno para las órdenes de cambio de \fontnomo{}.

\begin{example}
\begin{Large} 
Esto no es verdad, pero
qué diantres cabe esperar
en estos tiempos...\par
\end{Large}
\end{example}

Esto le ahorrará andar contando llaves.

\subsection{¡Atención: peligro!}

Como se comenta al principio de este capítulo, es peligroso sembrar el documento con órdenes explícitas como ésas, pues funcionan contra la idea básica de \LaTeX{}, que es separar la estructura de su documento del aspecto visual.  Esto significa que si usted usa la misma orden de cambio de \fontnomo{} en varios lugares para componer un tipo especial de información, debería usar \verb|\newcommand| para definir una ``orden lógica encubridora'' para la orden de cambio de \fontnomo{}.

\begin{example}
\newcommand{\ojo}[1]{%
 \textbf{#1}}
No \ojo{entre} en esta sala; está 
ocupada por \ojo{máquinas} de 
origen y propósito desconocidos.
\end{example}

Este enfoque tiene la ventaja de que usted puede decidir en una etapa posterior que quiere usar alguna representación visual de peligro distinta de \verb|\textbf|, sin tener que recorrer todo el documento identificando cada aparición de \verb|\textbf| y después deduciendo si ahí se usó para señalar un peligro o por alguna otra razón.


\subsection{Consejo}

Para concluir este viaje al mundo de las \fontsnomo{} y sus tamaños, acepte este humilde consejo:\nopagebreak

\begin{quote}
  \underline{\textbf{{\Huge¡}Recuerde\Huge!}} \textit{Cuantas}
  \textsf{M\textbf{\LARGE Á} \texttt{S}} \textsl{\fontsnomo{}} \Huge use
  \tiny en \footnotesize \textbf{un}  \small \texttt{documento},
  \large \textit{tanto} \normalsize más \textsc{legible} y
  \textsl{\textsf{lindo} \large s\Large e\LARGE r\huge á}.
\end{quote}

\section{Espaciado}
 
\subsection{Espacio entre renglones}

\index{espacio entre renglones} Si quiere usar mayor espacio entre renglones, puede cambiar su valor poniendo la orden
\begin{lscommand}
\ci{linespread}\verb|{|\emph{factor}\verb|}|
\end{lscommand}
en el preámbulo de su documento. Use \verb|\linespread{1.3}| para espaciado de ``uno y medio'' y \verb|\linespread{1.6}| para espaciado ``doble''.  Normalmente los renglones no se separan, así que el factor por omisión es~1.\index{doble espaciado de renglones}

Tenga en cuenta que el efecto de la orden \ci{linespread} es bastante drástico y no apropiado para publicar un trabajo.  Así que si tiene una buena razón para cambiar el espacio entre renglones quizá prefiera usar la orden:
\begin{lscommand}
\verb|\setlength{\baselineskip}{1.5\baselineskip}|
\end{lscommand}

\begin{example}
{\setlength{\baselineskip}%
           {1.5\baselineskip}
Este párrafo está compuesto con
el salto de línea base puesto a
1,5 de lo que era antes.  Fíjese
en la orden par al final del
párrafo.\par}

Este párrafo tiene un propósito
claro: mostrar que, una vez se
cierran las llaves, todo vuelve
a la normalidad.
\end{example}

\subsection{Formato de párrafo}\label{parsp}

En \LaTeX{}, hay dos parámetros que influyen en el aspecto del párrafo.  Poniendo una definición
\begin{code}
\ci{setlength}\verb|{|\ci{parindent}\verb|}{0pt}| \\
\verb|\setlength{|\ci{parskip}\verb|}{1ex plus 0.5ex minus 0.2ex}|
\end{code}
en el preámbulo del \filenomo{} de entrada, puede cambiar el aspecto de los párrafos.  Estas dos órdenes incrementan el espacio entre dos párrafos y establecen la sangría de párrafo a cero.

Las partes \texttt{plus} y \texttt{minus} de la longitud de arriba dicen a \TeX{} que puede comprimir y expandir el salto entre párrafos la cantidad indicada, si es necesario para ajustar apropiadamente los párrafos en la página.

En algunos países europeos los párrafos suelen separarse algo y no se sangran.  Pero tenga en cuenta que esto tiene su efecto en el índice general; sus renglones se espaciarán más en ese caso.  Para evitarlo, puede mover las dos órdenes del preámbulo a un lugar en su documento detrás de la orden \verb|\tableofcontents| o no usarlo en absoluto, porque verá que muchos libros profesionales usan sangría y no espacio para separar párrafos.

Si quiere sangrar un párrafo que no está sangrado, puede usar
\begin{lscommand}
\ci{indent}
\end{lscommand}
al principio del párrafo. Obviamente,  sólo tendrá efecto cuando \verb|\parindent| no valga cero. Para sangrar el primer párrafo tras cada título de sección, use el paquete \pai{indentfirst} del lote `tools'.

Para crear un párrafo no sangrado, puede usar
\begin{lscommand}
\ci{noindent}
\end{lscommand}
como primera orden del párrafo.  Puede ser útil si empieza un documento con texto de párrafo y no con una orden de sección.

\subsection{Espacio horizontal}
\label{sec:hspace}

\LaTeX{} determina los espacios entre palabras y automáticamente.  Para añadir espacio horizontal, use:
\index{horizontal!espacio}
\begin{lscommand}
\ci{hspace}\verb|{|\emph{longitud}\verb|}|
\end{lscommand}
Si dicho espacio debiera mantenerse incluso si cae al final o al principio de renglón, use \verb|\hspace*| en lugar de \verb|\hspace|. La \emph{longitud} en el caso más simple es sólo un número más una unidad.  Las unidades más importantes se listan en el cuadro~\ref{units}. \index{unidades}\index{dimensiones}

\begin{example}
Éste\hspace{1.5cm}es un espacio
de 1,5 cm. 
\end{example}
\suppressfloats
\begin{table}[tbp]
\caption{Unidades \TeX{}.} \label{units}\index{unidades}
\begin{lined}{9.5cm} 
\begin{tabular}{@{}ll@{}}
\texttt{mm} & milímetro $\approx 1/25$~pulgada \quad \demowidth{1mm} \\
\texttt{cm} & centímetro = 10~mm  \quad \demowidth{1cm}                     \\
\texttt{in} & pulgada $=$ 25,4~mm \quad \demowidth{1in}                    \\
\texttt{pt} & punto $\approx 1/72$~pulgada $\approx \frac{1}{3}$~mm  \quad\demowidth{1pt}\\
\texttt{em} & $\approx$ anchura de una `M' en la \fontnomo{} actual \quad \demowidth{1em}\\
\texttt{ex} & $\approx$ altura de una `x' en la \fontnomo{} actual \quad \demowidth{1ex}
\end{tabular}

\bigskip
\end{lined}
\end{table}

\label{cmd:stretch} 
La orden
\begin{lscommand}
\ci{stretch}\verb|{|\emph{n}\verb|}|
\end{lscommand} 
genera espacio especial, que se expande hasta llenar todo el espacio sobrante en un renglón.  Si dos órdenes \verb|\hspace{\stretch{|\emph{n}\verb|}}| tienen lugar en el mismo renglón, los espacios crecen proporcionalmente a sus argumentos.

\begin{example}
x\hspace{\stretch{1}}
x\hspace{\stretch{3}}x
\end{example}

Al usar espacio horizontal junto con texto, puede tener sentido hacer que el espacio ajuste su tamaño en relación con el tamaño de la \fontnomo{} actual.  Esto puede hacerse usando las unidades relativas a la \fontnomo{} \texttt{em} y \texttt{ex}:

\begin{example}
{\Large{}gran\hspace{1em}y}\\
{\tiny{}pequeña\hspace{1em}y}
\end{example}
 
\subsection{Espacio vertical}

\LaTeX{} determina automáticamente el espacio entre párrafos, secciones, subsecciones, etc. Si es necesario, puede añadirse espacio vertical adicional \emph{entre dos párrafos} con la orden:
\begin{lscommand}
\ci{vspace}\verb|{|\emph{longitud}\verb|}|
\end{lscommand}

Esta orden debería usarse normalmente entre dos renglones vacíos.  Si el espacio debe preservarse en lo alto o en lo bajo de la página, use la versión  de la orden con asterisco, \verb|\vspace*|, en lugar de \verb|\vspace|. \index{vertical!espacio}

La orden \verb|\stretch|, acompañada de \verb|\pagebreak|, puede usarse para escribir texto en el último renglón de una página, o para centrar texto verticalmente en una página.
\begin{code}
\begin{verbatim}
Algo de texto...

\vspace{\stretch{1}}
Esto va en la última línea de la página. \pagebreak
\end{verbatim}
\end{code}

Espacio adicional entre dos líneas del \emph{mismo} párrafo o dentro de una tabla se indica con la orden
\begin{lscommand}
\ci{\bs}\verb|[|\emph{longitud}\verb|]|
\end{lscommand}


Con  \ci{bigskip} y \ci{smallskip} puede saltar una cantidad predefinida de espacio vertical sin tener que preocuparse de números exactos.

\section{Composición de la página}

\begin{figure}[!hp]
\begin{center}
\makeatletter\@mylayout\makeatother
\end{center}
\vspace*{3ex}
\caption{Parámetros de composición de la página.}
\label{fig:layout}
\cih{footskip}
\cih{headheight}
\cih{headsep}
\cih{marginparpush}
\cih{marginparsep}
\cih{marginparwidth}
\cih{oddsidemargin}
\cih{paperheight}
\cih{paperwidth}
\cih{textheight}
\cih{textwidth}
\cih{topmargin}
\end{figure}
\index{página!composición}
\LaTeXe{} le permite indicar el \wi{tamaño del papel} en la orden \\ \verb|\documentclass|.  Después calcula los \wi{márgenes} adecuados, pero a veces usted no estará contento con los valores predefinidos. Naturalmente, puede cambiarlos. %no idea why this is needed here ...

\thispagestyle{fancyplain}
La figura~\ref{fig:layout} muestra todos los parámetros que pueden cambiarse.  La figura se creó con el paquete \pai{layout} del lote `tools'. \footnote{\CTANref|macros/latex/required/tools|}

\textbf{¡ESPERE!} Antes de lanzarse al frenesí de  ``Hagamos esa página estrecha un poco más ancha'', dedique unos segundos a pensar. Como muchas cosas en \LaTeX, hay una buena razón para que el aspecto de la página sea como es.

Por supuesto, comparada con su página recién salida de un paquete ofimático (como OpenOffice Writer o  MS Word), parece horrorosamente estrecha.  Pero eche un vistazo a su libro favorito\footnote{Me refiero a un libro real impreso y producido por una editorial con reputación.} y cuente el número de caracteres en una línea de texto normal.  Hallará que no hay más de en torno a 66 caracteres en cada renglón.  Ahora haga lo mismo con su página de \LaTeX{}; verá lo mismo.  La experiencia muestra que la lectura se vuelve difícil en cuanto hay más caracteres por renglón.  Es así porque a los ojos les resulta difícil moverse desde el final de un renglón al principio del siguiente.  Es la misma razón por la que los periódicos se componen en múltiples columnas. Así que si incrementa la anchura de su texto, tenga en cuenta que está haciendo la vida más difícil a los lectores de su documento.
 
Si de cualquier forma quiere hacerlo, \LaTeX{} proporciona dos órdenes para cambiar estos parámetros.  Se usan normalmente en el preámbulo del documento.

La primera orden asigna un valor fijo a cualquiera de los parámetros:
\begin{lscommand}
\ci{setlength}\verb|{|\emph{parámetro}\verb|}{|\emph{longitud}\verb|}|
\end{lscommand}

La segunda orden añade longitud a cualquier parámetro:
\begin{lscommand}
\ci{addtolength}\verb|{|\emph{parámetro}\verb|}{|\emph{longitud}\verb|}|
\end{lscommand} 

Esta segunda orden es de hecho más útil que la orden \ci{setlength}, pues puede usted así trabajar en relación a las valores establecidos. Para añadir un centímetro a la anchura total del texto, pongo las siguientes órdenes en el preámbulo del documento:
\begin{code}
\verb|\addtolength{\hoffset}{-0.5cm}|\\
\verb|\addtolength{\textwidth}{1cm}|
\end{code}

En este  contexto, quizá quiera mirar el paquete \pai{calc}.  Le permite usar operaciones aritméticas en el argumento de \ci{setlength} y en otros lugares donde puede introducir valores numéricos en argumentos de funciones.

\section{Más diversión con las longitudes}

Siempre que sea posible, evite usar longitudes absolutas en los documentos \LaTeX{}.  Intente basar las cosas en la anchura o altura de otros elementos de la página.  Para la anchura de una figura puede referirse a \verb|\textwidth| al componer la página.

Las siguientes 3 órdenes le permiten determinar la anchura, altura y profundidad de una cadena de texto.

\begin{lscommand}
\ci{settoheight}\verb|{|\emph{variable}\verb|}{|\emph{texto}\verb|}|\\
\ci{settodepth}\verb|{|\emph{variable}\verb|}{|\emph{texto}\verb|}|\\
\ci{settowidth}\verb|{|\emph{variable}\verb|}{|\emph{texto}\verb|}|
\end{lscommand}

El ejemplo siguiente muestra una posible aplicación de estas órdenes.

\begin{example}
\flushleft
\newenvironment{vardesc}[1]{%
  \settowidth{\parindent}{#1:\ }
  \makebox[0pt][r]{#1:\ }}{}

\begin{displaymath}
a^2+b^2=c^2
\end{displaymath}

\begin{vardesc}{Donde}$a$, 
$b$ -- son adyacentes al ángulo
recto de un triángulo rectángulo.

$c$ -- es la hipotenusa del
triángulo, y 

$d$ -- no sale aquí
en absoluto. 
\end{vardesc}
\end{example}

\section{Cajas}

\LaTeX{} construye sus páginas colocando cajas.  En principio, cada letra es una cajita, que se pega a otras letras para formar palabras. Éstas se pegan de nuevo a otras palabras, pero con un pegamento especial, que es tan elástico que una serie de palabras puede comprimirse o expandirse para rellenar exactamente un renglón de la página.

Esto es una simplificación de lo que realmente ocurre, pero realmente ocurre: \TeX{} trabaja con pegamento y cajas. Las letras no son las únicas cosas que son cajas.  Puede poner virtualmente cualquier cosa en una caja, incluso otras cajas.  Cada caja será manejada por \LaTeX{} como si fuera una simple letra.

En los capítulos anteriores ya ha encontrado algunas cajas, aunque no lo parezcan.  Los entornos \ei{tabular} e \ci{includegraphics}, por ejemplo, producen cajas.  Esto significa que puede usted fácilmente colocar dos tablas o imágenes una al lado de la otra.  Basta con asegurarse de que su anchura combinada no excede la anchura del texto.

Puede también empaquetar un párrafo de su elección en una caja con la orden

\begin{lscommand}
\ci{parbox}\verb|[|\emph{pos}\verb|]{|\emph{anchura}\verb|}{|\emph{texto}\verb|}|
\end{lscommand}

o el entorno

\begin{lscommand}
\verb|\begin{|\ei{minipage}\verb|}[|\emph{pos}\verb|]{|\emph{anchura}\verb|}| texto 
\verb|\end{|\ei{minipage}\verb|}|
\end{lscommand}

El parámetro \texttt{pos} puede tomar una de las letras \texttt{c, t} o \texttt{b} para controlar la alineación vertical de la caja, relativa a la línea base del texto que la rodea. \texttt{anchura} toma como argumento la longitud que indica la anchura de la caja.  La principal diferencia entre una \ei{minipage} y una \ci{parbox} es que usted no puede usar todas las órdenes y entornos dentro de una \ei{parbox}, mientras que casi todo es posible en una \ei{minipage}.

Mientras que \ci{parbox} empaqueta un párrafo entero partiendo renglones y todo, hay también una clase de órdenes encajonadoras que trabajan sólo con material alineado horizontalmente.  Ya conocemos una de ellas; se llama \ci{mbox}.  Simplemente empaqueta una serie de cajas en otra, y puede usarse para impedir a \LaTeX{} romper dos palabras.  Como puede  poner cajas dentro de cajas, estos empaquetadores de cajas horizontales le dan total flexibilidad.

La orden
\begin{lscommand}
\ci{makebox}\verb|[|\emph{anchura}\verb|][|\emph{pos}\verb|]{|\emph{texto}\verb|}|
\end{lscommand}

donde \texttt{anchura} define la anchura de la caja resultante vista desde fuera,\footnote{Esto significa que puede ser más pequeña que el material dentro de ella.  Usted puede incluso poner la anchura 0pt de forma que el texto de dentro de la caja se componga sin afectar a las cajas de alrededor.}  tiene un efecto parecido. Además de las expresiones de longitud, puede también usar \ci{width}, \ci{height}, \ci{depth} y \ci{totalheight} en el parámetro de anchura.  Se establecen a partir de valores obtenidos midiendo el \emph{texto} compuesto.  El parámetro \emph{pos} toma una letra como valor: \textbf{c}enter (centro), flush\textbf{l}eft (izquierda), flush\textbf{r}ight (derecha) o \textbf{s}pread (expandir el texto hasta llenar la caja).

La orden \ci{framebox} funciona exactamente igual que \ci{makebox}, pero dibuja una caja alrededor del texto.

El ejemplo siguiente le muestra algunas cosas que podría hacer con las órdenes \ci{makebox} y \ci{framebox}.

\begin{example}
\makebox[\textwidth]{%
    c e n t r a d o}\par
\makebox[\textwidth][s]{%
    e x p a n d i d o}\par
\framebox[1.1\width]{A la medida} \par
\framebox[0.8\width][r]{Muy ancho} \par
\framebox[1cm][l]{Y otro también...} 
¿Puede leer esto?
\end{example}

Ahora que controlamos lo horizontal, el siguiente paso obvio es ir por la vertical.\footnote{El control total sólo se obtiene controlando tanto lo horizontal como lo vertical\ldots}

La orden
\begin{lscommand}
\ci{raisebox}\verb|{|\emph{sube}\verb|}[|\emph{extiende-sobre-línea-base}\verb|][|\emph{extiende-bajo-línea-base}\verb|]{|\emph{texto}\verb|}|
\end{lscommand}
le permite definir las propiedades verticales de una caja. Puede usar \ci{width}, \ci{height}, \ci{depth} y \ci{totalheight} en los tres primeros parámtros, para afectar al tamaño de la caja dentro del argumento \emph{texto}.


\begin{example}
\raisebox{0pt}[0pt][0pt]{\Large%
\textbf{Aaaa\raisebox{-0.3ex}{a}%
\raisebox{-0.7ex}{aa}%
\raisebox{-1.2ex}{h}%
\raisebox{-2.2ex}{h}%
\raisebox{-4.5ex}{h}}}
---gritó, pero ni siquiera el más
próximo se dio cuenta de que 
algo terrible le había sucedido...
\end{example}

\section{Líneas y puntales}
\label{sec:rule}

Hace unas páginas puede haber visto la orden

\begin{lscommand}
\ci{rule}\verb|[|\emph{sube}\verb|]{|\emph{anchura}\verb|}{|\emph{altura}\verb|}|
\end{lscommand}

Usada normalmente produce simplemente una caja negra.

\begin{example}
\rule{3mm}{.1pt}%
\rule[-1mm]{5mm}{1cm}%
\rule{3mm}{.1pt}%
\rule[1mm]{1cm}{5mm}%
\rule{3mm}{.1pt}
\end{example}

Esto es útil para dibujar líneas verticales y horizontales. La línea de la página del título, por ejemplo, ha sido creada con una orden \ci{rule}.

Un caso especial es una línea sin anchura pero con cierta altura.  En composición profesional se llama \wi{puntal}.  Se usa para garantizar que un elemento de una página tiene una cierta altura mínima.  Podría usarlo en un entorno \texttt{tabular} para asegurarse de que una fila tiene cierta altura mínima.

\begin{example}
\begin{tabular}{|c|}
\hline
\rule{1pt}{4ex}Costeru...\\
\hline
\rule{0pt}{4ex}Puntal\\
\hline
\end{tabular}
\end{example}

\bigskip
{\flushright Fin.\par}

% Local Variables:
% TeX-master: "lshort2e"
% mode: latex
% mode: flyspell
% End: